\documentclass[11pt]{article}
\usepackage[T1]{fontenc}
\usepackage{lmodern}
\usepackage[margin=1.02in]{geometry}
\usepackage{amsmath,amssymb,amsthm,mathtools,mathrsfs}
\usepackage{bm,enumitem,booktabs,array,tabularx}
\usepackage{microtype,aliascnt,needspace}
\usepackage[colorlinks=true,linkcolor=blue,citecolor=blue,urlcolor=blue]{hyperref}
\usepackage[nameinlink,capitalise]{cleveref}
\setlength{\emergencystretch}{2em}
\raggedbottom
\numberwithin{equation}{section}
\newtheorem{theorem}{Theorem}[section]
\newaliascnt{proposition}{theorem}
\newtheorem{proposition}[proposition]{Proposition}
\aliascntresetthe{proposition}
\newaliascnt{lemma}{theorem}
\newtheorem{lemma}[lemma]{Lemma}
\aliascntresetthe{lemma}
\newaliascnt{corollary}{theorem}
\newtheorem{corollary}[corollary]{Corollary}
\aliascntresetthe{corollary}
\theoremstyle{definition}
\newaliascnt{definition}{theorem}
\newtheorem{definition}[definition]{Definition}
\aliascntresetthe{definition}
\newaliascnt{assumption}{theorem}
\newtheorem{assumption}[assumption]{Assumption}
\aliascntresetthe{assumption}
\newaliascnt{example}{theorem}
\newtheorem{example}[example]{Example}
\aliascntresetthe{example}
\theoremstyle{remark}
\newaliascnt{remark}{theorem}
\newtheorem{remark}[remark]{Remark}
\aliascntresetthe{remark}
\crefname{theorem}{Theorem}{Theorems}
\Crefname{theorem}{Theorem}{Theorems}
\crefname{proposition}{Proposition}{Propositions}
\Crefname{proposition}{Proposition}{Propositions}
\crefname{lemma}{Lemma}{Lemmas}
\Crefname{lemma}{Lemma}{Lemmas}
\crefname{corollary}{Corollary}{Corollaries}
\Crefname{corollary}{Corollary}{Corollaries}
\crefname{definition}{Definition}{Definitions}
\Crefname{definition}{Definition}{Definitions}
\crefname{assumption}{Assumption}{Assumptions}
\Crefname{assumption}{Assumption}{Assumptions}
\crefname{example}{Example}{Examples}
\Crefname{example}{Example}{Examples}
\crefname{remark}{Remark}{Remarks}
\Crefname{remark}{Remark}{Remarks}
\newcommand{\R}{\mathbb R}
\newcommand{\C}{\mathbb C}
\newcommand{\N}{\mathbb N}
\newcommand{\Z}{\mathbb Z}
\newcommand{\dd}{\mathrm d}
\newcommand{\ii}{\mathrm i}
\newcommand{\id}{\mathrm{id}}
\newcommand{\Tr}{\operatorname{Tr}}
\newcommand{\Rea}{\operatorname{Re}}
\newcommand{\supp}{\operatorname{supp}}
\newcommand{\rank}{\operatorname{rank}}
\newcommand{\Ker}{\operatorname{Ker}}
\newcommand{\Ran}{\operatorname{Ran}}
\newcommand{\diag}{\operatorname{diag}}
\newcommand{\Sym}{\operatorname{Sym}}
\newcommand{\End}{\operatorname{End}}
\newcommand{\SU}{\mathrm{SU}}
\newcommand{\U}{\mathrm U}
\newcommand{\HS}{\mathrm{HS}}
\newcommand{\Ric}{\operatorname{Ric}}
\newcommand{\Riem}{\operatorname{Riem}}
\newcommand{\Scal}{\operatorname{Scal}}
\newcommand{\vol}{\operatorname{vol}}
\newcommand{\norm}[1]{\lVert#1\rVert}
\newcommand{\abs}[1]{\lvert#1\rvert}
\newcommand{\ip}[2]{\langle#1,#2\rangle}
\newcommand{\Dtest}{\mathcal D}
\newcommand{\A}{\mathcal A}
\newcommand{\M}{\mathcal M}
\newcommand{\Sact}{\mathscr S}
\newcommand{\Vtest}{\mathscr V}
\newcommand{\Fcurv}{\mathsf F}
\newcommand{\Eins}{\mathsf G}
\newcommand{\Evar}{\mathcal E}
\newcommand{\bPsi}{\overline\Psi}
\newcommand{\dV}{\dd V}
\newcommand{\wt}{\widetilde}
\newcommand{\Iface}{\mathfrak I}
\newcolumntype{L}[1]{>{\raggedright\arraybackslash}p{#1}}


\title{\textbf{Variational Closure of Finite Einstein--Standard-Model Actions}}


\author{Aurelien Pelissier$^{\dagger}$\\
{\small $^{\dagger}$Correspondence: \href{mailto:aurelien.pelissier.38@gmail.com}{aurelien.pelissier.38@gmail.com}}}
\date{}


% Presentation revision V1. All technical material is included below.
% Author-side submission checks: affiliation, funding/competing-interest
% declarations, and any applicable account of assisted preparation must be
% completed from the author's own information before submission.
\hypersetup{pdftitle={Variational Closure of Finite Einstein--Standard-Model Actions},
 pdfauthor={Aurelien Pelissier},bookmarksnumbered=true}
\begin{document}
\maketitle

\begin{abstract}
We prove a variational closure theorem for finite Lorentzian gravity--matter systems with Standard-Model field content and a common action. Under compactness, first-variation consistency, and physically scaled stationarity, subsequences of reconstructed fields converge to distributional solutions of the classical Einstein--Standard-Model equations. The argument requires neither uniform coframe convergence nor strong spinor-derivative convergence, while compactness of the covariant Higgs gradient suffices to close the quartic sector. When bosonic compactness fails, the limit retains explicit stress defects whose vanishing is characterized by positive energy measures. For an explicit finite action, a gauge-invariant Wilson-transport criterion yields zero-defect closure at critical gauge regularity. On regular branches we also obtain quantitative convergence and an exact solution shadowing a single finite record from corrected initial data, without prescribing its future trajectory. Renewal Geometry provides one realization of the finite interface, but the closure theorem is independent of that construction.
\end{abstract}

\medskip
\noindent\textbf{Keywords.} Einstein--Standard-Model system; finite gravity--matter actions; variational closure; continuum limits; nonlinear compactness; stress defects; hyperbolic stability; Renewal Geometry.






% =============================================================================
\section{Introduction}
\label{sec:intro}
% =============================================================================

The Standard Model and general relativity organize two different kinds of
geometric structure.  General relativity identifies gravitation with the
Lorentzian geometry of spacetime, whereas the Standard Model introduces an
internal gauge structure, chiral fermion representations, a Higgs sector,
Yukawa couplings, and a three-generation matter multiplicity over an already
given spacetime.  In the usual continuum formulation these ingredients are
specified before the coupled field equations are written.

A natural finite-to-continuum question is whether a gravity--matter theory
specified at finite resolution can recover the classical
Einstein--Standard-Model system without inserting the limiting field equations
by construction.  This question is logically distinct from the microscopic
origin of the finite field content itself.  A finite model may carry the
correct Lorentzian signature, gauge group, matter representations, and common
gravity--matter action while nevertheless failing to possess a controlled
continuum limit.  Conversely, a convergent discretization of a prescribed
continuum theory need not explain why those particular structural data were
present in the first place.

The present paper isolates this analytic closure problem.  We formulate a
finite Lorentzian--Standard-Model interface containing the structural bundle
data, finite Dirac/Yukawa incidence, positive source geometry, and one common
gravity--matter action.  Reconstruction maps, physical variation maps,
compactness, and stationarity are imposed separately.  The principal question
is then whether these finite data have the classical
Einstein--Yang--Mills--Higgs--Dirac--Yukawa equations as a variational
continuum limit.

\paragraph{Geometric and variational context.}

There is a long history of describing Standard-Model structure in geometric
terms.  In noncommutative geometry, the Connes--Lott construction realizes
electroweak and Standard-Model data through a finite noncommutative geometry
\cite{ConnesLott1991}, finite spectral triples admit a diagrammatic
classification \cite{Krajewski1998}, and the spectral action produces
gravitational, Yang--Mills, Higgs, and fermionic terms from one spectral datum
\cite{ChamseddineConnes1997}.  The almost-commutative construction was later
extended to include neutrino mixing and Majorana masses
\cite{ChamseddineConnesMarcolli2007}; Lorentzian formulations of the
noncommutative Standard-Model geometry have also been investigated, including
the fermionic and doubling issues specific to indefinite signature
\cite{Barrett2007}.  Causal fermion systems provide a different finite or
microscopic entrance in which effective gauge and gravitational dynamics
arise from a causal action principle \cite{Finster2016}.

A complementary literature concerns the preservation of geometric and
variational structure under discretization.  Regge calculus initiated a
simplicial treatment of gravitational curvature \cite{Regge1961};
piecewise-flat curvature measures were subsequently shown to converge under
refinement \cite{CheegerMullerSchrader1984}, and more recent work has placed
linearized Regge calculus, densitized scalar curvature, and the Einstein
tensor in rigorous finite-element and distributional frameworks
\cite{ChristiansenRegge2011,ChristiansenRegge2024,
GawlikNeunteufelScalar2025,GawlikNeunteufelEinstein2026}.  Wilson lattice
gauge theory retains exact finite gauge invariance \cite{Wilson1974}, while
discrete variational field theories preserve action-derived multisymplectic
and Noether structures \cite{MarsdenPatrickShkoller1998}.  Finite element
exterior calculus provides stable discrete differential complexes
\cite{ArnoldFalkWinther2006,ArnoldFalkWinther2010}.  Gauge-invariant
Maxwell and Maxwell--Klein--Gordon discretizations give complementary
scheme-specific consistency and convergence results
\cite{ChristiansenHalvorsen2009,ChristiansenHalvorsen2011,
ChristiansenHalvorsenScheid2023}, and simplicial gauge actions can combine
finite gauge invariance, consistency, and discrete Noether identities
\cite{ChristiansenHalvorsen2012}.  In these approaches the continuum
geometric or gauge structure is specified before the discretization.  The
formulation used here instead isolates structural, compactness,
first-variation consistency and stationarity hypotheses that can be checked
for different finite gravity--matter realizations.  It is therefore
complementary to, rather than a replacement for, scheme-specific stability
and convergence analyses.

\paragraph{One upstream realization: Renewal Geometry.}

Renewal Geometry provides one realization of the finite interface considered
below.  Predictive-state constructions characterize state by the information
required for future prediction rather than by a separately postulated hidden
coordinate \cite{ShaliziCrutchfield2001}, while multitime operational
formalisms provide a related language for processes with memory
\cite{ChiribellaEtAl2009,PollockEtAl2018}.  In the Renewal Geometry
construction developed in
\cite{PelissierRenewalSpacetime2026,PelissierSpectral2026}, operational
histories are quotiented by equality of conditional future statistics, and the
resulting minimum persistent carrier
$Z_X^{\min}$, represented history algebra
$\mathcal A_X^{\mathrm{hist}}$, and positive word-Gram hierarchy are organized
into the \emph{Renewal Grand Tensor},
\begin{equation}
\boldsymbol{\mathbb T}_{\mathrm{Grand},X}
=
\left(
\mathcal A_X^{\mathrm{hist}},
\{\mathbb K_{X,r}\}_{r\ge0},
Z_X^{\min}
\right).
\label{eq:intro-grand-tensor}
\end{equation}
The terminology ``tensor'' refers to this finite predictive data structure,
not to a tensor field on a pre-existing manifold.

On certified branches, the corresponding Lorentzian construction provides a
relational carrier, Lorentzian signature, ADM-type data, and gravitational
normalization \cite{PelissierRenewalSpacetime2026}; the spectral construction
provides compatible differential and spin data
\cite{PelissierSpectral2026}; and the structural gravity--matter construction
provides a Standard-Model matter packet and common finite gravity--matter
action \cite{PelissierDuality2026}.  These constructions establish one
nonempty realization of the finite interface.  None of their microscopic
reconstruction arguments is used in place of the compactness or
first-variation estimates proved here.

\paragraph{The continuum obstruction.}

The main analytical difficulty is that neither convergence of fields nor
convergence of action values is sufficient to pass to the coupled classical
equations.  The Yang--Mills stress is quadratic in curvature, the Higgs stress
contains quadratic covariant derivatives and a quartic potential, and the
metric variation of the fermionic action also varies the spin geometry.
Weak convergence of the primitive fields can therefore fail to determine the
limit of the quantities that gravitate.

This issue is familiar from several directions.  Gauge compactness requires
careful control of both gauge choice and curvature \cite{Uhlenbeck1982}, and
at the critical four-dimensional Yang--Mills energy scale minimizing
sequences can develop curvature concentration and bubbling
\cite{Sedlacek1982}.  Critical Sobolev embeddings likewise permit
concentration phenomena \cite{Lions1985}.  More generally, weak limits of
nonlinear PDEs can retain oscillation or concentration information through
measure-valued and microlocal descriptions, including DiPerna--Majda
measures, generalized Young measures, H-measures and microlocal defect
measures \cite{DiPernaMajda1987,AlibertBouchitte1997,Tartar1990,Gerard1991}.
The covariance measures used below are deliberately coarser: they remain in
physical space and retain only the quadratic information needed by the
Yang--Mills and Higgs metric variations.  In general relativity,
high-frequency gravitational waves already produce an effective averaged
stress tensor in the Isaacson expansion \cite{Isaacson1968}; Burnett's weak
limit formulation and its modern developments give a rigorous continuation
of this effective-stress perspective \cite{Burnett1989,HuneauLuk2024}.
Spinors introduce an additional metric-dependence through the spin bundle and
Dirac operator, whose variation can be organized using the metric-spinor
constructions of Bourguignon--Gauduchon \cite{BourguignonGauduchon1992} and
the universal spinor-bundle framework of M\"uller--Nowaczyk
\cite{MullerNowaczyk2017}.  We work throughout on a torsion-free metric
branch; independently varying the Cartan connection would instead lead
generically to Einstein--Cartan dynamics in the presence of spin current
\cite{Kibble1961,HehlEtAl1976}.

For this reason the relevant convergence object is the
\emph{complete physical first variation}.  At cutoff $h$, let
$z_h^{\mathrm d}$ denote the finite configuration and
$z_h=\mathcal R_hz_h^{\mathrm d}$ its reconstructed field tuple.  The finite
first variation evaluated on a lifted physical test must approach the
continuum first variation evaluated on $z_h$, while the finite common-action
residual must vanish in the same physical normalization.  If the native
action covector has squared residual $R_h$, the lift of a unit physical test
has source norm at most $L_h$, and the identification error is $e_h$, then
\begin{equation}
\norm{\mathcal F_h}_{(\Vtest_K^r)^*}
\le
L_h\sqrt{R_h}+e_h.
\label{eq:intro-scaled-stationarity}
\end{equation}
Thus the natural stationarity condition is
$L_h\sqrt{R_h}+e_h\to0$, not merely $R_h\to0$.

\paragraph{Main results.}

We first prove a regulator-independent subsequential variational closure
theorem in a reduced action topology.  For any finite-interface sequence satisfying the stated reduced compactness,
first-variation consistency, and physically scaled stationarity hypotheses,
every cutoff sequence has a subsequence whose reconstructed fields solve
\begin{equation}
\Eins_{\mu\nu}(g)+\Lambda g_{\mu\nu}
=
\kappa T^{\mathrm{SM}}_{\mu\nu},
\qquad
\Evar_A=\Evar_H=\Evar_\Psi=\Evar_{\bPsi}=0
\label{eq:intro-ESM-system-new}
\end{equation}
in distributions.  No candidate continuum solution, limiting stress tensor,
or summable adjacent-cutoff transport is assumed.  It appears as
one additional condition that upgrades subsequential closure to a unique
route-independent limit; uniqueness in a prescribed Cauchy class gives an
alternative without summable transport.

The compactness assumptions can be obtained in several ways.  On a regular
branch, Sobolev bounds imply strong bosonic compactness, while Lorentzian
symmetric-hyperbolic stability gives strong spinor $H^1$ compactness.  At
lower regularity, compact screens control the nonlinear bosonic packets.
The reduced closure criterion further shows that strong compactness of the
covariant Higgs gradient, together with $A_h\to A$ in $L^4$, already closes
the critical Higgs endpoint, whereas bounded weak $H^1$ spinors suffice for
the complete bilinear fermionic first variation.  On the geometric side,
strong $H^1$ convergence of the coframe on a common compact nondegenerate
chart is sufficient for continuity of the complete first variation; no
strong $L^\infty$ convergence of the coframe is needed for zero-defect closure.

We also prove stronger quantitative results for an explicit finite local
gravity--Yang--Mills--Higgs--Dirac action.  A small finite Euler row,
controlled initial and harmonic-gauge mismatch, and measured high-frequency
leakage imply quantitative convergence of the actual reconstructed fields,
curvature, and complete Standard-Model stress.  On a regular data-centred
branch, positive full-record reader control upgrades the finite Euler row to
the differentiated physical forcing topology; a small composite initial
constraint residual is then corrected only at the initial slice, and the
physical Cauchy problem produces a unique exact Einstein--Standard-Model
solution shadowing the observed record on the certified slab.  Thus no exact
future comparison trajectory is supplied in advance.  The same-source
refinement shows how the required finite budgets can be obtained
simultaneously on one actually observed record without replacing that record
by an expectation or a filtered configuration.

Finally, when the bosonic compactness required by the classical theorem
fails, the weak limit retains the unresolved nonlinear information as
explicit Yang--Mills and Higgs stress defects.  The classical
Einstein--Standard-Model system is therefore identified as the zero-defect
branch of a broader weak-limit theory rather than being obtained by formally
discarding oscillation or concentration.  At the critical four-dimensional
Yang--Mills energy endpoint we obtain this defect limit modulo cutoff-dependent
local internal gauges, using an equivariant physical-test topology.  For the
explicit local action, a gauge-invariant Wilson translation screen of the
literal curvature and Higgs-gradient packets yields strong local $L^2$
packet compactness and hence the zero-defect branch without a separate
packet-energy uniform-integrability hypothesis.  For minimal coupling, the
timelike energy contraction independently characterizes vanishing of the
complete bosonic stress defect in terms of strong curvature and
covariant-gradient compactness.  For the explicit finite Higgs packet we also
isolate a strictly stronger higher-integrability certificate: the positive
companion kinetic density is uniformly comparable to $|K_h^H|^2$, so its
$L^{4/3}$ bound is equivalent to an $L^{8/3}$ bound for the covariant Higgs
difference.  Strong $L^2$ compactness alone yields only $L^1$ convergence of
this positive density and does not imply the stronger endpoint.

The results are classical and finite-to-continuum in scope.  They do not
address the construction of a renormalized quantum Standard Model coupled to
gravity, nor do they assert that every microscopic dynamics realizes the
finite interface or selects the controlled branch.  The analysis instead
provides a portable closure theorem: any finite gravity--matter realization
satisfying the stated structural and analytic hypotheses has the corresponding
classical Einstein--Standard-Model system as its variational continuum limit.

\paragraph{Organization.}

\Cref{sec:interface} specifies the finite interface, the continuum action and
its physical variations. \Cref{sec:limit} states the reduced closure theorem
before developing its two endpoint mechanisms and proof.
\Cref{sec:defects} passes from the joint weak-limit identity to positive
defects, critical gauge-quotient closure and the finite Wilson-screen
criterion. \Cref{sec:native-branch} gives the quantitative finite-action
realization and culminates in the data-centred exact-shadow theorem.
\Cref{sec:discussion} discusses interpretation and scope.
The appendices retain the supporting estimates and realization-specific
refinements; the guide preceding them identifies their dependencies.

\section{Finite Einstein--Standard-Model interface}
\label{sec:interface}

The analytic problem begins with finite data carrying one common
gravity--matter action. We first separate these structural data from the
reconstruction and convergence assumptions, then specify the continuum
variation that the finite action must approximate. Renewal Geometry supplies
one structural realization; its representation and source identities are
recorded in \Cref{app:structural-packet} and are not compactness hypotheses.

\begin{definition}[Finite classical Einstein--Standard-Model interface]
\label{def:finite-interface}
At cutoff $h$, a finite classical Einstein--Standard-Model interface $\Iface_h$ consists of the following data.
\begin{enumerate}[label=\textup{(I\arabic*)},leftmargin=3.2em]
\item A finite configuration space $\mathcal Q_h$ with a reconstructed oriented, time-oriented Lorentzian carrier and local coframe variables, together with a fixed structural realization on a smooth comparison cylinder $M$ whenever the reconstruction map below is applied.
\item A fixed finite-rank Standard-Model bundle type with
\begin{equation}
 G_{\mathrm{SM}}=S(\U(3)\times\U(2))
 \cong[\SU(3)\times\SU(2)\times\U(1)]/\Z_6,
 \label{eq:gauge-group}
\end{equation}
the chiral fermion representations of \Cref{tab:SM-representations}, one rank-three generation factor, the Higgs doublet, and any optional neutral/Majorana blocks declared as part of the branch.
\item A gauge-covariant finite Dirac/Yukawa incidence operator $D_{F,h}$ on that fixed structural carrier and a coefficient bank $\theta_h=(\kappa_h,\Lambda_h,g_{1,h},g_{2,h},g_{3,h},\lambda_{H,h},v_{H,h},\boldsymbol Y_h)$.
\item A differentiable common finite action
\[
 \Sact_h=\Sact_{\mathrm g,h}+\Sact_{\mathrm{SM},h}
\]
with one physical relative normalization between gravity and matter and with the corresponding finite gauge/relabeling covariance on the declared branch.
\item A positive source geometry for physical variations, including a faithful source quotient or equivalent positive Gram on the retained variation directions. This is used only through the physical test-lift estimates stated explicitly below; no algebraic commutant gap is treated as an analytic compactness estimate.
\end{enumerate}
No reconstruction-to-continuum map, compactness, stationarity, stress convergence, or field equation is included in this structural definition. Reconstruction maps and physical test lifts are introduced separately in \Cref{def:regulator}.
\end{definition}

\subsection{Fields, bundles and physical tests}

The fixed representation packet and its covariant Yukawa blocks are given
in \Cref{tab:SM-representations,lem:SM-descent}. The analytic arguments below
use their finite rank, unitarity of the internal representations and the
affine Higgs dependence of the Yukawa map, rather than the mechanism by
which that packet was obtained.

We work on an oriented, time-oriented spin cylinder
\begin{equation}
 M=(0,T)\times\Sigma,
 \qquad \partial\Sigma=\varnothing,
 \label{eq:cylinder}
\end{equation}
where $\Sigma$ is a compact smooth three-manifold. All conclusions are local on compact subcylinders. There is no physical timelike boundary in the main theorem, and test variations vanish near the initial and final time slices. A boundary extension requires its own flux and trace estimates; it is not implicit in \eqref{eq:cylinder}.

Choose a smooth auxiliary Riemannian metric $r$ and density $\dV_0$ on $M$. All $L^p$, Sobolev, and Hilbert-space norms refer to these positive comparison data and fixed smooth bundle metrics. Lorentzian contractions occur only in the signed action. On a compact nondegenerate coframe chart the different positive comparison norms are uniformly equivalent. No positivity of the Lorentzian action is assumed.

Let $P\to M$ be a fixed principal $G_{\mathrm{SM}}$ bundle of the type specified by \Cref{def:finite-interface}. The chiral fermion bundle, generation factor, and Higgs bundle $E_H$ are the corresponding associated bundles. We write
\begin{equation}
 z=(e,A,H,\Psi,\bPsi),\qquad
 g_{\mu\nu}=\eta_{ab}e^a_\mu e^b_\nu,
 \quad \eta=\diag(-1,1,1,1).
 \label{eq:fields}
\end{equation}
We use the curvature convention $[\nabla_\mu,\nabla_\nu]v^\rho=R^\rho{}_{\sigma\mu\nu}v^\sigma$ and $\Ric_{\sigma\nu}=R^\rho{}_{\sigma\rho\nu}$. The coframe and spinor expressions are local in a fixed finite atlas, with compatible frame changes on overlaps. A global coframe is not required. The fields $\Psi,\bPsi$ are classical realified spinor variables, varied independently before imposing the physical duality relation. They are not a Grassmann integration measure.

Put
\begin{equation}
 \Fcurv_A=\dd A+A\wedge A,\qquad
 K=D_AH=\dd H+\rho_H(A)H.
 \label{eq:curvature-higgs}
\end{equation}
The spin derivative $\nabla^{e,A}$ uses the Levi--Civita spin connection of $e$ and the retained gauge representation. The map $\mathscr M_{\boldsymbol Y}(H)$ is the finite Dirac/Yukawa coupling on the fixed fermion carrier. In real Higgs coordinates it is affine, with a possible constant Majorana block, and satisfies
\begin{equation}
 \norm{\mathscr M_{\boldsymbol Y}(H)}\le C_{\boldsymbol Y}(1+\abs H),
 \qquad
 \norm{D_H\mathscr M_{\boldsymbol Y}(H)[\eta_H]}
 \le C_{\boldsymbol Y}\abs{\eta_H}.
 \label{eq:yukawa-bound}
\end{equation}
Its covariance is part of the finite incidence interface. If the optional Majorana sector is retained, the fermion carrier includes its charge-conjugate blocks and the realified bilinear is understood on that doubled carrier with its fixed normalization.

\begin{definition}[Physical variation space]
\label{def:tests}
For a compact region $K\Subset M$, let $\Vtest_K$ consist of smooth compactly supported tuples
\[
 v=(k,a,\eta_H,\eta_\Psi,\eta_{\bPsi}),
 \qquad k\in\Dtest(M;\Sym^2 TM),\quad
 a\in\Dtest(M;T^*M\otimes\operatorname{ad}P),
\]
with the remaining entries in the indicated matter bundles. Thus $k^{\mu\nu}=\delta g^{\mu\nu}$. Fix an integer $r_0\ge4$ and give this space the sum of the $C^{r_0}$ norms in the chosen atlas. Its completion is denoted $\Vtest_K^{r_0}$. More generally, $\Vtest_K^r$ denotes the completion in the $C^r$ test norm for an integer $r\ge1$; $r_0$ remains the fixed regulator convention. A determining core means a subspace dense in this test space, with the stated uniform bounds extending to its completion.
\end{definition}

The symmetric coframe lift of $k$ is
\begin{equation}
 \dot e^a_\mu(k)=-\tfrac12e^a_\alpha g_{\mu\beta}k^{\alpha\beta}.
 \label{eq:metric-lift}
\end{equation}
It gives $\dot g_{\mu\nu}=-g_{\mu\alpha}g_{\nu\beta}k^{\alpha\beta}$. Spinor components are transported in the same local spin-frame identification. Antisymmetric coframe variations are local Lorentz gauge directions, not additional metric tests.

\subsection{The action and its complete metric variation}

Let $\theta=(\kappa,\Lambda,g_1,g_2,g_3,\lambda_H,v_H,\boldsymbol Y)$ be a fixed finite coefficient bank, with $\kappa,g_j,\lambda_H>0$. The matter action is
\begin{align}
 \Sact_{\mathrm{SM},\theta}(z)
 &=\int_M\mathcal L_{\mathrm{SM},\theta}(z)\,\dV_g,
 \label{eq:SM-action}\\
 \mathcal L_{\mathrm{SM},\theta}
 &=-\tfrac14\ip{\Fcurv_{\mu\nu}}{\Fcurv^{\mu\nu}}_{\boldsymbol g}
   -g^{\mu\nu}\Rea\ip{D_\mu H}{D_\nu H}
   -V(H)+\mathcal L_D,\nonumber\\
 V(H)&=\lambda_H(\abs H^2-v_H^2)^2,\nonumber\\
 \mathcal L_D
 &=\Rea\left\{\tfrac{\ii}{2}
 \bigl[\bPsi\gamma^\mu\nabla^{e,A}_\mu\Psi
       -(\nabla^{e,A}_\mu\bPsi)\gamma^\mu\Psi\bigr]
       -\bPsi\mathscr M_{\boldsymbol Y}(H)\Psi\right\}.
 \label{eq:dirac-density}
\end{align}
The positive invariant Lie-algebra metric in the first term has the three coefficients $g_j^{-2}$. Pairings in the fermionic expression include all retained chiral and generation indices. The gravitational action is
\begin{equation}
 \Sact_{\mathrm g,\theta}(g)
 =\frac{1}{2\kappa}\int_M(\Scal(g)-2\Lambda)\,\dV_g,
 \qquad \Sact_\theta=\Sact_{\mathrm g,\theta}+\Sact_{\mathrm{SM},\theta}.
 \label{eq:total-action}
\end{equation}
Only compactly supported variations are used. At limited regularity the gravitational first variation is defined by the first-derivative bulk expression obtained by integrating the torsion-free Palatini action by parts; \Cref{app:palatini} fixes this definition. It agrees with \eqref{eq:total-action} on smooth fields up to the appropriate boundary term.

\begin{remark}[Metric gravity versus independent Cartan gravity]
\label{rem:torsion}
In \eqref{eq:total-action} the spin connection is the torsion-free function $\omega(e)$, including in the matter action. This is a declared second-order branch. If instead $\omega$ is varied independently in the total Palatini--fermion action, its Euler equation contains the spin current. Eliminating that connection generally produces Einstein--Cartan interactions, not \eqref{eq:total-action}. The spinless Palatini elimination in the Lorentzian construction \cite{PelissierRenewalSpacetime2026} does not establish torsionlessness for a newly coupled spin-current sector. \Cref{app:palatini} retains the two alternatives explicitly.
\end{remark}

The Hilbert stress convention is
\begin{equation}
 D_g\Sact_{\mathrm{SM},\theta}(z)[k]
 =-\tfrac12\int_M T^{\mathrm{SM}}_{\mu\nu}(z)k^{\mu\nu}\,\dV_g.
 \label{eq:stress-definition}
\end{equation}
For smooth fields the gauge and Higgs contributions are
\begin{align}
 T^{\mathrm{YM}}_{\mu\nu}
 &=\ip{\Fcurv_{\mu\alpha}}{\Fcurv_\nu{}^\alpha}_{\boldsymbol g}
   -\tfrac14g_{\mu\nu}\ip{\Fcurv_{\alpha\beta}}{\Fcurv^{\alpha\beta}}_{\boldsymbol g},
 \label{eq:YM-stress}\\
 T^H_{\mu\nu}
 &=2\Rea\ip{D_\mu H}{D_\nu H}
   -g_{\mu\nu}\bigl(\abs{D_AH}_g^2+V(H)\bigr).
 \label{eq:H-stress}
\end{align}
The Dirac--Yukawa contribution is always taken from the complete metric variation of \eqref{eq:dirac-density}. In particular, the spin-connection variation is not dropped by holding a metric-dependent Dirac operator fixed.

At the reduced regularity, the gravitational distribution is fixed by the
local first-derivative representative
\[
 \frac{1}{2\kappa}\bigl[-\dd B(e)\wedge\omega(e)
       +B(e)\wedge\omega(e)\wedge\omega(e)\bigr]
       -\frac{\Lambda}{\kappa}\,\dV_g,
\]
where $B(e)$ is normalized by
$B_{ab}(e)\wedge R^{ab}(\omega(e))=\Scal(g)\dV_g$.
Its complete variation contains only bounded coframe coefficients, first
coframe derivatives and first derivatives of the metric test. The cutoff
construction and agreement with the smooth Einstein variation are established
in \Cref{app:palatini}. This specifies the distribution used in the main
theorem without assigning pointwise meaning to a rough curvature tensor.

\begin{remark}[Chirality, finite covariance, and ultraviolet modes]
\label{rem:chiral-scope}
The chiral labels specify the fixed continuum representation packet, not a
claim that the finite difference operator is a doubler-free quantum chiral
regulator.  The lattice-chirality obstruction is a spectral issue
\cite{NielsenNinomiya1981,Friedan1982}; describing the variational spinors as
classical does not remove it.  For example, a centered free derivative has
symbol $\ii\sin(h\xi_\mu)/h$, which is not uniformly coercive on the full
lattice spectrum.  The present argument instead imposes compactness, or the
explicit measured-tail bounds, on the \emph{complete reconstructed record}.
It does not discard ultraviolet modes or infer spinor compactness from a
small lattice Dirac row.  Nor does it construct a chiral determinant,
a fermionic integration measure, or quantum anomaly cancellation.  Exact
internal gauge covariance of the finite common action and consistency on
the controlled classical branch are the finite statements used below.
\end{remark}

\subsection{Finite reconstruction and physically scaled stationarity}

\begin{definition}[A finite-interface regulator sequence]
\label{def:regulator}
A finite-interface regulator sequence is a cofinal family $(\Iface_h)_{h\downarrow0}$ of structural interfaces from \Cref{def:finite-interface}, with finite configurations $z_h^{\mathrm d}\in\mathcal Q_h$, together with analytic reconstruction maps giving smooth local fields
\[
 z_h=\mathcal R_h(z_h^{\mathrm d})=(e_h,A_h,H_h,\Psi_h,\bPsi_h)
\]
on the same bundles over $M$, with the coefficient bank $\theta_h$ carried by $\Iface_h$. For every $K\Subset M$ there are linear physical test lifts $\mathcal I_h:\Vtest_K\to T_{z_h^{\mathrm d}}\mathcal Q_h$, supported in a fixed slightly larger region. Define
\begin{align}
 c_h(K)&=\sum_{b\in\{\mathrm g,\mathrm{SM}\}}\;\sup_{\norm v_{C^{r_0}}\le1}
 \abs{D\Sact_{b,h}(z_h^{\mathrm d})[\mathcal I_hv]
       -D\Sact_{b,\theta_h}(z_h)[v]},
 \label{eq:C1-consistency}\\
 \varepsilon_h(K)&=\sup_{\norm v_{C^{r_0}}\le1}
 \abs{D\Sact_h(z_h^{\mathrm d})[\mathcal I_hv]}.
 \label{eq:stationarity}
\end{align}
The supremum is over $\Vtest_K$. First-variation consistency means $c_h(K)\to0$. Physical common-action stationarity means $\varepsilon_h(K)\to0$.
\end{definition}

The comparison in \eqref{eq:C1-consistency} is with the action evaluated on the \emph{finite reconstructed fields}, not on their prospective limit. It can be checked by local interpolation, quadrature, and differentiated stencil estimates. It therefore does not assume the continuum Euler equation. \Cref{app:reconstruction} supplies a concrete smooth comparison realization. The existence of such estimates on a prescribed mesh must not be confused with their validity on every microscopic renewal complex.

Convergence of action values alone cannot replace \eqref{eq:C1-consistency}: $s_n(x)=n^{-1}\sin(nx)$ converges uniformly to zero, whereas $s_n'(0)=1$. Likewise, matching only normalized probabilities loses the affine action normalization needed for \eqref{eq:stationarity}.

\begin{proposition}[Scaled source control of physical stationarity]
\label{prop:source-bound}
Let $a_h$ be a finite action covector on a positive source space with Gram $G_h$. Suppose
\[
 \norm{a_h}_{G_h^{-1}}^2\le R_h,
 \qquad \norm{v_h(v)}_{G_h}\le L_h\norm v_{C^{r_0}},
\]
and
\[
 \abs{D\Sact_h[\mathcal I_hv]-a_h[v_h(v)]}
 \le e_h\norm v_{C^{r_0}}.
\]
Then $\varepsilon_h(K)\le L_h\sqrt{R_h}+e_h$. In particular, the sufficient vanishing condition is $L_h\sqrt{R_h}+e_h\to0$, not merely $R_h\to0$.
\end{proposition}
\begin{proof}
Write the source pairing in the $G_h$ metric and apply Cauchy--Schwarz:
\[
 \abs{a_h[v_h(v)]}
 \le\norm{a_h}_{G_h^{-1}}\norm{v_h(v)}_{G_h}
 \le L_h\sqrt{R_h}\norm v_{C^{r_0}}.
\]
Adding the identification error and taking the supremum proves the claim.
\end{proof}

The same estimate holds for an averaged source after Cauchy--Schwarz in the probability variable. An averaged Euler identity, however, is not a fieldwise equation for one limiting metric.  Similarly, stationarity restricted to a constraint surface does not control a missing normal action row; \Cref{app:complete-source} gives the precise decomposition. The reduced closure theorem is deterministic; \Cref{cor:random-realization} states the additional requirements for a realized probabilistic version.

\subsection{The local finite realization}
\label{subsec:local-action-conventions}

The regulator-independent theorem does not prescribe a stencil. The
finite Wilson criterion and the quantitative results use the particular
link/plaquette action $S_h^{\rm loc}$ specified in \Cref{app:native-action}.
Its Cartan and spin links share the same coframe-derived torsion-free
connection, and all sectors retain the relative normalization of
\eqref{eq:total-action}.

For that realization fix a periodic comparison box $\mathscr Q$ of period
$2\pi$ in each coordinate, an odd integer $n$, and $h=2\pi/n$.
The physical cylinder lies strictly inside a buffered coordinate region;
the auxiliary temporal seam imposes no physical periodic-time condition.
The notation $\mathsf R_h^0$ denotes the piecewise-constant reconstruction,
$\mathcal I_h^{\rm trig}$ the full trigonometric reconstruction, and
$D^+_{\mu,h}$ the ordinary forward difference. Their comparison is proved in
\Cref{lem:native-reconstruction-identification}.
For internal links $U_\mu$, let $U_{\mu\nu}$ be the ordered plaquette
product in \eqref{eq:native-plaquettes}. The literal packets are
\[
 F^h_{\mu\nu}=h^{-2}\operatorname{Log}U_{\mu\nu},\qquad
 K^h_\mu=\frac{\rho_H(U_\mu(x))H(x+he_\mu)-H(x)}h.
\]
Their logarithms are always taken on the declared admissible chart.
The complete physical covector uses the positive mass metric with link
tangent $h^{-1}\dot U_\mu U_\mu^{-1}$; its norm, the action and the packet
norms are preserved by actual site gauges. These conventions do not filter
or replace the finite record.

\section{Variational closure}
\label{sec:limit}

We now state the principal closure theorem. Its hypotheses are expressed
in the topology of the complete first variation, rather than the topology
of the primitive fields alone. The proof has three steps: extract the
nonlinear packets, establish continuity of the complete first variation,
and pass physically scaled finite stationarity to the limit. The two
sector-specific steps are proved in \Cref{subsec:endpoint-mechanisms}; the
elementary product and screen estimates are collected in
\Cref{app:analytic-preliminaries}.

\subsection{Reduced topology and the main theorem}
\label{subsec:reduced-main}

The following convergence class records the different roles of first-order
bilinear and second-order quadratic sectors.  It is a sufficient action
topology; no claim of optimal regularity for the Cauchy problem is intended.

\begin{definition}[Reduced action convergence]
\label{def:reduced-topology}
On each compact smooth chart $K$, reduced action convergence means
\begin{align}
 e_h&\to e\quad\text{strongly in }H^1,
 &e_h(x),e(x)&\in\mathcal K_e\quad\text{a.e.},\notag\\
 A_h&\to A\quad\text{in }L^4,
 &\Fcurv_{A_h}&\to\Fcurv_A\quad\text{in }L^2,\notag\\
 H_h&\to H\quad\text{in }L^2,
 &D_{A_h}H_h&\to D_AH\quad\text{in }L^2,\notag\\
 \Psi_h&\rightharpoonup\Psi\quad\text{in }H^1,
 &\bPsi_h&\rightharpoonup\bPsi\quad\text{in }H^1,
 \label{eq:reduced-topology}
\end{align}
where $\mathcal K_e$ is a fixed compact subset of one oriented,
time-oriented, nondegenerate coframe chart, together with convergence of the
coefficient bank in a compact physical parameter set.  All bundle
identifications are fixed.  Strong $H^1\cap L^4$ convergence of $H_h$
follows from \Cref{prop:covariant-higgs-endpoint}; it is not an independent
condition.  The geometric condition is strong first-jet convergence on a
bounded chart and does not assert $\norm{e_h-e}_\infty\to0$.
\end{definition}

A common compact screen means the approximation property in
\Cref{lem:screen}; \Cref{prop:time-compactness} gives a sufficient
spatial-screen and time-regularity criterion. The next definition turns
those estimates into a compactness certificate for the reduced topology.

\begin{definition}[Reduced compactness certificate]
\label{def:reduced-certificate}
A regulator sequence satisfies the reduced compactness certificate on
$K\Subset M$ when, after fixed local gauge and spin-frame identifications,
\begin{enumerate}[label=\textup{(R\arabic*)},leftmargin=3.2em]
\item the coframe values lie almost everywhere in one fixed compact
      nondegenerate coframe chart, the coframes are bounded in $H^1$, and
      their first-derivative packet has a common compact screen in $L^2$;
\item the connections are precompact in $L^4$, and the curvature packet is
      bounded and has a common compact screen in $L^2$;
\item the Higgs fields are bounded in $L^2$, and the covariant-gradient
      packet is bounded and has a common compact screen in $L^2$;
\item $\Psi_h$ and $\bPsi_h$ are bounded in $H^1$; no spinor first-jet
      screen is required;
\item the coefficient banks remain in a compact subset of the physical
      parameter region.
\end{enumerate}
The first item deliberately requires only boundedness of the coframe values,
not their precompactness in $L^\infty$.  Rellich compactness supplies the
zeroth-order $L^2$ compactness, while the screen supplies strong compactness
of the first derivatives.
\end{definition}

\Needspace{8\baselineskip}
\begin{theorem}[Reduced Einstein--Standard-Model variational closure]
\label{thm:reduced-closure}
Let $(\Iface_h)$ be a finite-interface regulator sequence satisfying
\Cref{def:reduced-certificate} on every compact region, first-variation
consistency $c_h(K)\to0$, and scaled common-action stationarity
$L_h\sqrt{R_h}+e_h\to0$.  Then every cutoff sequence has a subsequence for
which
\begin{align}
 e_h&\to e &&\text{strongly in }H^1,\notag\\
 A_h&\to A &&\text{in }L^4,\qquad
 \Fcurv_{A_h}\to\Fcurv_A &&\text{in }L^2,\notag\\
 H_h&\to H &&\text{in }H^1\cap L^4,\qquad
 D_{A_h}H_h\to D_AH &&\text{in }L^2,\notag\\
 \Psi_h&\rightharpoonup\Psi &&\text{in }H^1,\qquad
 \bPsi_h\rightharpoonup\bPsi &&\text{in }H^1.
 \label{eq:reduced-limits}
\end{align}
The coframe values of the subsequence and its limit remain in the same
compact nondegenerate chart almost everywhere; uniform convergence of the
coframes is not asserted.  Along this subsequence the complete reconstructed
first variation converges in $(\Vtest_K^{r_0})^*$ and the limit is a
distributional solution of all classical Einstein--Standard-Model Euler
equations.  In particular, with the gravitational distribution defined by
\Cref{app:palatini},
\begin{equation}
 (\Eins(g)+\Lambda g)\dV_g=\kappa T^{\rm SM}(z)\dV_g.
 \label{eq:reduced-Einstein-limit}
\end{equation}
Thus neither strong $L^\infty$ coframe convergence, an independent Higgs
quartic compactness condition, nor strong spinor first-jet compactness is
part of this closure criterion.
\end{theorem}

The conclusion is distributional stationarity of the classical,
non-quantized field theory. It does not by itself assert smoothness or
well-posedness of a Cauchy problem. In particular, the theorem needs strong
compactness of the bosonic quantities entering quadratic stresses, but
only weak first-jet control of the bilinear fermionic sector.

\subsection{The two endpoint mechanisms}
\label{subsec:endpoint-mechanisms}

The apparent critical Higgs obstruction is removed by using the covariant
gradient itself. Once the connections are compact in $L^4$, that gradient
controls the ordinary first derivative and closes the quartic potential.

\begin{proposition}[Covariant-gradient compactness closes the Higgs endpoint]
\label{prop:covariant-higgs-endpoint}
Let $K\Subset M$ be a bounded smooth chart and let $H_h$ be the smooth reconstructed Higgs fields.  Suppose
\[
 A_h\to A\quad\text{in }L^4(K),\qquad
 \sup_h\norm{H_h}_{L^2(K)}<\infty,
 \qquad D_{A_h}H_h\to Y\quad\text{in }L^2(K).
\]
Then $(H_h)$ is precompact in $H^1(K)$.  Every subsequential limit $H$ satisfies $Y=D_AH$, and the convergence along that subsequence is strong in $H^1(K)$ and hence in $L^4(K)$.  If $H_h\to H$ in $L^2(K)$ is already known, the whole sequence converges strongly in these spaces.  In particular the quartic potential and the complete Higgs stress converge in $L^1$ without an independent Orlicz or $L^{4+\delta}$ hypothesis.
\end{proposition}
\begin{proof}
The closure of $\{A_h\}\cup\{A\}$ is norm compact in $L^4(K)$.  Multiplication by a fixed $L^4$ coefficient is compact $H^1\to L^2$: approximate the coefficient in $L^4$ by bounded functions and compose multiplication by the bounded approximation with the compact embedding $H^1\Subset L^2$.  Uniformity on the compact coefficient set therefore gives, for every $\epsilon>0$,
\begin{equation}
 \norm{BH}_2\le \epsilon\norm H_{H^1}+C_\epsilon\norm H_2
 \qquad(B\in\overline{\{A_h\}\cup\{A\}}).
 \label{eq:relative-multiplier}
\end{equation}
Since $\dd H_h=D_{A_h}H_h-\rho_H(A_h)H_h$, the ordinary $H^1$ norm and \eqref{eq:relative-multiplier} give a uniform $H^1$ bound after the first term on the right is absorbed.  Extract $H_h\rightharpoonup H$ in $H^1$ and $H_h\to H$ in $L^2$.

The product identity
\[
 A_hH_h-AH=(A_h-A)H_h+A(H_h-H)
\]
converges to zero in $L^2$: the first term uses boundedness of $H_h$ in $L^4$, and the second uses compact multiplication by the fixed $A\in L^4$.  Passing to distributions in the formula for $\dd H_h$ identifies $Y=D_AH$.  Applying the same graph estimate to $H_h-H$ gives strong $H^1$ convergence.  The continuous embedding $H^1(K)\hookrightarrow L^4(K)$ in four dimensions gives strong $L^4$ convergence, and the quadratic and quartic stress terms then converge by the product estimates of \Cref{lem:products}.
\end{proof}

The fermionic mechanism is different. Each differentiated term in the
Dirac action has an undifferentiated spinor factor, so weak--strong pairing
can pass the first variation. The metric dependence of the spin connection
must be included in that argument, as in the following proposition.

\begin{proposition}[Weak-$H^1$ continuity of the complete fermionic first variation]
\label{prop:weak-fermion}
On $K\Subset M$, suppose the coframes $e_h,e$ take values almost everywhere
in one compact nondegenerate coframe chart and
\begin{align*}
 e_h&\to e\quad\text{strongly in }H^1,
 &A_h&\to A\quad\text{in }L^4,\\
 H_h&\to H\quad\text{in }L^4,
 &\Psi_h&\rightharpoonup\Psi\quad\text{in }H^1,\\
 &&\bPsi_h&\rightharpoonup\bPsi\quad\text{in }H^1,
\end{align*}
and let the finite Yukawa/Majorana coefficients converge.  Then the complete
first variation of the Dirac--Yukawa action \eqref{eq:dirac-density},
including its metric and spin-connection variation, converges against every
physical test tuple.  For every integer $r\ge2$,
\begin{equation}
 \norm{D\Sact_{D,h}-D\Sact_D}_{(\Vtest_K^r)^*}\longrightarrow0.
 \label{eq:weak-fermion-variation}
\end{equation}
The same conclusion holds if either or both of the gauge and Higgs
hypotheses are weakened to strong $L^2$ convergence with uniform $L^4$
bounds. In particular $A_h\to A$ strongly in $L^2$ with
$\sup_h\norm{A_h}_4<\infty$ suffices for the fermionic sector.
No strong convergence of
the spinor first derivatives, and no uniform convergence of the coframes,
is required.  In particular the fermionic Hilbert stress has no independent
defect in this class.
\end{proposition}
\begin{proof}
Weak $H^1$ convergence and Rellich compactness give strong $L^2$
convergence of both undifferentiated spinors and uniform $L^4$ bounds.
Their bilinears converge strongly in $L^1$ and remain bounded in $L^2$,
hence converge weakly in $L^2$.  A differentiated Dirac term pairs one
weakly convergent $L^2$ derivative with one strongly convergent $L^2$
undifferentiated spinor.  By \Cref{lem:products}, the coframe Clifford
coefficients are bounded, converge in measure, and preserve this strong
$L^2$ factor.  The spin connection and its metric variation converge
strongly in $L^2$ by \eqref{eq:spin-connection-convergence}.  Gauge
coefficients need only converge strongly in $L^2$, as does the affine
Yukawa coefficient. These coefficients pair with weakly convergent $L^2$
spinor bilinears. In a spinor variation a bilinear is replaced by one
spinor and one bounded test; this factor is strongly convergent in $L^2$.
The $L^4$ bounds retain uniform control of the covariant products. This
proves both relaxed hypotheses.  Thus every term in the complete
first variation converges for a fixed smooth test, including the complete
metric--spin-connection chain.

The same expansions give a cutoff-independent bound by
$C\norm v_{C^1}$ on bounded $H^1$ spinor sets and on the fixed compact
coframe chart.  The $C^r$ unit ball supported in a fixed compact set is
totally bounded in $C^1$ for $r\ge2$.  A finite-net argument upgrades
pointwise convergence to the dual norm \eqref{eq:weak-fermion-variation}.
\end{proof}

\subsection{Continuity and proof of closure}
\label{subsec:closure-proof}

The two endpoint propositions supply precisely the sectorwise convergence
summarized in \Cref{tab:action-topology}. The remaining gravity and
Yang--Mills terms are bounded algebraic coefficients multiplying strongly
convergent quadratic packets. No limiting stress tensor is supplied as
independent data.

\begin{table}[htbp]
\centering\small
\caption{Sectorwise hypotheses for reduced variational closure.  All norms
are local positive comparison norms.  The strong curvature and gradient
requirements concern nonlinear packets, not independent tensor variables.}
\label{tab:action-topology}
\begin{tabularx}{\textwidth}{@{}l>{\raggedright\arraybackslash}X>{\raggedright\arraybackslash}X@{}}
\toprule
Sector & Convergence used & Consequence\\
\midrule
Gravity & $e_h\to e$ strongly in $H^1$ on one compact nondegenerate coframe chart &
Continuity of the distributional Einstein first variation.\\[3pt]
Yang--Mills & $A_h\to A$ in $L^4$, $F_{A_h}\to F_A$ in $L^2$ &
Identification of curvature and strong $L^1$ convergence of gauge stress.\\[3pt]
Higgs & $H_h\to H$ in $L^2$, $D_{A_h}H_h\to D_AH$ in $L^2$ &
Strong $H^1\cap L^4$ convergence and closure of the quartic potential.\\[3pt]
Fermions & Weak $H^1$ convergence of both spinor variables &
Continuity of the complete bilinear first variation; no independent fermionic stress defect.\\
\bottomrule
\end{tabularx}
\end{table}

\begin{proposition}[Continuity in the reduced action topology]
\label{prop:reduced-continuity}
Under \Cref{def:reduced-topology}, for every integer $r\ge2$,
\begin{equation}
 \norm{D\Sact_{\theta_h}(z_h)-D\Sact_\theta(z)}_{(\Vtest_K^r)^*}\longrightarrow0.
 \label{eq:reduced-variation-continuity}
\end{equation}
The conclusion holds separately for gravity, Yang--Mills, Higgs and the
complete Dirac--Yukawa contribution, including the metric--spin connection
chain.  In particular the limiting matter metric variation is the Hilbert
stress of the limiting fields, not an independently prescribed tensor.
\end{proposition}
\begin{proof}
\Cref{prop:covariant-higgs-endpoint} gives strong $L^4$ convergence of
$H_h$.  The curvature, covariant-gradient and quartic products then converge
in $L^1$ by \Cref{lem:products} and \eqref{eq:quartic-continuity}.
The first-order gravitational representative
\eqref{eq:palatini-first-order} is a finite sum of bounded smooth coframe
coefficients multiplying products of first coframe derivatives, plus terms
linear in a first derivative of the metric test.  The derivative packets
converge strongly in $L^2$, while the coefficients remain uniformly bounded
and converge in measure; \Cref{lem:products} therefore passes every term
uniformly on bounded $C^1$ test sets.  The Yang--Mills and Higgs metric rows
have the same bounded-coefficient/strong-product structure.  The complete
fermionic contribution follows from \Cref{prop:weak-fermion}.  A finite-net
argument on the $C^r$ unit ball yields
\eqref{eq:reduced-variation-continuity}.
\end{proof}

It remains to combine this continuity with compactness and finite
stationarity. The abstract triangle inequality is recorded in
\Cref{thm:abstract-closure}; the following proof applies it to the complete
coupled action.

\begin{proof}[Proof of \Cref{thm:reduced-closure}]
On a compact chart, the bounded $H^1$ coframes are precompact in $L^2$.
The common screen lemma gives, after extraction, strong convergence of their
first derivatives.  Hence $e_h\to e$ strongly in $H^1$.  After a further
subsequence the coframes converge almost everywhere, so the closed compact
coframe range is preserved by the limit and all algebraic coefficient and
inverse bounds remain uniform.  The same screen argument gives strong
$L^2$ convergence of the curvature and covariant Higgs-gradient packets,
while the assumed zeroth-order compactness gives $A_h\to A$ in $L^4$.
The identity
$\dd A_h=\Fcurv_{A_h}-A_h\wedge A_h$ identifies the curvature limit with
$\Fcurv_A$ in distributions.

By \Cref{prop:covariant-higgs-endpoint}, the bounded $L^2$ Higgs family with
strong covariant gradients has a subsequence converging strongly in
$H^1\cap L^4$, and its limiting gradient is $D_AH$.  The two spinor families
have weakly convergent $H^1$ subsequences, with strong $L^2$ convergence by
Rellich compactness.  Extract the finite parameter bank as well and use a
diagonal argument over the fixed atlas and a countable compact exhaustion.
This gives \eqref{eq:reduced-limits}.

Let
\[
 \omega_h(K)=\norm{D\Sact_{\theta_h}(z_h)-D\Sact_\theta(z)}_
 {(\Vtest_K^{r_0})^*}.
\]
It tends to zero by \Cref{prop:reduced-continuity}.  The abstract closure
estimate and \Cref{prop:source-bound} give
\begin{equation}
 \norm{D\Sact_\theta(z)}_{(\Vtest_K^{r_0})^*}
 \le L_h\sqrt{R_h}+e_h+c_h(K)+\omega_h(K)\longrightarrow0.
 \label{eq:reduced-residual-bound}
\end{equation}
Independent gauge, Higgs, spinor and dual-spinor tests give their Euler
identities.  An inverse-metric test gives the gravitational variation
$(\Eins(g)+\Lambda g)\dV_g/(2\kappa)$ and the matter variation
$-T^{\rm SM}\dV_g/2$, proving \eqref{eq:reduced-Einstein-limit}.  The
first-order gravitational representative fixes this identity at the stated
metric regularity; no multiplication of an arbitrary curvature distribution
by an uncontrolled coefficient is used.
\end{proof}

\begin{remark}[Scope of the weaker geometric topology]
\label{rem:weaker-geometry-scope}
The preceding weakening concerns the zero-defect first-variation passage.
It is not used to strengthen a theorem with a nonzero concentrating stress
measure.  A uniformly bounded coefficient that converges only in measure
need not commute with a singular weak-star measure.  Accordingly the defect
results below retain the stronger geometric hypotheses stated there.  This
separates the regularity needed for the classical zero-defect branch from
that needed to identify an unresolved measure-valued source.
\end{remark}

\subsection{Stress and uniqueness consequences}

\begin{corollary}[Reconstructed and operational stress in the reduced class]
\label{cor:reduced-stress}
Along a subsequence of \Cref{thm:reduced-closure}, the Yang--Mills and Higgs
stress densities converge strongly in $L^1$.  The complete matter metric
first variation converges in the dual $C^r$ test norm for $r\ge2$, hence in
$H^{-m}$ for every integer $m>4$ on a fixed compact test region.  For the
operational stress
\[
 \langle\mathsf T_h,k\rangle
 =-2D\Sact_{\mathrm{SM},h}[\mathcal I_h(k,0,0,0,0)],
\]
sectorwise consistency in \eqref{eq:C1-consistency} gives convergence to
$T^{\rm SM}\dV_g$ in $(\Vtest_K^{r_0})^*$ and in $H^{-m}$ for $m>r_0+2$.
Strong $L^1$ convergence of the fermionic stress is not asserted under weak
spinor $H^1$ convergence.
\end{corollary}
\begin{proof}
The first assertion follows from the quadratic and quartic product
estimates.  Apply \Cref{prop:reduced-continuity} to metric tests and use
$H_0^m(K)\hookrightarrow C^r(K)$ for $m>r+2$ in four dimensions.
Taking $r=2$ proves the reconstructed statement.  The operational statement
adds the separately controlled matter consistency error, which is measured
in the $C^{r_0}$ test norm.
\end{proof}

Subsequential closure and selection of a unique continuum trajectory are
separate assertions. A unique Cauchy class removes subsequence ambiguity
provided its initial traces are identified independently.

\begin{proposition}[Uniqueness without summable adjacent-cutoff transport]
\label{prop:reduced-unique}
Suppose the hypotheses of \Cref{thm:reduced-closure} hold for the full
regulator family with one limiting coefficient bank.  Suppose further that
all accumulation points, in fixed gauge and frame identifications, belong
to a solution class in which the limiting Cauchy data determine a unique
solution $z_*$.  Assume convergence of the relevant Cauchy traces separately
in that class.  Then the entire family converges to $z_*$ in reduced action
convergence.  No summability of adjacent-cutoff errors is needed.
\end{proposition}
\begin{proof}
On bounded $H^1$ sets the weak spinor topology is metrizable because $H^1$
is a separable Hilbert space; the other factors in
\eqref{eq:reduced-topology} are norm topologies.  A countable exhaustion
therefore metrizes the relevant local product topology on the bounded
family.  If convergence failed, a sequence outside a fixed neighborhood of
$z_*$ would have, by reduced compactness, a convergent subsequence.  Its
limit is a solution with the prescribed traces and hence equals $z_*$,
a contradiction.  Trace convergence is an additional hypothesis:
spacetime compactness alone does not identify data on an initial slice.
\end{proof}

The conditional Sobolev upgrade is \Cref{cor:strong-solution-upgrade}.
Strong spinor first-jet convergence, regular Sobolev criteria, summable
cofinal transport and realized-random variants are retained in
\Cref{app:strong-packet}. They strengthen the conclusions on their stated
branches without becoming hypotheses of the reduced theorem.

\section{Weak limits and zero-defect closure}
\label{sec:defects}

The main theorem identifies the classical limit when the nonlinear
bosonic packets are compact. We next retain, rather than discard, the
information lost when those packets converge only weakly. The order of the
argument is: identify the joint defect source, characterize its vanishing
by positive energies, pass to cutoff-dependent gauge representatives, and
finally give a finite criterion that restores strong compactness.

\subsection{The joint weak-limit alternative}

On the weak-limit branch, assume compatible coframes converge strongly
in $H^1$ and uniformly to continuous nondegenerate coframes, with uniform
inverse bounds. Let $A_h\to A$ in $L^4$, let the Higgs fields converge weakly
in $H^1$, and identify the weak $L^2$ limits of $F_{A_h}$ and $D_{A_h}H_h$
with $F_A$ and $D_AH$. Both spinor variables remain weakly bounded in
$H^1$, and the parameters converge in the physical region. Uniform
geometric convergence is retained here because a singular stress measure
must be multiplied by a continuous limiting coefficient.

\Cref{thm:higgs-defect,prop:YM-defect} construct the measure limits
\[
 T_h^{\rm YM}\dV_{g_h}\rightharpoonup^*
       T^{\rm YM}\dV_g+\mathfrak S_{\rm YM},\qquad
 T_h^H\dV_{g_h}\rightharpoonup^*
       T^H\dV_g+\mathfrak S_H,
\]
where $\mathfrak S_H=\mathfrak K_H-\lambda_Hg\nu_H$:
$\mathfrak K_H$ is the kinetic stress defect and $\nu_H\ge0$ the quartic
concentration measure. Their construction and the matter-equation passage
are proved in \Cref{app:weak-limit-estimates}. They assemble into the
following limit identity.

\begin{theorem}[Joint matter-defect Einstein limit]
\label{thm:joint-defect}
Assume the critical Higgs hypotheses of \Cref{thm:higgs-defect}, the
Yang--Mills hypotheses of \Cref{prop:YM-defect}, $A_h\to A$ strongly in $L^4$,
the compatible coframe and weak-$H^1$ spinor hypotheses of
\Cref{prop:weak-matter-equations}, and convergence of the gravitational first
variation.  The latter follows in particular from the strong first-coframe-jet hypothesis of the reduced theorem.  If the operational metric residual and its reconstruction-consistency error vanish, then
\begin{equation}
 \bigl(\Eins(g)+\Lambda g\bigr)\dV_g
 =\kappa\Bigl(T^{\mathrm{SM}}\dV_g
       +\mathfrak S_{\mathrm{YM}}+\mathfrak S_H\Bigr).
 \label{eq:joint-defect-Einstein}
\end{equation}
If additionally $g_h\to g$ in $C^1$, the total reconstructed stress
has the uniform measure bound and vanishing conservation residual of
\Cref{lem:defect-conservation-transport}, and the regular limiting matter
stress satisfies its Noether identity, then
$\mathfrak S_{\mathrm{YM}}+\mathfrak S_H$ is distributionally conserved.
These additional conditions concern conservation only, not
\eqref{eq:joint-defect-Einstein}. In particular, the classical
Einstein--Standard-Model equation is the zero-defect branch of the weak-limit
identity.
\end{theorem}
\begin{proof}
Test the finite metric equation against a compactly supported smooth symmetric tensor and pass the gravitational, regular fermionic, Yang--Mills, and Higgs terms separately using \Cref{prop:YM-defect,thm:higgs-defect} and the complete fermionic continuity in \Cref{prop:weak-matter-equations}.  This gives \eqref{eq:joint-defect-Einstein}.  Under the additional $C^1$ hypothesis,
\Cref{lem:defect-conservation-transport} passes the finite conservation
identity; subtraction of the separately valid regular Noether identity gives
conservation of the sum of the defects. Neither conservation nor the Ward
identity forces that sum to vanish.
\end{proof}

The matter equations and the metric equation have different weak-limit
requirements. The former pass by weak--strong pairing, whereas metric
variation retains two weak bosonic factors. The critical bubble and null-wave
examples in \Cref{ex:bubble,ex:null-wave} illustrate respectively
concentration and oscillatory kinetic stress. Neither matter stationarity
nor conservation alone removes either phenomenon.

\subsection{Positive defects and exact zero-defect characterization}
\label{subsec:positive-defects}

The signed tensor defect is controlled by positive covariance measures.
For a weak $L^2$ packet $Y_h\rightharpoonup Y$, write
\[
 Y_h\otimes Y_h\,\dV_0\rightharpoonup^*
 Y\otimes Y\,\dV_0+\mathsf Q_Y,
 \qquad \mu_Y=\operatorname{tr}\mathsf Q_Y\ge0.
\]
The covariance measure is positive semidefinite, and on a compact domain
including its boundary, $\mu_Y=0$ is equivalent to strong $L^2$ convergence.
\Cref{lem:positive-packet-defect} proves this assertion and its
bounded-contraction estimate. Applied to curvature and covariant Higgs
gradients, it yields the measures $\mu_{\rm YM}$ and $\mu_{\nabla H}$.

Apply this lemma to the realified packets
$Y_h^{\rm YM}=\Fcurv_{A_h}$ and $Y_h^H=D_{A_h}H_h$, using the fixed
positive comparison metrics.  Write their trace defects as
$\mu_{\rm YM}$ and $\mu_{\nabla H}$.  The volume factors and Lorentzian
contractions in the Hilbert stress are included in $B_h$ in
\eqref{eq:defect-contraction}.  On each compact nondegenerate chart,
\begin{equation}
 |\mathfrak S_{\rm YM}|\le C\mu_{\rm YM},\qquad
 |\mathfrak K_H|\le C\mu_{\nabla H},\qquad
 |\mathfrak S_H|\le C(\mu_{\nabla H}+\nu_H).
 \label{eq:stress-defect-domination}
\end{equation}
Thus positive packet defects control every component of the stress defects.
The local formulation uses compactly supported tests on a larger chart;
vanishing there yields strong convergence on smaller compact sets.  Including
the boundary in the compact-domain formulation prevents undetected loss of
mass at the edge of that domain.

\begin{theorem}[Timelike positivity and exact zero-defect equivalence]
\label{thm:zero-defect-characterization}
Assume the weak Yang--Mills and Higgs hypotheses of
\Cref{thm:joint-defect}, with positive limiting gauge coefficients and
$\lambda_H>0$, and let $n$ be a continuous future-directed unit timelike
vector field for $g$ on the compact chart.  There are $c,C>0$ such that,
as scalar Radon measures,
\begin{equation}
 c(\mu_{\rm YM}+\mu_{\nabla H}+\nu_H)
 \le (\mathfrak S_{\rm YM}+\mathfrak S_H)(n,n)
 \le C(\mu_{\rm YM}+\mu_{\nabla H}+\nu_H).
 \label{eq:timelike-defect-coercivity}
\end{equation}
Consequently the following conditions are equivalent along the extracted
sequence on the compact domain:
\begin{align}
 &\mathfrak S_{\rm YM}+\mathfrak S_H=0;\notag\\
 &(\mathfrak S_{\rm YM}+\mathfrak S_H)(n,n)=0;\notag\\
 &\mu_{\rm YM}=\mu_{\nabla H}=\nu_H=0;\notag\\
 &\Fcurv_{A_h}\to\Fcurv_A\text{ in }L^2,\quad
 D_{A_h}H_h\to D_AH\text{ in }L^2,\quad
 H_h\to H\text{ in }L^4.
 \label{eq:complete-zero-defect-equivalence}
\end{align}
In particular, independent Yang--Mills and Higgs stress defects cannot cancel
to produce a zero total bosonic defect in this minimally coupled class.
\end{theorem}
\begin{proof}
Choose a local orthonormal frame $(n,e_1,e_2,e_3)$ and put
$E_i=F(n,e_i)$ and $B_i=\frac12\varepsilon_{ijk}F(e_j,e_k)$.
Direct substitution into \eqref{eq:YM-stress} gives
\begin{equation}
 T^{\rm YM}(n,n)=\tfrac12\sum_i
 (|E_i|_{\boldsymbol g}^2+|B_i|_{\boldsymbol g}^2).
 \label{eq:YM-timelike-energy}
\end{equation}
This is positive definite on the complete six-component two-form packet.
Likewise, the kinetic part of \eqref{eq:H-stress} obeys
\begin{equation}
 T^{H,\rm kin}(n,n)=|D_nH|^2+\sum_i|D_{e_i}H|^2.
 \label{eq:H-timelike-energy}
\end{equation}
Both forms, multiplied by the positive volume density, are uniformly
comparable on the compact chart to the respective reference squared norms.
Contracting their covariance measures therefore gives nonnegative measures
bounded above and below by the corresponding trace defects.  The potential
defect is $-\lambda_Hg\nu_H$, whose contraction with $n\otimes n$ is
$\lambda_H\nu_H$.  The lower-degree potential terms converge strongly and
contribute no defect.  Uniform convergence of $g_h$, its inverse and the
parameters justifies all coefficient replacements by
\Cref{lem:positive-packet-defect}.  A finite frame cover gives
\eqref{eq:timelike-defect-coercivity}.

Vanishing of the full tensor implies vanishing of its timelike contraction.
The lower bound then forces all three positive measures to vanish.
\Cref{lem:positive-packet-defect,lem:critical-cubic} identify this with the
three strong convergences in \eqref{eq:complete-zero-defect-equivalence}.
Those convergences, by the quadratic and quartic estimates, make the tensor
defects vanish, closing the cycle.
\end{proof}

\begin{corollary}[Independent compactness certificates]
\label{cor:independent-bosonic-certificates}
Under the same hypotheses, $\mu_{\nabla H}=0$ already implies
$H_h\to H$ in $H^1\cap L^4$ and hence $\nu_H=0$.
Thus $\mu_{\rm YM}=\mu_{\nabla H}=0$ is sufficient and necessary for the
zero-defect branch in \Cref{thm:zero-defect-characterization}; a separate
quartic certificate is redundant once the covariant gradient is compact.
\end{corollary}
\begin{proof}
Rellich compactness provides $H_h\to H$ in $L^2$ along the weak $H^1$
subsequence.  Trace-defect zero gives strong convergence of the covariant
gradient.  Apply \Cref{prop:covariant-higgs-endpoint}, followed by
\Cref{lem:critical-cubic}.
\end{proof}

\begin{remark}[Scope of the converse]
The converse uses the coercive \emph{timelike Hilbert-energy contraction},
not the Lorentzian trace of the stress or the signed action density.
For example the four-dimensional Yang--Mills stress is trace free even when
$\mu_{\rm YM}\ne0$.  Conservation alone also does not imply vanishing.
The conclusion does not include nonminimal improvement terms, a weakly
converging gravitational first jet, or an additional fermionic defect outside
the weak-$H^1$ class.  Those cases require their own defect analysis.
\end{remark}

Thus the minimal theory has a genuine two-sided criterion: zero total
bosonic stress defect means that the unresolved positive packet energies
vanish. The nonminimal improvement and the stronger Higgs
$L^{8/3}$ reserve are separate refinements, treated in
\Cref{app:improvement,app:positive-Higgs-envelope}.

\subsection{Critical gauge-quotient closure}
\label{subsec:critical-quotient}

The preceding defect identities were formulated after choosing local bundle
representatives.  At the four-dimensional Yang--Mills energy endpoint it is
preferable to state the compactness and first-variation hypotheses modulo the
internal gauge group itself.  Two additional ingredients are needed.  First,
cutoff-dependent gauge normalizations must be paired with a test topology that
is exactly equivariant.  Second, at critical curvature regularity one must
separate concentration control from the stronger compactness needed to remove
quadratic stress oscillations.

Fix on each local spinor--generation bundle a smooth positive-metric reference
connection $\nabla^{\rm ref}$ acting trivially on the internal gauge factor and
write
\begin{equation}
 \nabla^{{\rm ref},A}\Psi
 :=\nabla^{\rm ref}\Psi+\rho(A)\Psi.
 \label{eq:reference-gauge-spin-derivative}
\end{equation}
For a compactly supported physical test
$v=(k,a,\eta_H,\eta_\Psi,\eta_{\bPsi})$ define
\begin{align}
 \mathfrak n_z(v)
 :={}&\norm{k}_{W^{1,\infty}}
 +\norm a_{L^\infty}+\norm{D_Aa}_{L^2}\notag\\
 &+\norm{\eta_H}_{L^\infty}+\norm{D_A\eta_H}_{L^2}\notag\\
 &+\norm{\eta_\Psi}_{L^\infty}
   +\norm{\nabla^{{\rm ref},A}\eta_\Psi}_{L^2}
 +\norm{\eta_{\bPsi}}_{L^\infty}
   +\norm{\nabla^{{\rm ref},A}\eta_{\bPsi}}_{L^2}.
 \label{eq:equivariant-test-size}
\end{align}
All norms use the fixed positive comparison geometry.

The test size is exactly invariant under internal gauge transformations:
$\mathfrak n_{u\cdot z}(u\cdot v)=\mathfrak n_z(v)$, and the complete
first variation is bounded by $C_K\mathfrak n_z(v)$ on bounded weak packets.
The verification is \Cref{lem:equivariant-tests}. These facts allow tests in
a cutoff-dependent normalized gauge to be pulled back without introducing
derivatives of the normalizing gauge into the finite budget.

We say that a finite regulator has \emph{equivariant consistency and
stationarity budgets} $c_h^{\rm eq}(K)$ and
$\varepsilon_h^{\rm eq}(K)$ if
\begin{align}
 \abs{D\Sact_h(z_h^{\rm d})[\mathcal I_hv]
      -D\Sact_{\theta_h}(z_h)[v]}
 &\le c_h^{\rm eq}(K)\,\mathfrak n_{z_h}(v),
 \label{eq:equivariant-consistency}\\
 \abs{D\Sact_h(z_h^{\rm d})[\mathcal I_hv]}
 &\le\varepsilon_h^{\rm eq}(K)\,\mathfrak n_{z_h}(v)
 \label{eq:equivariant-stationarity}
\end{align}
for every smooth compactly supported physical test.  Because of
\eqref{eq:equivariant-test-invariance}, these inequalities can be applied to a
fixed test in a cutoff-dependent normalized gauge by pulling that test back to
the original finite record; no derivative of the normalizing gauge enters the
budget.

\begin{definition}[Uniformly integrable critical curvature energy]
\label{def:critical-curvature-ui}
A family of reconstructed connections has uniformly integrable critical
curvature energy on a compact chart $K$ if
\begin{equation}
 \lim_{\delta\downarrow0}\sup_h\ \sup_{\substack{E\subset K\ \text{measurable}\\
                         \operatorname{vol}_0(E)\le\delta}}
 \int_E\abs{F_{A_h}}_r^2\,\dV_0=0.
 \label{eq:critical-curvature-ui}
\end{equation}
This property is internal-gauge invariant.
\end{definition}

Uniform integrability of critical curvature energy supplies a common
small-energy cover and local Coulomb representatives by
\Cref{prop:critical-uhlenbeck}.  This supplements the classical critical
four-dimensional Yang--Mills compactness picture, where curvature
concentration and bubbling can occur even for bounded-energy sequences
\cite{Uhlenbeck1982,Sedlacek1982}: uniform integrability rules out such
concentration on the chosen compact chart, while a separate compactness
mechanism is still needed to remove quadratic stress oscillation.  Positive
covariant graph bounds then give ordinary $H^1$ bounds for the matter
variables in those representatives (\Cref{lem:critical-positive-graph}). The
resulting theorem concerns the gauge quotient; it does not select a canonical
representative.

\begin{theorem}[Critical-energy quotient defect closure]
\label{thm:critical-quotient-defect}
Let $(\Iface_h)$ be a finite-interface regulator sequence.  On every compact
chart assume:
\begin{enumerate}[label=\textup{(Q\arabic*)},leftmargin=3.2em]
\item the coframes are precompact in $L^\infty\cap H^1$, uniformly
      nondegenerate, and the coefficient banks lie in a compact physical set;
\item $\sup_h\norm{F_{A_h}}_2<\infty$ and the curvature energies satisfy
      \eqref{eq:critical-curvature-ui};
\item with $K_h=D_{A_h}H_h$,
      $\sup_h(\norm{H_h}_2+\norm{K_h}_2)<\infty$;
\item the spinors satisfy the gauge-invariant comparison bounds
      \begin{equation}
       \sup_h\bigl(
       \norm{\Psi_h}_2+\norm{\nabla^{{\rm ref},A_h}\Psi_h}_2
       +\norm{\bPsi_h}_2+\norm{\nabla^{{\rm ref},A_h}\bPsi_h}_2
       \bigr)<\infty;
       \label{eq:critical-spinor-bound}
      \end{equation}
\item $c_h^{\rm eq}(K)+\varepsilon_h^{\rm eq}(K)\to0$.
\end{enumerate}
Then every cutoff sequence has a subsequence and cutoff-dependent local
internal gauges such that, on smaller compact charts,
\begin{align}
 e_h&\to e &&\text{in }L^\infty\cap H^1,\notag\\
 A_h&\rightharpoonup A &&\text{in }W^{1,2},
 &A_h&\to A &&\text{in every }L^q,
   \quad q<4,\notag\\
 F_{A_h}&\rightharpoonup F_A &&\text{in }L^2,\notag\\
 H_h&\rightharpoonup H &&\text{in }H^1,
 &H_h&\to H &&\text{in every }L^q,
   \quad q<4,\notag\\
 D_{A_h}H_h&\rightharpoonup D_AH &&\text{in }L^2,\notag\\
 \Psi_h&\rightharpoonup\Psi,
 \quad\bPsi_h\rightharpoonup\bPsi
 &&\text{in }H^1.
 \label{eq:critical-quotient-convergence}
\end{align}
The Yang--Mills, Higgs, Dirac and dual-Dirac equations hold distributionally,
and the metric equation is
\begin{equation}
 (\Eins(g)+\Lambda g)\dV_g
 =\kappa\left(T^{\rm SM}(z)\dV_g+\mathfrak S_{\rm YM}+\mathfrak S_H\right).
 \label{eq:critical-quotient-einstein}
\end{equation}
If, in addition, $|D_{A_h}H_h|^2$ and $|H_h|^4$ are uniformly integrable,
then the complete bosonic defect is an $L^1$ tensor density and the quartic
concentration measure vanishes.  The conclusion is a quotient statement: no
fixed gauge representative is selected canonically.
\end{theorem}

\begin{proof}
Extract the geometric and parameter subsequence and apply
\Cref{prop:critical-uhlenbeck} on a finite cover of each smaller compact
chart. The covariant graph estimates give weak $H^1$ matter compactness,
and strong subcritical convergence identifies the curvature and covariant
derivatives. Pulling back fixed normalized tests preserves their
$\mathfrak n$-size, so the equivariant budgets pass the matter Euler rows.
The relaxed fermionic continuity theorem supplies the complete fermionic
metric limit. Only the quadratic bosonic metric terms and the quartic
potential retain defects, giving \eqref{eq:critical-quotient-einstein}.
Uniform integrability makes the corresponding covariance measures
absolutely continuous, and Vitali removes the quartic measure under the
additional hypothesis. The complete sectorwise passage is given in
\Cref{app:critical-quotient-proof}.
\end{proof}

\subsection{The finite Wilson-screen criterion}
\label{subsec:finite-Wilson}

The quotient theorem permits a defect. To exclude it directly from the
explicit local finite action, we now screen the literal packets rather
than a continuum surrogate. The crucial distinction is between scalar
magnitude compactness, which follows from unitarity alone, and vector
packet compactness, which also uses a controlled gauge representative.

For the explicit local action one can recognize the missing strong packet
compactness directly on the same finite links that enter the action.  Let
$U_\mu(x)$ be the internal link of \Cref{app:native-action}.  For an integer
$m>0$ let
\begin{equation}
 \mathcal U_{\mu,m}(x)
 :=U_\mu(x)U_\mu(x+h e_\mu)\cdots
   U_\mu(x+(m-1)h e_\mu),
 \label{eq:open-finite-link}
\end{equation}
and use the inverse ordered product when $m<0$.  For a finite packet $Y_h^{\rm d}$
in a unitary representation define the Wilson shift
\begin{equation}
 \mathsf W^U_{\mu,m}Y_h^{\rm d}(x)
 :=\rho(\mathcal U_{\mu,m}(x))Y_h^{\rm d}(x+mhe_\mu),
 \label{eq:finite-Wilson-shift}
\end{equation}
with $\rho=\operatorname{Ad}$ for curvature.  Put
\begin{equation}
 \Omega_h(Y^{\rm d};\varrho)
 :=\max_\mu\max_{0<|m|h\le\varrho}
 \left(
 h^4\sum_{x\in Q'_{h,\mu,m}}
 \abs{\mathsf W^U_{\mu,m}Y_h^{\rm d}(x)-Y_h^{\rm d}(x)}^2
 \right)^{1/2}.
 \label{eq:finite-Wilson-modulus}
\end{equation}
The endpoint gauges telescope in \eqref{eq:open-finite-link}, so
\eqref{eq:finite-Wilson-modulus} is exactly site-gauge invariant.
Throughout the native finite-packet statements below we use
\begin{equation}
 \norm{Y_h}_{p,h}^p=h^4\sum_x|Y_h(x)|^p
 \label{eq:native-lp-main}
\end{equation}
for finite $p$, with the usual maximum norm for $p=\infty$.

Fix, independently of the record, a rooted spanning tree $\mathfrak T_h$ of
the periodic comparison grid.  Spanning- or maximal-tree gauges are standard
tools for eliminating redundant lattice link variables
\cite{LigterinkWaletBishop2000}.  The role of the tree here is more specific:
it produces a gauge-invariant same-record logarithmic seed whose smallness can
be fed into the exact finite Coulomb normalization below.  For a vertex $x$,
let $C_x(U)$ be the ordered product of the actual internal links along the tree
path from the root to $x$, with inverses on reversed edges, and define the
root-based closed words
\begin{equation}
 W^{\mathfrak T}_\mu(x)
 =C_x(U)U_\mu(x)C_{x+h e_\mu}(U)^{-1},
 \qquad
 B^{\mathfrak T}_{\mu,h}(x)
 =h^{-1}\operatorname{Log}W^{\mathfrak T}_\mu(x).
 \label{eq:rooted-Wilson-seed}
\end{equation}
The seed is declared inadmissible when one of the required logarithms leaves
the fixed conjugation-invariant logarithm chart.  On the admissible set put
\begin{equation}
 \mathfrak c_h^{\mathfrak T}(U)
 :=\norm{B_h^{\mathfrak T}}_{4,h}+\norm{\mathbb F_h}_{2,h}.
 \label{eq:rooted-Wilson-certificate}
\end{equation}
Endpoint cancellation shows that this quantity is invariant under arbitrary
finite site gauges, up to the irrelevant constant conjugation at the root.
The construction and the associated exact finite Coulomb normalization are
proved in \Cref{app:finite-gauge-normalization}.

\begin{theorem}[Finite Wilson-screen zero-defect criterion]
\label{thm:finite-Wilson-zero-defect}
Consider the explicit local action of \Cref{app:native-action} on the
periodic buffered comparison grid, with conclusions restricted to an
interior physical cylinder.  Let
\begin{equation}
 \mathbb F_h=\{F^h_{\alpha\beta}\}_{\alpha<\beta},
 \qquad
 \mathbb K_h=\{K^h_\alpha\}_{\alpha=0}^3
 \label{eq:literal-finite-bosonic-packets}
\end{equation}
be its literal curvature and Higgs-gradient packets.  In this theorem all
packet norms and Wilson translation sums are taken over the entire periodic
comparison box, including periodic shifts; only the final Euler tests are
restricted to the interior physical cylinder.  Assume:
\begin{enumerate}[label=\textup{(F\arabic*)},leftmargin=3.2em]
\item the nodal coframes $E_h$ take values in one fixed compact
      nondegenerate coframe chart, the families
      $\mathsf R_h^0E_h$ and $\mathsf R_h^0D_h^+E_h$ are precompact in
      $L^2$, the scaled connection margin
      \eqref{eq:native-gravity-log-margin} holds, and the complete unfiltered
      reconstructed coframes remain in the same compact chart;
\item for one predeclared rooted tree the same-record certificate
      \eqref{eq:rooted-Wilson-certificate} obeys
      $\mathfrak c_h^{\mathfrak T}(U_h)\le\epsilon_c$, where $\epsilon_c$
      is the cutoff-independent threshold in
      \Cref{thm:finite-Coulomb-normalization};
\item $\norm{\mathbb F_h}_{2,h}+\norm{\mathbb K_h}_{2,h}
      +\norm{H_h}_{2,h}$ is bounded and
      \begin{equation}
       \lim_{\varrho\downarrow0}\limsup_{h\downarrow0}
       \left[\Omega_h(\mathbb F_h;\varrho)
             +\Omega_h(\mathbb K_h;\varrho)\right]=0;
       \label{eq:two-finite-Wilson-screens}
      \end{equation}
\item the spinor and dual-spinor blocks have uniformly bounded positive
      internal-link graph norm \eqref{eq:native-spinor-graph}, and the
      coefficient banks remain in a compact physical parameter set;
\item the norm $\sigma_h^{\rm phys}$ of the complete finite action covector
      in the positive physical link/matter mass metric tends to zero.
\end{enumerate}
Then there are site gauges acting on the same finite records for which the
normalized logarithmic connections are in exact discrete Coulomb gauge and,
after extraction,
\begin{align}
 e_h&\to e &&\text{strongly in }H^1_{\rm loc},\notag\\
 A_h&\to A &&\text{strongly in }L^4_{\rm loc},\notag\\
 F_{A_h}&\to F_A &&\text{strongly in }L^2_{\rm loc},\notag\\
 H_h&\to H &&\text{strongly in }H^1_{\rm loc}\cap L^4_{\rm loc},\notag\\
 D_{A_h}H_h&\to D_AH &&\text{strongly in }L^2_{\rm loc},\notag\\
 \Psi_h&\rightharpoonup\Psi,
 \quad\bPsi_h\rightharpoonup\bPsi &&\text{weakly in }H^1_{\rm loc}.
 \label{eq:Wilson-strong-convergence}
\end{align}
The complete sectorwise first-variation consistency errors relative to the
unfiltered reconstructed fields tend to zero.  Every subsequential limit
therefore satisfies the ordinary distributional Einstein--Standard-Model
equations, the bosonic stress densities converge strongly in $L^1_{\rm loc}$,
and the complete fermionic metric variation converges in the dual $C^r$ test
norm for every $r\ge2$.

No growing-band or high-derivative reserve, independently assumed matter
first-variation consistency, separate quartic compactness condition, strong
spinor first-jet compactness, or independent packet-energy
uniform-integrability hypothesis is used.  The small rooted certificate is a
substantive finite non-Abelian gauge-normalization condition; it is not
inferred from the curvature bound alone.
\end{theorem}
\begin{proof}
By \Cref{prop:rooted-gauge-certificate}, the rooted words in
\eqref{eq:rooted-Wilson-seed} are an exact tree representative of the
original links, and the small certificate constructs an exact finite site
gauge in which the logarithmic connection obeys
$\delta_hA_h=0$ and $\norm{A_h}_{4,h}\le\epsilon_*$.  The same-record gauge
transformation leaves the complete action, the physical action-covector
norm, the literal curvature and Higgs-gradient norms, their Wilson moduli,
and the positive internal-link spinor graph norms unchanged.

The normalized $L_h^4$ bound gives a uniform $L_h^2$ connection bound on
the fixed-volume box.  Apply \Cref{thm:native-Wilson-compactness} separately
to $\mathbb F_h$ and $\mathbb K_h$.  Their raw reconstructions are therefore
strongly precompact in $L^2$, and the associated positive packet energies are
uniformly integrable.  The critical discrete Coulomb absorption theorem
\Cref{thm:native-discrete-Coulomb} converts strong literal curvature
compactness into strong $L^4$ connection convergence and strong convergence
of the full ordinary connection first differences.  The native Higgs graph
argument \Cref{prop:native-Higgs-compactness} then gives the Higgs statements
in \eqref{eq:Wilson-strong-convergence}; the positive spinor graph gives the
weak $H^1$ spinor subsequence.

The coframe precompactness gives, after extraction, strong $L^2$ convergence
of both values and first differences; discrete summation by parts identifies
the latter limit with the derivative of the former.  Hence
\Cref{prop:native-gravity-firstjet} gives strong $H^1$ geometric convergence
without uniform coframe convergence.  The all-sector native theorem \Cref{thm:native-firstvariation-no-band} now shows
that the unchanged local finite action has the complete classical first
variation as its limit, uniformly on bounded $C^r$ physical tests.  Because
the gauge normalization is an isometry for the positive physical tangent
metric, assumption \textup{(F5)} makes this limiting first variation vanish.
Independent metric, gauge, Higgs, spinor and dual-spinor tests give the full
Euler system.  Strong $L^2$ convergence of the curvature and covariant Higgs
gradient, together with strong $L^4$ Higgs convergence, gives strongly
convergent quadratic and quartic densities.  Their bounded metric
coefficients converge in measure, so \Cref{lem:products} gives strong
$L^1$ convergence of the bosonic stress densities directly.  This step does
not apply a merely measure-convergent coefficient to a nonzero defect
measure.  The complete fermionic metric variation is covered by
\Cref{prop:weak-fermion}.
\end{proof}

The proof therefore has a fixed direction: actual site-gauge normalization,
compactness of the literal packets, identification of the reconstructed
fields, and only then use of the finite action covector. The small rooted
certificate is not inferred from a curvature bound. The
\emph{determinant-resolved} extension in
\Cref{cor:determinant-resolved-closure} restricts the smallness requirement
to the semisimple factor on the stated zero-flux branch; its complete
construction is in \Cref{app:determinant-sector}.

The Wilson criterion is qualitative and requires no high-order Cauchy
reserve. Quantitative control of the fields, curvature and complete stress
belongs to a different, regular branch, developed next.

\section{Quantitative finite-action realization}
\label{sec:native-branch}

We now strengthen qualitative closure to quantitative comparison on a
regular slab. The first result compares the explicit finite action with a
specified regular solution. The second constructs the comparison future
from a corrected initial datum of one finite record. The high-order
forcing estimates and the constraint solver are kept separate: a small
finite Euler row alone supplies neither a regular Cauchy solution nor
compatible initial data.

Throughout this section the action is the unchanged $S_h^{\rm loc}$ of
\Cref{app:native-action}; Fourier decompositions are estimates on its
complete record. \Cref{app:quantitative-residuals} proves the consistency
and forcing estimates, and \Cref{app:coupled-stability} proves the
coupled actual-jet stability theorem used in both comparisons.

\subsection{Finite residual and leakage budgets}

Fix buffered slabs $Q\Subset Q'$ inside a periodic auxiliary coordinate box $\mathscr Q$ of period $2\pi$ in each coordinate, with the auxiliary temporal seam outside $Q'$.  Let $n$ be odd, $h=2\pi/n$, and
\[
 \Lambda_h=\{\ell\in\Z^4:|\ell_\mu|\le(n-1)/2\}.
\]
Let $u_h$ be the calibrated nodal physical-coordinate record $(e,A,H,\Psi,\bPsi)$ in the open nondegenerate chart of \Cref{app:native-action}, and let $z_h=\mathcal I_h^{\rm trig}u_h$ be its trigonometric reconstruction.  For $1\le K<(n-1)/2$ define
\begin{equation}
 \tau_{h,j}(K)=\sum_{\substack{\ell\in\Lambda_h\\|\ell|_\infty>K}}
 \left(1+\frac{|\ell|_1}{K}\right)^j|\widehat u_h(\ell)|,
 \qquad j\ge0.
 \label{eq:native-tail}
\end{equation}
These are measured tails of the existing record; no Fourier coefficient is deleted.

For the finite action $S_h^{\rm loc}$ of \eqref{eq:native-local-action}, define its mass-normalized interior physical Euler row $E_h^{\rm raw}$ by
\begin{equation}
 D S_h^{\rm loc}(u_h)[v_h]
 =h^4\sum_{x\in Q'_h}\ip{E_h^{\rm raw}(u_h)(x)}{v_h(x)}
 \label{eq:raw-euler-row}
\end{equation}
for every interior nodal physical variation, with the fixed finite coefficient and representation metrics.  Write
\begin{equation}
 \norm{E_h^{\rm raw}(u_h)}_{0,h;Q'}\le \sigma_h.
 \label{eq:native-sigma}
\end{equation}
The quantity $\sigma_h$ is an actual finite-action covector norm, not a continuum residual.

The actual-jet state $\mathcal U(z)$ contains the reconstructed metric
and its first derivatives, the connection and curvature, the Higgs field
and covariant derivative, and both spinors with their tangential covariant
first derivatives. Its exact components and normal-spinor elimination
are given in \eqref{eq:actual-jet-state}. No independent jet variables
are supplied to the comparison theorem.

The finite consistency estimate is $O(hK^3)$ on the declared growing
smoothness scale (\Cref{prop:native-consistency}). The leakage transfer
retains the full tail and gives the following differentiated forcing budget.

Set $m=k+2$ for an integer $k\ge4$ and define
\begin{align}
 \varepsilon_{0,h}&=\sigma_h+hK^3,
 &\varepsilon_{c,h}&=\varepsilon_{0,h}+K^2\tau_{h,2}(K),\notag\\
 L_{h,k}&=1+\log\left(2+\frac{C_kK^{k+5}}{\varepsilon_{c,h}}\right),
 &F_{h,k}&=\varepsilon_{c,h}K^{k+1}L_{h,k}^{k+1}
             +K^{k+2}\tau_{h,m}(K).
 \label{eq:native-forcing-budget}
\end{align}

Under the calibrated field-value and inverse-metric margins,
$hK\le c_{\rm res}$, $\tau_{h,m}\le\tau_*$ and
$\varepsilon_{c,h}\le1$, \Cref{thm:native-source} proves
\[
 \mathcal Y_k(z_h):=
 \|\mathcal R_{\rm B}(z_h)\|_{L_t^2H_x^k}
 +\|\mathcal R_{\rm D}(z_h)\|_{L_t^2H_x^{k+1}}
 \le C F_{h,k}.
\]
Here $\mathcal R_{\rm B}$ collects the physical Einstein, Yang--Mills and
Higgs rows, and $\mathcal R_{\rm D}$ the Dirac and dual-Dirac rows.
The additional spatial derivative in the Dirac source is part of the
coupled prolongation estimate, not a freely interchangeable norm.

\subsection{Reference-centred quantitative closure}

Let $z_*$ be a regular torsion-free Einstein--Standard-Model solution on $Q=[0,T]\times\mathbb T^3$, written in harmonic coordinates and temporal internal gauge, and let $\mathcal U_*$ denote its actual-jet state from \Cref{app:coupled-stability}.  For an actual reconstructed record put
\begin{equation}
 i_{h,k}=\norm{\mathcal U_h(0)-\mathcal U_*(0)}_{H^k},
 \qquad
 \gamma_{h,k}=\norm{C(g_h)}_{L^2_tH_x^{k+1}}
              +\norm{\partial_tC(g_h)}_{L^2_tH_x^k},
 \label{eq:native-initial-gauge}
\end{equation}
where $C^\mu(g)=g^{\alpha\beta}\Gamma^\mu_{\alpha\beta}(g)$.  The state $\mathcal U_h$ is built from the actual derivatives and curvatures of $z_h$; no independent jet variables are supplied.

\begin{theorem}[Quantitative finite-action stability and Einstein--Standard-Model closure]
\label{thm:native-closure}
Let $k\ge4$.  Suppose the complete finite records exist on $Q$, are in temporal internal gauge, and satisfy the calibrated field-value, inverse-metric and buffered comparison hypotheses of \Cref{thm:native-source}.  No all-time Sobolev bound for their actual-jet states is imposed.  Suppose
\[
 K_h\to\infty,\qquad hK_h\to0,\qquad
 \tau_{h,k+2}(K_h)\le\tau_*,
\]
together with the finite-action row bound \eqref{eq:native-sigma}.  Define
\begin{equation}
 D_{h,k}^{\rm full}=i_{h,k}+\gamma_{h,k}+F_{h,k}.
 \label{eq:native-total}
\end{equation}
If $D_{h,k}^{\rm full}\to0$, then the all-time $C_tH_x^k$ bound for the actual state is a conclusion, and
\begin{align}
 &\norm{\mathcal U_h-\mathcal U_*}_{C_tH_x^k}
 +\norm{\Riem(g_h)-\Riem(g_*)}_{L^2_tH_x^{k-1}}\notag\\
 &\hspace{4em}
 +\norm{T^{\rm SM}(z_h)-T^{\rm SM}(z_*)}_{L^2_tH_x^k}
 \le C D_{h,k}^{\rm full}.
 \label{eq:native-geometry}
\end{align}
In particular the compactness, stress convergence and zero-defect conclusions needed for the classical limit are derived on this branch from the finite physical row, initial/gauge mismatch and measured leakage.  Existence of the finite record on the slab is a premise; no existence or selection theorem for an unspecified microscopic rule is inferred.
\end{theorem}
\begin{proof}
\Cref{thm:native-source} provides the physical forcing of the complete record before any comparison with $z_*$.  The gauge-defect-aware equation \eqref{eq:actual-jet-writer} applies to
these actual records even when $C(g_h)\ne0$; its added source is controlled
by the measured $\gamma_{h,k}$. Apply the one-sided coupled actual-jet
stability estimate \Cref{prop:coupled-bootstrap} with
$d=i_{h,k}+\gamma_{h,k}+\mathcal Y_k(z_h)$.  A first-exit argument makes the temporary bounded-state hypothesis self-improving and yields \eqref{eq:native-geometry}.  The stress estimate follows because the complete torsion-free Standard-Model stress is a smooth algebraic function of the controlled actual jets, with the normal spinor derivative recovered from the Dirac equation.  Thus no separately convergent stress tensor or curvature sequence is assumed.
\end{proof}

This theorem derives the all-time state bound by a first-exit argument.
The explicit rate and calibration consequences are
\Cref{cor:native-rate,cor:local-calibration-nonempty}. They remain
conditional on the stated finite-record budgets; no microscopic existence
or selection law is inferred.

\subsection{A posteriori exact physical shadows}
\label{subsec:aposteriori-shadow}

The preceding theorem compares a finite record with a regular exact solution
that is supplied as comparison data.  On a regular physical Cauchy chart that
external future trajectory can be removed.  The finite record itself supplies
the approximate forcing; only its initial datum is retracted to the exact
constraint manifold, after which the exact future is generated by the
Einstein--Standard-Model Cauchy problem.

Let $\mathcal X_k$ be the declared initial actual-jet chart and let
$\mathfrak C_k:\mathcal X_k\to\mathcal Y_k^{\rm con}$ be its complete
continuum constraint map, with the Sobolev orders of the target chosen for
the differential orders of the constraints. It includes the Einstein normal
constraints, internal Gauss constraints, defining-jet identities and the
chosen initial gauge completion. Put
\begin{equation}
 \mathcal C_k=\{d\in\mathcal X_k:\mathfrak C_k(d)=0\}.
 \label{eq:exact-physical-initial-manifold}
\end{equation}
The word ``exact'' in this subsection refers to this zero set, not to the
zero set of a sampled or projected constraint operator.

\begin{assumption}[Verified physical Cauchy tube]
\label{ass:physical-Cauchy-tube}
Fix $k\ge4$ and a harmonic-coordinate/temporal-internal-gauge actual-jet
chart. There is a declared compatible regular family
$\mathcal C_k^{\rm reg}\subset\mathcal C_k\cap H^{k+1}$, carrying a fixed
higher-order bound, and an open neighborhood $\mathcal O_k$ in the initial
$H^k$ topology such that every datum in
$\mathcal C_k^{\rm reg}\cap\mathcal O_k$ generates a unique physical
Einstein--Standard-Model solution on the common slab
$Q=[0,T]\times\mathbb T^3$. These solutions remain in one compact
nondegenerate chart, with a common spacelike foliation and a uniform
$C_tH_x^{k+1}$ bound for their actual-jet states. The constants in
\Cref{prop:coupled-bootstrap} are uniform throughout this family.
\end{assumption}

The regular family is part of the hypothesis: a bounded open $H^k$ ball is
not asserted to acquire an extra derivative under hyperbolic evolution.
Only the exact initial datum, not its future trajectory, will be supplied to
this Cauchy theorem.

\begin{definition}[Exact finite-dimensional constraint reduction]
\label{def:exact-initial-reduction}
Let $\widehat z_h$ denote the record after the regular normalization used in
\Cref{thm:aposteriori-physical-shadow}; put $\widehat z_h=z_h$ when no
normalization is needed. An exact reduction at its normalized initial
record $x_h$ consists of a $C^1$ correction family
$\iota_h(x_h,\lambda)\in\mathcal X_k$, $\lambda\in\R^{m_h}$,
$|\lambda|\le r_*$, a fixed injective linear map
$\jmath_h:\R^{m_h}\to\mathcal Y_k^{\rm con}$, and a $C^1$ map
$\Phi_h(x_h,\lambda)\in\R^{m_h}$ such that
\begin{align}
 \iota_h(x_h,0)&=\mathcal U(\widehat z_h)(0),
       \label{eq:exact-initial-basepoint}\\
 \mathfrak C_k(\iota_h(x_h,\lambda))
       &=\jmath_h\Phi_h(x_h,\lambda)
       \qquad(|\lambda|\le r_*).
       \label{eq:exact-initial-factorization}
\end{align}
Thus all complementary field-valued constraints are solved exactly
throughout the correction family, not just at its base point. Equivalently,
a splitting of the constraint target has identically zero complementary
component on this family and retained component $\Phi_h$. We require
\begin{equation}
 \norm{\iota_h(x_h,\lambda)-\iota_h(x_h,\mu)}_{H^k}
       \le B_*|\lambda-\mu|.
 \label{eq:exact-initial-reader-Lipschitz}
\end{equation}
The reduced coordinate norms and the constants in the ensuing estimates
are fixed in the physical normalization.
\end{definition}

The subscript $h$ labels the record and its correction family; it does not
make $\Phi_h$ a discretization of the remaining continuum constraints.
Constructing $\iota_h$ may itself use an independently verified elliptic
complement solver. A finite list of nodal constraint equations alone does
not establish \eqref{eq:exact-initial-factorization}.

The exact factorization is the substantive constraint hypothesis.
\Cref{prop:initial-constraint-retraction} shows that a retained residual
bounded by $\beta_h$, a uniform inverse margin $\gamma_*$ and correction
Lipschitz constant $B_*$ produce an exactly compatible datum at distance
at most $2B_*\gamma_*^{-1}\beta_h$. Validated evaluation errors and
complementary corrections are charged separately in
\Cref{cor:validated-initial-retraction,rem:initial-correction-bookkeeping}.

The positive reader criterion of \Cref{cor:native-reader-sobolev} has a
particularly useful regular specialization.  Fix integers
\begin{equation}
 q>s+4,
 \qquad s>k+1,
 \qquad k\ge4.
 \label{eq:shadow-regularity-indices}
\end{equation}
If one selected record has $V_h\le R^2$ and the reader LMI of
\eqref{eq:source-reader-LMI}, then
\begin{equation}
 \norm{z_h}_{H^q(Q)}\le C_R,
 \qquad
 \max_{|\alpha|\le s+2}\norm{\partial^\alpha z_h}_{L^\infty(Q)}\le C_R.
 \label{eq:shadow-reader-reserve}
\end{equation}
Thus no growing Fourier scale is needed on this branch.

The fixed reserve converts a small complete finite Euler row
$\sigma_h\le s_h$ into a strong physical forcing estimate.
\Cref{prop:shadow-residual-upgrade} proves
\[
 \varepsilon_h:=\|\mathcal R_{\rm B}(z_h)\|_{L^2}
       +\|(r_D,\bar r_D)(z_h)\|_{L^2}\le C_R(s_h+h),
 \qquad
 \eta_{h,k}:=C_R\bigl(\varepsilon_h^{1-k/s}
       +\varepsilon_h^{1-(k+1)/s}\bigr),
\]
with bosonic forcing bounded in $H^k$ and both spinor rows in $H^{k+1}$
by $\eta_{h,k}$. The comparison future can now be generated from the
corrected datum rather than supplied in advance.

\begin{theorem}[A posteriori exact Einstein--Standard-Model shadow]
\label{thm:aposteriori-physical-shadow}
Assume \Cref{ass:physical-Cauchy-tube}.  Let $z_h$ be one complete finite
record satisfying the reader hypothesis leading to
\eqref{eq:shadow-reader-reserve}, the finite Euler bound
$\sigma_h\le s_h$, and the exact initial reduction of \Cref{def:exact-initial-reduction} and the
retraction hypotheses of \Cref{prop:initial-constraint-retraction} with
$c_h\le\beta_h$.  Suppose the
record is already in the harmonic/temporal chart of
\Cref{ass:physical-Cauchy-tube}, or admits a certified regular normalization
$z_h\mapsto\widehat z_h$ that fixes its initial embedding and first
differential and changes the residual norms in
\eqref{eq:shadow-strong-residual} by at most a cutoff-independent factor.
Define
\begin{equation}
 \mathfrak d_h
 :=2B_*\gamma_*^{-1}\beta_h
   +C_R\sqrt T\,\eta_{h,k}.
 \label{eq:aposteriori-shadow-budget}
\end{equation}
Assume also that the initial retraction stays in the verified initial
tube, so $d_h\in\mathcal C_k^{\rm reg}\cap\mathcal O_k$ with its stated regularity
bounds.  If $\mathfrak d_h$ lies below the uniform stability radius of that
tube, then there exists a unique exact physical
Einstein--Standard-Model solution $z_{*,h}$ on $Q$ with initial datum $d_h$
from \Cref{prop:initial-constraint-retraction}, and
\begin{align}
 &\norm{\mathcal U(\widehat z_h)-\mathcal U(z_{*,h})}_{C_tH_x^k}
 +\norm{\Riem(\widehat g_h)-\Riem(g_{*,h})}_{L_t^2H_x^{k-1}}
 \notag\\
 &\qquad
 +\norm{T^{\rm SM}(\widehat z_h)-T^{\rm SM}(z_{*,h})}_{L_t^2H_x^k}
 \le C_R\mathfrak d_h.
 \label{eq:aposteriori-physical-shadow}
\end{align}
Uniqueness here is for the corrected datum $d_h$ in the declared gauge
class, not among all exact solutions lying near the record. In particular,
the exact comparison future is a conclusion of the finite record certificate
and the verified Cauchy tube; it is not supplied in advance.  If the corrected initial
data $d_h$ converge in a Cauchy class for which the physical solution is
unique, the shadows $z_{*,h}$ converge to that unique exact solution and the
certified finite records have the same continuum limit.
\end{theorem}
\begin{proof}
By \Cref{prop:shadow-residual-upgrade}, the complete normalized record has
bosonic forcing in $L^2_tH_x^k$ and Dirac forcing in
$L^2_tH_x^{k+1}$ with total size at most $C_R\eta_{h,k}$.  Proposition
\ref{prop:initial-constraint-retraction} supplies a compatible exact initial
datum whose $H^k$ distance from the normalized actual initial state is at most
$2B_*\gamma_*^{-1}\beta_h$.  Assumption
\ref{ass:physical-Cauchy-tube} therefore constructs the exact physical
solution $z_{*,h}$ from that datum on the same slab.

Apply \Cref{prop:coupled-bootstrap} with reference solution $z_{*,h}$.
The stated harmonic normalization makes $C(\widehat g_h)=0$ on this branch;
otherwise the additional budget of \Cref{cor:shadow-harmonic-defect} is used.
Cauchy--Schwarz in time bounds the forcing contribution to its integrated
energy estimate by $C_R\sqrt T\,\eta_{h,k}$.  The sum of this contribution and
the initial mismatch is \eqref{eq:aposteriori-shadow-budget}, and the assumed
tube smallness closes the first-exit argument.  The state, curvature and
complete stress estimates are exactly
\eqref{eq:aposteriori-physical-shadow}.  The final assertion follows from
continuous dependence and uniqueness in the stated Cauchy class.
\end{proof}

The harmonic-defect version and a transparent $O_R(h^{1/2})$ rate are
\Cref{cor:shadow-harmonic-defect,cor:aposteriori-shadow-rate}. The correction
acts only on the initial comparison datum; the observed spacetime record
is unchanged. Uniqueness is attached to that corrected datum, not to an
unspecified collection of nearby exact trajectories.

\paragraph{Obtaining the budgets on one record.}
\label{subsec:quantitative-source-bridge}
\Cref{sec:native-selection} gives a source-level entrance to the deterministic
results. Signed action drift, positive reader energy and initial alignment
or constraint control are tested simultaneously on one actually observed
record. The corresponding selection theorems retain the original failure
probabilities and do not replace the field by an expectation.
The reader certificates, alternative prefix-balance entrance and generated
comparison families are recorded in
\Cref{app:native-source-selection,app:prefix-stationarity,app:generated-dynamics}.
These are sufficient realization mechanisms, not additional premises of
reduced or Wilson-screen closure.

\section{Interpretation and discussion}
\label{sec:discussion}

The central result of this paper is a finite-to-continuum statement rather
than a reconstruction theorem for the microscopic origin of the Standard
Model.  We begin from a finite Lorentzian--Standard-Model interface carrying
one gravity--matter action and ask under which analytic conditions its
refinement has the classical coupled field equations as a continuum limit.
The structural origin of the finite carrier, the analytic control required
for the limit, and the eventual uniqueness of that limit are logically
distinct questions.

This separation is important because the presence of the correct finite
field content and a common variational principle does not by itself imply
classical dynamics.  Nonlinear products must converge in the topologies in
which they enter the complete first variation, finite stationarity must remain
small after being tested in physical directions, and gauge, spin, and
Lorentzian geometric structures must remain compatible across the refining
family.  The Einstein--Standard-Model equations arise here from the
simultaneous control of these ingredients rather than being built into the
definition of the finite interface.

Several neighboring convergence theories should therefore be distinguished
from the present statement.  Rigorous Regge and finite-element results control
distributional scalar or Einstein curvature for specified approximating
metrics \cite{CheegerMullerSchrader1984,ChristiansenRegge2024,
GawlikNeunteufelScalar2025,GawlikNeunteufelEinstein2026}, while
gauge-invariant Maxwell--Klein--Gordon schemes establish convergence for
specific discretizations by combining discrete symmetry, energy estimates and
compactness \cite{ChristiansenHalvorsen2011,
ChristiansenHalvorsenScheid2023}.  Those results provide stronger
scheme-specific stability information than is assumed abstractly here.  The
present closure theorem addresses a different question: conditional on
compactness, first-variation consistency and physically scaled stationarity,
it passes the complete coupled Lorentzian gravity--gauge--Higgs--spinor
variation to the limit and records any unresolved bosonic quadratic products
as explicit stress defects.  Thus regulator independence is obtained by
stating the analytic hypotheses explicitly, not by asserting that every
finite regulator satisfies them.

\subsection{Closure, quantitative comparison and uniqueness}

The natural convergence object is the complete physical first variation.
Convergence of action values does not identify critical configurations, and
a small source-coordinate residual must still be tested against the
physical lift. This is why the condition is
\begin{equation}
L_h\sqrt{R_h}+e_h\longrightarrow0.
\label{eq:discussion-scaled-stationarity}
\end{equation}

rather than $R_h\to0$ alone. The lift cost is part of the
finite-to-continuum dictionary.

Under \Cref{thm:reduced-closure}, controlled sequences have accumulation
points satisfying
\begin{equation}
\Eins_{\mu\nu}(g)+\Lambda g_{\mu\nu}
=
\kappa T^{\mathrm{SM}}_{\mu\nu},
\qquad
\Evar_A=\Evar_H=\Evar_\Psi=\Evar_{\bPsi}=0
\label{eq:discussion-ESM}
\end{equation}

in distributions. The compactness step produces the candidate fields;
continuity and stationarity identify their equations. Cofinal uniqueness
requires an additional mechanism, either summable transport or uniqueness
in a fixed Cauchy class with independently convergent traces.

The quantitative theorems answer a stronger question on a regular branch.
The reference-centred result bounds the actual fields, curvature and stress
in terms of a finite action row, initial/gauge mismatch and measured leakage.
The data-centred theorem instead uses a positive reader reserve, an exact
initial-constraint reduction and a verified physical Cauchy tube to construct
the comparison future. It does not assume that a bounded $H^k$ neighborhood
acquires an extra derivative or that finitely many sampled constraints solve
the full continuum constraint system.

The source-level theorems supply a further, conditional entrance to these
budgets on one observed record. Their probabilistic estimates retain failure
events and do not define a mean-field geometry. Whether a particular
microscopic dynamics satisfies their drift, reader and legal-domain
hypotheses remains a separate selection question.

\subsection{Nonlinear compactness and the defect alternative}

The distinction between primitive fields and nonlinear packets is visible
in the weak-limit equation
\begin{equation}
\bigl(\Eins(g)+\Lambda g\bigr)\dV_g
=
\kappa\Bigl(
T^{\mathrm{SM}}\dV_g
+\mathfrak S_{\mathrm{YM}}
+\mathfrak S_H
\Bigr).
\label{eq:discussion-defect}
\end{equation}

The defects retain information invisible in the weak field limits but
visible in the metric variation. In the minimally coupled theory,
\Cref{thm:zero-defect-characterization} identifies their vanishing with the
vanishing of positive curvature and covariant-gradient energy defects;
Higgs quartic compactness then follows. Conservation alone is not such a
criterion. Nonminimal improvement terms and weaker geometric first-jet
limits require their own defect analysis.

The geometry needed for zero-defect closure is weaker than the geometry
used to multiply a nonzero singular stress measure. Strong $H^1$ coframe
convergence on a common compact nondegenerate chart suffices for the former,
while the defect identities retain uniform coefficient convergence to
continuous representatives. The two hypotheses are not interchangeable.

At critical gauge regularity, equivariant tests make quotient closure
compatible with cutoff-dependent local gauges. For the explicit action,
Wilson translation control first makes scalar packet magnitudes compact;
a same-record normalization then gives vector compactness and identifies
the limiting fields. The finite rooted certificate remains a substantive
non-Abelian condition, and zero determinant flux is a separate topological
restriction of the global periodic normalization.

The stronger Higgs diagnostic is distinct from this qualitative closure:
$e_h^{H,+}\asymp|K_h^H|^2$ gives equivalence between an $L^{4/3}$ bound
for the positive kinetic density and an $L^{8/3}$ covariant-gradient bound.
Strong $L^2$ gradient compactness yields only $L^1$ convergence of that
density. The additional integrability is retained as an optional reserve,
not silently inserted into the zero-defect hypotheses.

The fermionic first variation behaves differently because the action is
first order and bilinear. Weak $H^1$ control suffices for its complete
weak--strong passage, including the spin-connection variation. On a regular
branch, stronger spinor convergence follows from Lorentzian
symmetric-hyperbolic stability. Variational stationarity and a controlled
Cauchy evolution therefore remain distinct conclusions.

\subsection{Structural provenance and classical scope}

Renewal Geometry supplies one upstream realization through its predictive
carrier, compatible differential/spin structure and structural matter packet
\cite{PelissierRenewalSpacetime2026,PelissierSpectral2026,PelissierDuality2026}.
The present analysis begins after that interface has been obtained and is
portable to another realization with the same stated hypotheses. The
logical separation is
\begin{equation}
\begin{aligned}
\text{structural realization:}\quad&
\text{a finite Lorentzian--Standard-Model interface is obtained},\\
\text{analytic closure:}\quad&
\text{a controlled refinement satisfies the classical equations},\\
\text{microscopic selection:}\quad&
\text{a microscopic dynamics naturally enters that controlled class}.
\end{aligned}
\label{eq:discussion-three-levels}
\end{equation}

The second is the principal subject of this paper; the source-level
criteria give conditional routes toward the third, not a universality
claim for microscopic dynamics.

The endpoint established here is classical.  The spinor variables are
classical variational fields, and the limiting equations are the classical
Einstein--Yang--Mills--Higgs--Dirac--Yukawa equations.  The analysis does not
construct a fermionic path-integral measure, a renormalized quantum
Standard Model coupled to gravity, or a quantum gravitational continuum
measure.  Operational probabilities in a microscopic predictive theory and a
quantum field-theoretic measure play different mathematical roles and should
not be identified.

Likewise, the closure theorem does not determine the numerical Standard-Model
parameters.  Gauge couplings, the Higgs potential parameters, Yukawa
eigenvalues and mixing data, and their renormalization-group evolution remain
independent inputs or dynamical data on the finite structural carrier.

\paragraph{Outlook.}

The main remaining questions are consequently more sharply separated.
On the analytic side, one may seek weaker compactness hypotheses, a more
complete classification of admissible stress defects, boundary versions of
the closure theorem, and broader uniqueness criteria for cofinal refinement.
On the microscopic side, the central problem is to identify natural classes
of finite dynamics for which the interface, stationarity, reader, and
compactness hypotheses are stable or typical rather than externally imposed.  Constraint
compatibility of a principal operator alone does not control the complete
state-dependent source or supply a spatially uniform stationary family.

Within the scope proved here, the central conclusion is that classical
Einstein--Standard-Model dynamics can be characterized as the variational
closure of a finite gravity--matter theory once the physically relevant
nonlinear packets and variations are controlled.  When that control is
strong enough, the ordinary coupled field equations emerge; when it is not,
the unresolved finite structure remains visible through explicit effective
stress defects rather than being discarded by definition.

\section*{Data and formalization availability}
Selected finite algebraic components of the construction have been formalized in
\texttt{Lean}.  Source code is available at
\url{https://github.com/AI-MathPhys/renewal-geometry}.

\clearpage
\appendix
\crefalias{section}{appendix}
\crefname{appendix}{Appendix}{Appendices}
\Crefname{appendix}{Appendix}{Appendices}
\section*{Guide to the appendices}
\addcontentsline{toc}{section}{Guide to the appendices}

The appendices are part of the same manuscript and contain the complete
supporting proofs. They are organized by their use in the main argument:
\Cref{app:analytic-preliminaries,app:strong-packet} support reduced and
strong-packet closure; \Cref{app:weak-limit-estimates} supports the weak-limit
and quotient results; \Cref{app:local-realization,app:native-critical-closure,app:finite-gauge-normalization,app:determinant-transfer}
specify the finite action and prove its critical packet and gauge estimates;
\Cref{app:quantitative-residuals,app:initial-shadow-refinements,app:stability-package}
supply quantitative forcing, constraint correction and coupled stability.
\Cref{app:comparison-families,sec:native-selection,app:source-proof-package}
retain the comparison constructions and source-level entrances. No external
supplement is needed to follow the proof dependencies.

\section{Continuity, compactness and metric-variation estimates}
\label{app:analytic-preliminaries}

The estimates in this appendix support the compactness and continuity
steps of \Cref{thm:reduced-closure}. The abstract closure argument is
independent of the field theory; the product and coframe estimates identify
the complete first variation in the reduced topology.

\subsection{An abstract first-variation closure principle}

The passage from finite stationarity to a limiting Euler equation is a
first-variation statement.  Its abstract form does not require positivity,
convexity, minimization, or a prescribed evolution of the finite records.
The sector-dependent continuity needed to apply it is proved below.

\begin{lemma}[First-variation closure]
\label{thm:abstract-closure}
Let $\mathscr X$ be a class of reconstructed configurations, with a specified
sequential convergence, and let $\mathscr V_K$ be a Banach space of physical
tests.  Let $S$ have first variations $DS(z)\in\mathscr V_K^*$ on this class.
Suppose $z_h\to z$ and
\begin{equation}
 \omega_h(K):=\norm{DS(z_h)-DS(z)}_{\mathscr V_K^*}\longrightarrow0.
 \label{eq:abstract-continuity}
\end{equation}
Let $S_h$ be differentiable finite actions, $z_h=\mathcal R_hz_h^{\rm d}$,
and $I_h$ linear test lifts.  Assume
\begin{align}
 \sup_{\norm v\le1}|DS_h(z_h^{\rm d})[I_hv]-DS(z_h)[v]|&\le c_h(K),\notag\\
 \sup_{\norm v\le1}|DS_h(z_h^{\rm d})[I_hv]|&\le\varepsilon_h(K).
 \label{eq:abstract-budgets}
\end{align}
Then
\begin{equation}
 \norm{DS(z)}_{\mathscr V_K^*}
 \le\varepsilon_h(K)+c_h(K)+\omega_h(K).
 \label{eq:abstract-closure}
\end{equation}
In particular, if $c_h(K)+\varepsilon_h(K)\to0$, then $DS(z)=0$.
If the reconstructed family is sequentially precompact and these hypotheses
hold along every convergent subsequence, all its accumulation points are
stationary configurations of $S$.
\end{lemma}
\begin{proof}
For a unit test, insert and subtract $DS(z_h)[v]$ and
$DS_h(z_h^{\rm d})[I_hv]$.  The three differences are bounded by the three
terms in \eqref{eq:abstract-closure}.  Taking the supremum and then the limit
proves stationarity.  Apply the same argument to each convergent subsequence
for the last assertion.
\end{proof}

This principle concerns convergence of critical configurations, rather than
only convergence of minimum problems.  Plain $\Gamma$-convergence of the
action functionals \cite{DalMaso1993,Braides2002} does not by itself supply
\eqref{eq:abstract-continuity} or \eqref{eq:abstract-budgets}.  There are
stronger variational frameworks for critical points: for example,
Jerrard--Sternberg obtain critical points of $\Gamma$-converging functionals
near suitable nondegenerate critical points \cite{JerrardSternberg2009}, and
Mosco/Attouch-type theories relate variational convergence to convergence of
subdifferentials in convex settings \cite{Attouch1984}.  Those results use
additional structure not assumed for the nonconvex Lorentzian gauge system
here.  The elementary oscillatory example following \Cref{def:regulator}
already separates uniform action convergence from first-variation
convergence.  Variational convergence of gradient flows also requires
additional control of slopes and dissipation \cite{SandierSerfaty2004}; no
gradient-flow law for the Lorentzian common action is assumed here.

\subsection{Products and bounded geometric coefficients}

\begin{lemma}[Local product and bounded-coefficient continuity]
\label{lem:products}
If $u_h\to u$ in $L^p$, $v_h\to v$ in $L^q$, and
$1/p+1/q=1/r\le1$, then $u_hv_h\to uv$ in $L^r$.  In particular,
\begin{equation}
 \norm{B(Y_h,Y_h)-B(Y,Y)}_1
 \le C\norm{Y_h-Y}_2(\norm{Y_h}_2+\norm Y_2)
 \label{eq:quadratic-product}
\end{equation}
for a bounded bilinear contraction $B$.

Suppose in addition that $B_h$ are finite matrix-valued functions with
$\sup_h\norm{B_h}_\infty<\infty$ and $B_h\to B$ in measure.  If
$u_h\to u$ in $L^p$, $1\le p<\infty$, then
\begin{equation}
 B_hu_h\longrightarrow Bu\quad\text{in }L^p.
 \label{eq:bounded-coefficient-product}
\end{equation}
If $u_h\to u$ and $v_h\to v$ in $L^2$, then
$B_h(u_h\otimes v_h)\to B(u\otimes v)$ in $L^1$.

Finally, let $\mathcal K_e$ be a compact subset of one oriented,
time-oriented, nondegenerate coframe chart and suppose
$e_h,e\in\mathcal K_e$ almost everywhere and $e_h\to e$ strongly in
$H^1$.  Then every smooth algebraic coefficient $F(e,e^{-1})$ satisfies
\begin{equation}
 F(e_h,e_h^{-1})\to F(e,e^{-1})
 \quad\text{in }H^1\text{ and in every finite }L^p,
 \qquad \sup_h\norm{F(e_h,e_h^{-1})}_\infty<\infty.
 \label{eq:bounded-chart-composition}
\end{equation}
In particular,
\begin{equation}
 \omega(e_h)\to\omega(e)\quad\text{in }L^2,
 \qquad
 \dot\omega(e_h)[\dot e_h(k)]\to
 \dot\omega(e)[\dot e(k)]\quad\text{in }L^2,
 \label{eq:spin-connection-convergence}
\end{equation}
uniformly for $k$ in bounded $C^1$ test sets.  No strong
$L^\infty$ convergence of $e_h$ is required.
\end{lemma}
\begin{proof}
Subtract products one factor at a time and apply H\"older.  For
\eqref{eq:bounded-coefficient-product}, split
$B_hu_h-Bu=B_h(u_h-u)+(B_h-B)u$.  The first term tends to zero in
$L^p$.  Every subsequence of the second has a further subsequence on which
$B_h\to B$ almost everywhere; dominated convergence applies because the
coefficients are uniformly bounded.  The subsequence criterion proves the
claim for the full sequence.  The tensor-product assertion follows from
$u_h\otimes v_h\to u\otimes v$ in $L^1$ and the same argument.

Strong $H^1$ convergence of $e_h$ implies convergence in measure.  Smooth
coefficients on the compact chart are uniformly bounded with uniformly
bounded derivatives, hence converge in measure and in every finite $L^p$.
The chain rule gives
\[
 \partial F(e_h)-\partial F(e)
 =DF(e_h)(\partial e_h-\partial e)
  +(DF(e_h)-DF(e))\partial e,
\]
and the second term is controlled by
\eqref{eq:bounded-coefficient-product}.  This proves
\eqref{eq:bounded-chart-composition}.  In local frames the Levi--Civita spin
connection is a finite sum $P(e)\partial e$, while its physical metric
variation has the form
$P_1(e)k\,\partial e+P_2(e)\partial k$.  Applying the same
bounded-coefficient argument proves \eqref{eq:spin-connection-convergence}
and its uniformity on bounded $C^1$ test sets.
\end{proof}

\subsection{Compact screens for the quantities that enter the action}

\begin{lemma}[Common compact screen]
\label{lem:screen}
Let $Y_h\rightharpoonup Y$ in a fixed Hilbert space $\mathcal Y$. Suppose that, for every $R$, there are bounded $S_{h,R}$ and compact $S_R$ with
\begin{equation}
 \norm{S_{h,R}-S_R}_{\mathrm{op}}\longrightarrow0,
 \qquad
 \lim_{R\to\infty}\sup_h\norm{(I-S_{h,R})Y_h}=0,
 \qquad S_RY\longrightarrow Y.
 \label{eq:screen}
\end{equation}
Then $Y_h\to Y$ strongly. The same conclusion holds for a whitened packet $Z_h=\mathbb A_h^{1/2}Y_h$ provided $Z_h\rightharpoonup\mathbb A^{1/2}Y$, the screen conditions hold for $Z_h$, and
\[
 mI\preceq\mathbb A_h\preceq MI,
 \qquad \mathbb A_h^{-1/2}\to\mathbb A^{-1/2}\text{ strongly},
 \quad 0<m\le M<\infty.
\]
\end{lemma}
\begin{proof}
For fixed $R$, compactness gives $S_RY_h\to S_RY$, and norm convergence of $S_{h,R}$ gives $(S_{h,R}-S_R)Y_h\to0$. Decompose
\[
 Y_h-Y=(I-S_{h,R})Y_h+(S_{h,R}Y_h-S_RY)+(S_R-I)Y.
\]
Take the upper limit in $h$ and then let $R\to\infty$. This proves the first assertion. For the second, apply the first assertion to $Z_h$ and use
\[
 \norm{Y_h-Y}\le m^{-1/2}\norm{Z_h-Z}
   +\norm{(\mathbb A_h^{-1/2}-\mathbb A^{-1/2})Z}.
\]
\end{proof}

A screen for the action must include the full quadratic packet. For example, $\Fcurv_{A_h}$, $D_{A_h}H_h$, and the first spinor jets are different components, even when all are represented on the same finite history. A gap in the multiplicity commutant does not control their spatial or temporal frequencies.

\begin{proposition}[From spatial screens to spacetime compactness]
\label{prop:time-compactness}
Let $\mathcal H=L^2(\Sigma;E)$ and let $P_R$ be finite-rank orthogonal projections with smooth ranges and $P_R\to I$ strongly. Suppose $Y_h$ is bounded in $L^2(0,T;\mathcal H)$,
\[
 \lim_{R\to\infty}\sup_h\norm{(I-P_R)Y_h}_{L^2_t\mathcal H}=0,
\]
and, for each fixed $R$, $P_RY_h$ is bounded in $H^1(0,T;P_R\mathcal H)$. Then $(Y_h)$ is precompact in $L^2(0,T;\mathcal H)$. The last condition follows, for example, from a bound on $\partial_tY_h$ in $L^2(0,T;H^{-m}(\Sigma;E))$ for a fixed $m$.
\end{proposition}
\begin{proof}
For fixed $R$, the range is finite dimensional, so one-dimensional compact Sobolev embedding makes $(P_RY_h)$ precompact in $L^2_t\mathcal H$. The tail estimate gives a finite $\epsilon$-net for $(Y_h)$ by first choosing $R$ and then a finite net for the projected sequence. For the final assertion, pair $\partial_tY_h$ with a smooth orthonormal basis of $P_R\mathcal H$; each coefficient has a bounded $L^2$ time derivative.
\end{proof}

This elementary special case of spacetime compactness is consistent with the general theory of Simon \cite{Simon1986}. The time condition is indispensable: $Y_h(t,x)=\sin(t/h)v(x)$ remains in a single spatial mode but has no strongly convergent subsequence in spacetime $L^2$. Thus compact spatial Hodge screens alone do not establish the required spacetime result.

\subsection{Sectorwise first-variation estimates}
\label{app:sectorwise-variation}

Let $d_K(z_h,z)$ denote the sum of the difference norms in \eqref{eq:strong-geometry}--\eqref{eq:strong-spinors}, together with $\abs{\theta_h-\theta}$.  The same formula defines $d_K(z_h,z_j)$ for two reconstructed cutoffs on the common local bundle realization.  When necessary, take these norms on a slightly larger compact region containing the supports of the test lifts.

\begin{proposition}[First-variation continuity from primitive fields]
\label{prop:variation-continuity}
On a uniformly bounded classical strong packet,
\begin{equation}
 \norm{D\Sact_{\theta_h}(z_h)-D\Sact_\theta(z)}_{(\Vtest_K^{r_0})^*}
 \le C_K d_K(z_h,z).
 \label{eq:variation-bound}
\end{equation}
The same estimate holds separately for the gravitational and matter terms. The matter metric variation has the limit prescribed by \eqref{eq:stress-definition}; it is not an independent convergence assumption.
\end{proposition}

\begin{proof}
All arguments take place in one compact region and a finite bundle atlas. A partition of unity introduces only fixed smooth coefficients. We give the separate sector estimates, including the terms specific to metric variation.

\emph{Yang--Mills sector.} Write its local density as a bounded coframe coefficient contracted with two curvatures. A gauge variation gives
\[
 \dot\Fcurv_{A_h}=\dd a+[A_h,a],
\]
which converges in $L^2$ uniformly for bounded smooth $a$, by $A_h\to A$ in $L^4$. The pairing of $\dot\Fcurv_{A_h}$ and $\Fcurv_{A_h}$ therefore converges in $L^1$. A metric variation changes the coefficient and has the form
\begin{equation}
 B'_h[k](\Fcurv_{A_h},\Fcurv_{A_h}).
 \label{eq:metric-gauge-quadratic}
\end{equation}
The coefficient converges in $L^\infty$ and both curvature factors converge in $L^2$. Equation \eqref{eq:quadratic-product} proves convergence with the stated norm bound. If curvature were only weakly convergent, \eqref{eq:metric-gauge-quadratic} would contain two weak factors; a strongly convergent Hodge coefficient would not repair that loss.

\emph{Higgs kinetic sector.} With $K_h=D_{A_h}H_h$,
\[
 \dot K_h=D_{A_h}\eta_H+\rho_H(a)H_h.
\]
This converges strongly in $L^2$. The field variation is a pairing of $\dot K_h$ and $K_h$; the metric variation is a bounded coefficient applied to $K_h\otimes K_h$. Both follow from \Cref{lem:products}.

\emph{Higgs potential.} On a bounded $L^4$ set,
\begin{align}
 \norm{\abs{H_h}^4-\abs H^4}_1
 &\le C(\norm{H_h}_4+\norm H_4)^3\norm{H_h-H}_4,\notag\\
 \norm{\abs{H_h}^2H_h-\abs H^2H}_{4/3}
 &\le C(\norm{H_h}_4+\norm H_4)^2\norm{H_h-H}_4.
 \label{eq:quartic-continuity}
\end{align}
The lower-degree terms satisfy the analogous estimates. These inequalities pass both the volume variation of the potential and its Higgs variation, including the finite parameter limits.

\emph{Dirac--Yukawa sector.} Sobolev embedding on the four-dimensional compact chart gives strong $L^4$ convergence of both spinor variables from \eqref{eq:strong-spinors}. The derivative terms have the form
\[
 C(e_h)\bPsi_h\partial\Psi_h,
 \qquad C(e_h)(\partial\bPsi_h)\Psi_h.
\]
They converge in $L^1$, using the $L^2$ bounds for spinors and first derivatives. Spin-connection terms contain one $L^2$ coefficient $\omega(e_h)$ and the $L^2$ product $\bPsi_h\Psi_h$; they converge in $L^1$ by \eqref{eq:spin-connection-convergence}. Gauge and Yukawa terms contain, respectively, $A_h\bPsi_h\Psi_h$ and an affine function of $H_h$ multiplying $\bPsi_h\Psi_h$. H\"older with three $L^4$ factors gives convergence in $L^{4/3}$, hence in $L^1$ on the compact chart.

Varying $\Psi,\bPsi,A,H$ replaces one of these factors by its corresponding smooth test or derivative. Varying the metric differentiates $C(e_h)$ and $\omega(e_h)$. The additional spin term is $\dot\omega(e_h)[\dot e_h(k)]\bPsi_h\Psi_h$, again an $L^2$--$L^2$ product. Thus the full metric variation, including the spin connection, converges uniformly on bounded $C^1$ test sets. This proves the matter assertion without omitting a fermionic stress term.

\emph{Gravity.} Let $B(e)$ be the coframe bivector density of \Cref{app:palatini}. On compactly supported variations, integration by parts changes the torsion-free Palatini density into
\begin{equation}
 -\dd B(e)\wedge\omega(e)+B(e)\wedge\omega(e)\wedge\omega(e),
 \label{eq:palatini-first-order}
\end{equation}
with the chosen Einstein--Hilbert normalization. The volume term is algebraic in $e$. In the first variation of \eqref{eq:palatini-first-order}, each differentiated factor is in $L^2$, each algebraic coefficient is in $L^\infty$, and the variation introduces at most $k$ and $\partial k$. The coefficient and connection estimates of \Cref{lem:products} therefore apply to every term. They also define the gravitational first variation at the limiting $H^1\cap L^\infty$ metric. On smooth metrics it is the ordinary Einstein first variation, and its distributional extension is fixed by the integration-by-parts formula, as explained in \Cref{app:palatini}.

Subtracting each finite multilinear expression one factor at a time gives a bound linear in $d_K$. The other factors are uniformly bounded by the packet hypotheses. A finite sum of these estimates proves \eqref{eq:variation-bound}.
\end{proof}

\begin{remark}[Weak--strong gravity is a different route]
The geometric condition \eqref{eq:strong-geometry} is a strong first-jet route, not a requirement of strong second derivatives. \Cref{prop:weak-palatini} gives an alternative Palatini passage from strong coframe products and weak, identified curvature. That alternative does not weaken the strong curvature condition in the \emph{Yang--Mills stress}. The gravitational Palatini term is linear in curvature; the Yang--Mills metric variation is quadratic in curvature.
\end{remark}

\subsection{Palatini products, metric variations, and the spin-current alternative}
\label{app:palatini}

\subsubsection{The first-order gravitational distribution}

Fix the orientation and curvature sign so that, for the torsion-free connection,
\begin{equation}
 B_{ab}(e)\wedge R^{ab}(\omega(e))=\Scal(g)\dV_g,
 \label{eq:B-normalization}
\end{equation}
where $B(e)$ is the appropriately normalized internal dual of $e\wedge e$. For a cutoff function $\chi$ equal to one near the support of a physical test,
\begin{align}
 \int\chi B\wedge R(\omega)
 =\int\bigl[-\dd(\chi B)\wedge\omega
              +\chi B\wedge\omega\wedge\omega\bigr].
 \label{eq:Palatini-IBP}
\end{align}
Exterior Stokes proves the identity for smooth compactly supported data. The terms involving $\dd\chi$ do not contribute to a variation supported where $\chi=1$.

For $e,e^{-1}\in L^\infty$ with $e\in H^1$, the Levi--Civita connection $\omega(e)$ is in $L^2$. Every term on the right of \eqref{eq:Palatini-IBP} is integrable, and its derivative in the lift \eqref{eq:metric-lift} is bounded by a constant times $\norm k_{C^1}$. This defines the local Einstein first-variation distribution. Different choices of $\chi$ give the same functional on tests supported in their common unit region. On smooth metrics the distribution is $(\Eins(g)+\Lambda g)\dV_g/(2\kappa)$.

This definition keeps the derivative of the connection in divergence form.
The need to specify a regularity class in which nonlinear curvature is a
well-defined distribution is classical in general relativity
\cite{GerochTraschen1987}; the representative used here is consistent with
the stability framework for locally bounded metrics with locally
square-integrable first derivatives and bounded inverse developed by
LeFloch--Mardare \cite{LeFlochMardare2007}.  It does not assert a pointwise
curvature tensor or an unqualified product of a rough coefficient with an
arbitrary distribution.

\subsubsection{A weaker curvature passage}

\begin{proposition}[Weak-curvature Palatini passage]
\label{prop:weak-palatini}
On a compact four-dimensional region suppose
\begin{equation}
 e_h\to e\text{ in }L^2,\qquad
 \sup_h\norm{e_h}_6<\infty,\qquad
 R_h\rightharpoonup R\text{ in }L^p,
 \quad p>\tfrac32.
 \label{eq:Palatini-hypotheses}
\end{equation}
Assume $R$ is identified with the curvature of the limiting represented connection. Then $B(e_h)\to B(e)$ in $L^{p'}$, $p'=p/(p-1)$, and the Palatini pairings converge. The volume densities converge in $L^1$. If the represented coframe variations converge strongly in $L^2$, are uniformly bounded in $L^6$, and have compatible supports, their coframe first variations also converge. For a full metric variation, any connection-Euler pairing introduced by the lift must additionally tend to zero, or be retained as part of the limiting connection equation.
\end{proposition}
\begin{proof}
Interpolation between strong $L^2$ and bounded $L^6$ gives strong $L^q$ convergence for every $2\le q<6$; the lower exponents follow by finite volume. Since $p>3/2$, $2p'<6$. The product estimate gives
\[
 \norm{e_h\wedge e_h-e\wedge e}_{p'}
 \le C\norm{e_h-e}_{2p'}(\norm{e_h}_{2p'}+\norm e_{2p'})\to0.
\]
Pair with weak $L^p$ curvature. Taking $q=4$ gives convergence of the four-coframe volume product in $L^1$ by a four-factor telescoping estimate. The same argument applies to $\dot B_h$, a sum of one coframe and one coframe-variation factor, and to the volume variation. A connection variation produces its own Euler pairing after integration by parts. The stated extra condition, rather than coframe compactness alone, is what controls that term.
\end{proof}

The hypothesis identifying $R$ is substantive. A weak limit of matrix-valued two-forms is not automatically the curvature of the limiting connection. Nor does the Palatini lemma apply to a quadratic gauge stress by replacing one factor with a Hodge coefficient.

\subsubsection{Stationary elimination is not a universal torsion argument}

The independent connection Euler equation for a Palatini--Holst action has the form
\begin{equation}
 \mathcal C_e(T)=J_{\mathrm{spin}},
 \label{eq:Cartan-source}
\end{equation}
where $\mathcal C_e$ is the Cartan map determined by the coframe and the Palatini--Holst coefficients. On a nondegenerate chart with a uniformly bounded inverse of that map,
\[
 \norm T\le C\norm{J_{\mathrm{spin}}+r_{\mathrm{conn}}}.
\]
Thus zero spin current and vanishing connection residual imply torsionlessness. With nonzero current, the inverse instead reconstructs contorsion. Substitution into the total action retains the resulting spin-current interaction \cite{Kibble1961,HehlEtAl1976}. This is compatible with the connection alternative discussed in the Lorentzian construction \cite{PelissierRenewalSpacetime2026}, but differs from the metric action chosen in the main theorem.

There is also a separate normalization condition. A nonzero Holst coefficient alone does not imply a nonzero Einstein coefficient: on a torsion-free connection the Holst coframe equation reduces to the first Bianchi identity. The main action fixes a nonzero Palatini/Einstein coefficient $1/(2\kappa)$ before taking the limit. Pure Holst stationarity therefore cannot be used to infer \eqref{eq:Einstein-SM}.

\section{Strong-packet and regularity refinements}
\label{app:strong-packet}

These results strengthen reduced closure on additional regular branches.
Independent quartic integrability and spinor first-jet screens are optional
sufficient conditions here, not premises of \Cref{thm:reduced-closure}.
The sectorwise estimates used below are in
\Cref{app:sectorwise-variation}.

\begin{remark}[Meaning of the classical limit]
\label{rem:classical-distributional}
Throughout, ``classical'' distinguishes the non-quantized field theory from
a quantum measure.  At the reduced regularity the conclusion is a
\emph{distributional solution of the classical Euler system}, with the
Einstein first variation defined by \Cref{app:palatini}.  Neither smoothness
nor well-posedness of an initial-value problem follows from distributional
stationarity alone.  The regularity and evolution statements are given
separately in \Cref{cor:strong-solution-upgrade,thm:native-closure}.
\end{remark}

The Higgs potential requires a different endpoint norm. On a four-dimensional bounded region, an $H^1$ bound controls $L^4$ but does not make its unit ball compact in $L^4$.

\begin{proposition}[A local quartic compactness certificate]
\label{prop:orlicz}
Let $H_h\rightharpoonup H$ in $H^1(K)$ on a bounded smooth four-dimensional region. Suppose $\Phi:[0,\infty)\to[0,\infty)$ is increasing, satisfies $\Phi(t)/t\to\infty$, and
\begin{equation}
 \sup_h\int_K\Phi(\abs{H_h}^4)\,\dV_0<\infty.
 \label{eq:orlicz}
\end{equation}
Then $H_h\to H$ strongly in $L^4(K)$. A uniform $L^{4+\delta}$ bound, $\delta>0$, is sufficient.
\end{proposition}
\begin{proof}
Rellich compactness gives strong $L^2$ convergence and convergence in measure. For any $B>0$,
\[
 \int_{\{\abs{H_h}^4>B\}}\abs{H_h}^4
 \le\sup_{t>B}\frac{t}{\Phi(t)}\int_K\Phi(\abs{H_h}^4).
\]
The right side tends to zero uniformly in $h$. Hence the fourth powers are uniformly integrable. Convergence in measure and this uniform integrability give $\norm{H_h-H}_4\to0$ by the elementary Vitali argument, applied after using $\abs{H_h-H}^4\le8(\abs{H_h}^4+\abs H^4)$. Take $\Phi(t)=t^{1+\delta/4}$ for the last assertion.
\end{proof}

\subsection{A-priori routes to the action topology}

The compact-screen formulation is useful at low regularity, but it should not be read as the only way to obtain the nonlinear action topology. On a regular branch, ordinary a-priori Sobolev estimates imply the required bosonic compactness directly.

\begin{proposition}[A-priori Sobolev bounds imply bosonic strong compactness]
\label{prop:sobolev-bosonic}
Let $K\Subset M$ be a bounded smooth four-dimensional region. Suppose that for some $\sigma>0$, after fixed local gauge and frame choices,
\begin{equation}
 \sup_h\Bigl(
 \norm{e_h}_{H^{2+\sigma}(K)}
 +\norm{A_h}_{H^{1+\sigma}(K)}
 +\norm{H_h}_{H^{1+\sigma}(K)}
 \Bigr)<\infty,
 \qquad
 \sup_h\norm{e_h^{-1}}_{L^\infty(K)}<\infty.
 \label{eq:apriori-bosonic}
\end{equation}
Assume also that the coefficient banks remain in a compact subset of the physical parameter region. Then every cutoff sequence has a subsequence such that
\begin{align}
 e_h&\to e &&\text{strongly in }H^1(K)\cap L^\infty(K),\notag\\
 A_h&\to A &&\text{strongly in }H^1(K)\text{ and hence in }L^4(K),\notag\\
 H_h&\to H &&\text{strongly in }H^1(K)\cap L^4(K),\notag\\
 \Fcurv_{A_h}&\to\Fcurv_A &&\text{strongly in }L^2(K),\notag\\
 D_{A_h}H_h&\to D_AH &&\text{strongly in }L^2(K).
 \label{eq:apriori-bosonic-convergence}
\end{align}
Moreover $H_h$ is uniformly bounded in $L^{4+\delta}(K)$ for some $\delta>0$ depending only on $\sigma$ and the dimension, so quartic uniform integrability follows automatically. Thus, on this regular branch, the bosonic part of the strong packet is obtained from a-priori boundedness rather than assumed precompactness or screen tails.
\end{proposition}
\begin{proof}
Rellich--Kondrachov gives compact embeddings $H^{2+\sigma}(K)\Subset H^1(K)$ and $H^{1+\sigma}(K)\Subset H^1(K)$. Since $2+\sigma>2=d/2$, the first family is also precompact in $C^0(K)$ after lowering the H\"older exponent, hence in $L^\infty(K)$. The uniform inverse bound preserves the nondegenerate signature chart. Strong $H^1$ convergence of $A_h$ gives strong $L^4$ convergence in four dimensions. Therefore
\[
 \dd A_h\to\dd A\quad\text{in }L^2,
 \qquad
 A_h\wedge A_h\to A\wedge A\quad\text{in }L^2,
\]
which proves the curvature convergence. The same argument, using
\[
 D_{A_h}H_h=\dd H_h+\rho_H(A_h)H_h,
\]
gives strong $L^2$ convergence of the Higgs covariant derivative. Finally $H^{1+\sigma}$ embeds continuously into some $L^{4+\delta}$ with $\delta>0$, yielding quartic uniform integrability. Finite-dimensional compactness handles the coefficients.
\end{proof}

\begin{proposition}[Lorentzian Dirac stability produces spinor $H^1$ compactness]
\label{prop:dirac-stability}
Let $I\times\Sigma$ be a compact subslab with $\Sigma$ closed. In a common local spin trivialization suppose the reconstructed spinors satisfy symmetric-hyperbolic systems
\begin{equation}
 A_h^0\partial_t\Psi_h+A_h^j\partial_j\Psi_h+B_h\Psi_h=r_h,
 \label{eq:dirac-hyperbolic}
\end{equation}
with the analogous dual system for $\bPsi_h$. Assume the matrices $A_h^0$ are uniformly positive definite, the systems are uniformly symmetric hyperbolic, and
\begin{equation}
 \sup_h\bigl(
 \norm{A_h^\mu}_{W^{1,\infty}(I\times\Sigma)}
 +\norm{B_h}_{L^\infty_{t,x}}
 +\norm{\Psi_h}_{L^\infty_tH^2_x}
 +\norm{\bPsi_h}_{L^\infty_tH^2_x}
 \bigr)<\infty.
 \label{eq:dirac-uniform}
\end{equation}
Suppose the coefficients are Cauchy in $L^\infty(I\times\Sigma)$, the spinor and dual-spinor initial data are Cauchy in $L^2(\Sigma)$, and the residuals $r_h,\bar r_h$ are Cauchy in $L^2_tL^2_x$. Then $\Psi_h$ and $\bPsi_h$ are Cauchy in $H^1(I\times\Sigma)$. In particular, if $r_h,\bar r_h\to0$, the limiting fields satisfy the Lorentzian Dirac--Yukawa and dual equations determined by the limiting coefficients.
\end{proposition}
\begin{proof}
For $w=\Psi_h-\Psi_j$, subtract the two first-order systems and use the $h$-system as the principal symmetric-hyperbolic operator. Uniform positivity of $A_h^0$ and the coefficient bounds in \eqref{eq:dirac-uniform} give the standard $L^2$ energy inequality
\begin{equation}
 \norm{w(t)}_{L^2_x}
 \le C_T\left[
 \norm{w(0)}_{L^2_x}
 +\norm{r_h-r_j}_{L^1_tL^2_x}
 +\delta_{h,j}
 \right],
 \label{eq:dirac-L2-stability}
\end{equation}
where $\delta_{h,j}\to0$ is the coefficient distance multiplied by uniform $L^2$ bounds for the first derivatives of $\Psi_j$. The spatial derivative bounds follow from the uniform $H^2_x$ estimate; the normal derivative is bounded by solving \eqref{eq:dirac-hyperbolic} for $\partial_t\Psi_j$, using uniform invertibility of $A_j^0$ and boundedness of the residuals. Hence $\Psi_h$ is Cauchy in $C_tL^2_x$.

Interpolation on the closed three-manifold gives, uniformly in $t$,
\[
 \norm{w(t)}_{H^1_x}
 \le C\norm{w(t)}_{L^2_x}^{1/2}
        \norm{w(t)}_{H^2_x}^{1/2},
\]
so $\Psi_h$ is Cauchy in $C_tH^1_x$. Finally solve \eqref{eq:dirac-hyperbolic} for $\partial_t\Psi_h$. Strong convergence of the spatial $H^1$ terms and coefficients, together with $r_h$ Cauchy in $L^2_tL^2_x$, gives strong convergence of $\partial_t\Psi_h$ in $L^2_tL^2_x$. This is strong $H^1(I\times\Sigma)$ convergence. The dual system is identical. Passing to the limit in \eqref{eq:dirac-hyperbolic} proves the final assertion.
\end{proof}

\subsection{Primitive compactness certificates}

The strong packet used below is a conclusion of a concrete compactness mechanism once the quantities entering the nonlinear first variation are controlled in their own topologies. The low-regularity formulation is deliberately modular: direct spacetime compact screens may be used, or spatial screens may be combined with the time regularity of \Cref{prop:time-compactness}. For spinors one may instead use the Lorentzian evolution route of \Cref{prop:dirac-stability}; the regular Sobolev branch of \Cref{thm:regular-branch} avoids compact screens altogether.

\begin{definition}[Classical compactness certificate]
\label{def:compactness-certificate}
A finite-interface regulator sequence satisfies the classical compactness certificate on $K\Subset M$ if, after fixed local gauge and spin-frame identifications,
\begin{enumerate}[label=\textup{(C\arabic*)},leftmargin=3.2em]
\item the coframes are precompact in $L^\infty(K)$, uniformly nondegenerate, and bounded in $H^1(K)$; their first-derivative packet has a common compact screen in $L^2(K)$;
\item the connections are precompact in $L^4(K)$, the curvatures are bounded in $L^2(K)$, and the curvature packet has a common compact screen in $L^2(K)$;
\item the Higgs fields are bounded in $H^1(K)$, the covariant gradients are bounded in $L^2(K)$ and have a common compact screen, and there is an increasing $\Phi$ with $\Phi(t)/t\to\infty$ such that
\[
 \sup_h\int_K\Phi(\abs{H_h}^4)\,\dV_0<\infty;
\]
\item for the spinors, either
      \begin{enumerate}[label=\textup{(C4\alph*)},leftmargin=3.0em]
      \item $\Psi_h,\bPsi_h$ are precompact in $L^2(K)$ and bounded in $H^1(K)$, while their ordinary first-jet packets in the fixed local frames have common compact screens in $L^2(K)$; or
      \item on a compact subslab containing $K$, the reconstructed Dirac and dual-Dirac systems satisfy the hypotheses of \Cref{prop:dirac-stability};
      \end{enumerate}
\item the coefficient banks remain in a compact subset of the region $\kappa,g_j,\lambda_H>0$.
\end{enumerate}
A ``common compact screen'' means the hypotheses of \Cref{lem:screen}; it may instead be verified by the spatial-tail and time-regularity conditions of \Cref{prop:time-compactness}.
\end{definition}

\subsection{The primitive convergence class}

\begin{definition}[Classical strong packet]
\label{def:strong-packet}
A finite-interface regulator sequence has a classical strong packet with limit $z=(e,A,H,\Psi,\bPsi)$ if, on every $K\Subset M$ and in the fixed local bundle identifications,
\begin{align}
 &e_h\to e\text{ in }L^\infty(K)\cap H^1(K),
 \quad \sup_h(\norm{e_h}_\infty+\norm{e_h^{-1}}_\infty)<\infty,
 \label{eq:strong-geometry}\\
 &A_h\to A\text{ in }L^4(K),
 \qquad \Fcurv_{A_h}\to\Fcurv_A\text{ in }L^2(K),
 \label{eq:strong-gauge}\\
 &H_h\to H\text{ in }L^4(K),
 \qquad D_{A_h}H_h\to D_AH\text{ in }L^2(K),
 \label{eq:strong-Higgs}\\
 &\Psi_h\to\Psi,\quad\bPsi_h\to\bPsi
       \text{ in }H^1(K),
 \label{eq:strong-spinors}\\
 &\theta_h\to\theta,
 \quad \kappa_h,g_{j,h},\lambda_{H,h}\text{ remain positive and bounded away from zero}.
 \label{eq:strong-parameters}
\end{align}
The limiting coframe is in the same orientation and signature chart. The curvatures in \eqref{eq:strong-gauge} are identified with the reconstructed connections, not independent weak tensor limits.
\end{definition}

These sufficient conditions use screens for the actual curvature and
covariant-gradient packets, the coframe noncollapse bound and, on the
stated strong-packet branch, \Cref{prop:orlicz} for the Higgs amplitude.
No optimality is claimed.

The Higgs conditions imply strong ordinary $H^1$ convergence because
\[
 \dd H_h=D_{A_h}H_h-\rho_H(A_h)H_h
\]
and the second term converges strongly in $L^2$ by H\"older. For spinors, route \textup{(C4a)} is a direct compactness certificate, whereas route \textup{(C4b)} derives the same $H^1$ compactness from the Lorentzian Dirac evolution. No elliptic coercivity of a Lorentzian Dirac graph norm is used.

\begin{theorem}[Finite certificates produce the classical strong packet]
\label{thm:certificate-packet}
If a finite-interface regulator sequence satisfies \Cref{def:compactness-certificate} on every $K\Subset M$, then every cutoff sequence has a subsequence and a field tuple $z=(e,A,H,\Psi,\bPsi)$ for which the convergences of \Cref{def:strong-packet} hold locally.  The limiting curvature is $\Fcurv_A$ and the limiting Higgs derivative is $D_AH$; they are not independent tensor limits.
\end{theorem}
\begin{proof}
On each compact chart, extract the uniformly convergent coframes and the
$L^4$-convergent connections. The common screens give strong $L^2$
convergence of the coframe derivatives, curvatures and covariant Higgs
gradients. Passing to distributions in $\dd A_h=F_{A_h}-A_h\wedge A_h$
identifies the curvature. Rellich compactness and \Cref{prop:orlicz} give
strong $L^4$ Higgs convergence; the identity
$\dd H_h=D_{A_h}H_h-\rho_H(A_h)H_h$ identifies the gradient and gives
strong $H^1$ convergence. Under \textup{(C4a)}, value compactness and the
first-jet screens give strong $H^1$ spinor convergence; under
\textup{(C4b)}, apply \Cref{prop:dirac-stability}. Extract the coefficient
bank and diagonalize over the atlas and compact exhaustion. Uniform inverse
bounds preserve the signature chart.
\end{proof}

\Needspace{8\baselineskip}
\subsection{Subsequential variational closure and cofinal uniqueness}

\begin{corollary}[Strong-packet refinement of variational closure]
\label{thm:main-limit}
Let $(\Iface_h)_{h\downarrow0}$ be any finite-interface regulator sequence in the sense of \Cref{def:regulator}. Assume it satisfies the classical compactness certificate of \Cref{def:compactness-certificate} on every $K\Subset M$, the finite reconstruction is first-variation consistent,
\begin{equation}
 c_h(K)\longrightarrow0,
 \label{eq:main-consistency}
\end{equation}
and the physical common-action residual is controlled by \Cref{prop:source-bound} with
\begin{equation}
 L_h\sqrt{R_h}+e_h\longrightarrow0.
 \label{eq:main-source-vanishing}
\end{equation}
Then every cutoff sequence $h_n\downarrow0$ has a subsequence $h_{n_j}$ and a field tuple $z=(e,A,H,\Psi,\bPsi)$ such that $z_{h_{n_j}}$ converges to $z$ in the classical strong-packet topology on compact subsets. Along this subsequence the reconstructed first variations converge to $D\Sact_\theta(z)$ and
\begin{equation}
 D\Sact_\theta(z)[v]=0\qquad(v\in\Vtest_K).
 \label{eq:main-stationary}
\end{equation}
Consequently the limiting fields solve the classical Yang--Mills, Higgs, and Dirac--Yukawa equations, and
\begin{equation}
 \Eins_{\mu\nu}(g)+\Lambda g_{\mu\nu}
 =\kappa T^{\mathrm{SM}}_{\mu\nu}(z)
 \quad\text{in }\Dtest'(M).
 \label{eq:Einstein-SM}
\end{equation}
More precisely, for the selected subsequence,
\begin{equation}
 \norm{D\Sact_\theta(z)}_{(\Vtest_K^{r_0})^*}
 \le L_{h_{n_j}}\sqrt{R_{h_{n_j}}}+e_{h_{n_j}}
      +c_{h_{n_j}}(K)+C_Kd_K(z_{h_{n_j}},z).
 \label{eq:main-residual-bound}
\end{equation}
No summable adjacent-cutoff transport, continuum stress tensor, Euler equation, or candidate strong limit is assumed. Compactness produces accumulation points; consistency and stationarity identify every such strong-packet accumulation point as a classical Einstein--Standard-Model solution.
\end{corollary}

\begin{proof}
\Cref{thm:certificate-packet} supplies a strong-packet subsequence.
\Cref{prop:variation-continuity} bounds its continuum first-variation
difference by $C_Kd_K$; \Cref{prop:source-bound} bounds the finite
stationarity term by $L_h\sqrt{R_h}+e_h$. Adding the consistency error
gives \eqref{eq:main-residual-bound}. All three contributions vanish.
Independent matter tests yield their Euler rows, and the inverse-metric
test with \eqref{eq:stress-definition} gives \eqref{eq:Einstein-SM}.
The gravitational distribution is the first-order representative of
\Cref{app:palatini}.
\end{proof}

\begin{proposition}[Summable cofinal transport]
\label{prop:cofinal-transport}
Let $h_n\downarrow0$ be a cofinal sequence satisfying the classical compactness certificate on every compact region, with adjacent cutoffs represented on the same local bundles.  Suppose
\begin{equation}
 d_K(z_{h_{n+1}},z_{h_n})\le\eta_{n,K},
 \qquad \sum_{n=1}^\infty\eta_{n,K}<\infty
 \label{eq:cofinal-transport}
\end{equation}
for every $K\Subset M$.  Then $(z_{h_n})$ is Cauchy in each norm entering $d_K$ and converges to a unique local strong-packet limit.  If two admissible cofinal routes admit direct comparisons whose $d_K$ error tends to zero after both cutoffs pass a common scale, their limits agree.  In particular, if \Cref{thm:certificate-packet} supplies subsequential strong-packet limits and \eqref{eq:cofinal-transport} holds, no subsequence choice remains.
\end{proposition}
\begin{proof}
For $m>n$, the triangle inequality bounds
$d_K(z_{h_m},z_{h_n})$ by $\sum_{j=n}^{m-1}\eta_{j,K}\to0$.
Completeness gives a unique norm limit; a subsequential strong packet from
\Cref{thm:certificate-packet} supplies its curvature and covariant-derivative
identifications. Passing the direct comparison between two routes to their
limits proves route independence.
\end{proof}

\begin{corollary}[Cofinal uniqueness and route independence]
\label{cor:cofinal-unique}
Under the hypotheses of \Cref{thm:main-limit}, let a cofinal sequence additionally satisfy the summable transport estimate \eqref{eq:cofinal-transport}. Then the entire cofinal sequence converges in the strong-packet topology to a single classical Einstein--Standard-Model solution. If two admissible cofinal routes satisfy the direct-comparison hypothesis of \Cref{prop:cofinal-transport}, their limits agree.
\end{corollary}
\begin{proof}
Apply \Cref{prop:cofinal-transport} and identify the unique accumulation
point by \Cref{thm:main-limit}.
\end{proof}

\begin{theorem}[Regularity criterion for variational closure]
\label{thm:regular-branch}
This is a sufficient regularity criterion.
Let $(\Iface_h)$ be a finite-interface regulator sequence and let $Q=I\times\Sigma\Subset M$ be a compact subslab. Suppose that for some $\sigma>0$, in fixed local gauges and frames,
\begin{equation}
 \sup_h\left(
 \norm{e_h}_{H^{3+\sigma}(Q)}
 +\norm{A_h}_{H^{3+\sigma}(Q)}
 +\norm{H_h}_{H^{3+\sigma}(Q)}
 \right)<\infty,
 \qquad
 \sup_h\norm{e_h^{-1}}_{L^\infty(Q)}<\infty,
 \label{eq:regular-bosonic-bound}
\end{equation}
and the coefficient banks remain in a compact physical set. Suppose the reconstructed spinors and dual spinors obey the Lorentzian systems of \Cref{prop:dirac-stability}, with uniformly bounded $L^\infty_tH^2_x$ norms, precompact initial data in $L^2(\Sigma)$, and residuals tending to zero in $L^2_tL^2_x$. Finally assume first-variation consistency and scaled stationarity, \eqref{eq:main-consistency} and \eqref{eq:main-source-vanishing}.

Then every cutoff sequence has a subsequence converging in the full classical strong packet on $Q$, and every such limit solves the classical Einstein--Standard-Model equations there. No compact-screen hypothesis and no a-priori spinor first-jet compactness is required on this regular branch.
\end{theorem}
\begin{proof}
\Cref{prop:sobolev-bosonic} gives strong bosonic compactness.
The stronger bounds \eqref{eq:regular-bosonic-bound} also give uniform
coefficient convergence and the bounds required by
\Cref{prop:dirac-stability}. Extract Cauchy initial spinors and use the
vanishing residuals to obtain their strong spacetime $H^1$ convergence.
The resulting strong packet, consistency and scaled stationarity imply
closure by \Cref{thm:main-limit}.
\end{proof}

\begin{corollary}[Renewal Geometry realization]
\label{cor:renewal-realization}
On the source-native Standard-Model branch of
\Cref{prop:renewal-interface}, the compactness, consistency and scaled
stationarity hypotheses of \Cref{thm:main-limit} give a subsequential
Einstein--Standard-Model limit. The additional hypotheses of
\Cref{cor:cofinal-unique} give a unique route-independent limit.
For the explicit local action, the finite physical-row, initial/gauge and
leakage budgets of \Cref{thm:native-closure} instead yield quantitative
regularity and stress convergence, with cofinal summability on the rate
family of \Cref{cor:native-rate}.
\end{corollary}
\begin{proof}
Combine the structural realization \Cref{prop:renewal-interface} with
\Cref{thm:main-limit}, and use \Cref{cor:cofinal-unique} for the optional
transport conclusion. The quantitative statements follow from
\Cref{thm:native-closure,cor:native-rate}.
\end{proof}

\begin{corollary}[Exact finite critical branch]
\label{cor:exact-critical}
Under the compactness and consistency hypotheses of \Cref{thm:main-limit}, suppose the finite configurations are exact critical points in all lifted physical directions,
\[
 D\Sact_h(z_h^{\mathrm d})[\mathcal I_hv]=0
 \qquad(v\in\Vtest_K).
\]
Then every strong-packet accumulation point satisfies \eqref{eq:Einstein-SM} and all classical matter Euler equations without a separate source-residual hypothesis.
\end{corollary}
\begin{proof}
Exact criticality gives $\varepsilon_h(K)=0$ in \eqref{eq:stationarity}.  The proof of \Cref{thm:main-limit} then uses only $c_h(K)\to0$ and $d_K(z_h,z)\to0$.
\end{proof}

\begin{corollary}[Distributional and negative-Sobolev stress convergence]
\label{cor:stress-topology}
Assume the sectorwise consistency in \eqref{eq:C1-consistency}. Define the operational stress density by
\[
 \langle\mathsf T_h,k\rangle
 =-2D\Sact_{\mathrm{SM},h}[\mathcal I_h(k,0,0,0,0)].
\]
Along every strong-packet subsequence furnished by \Cref{thm:main-limit}, $\mathsf T_h\to T^{\mathrm{SM}}\dV_g$ in distributions. On each fixed compact test region it also converges in $H^{-m}$ for every integer $m>r_0+2$, after using the fixed reference density and bundle frames.
\end{corollary}
\begin{proof}
Sectorwise consistency and \Cref{prop:variation-continuity} give convergence
in the dual $C^{r_0}$ norm. Restrict the test functional to $H_0^m$ and use
$H_0^m\hookrightarrow C^{r_0}$ for $m>r_0+2$ in four dimensions.
\end{proof}

\begin{corollary}[A realized random branch]
\label{cor:random-realization}
On a common probability space suppose the reconstructed fields satisfy the compactness and consistency hypotheses of \Cref{thm:main-limit} and the cofinal-transport hypotheses of \Cref{cor:cofinal-unique} almost surely, with limiting field $z(\omega)$.  Let the source estimate of \Cref{prop:source-bound} hold pathwise with deterministic $L_h$ and $e_h\to0$, and assume
\[
 \sum_h L_h^2\,\mathbb E\norm{a_h}_{G_h^{-1}}^2<\infty.
\]
Then the limiting fields satisfy \eqref{eq:Einstein-SM} and the matter Euler equations almost surely. It is enough to formulate the hypotheses on a countable dense test core with the displayed uniform bounds.
\end{corollary}
\begin{proof}
Tonelli gives $\sum_hL_h^2\|a_h\|_{G_h^{-1}}^2<\infty$ almost surely.
Thus \Cref{prop:source-bound} makes the physical residual vanish.
Apply \Cref{thm:main-limit,cor:cofinal-unique} on a countable determining
core and extend by continuity. Independence is unnecessary.
\end{proof}

\subsection{Conditional regularity upgrade}

\begin{corollary}[Regularity upgrade on a common slab]
\label{cor:strong-solution-upgrade}
Let a limit furnished by \Cref{thm:reduced-closure} additionally satisfy,
in fixed gauges on $I\times\Sigma$,
\[
 (g,A,H)\in C_tH_x^s\cap C_t^1H_x^{s-1}\cap C_t^2H_x^{s-2},
 \qquad
 (\Psi,\bPsi)\in C_tH_x^{s-1}\cap C_t^1H_x^{s-2},\quad s\ge5,
\]
with a common nondegenerate Lorentzian chart.  Then its distributional Euler
identities are strong Sobolev field equations.  In particular all second-order
bosonic and first-order spinor rows vanish in $C_tH_x^{s-2}$.
If the fields are smooth, these are the pointwise Einstein--Standard-Model
equations.
\end{corollary}
\begin{proof}
Sobolev multiplication on the three-dimensional closed slices makes all
coefficients and products in the Euler rows continuous in the indicated
Sobolev spaces.  The distribution represented by each such continuous row
is zero by \Cref{thm:reduced-closure}; therefore the row is zero in that
space.  This is a regularity upgrade conditional on the displayed bounds,
not a smoothing conclusion for a hyperbolic equation.
\end{proof}

\begin{remark}[Stronger compactness and applications]
The strong-packet criteria, explicit spinor evolution estimate, summable
cofinal transport, exact-critical-point and realized-random refinements are
proved in \Cref{app:strong-packet}.  They refine the reduced theorem when
norm convergence of spinor first jets or a pathwise cofinal limit is needed.
The structural realization \Cref{prop:renewal-interface} may be combined with
\Cref{thm:reduced-closure} once its analytic hypotheses are verified; that
combination does not infer those hypotheses from the representation packet.
\end{remark}

\section{Weak-limit estimates and diagnostic examples}
\label{app:weak-limit-estimates}

This appendix proves the sectorwise passages used in the joint defect
identity and the critical gauge-quotient theorem. Uniform coefficient
convergence is used for measure-valued stresses; conservation has the
additional connection and Noether hypotheses stated explicitly below.

\subsection{Critical Higgs and Yang--Mills stress measures}

We first weaken the matter compactness rather than silently deleting the resulting source. Work on a compact chart with smooth $g_h$ converging uniformly to a continuous Lorentzian metric $g$, uniform inverse bounds and fixed signature, $A_h\to A$ in $L^4$, and
\begin{equation}
 H_h\rightharpoonup H\text{ in }H^1,
 \qquad H_h\to H\text{ almost everywhere and in }L^q\ (q<4).
 \label{eq:critical-Higgs}
\end{equation}
The latter convergence follows along a subsequence from Rellich compactness. These are positive comparison norms even though the action is Lorentzian.
Uniform convergence of the metric, its inverse and the positive volume density
is sufficient for all algebraic stress-measure passages in this subsection.
The limiting coefficients have their continuous representatives, so their
products with singular Radon measures are well defined. No convergence of
metric derivatives is used for those passages. A separate $C^1$ hypothesis
will be imposed when a covariant conservation identity is passed to the limit.

\begin{lemma}[Cubic Euler terms and quartic concentration]
\label{lem:critical-cubic}
Under \eqref{eq:critical-Higgs},
\begin{equation}
 \abs{H_h}^2H_h\rightharpoonup\abs H^2H\text{ in }L^{4/3},
 \qquad D_{A_h}H_h\rightharpoonup D_AH\text{ in }L^2.
 \label{eq:critical-weak}
\end{equation}
After extraction on a fixed compact domain, there is a nonnegative Radon measure $\nu_H$ with
\begin{equation}
 \abs{H_h}^4\dV_{g_h}\rightharpoonup^*
 \abs H^4\dV_g+\nu_H.
 \label{eq:quartic-defect}
\end{equation}
On that compact domain, including its boundary, $\nu_H=0$ is equivalent to strong $L^4$ convergence along the extracted sequence. The local version applies on smaller compact regions, without excluding escape through the edge of a larger chart.
\end{lemma}
\begin{proof}
The cubic sequence is bounded in the reflexive space $L^{4/3}$ and converges almost everywhere; boundedness and convergence in measure identify every weak subsequential limit with the displayed cubic. For the covariant derivative, write
\[
 A_hH_h-AH=(A_h-A)H_h+A(H_h-H).
\]
The first term tends to zero in $L^2$. For any $L^2$ test $v$, the second term pairs with $A^*v\in L^{4/3}$ and tends to zero by weak $L^4$ convergence of $H_h$. The ordinary derivative converges weakly in $L^2$.

The fourth-power densities have bounded mass. Weak-star compactness and Fatou's lemma against nonnegative continuous tests give \eqref{eq:quartic-defect}. On a compact domain, testing with one gives convergence of fourth-power integrals if $\nu_H=0$. Weak $L^4$ convergence and uniform convexity then imply strong convergence. The converse follows from the product estimate \eqref{eq:quartic-continuity} and uniform convergence of the volume density.
\end{proof}

Let $\mathfrak K_H$ be the tensor-valued measure defect in the complete kinetic stress:
\begin{align}
 &\left(2\Rea\ip{D_\mu H_h}{D_\nu H_h}
      -g_{h,\mu\nu}\abs{D_{A_h}H_h}_{g_h}^2\right)\dV_{g_h}
 \rightharpoonup^*\nonumber\\[-0.3em]
 &\hspace{2em}
 \left(2\Rea\ip{D_\mu H}{D_\nu H}
      -g_{\mu\nu}\abs{D_AH}_g^2\right)\dV_g+\mathfrak K_{H,\mu\nu}.
 \label{eq:kinetic-defect}
\end{align}
Such a subsequence exists by the positive $L^2$ bound on the covariant derivatives.
No positivity of the Lorentzian tensor $\mathfrak K_H$ is claimed.
Once the weak covariant-gradient limit has been identified, this stress
passage uses no further convergence of the connection. In particular it
applies to the critical quotient class below, where the connection is only
strongly compact in subcritical $L^q$.

\begin{theorem}[Critical Higgs stress-defect identity]
\label{thm:higgs-defect}
For the minimal action \eqref{eq:SM-action}, assume uniform metric convergence
with continuous limit and uniform inverse bounds, the Higgs convergence
\eqref{eq:critical-Higgs}, an identified weak gradient limit
$D_{A_h}H_h\rightharpoonup D_AH$ in $L^2$, and convergent parameters. Then
\begin{equation}
 T^H_h\dV_{g_h}\rightharpoonup^*
 T^H\dV_g+\mathfrak S_H,
 \qquad
 \mathfrak S_{H,\mu\nu}
 =\mathfrak K_{H,\mu\nu}-\lambda_Hg_{\mu\nu}\nu_H.
 \label{eq:H-defect}
\end{equation}
If the remaining stress contributions and the gravitational first variation have their intended limits, and the operational metric residual and consistency error vanish, then
\begin{equation}
 \bigl(\Eins(g)+\Lambda g\bigr)\dV_g
 =\kappa\bigl(T^{\mathrm{SM}}\dV_g+\mathfrak S_H\bigr).
 \label{eq:defect-Einstein}
\end{equation}
For the conservation assertion assume additionally $g_h\to g$ in $C^1$,
that the total reconstructed stress densities are uniformly bounded as
Radon measures and converge weakly star, and that their covariant divergences
vanish distributionally in the limit. Then their limiting stress is
conserved. If the regular limiting matter stress also satisfies its Noether
identity (in particular for smooth limiting matter solutions),
$\nabla^\mu\mathfrak S_{H,\mu\nu}=0$. Distributional matter stationarity
alone at the reduced regularity is not used to assert this extra Noether
identity. Conservation does not imply $\mathfrak S_H=0$.
\end{theorem}
\begin{proof}
The kinetic term is \eqref{eq:kinetic-defect}. In the potential, the quadratic density $\abs{H_h}^2$ converges in $L^1$, and the constant term converges with the volume density. The only additional potential contribution is the quartic measure \eqref{eq:quartic-defect}, with sign fixed by \eqref{eq:H-stress}. This proves \eqref{eq:H-defect}. Test the metric equation with a compactly supported smooth $k$ and pass each term to obtain \eqref{eq:defect-Einstein}.

The conservation statement follows from the separate transport lemma below,
using the additional $C^1$ metric convergence. Subtract the limiting regular
Noether identity only when it is defined and holds. A nonzero constant
symmetric tensor on a flat chart already disproves the implication from
vanishing divergence to a vanishing tensor.
\end{proof}

\begin{lemma}[Covariant conservation under a strong connection limit]
\label{lem:defect-conservation-transport}
Let smooth metrics $g_h$ converge to $g$ in $C^1$ on a compact chart, with
uniform inverse bounds. Let symmetric covariant tensor densities
$\mathsf T_h$ be uniformly bounded Radon measures and converge weakly star
to $\mathsf T$. If, for every compactly supported smooth vector field $X$,
\begin{equation}
 \int \nabla_{g_h}^{(\mu}X^{\nu)}\,\dd\mathsf T_{h,\mu\nu}\longrightarrow0,
 \label{eq:defect-conservation-tested}
\end{equation}
then $\int\nabla_g^{(\mu}X^{\nu)}\,\dd\mathsf T_{\mu\nu}=0$.
If $\mathsf T=T(z)\dV_g+\mathfrak S$ and $T(z)\dV_g$ separately has this
identity, then $\mathfrak S$ has it as well.
\end{lemma}
\begin{proof}
Put $B_h(X)=\nabla_{g_h}^{(\mu}X^{\nu)}$. The coordinate formula for the
covariant derivative, the inverse-metric bounds and $C^1$ convergence give
$\norm{B_h(X)-B(X)}_\infty\to0$, with continuous $B(X)$. Thus
\[
 \left|\int(B_h(X)-B(X))\,\dd\mathsf T_h\right|
 \le\norm{B_h(X)-B(X)}_\infty\,|\mathsf T_h|(K)\to0.
\]
The remaining fixed continuous coefficient passes through weak-star
convergence. Subtraction proves the last assertion. Merely uniform metric
convergence would not control $B_h(X)$ and is not sufficient for this step.
\end{proof}

\begin{proposition}[Yang--Mills quadratic stress defect]
\label{prop:YM-defect}
Work on a compact chart with smooth $g_h$ converging uniformly to a continuous
Lorentzian metric $g$, uniform inverse bounds and fixed signature, and assume
an identified curvature limit
\[
 \Fcurv_{A_h}\rightharpoonup\Fcurv_A\quad\text{in }L^2.
\]
After extraction there is a symmetric tensor-valued finite Radon measure $\mathfrak S_{\mathrm{YM}}$ such that
\begin{equation}
 T^{\mathrm{YM}}_h\dV_{g_h}
 \rightharpoonup^*
 T^{\mathrm{YM}}\dV_g+\mathfrak S_{\mathrm{YM}}.
 \label{eq:YM-defect}
\end{equation}
Strong $L^2$ convergence of $\Fcurv_{A_h}$ is sufficient for
$\mathfrak S_{\mathrm{YM}}=0$. No strong $L^4$ convergence of $A_h$ is
needed for this algebraic stress-measure statement once its weak curvature
limit is known to be $F_A$.
\end{proposition}
\begin{proof}
The positive reference $L^2$ bound on the curvature gives a uniform $L^1$ bound for every component of the quadratic stress density because the metric coefficients and volume densities converge uniformly.  Weak-star compactness of finite measures therefore gives a subsequential tensor-valued limit.  Subtracting the regular quadratic tensor built from $\Fcurv_A$ defines $\mathfrak S_{\mathrm{YM}}$.  If the curvatures converge strongly in $L^2$, the quadratic product estimate \eqref{eq:quadratic-product} makes the difference of the stress densities tend to zero in $L^1$.
\end{proof}

\begin{proposition}[Matter equations survive the bosonic weak limit]
\label{prop:weak-matter-equations}
In the weak Yang--Mills and Higgs class above, suppose additionally that
$e_h\to e$ in $L^\infty\cap H^1$ in compatible spin frames, with uniformly
bounded inverse, and that $\Psi_h,\bPsi_h$ converge weakly in $H^1$.
If first-variation consistency and physical stationarity hold for all matter
tests, the limiting Yang--Mills, Higgs and Dirac--Yukawa equations hold in
distributions.  The complete fermionic metric first variation has its
intended limit without a separate stress-convergence assumption.
\end{proposition}
\begin{proof}
In a gauge variation, $\dd a+[A_h,a]$ converges strongly in $L^2$ and pairs
with the weakly convergent curvature.  In the Higgs kinetic variation,
$D_{A_h}\eta_H$ converges strongly in $L^2$.  The gauge-current contribution
pairs $D_{A_h}H_h$ with $\rho_H(a)H_h$; the latter converges strongly in
$L^2$ by Rellich compactness.  The cubic potential force converges weakly
in $L^{4/3}$ by \Cref{lem:critical-cubic}.  Thus all bosonic matter
variations converge against fixed smooth tests, even though their metric
variations may retain quadratic defects.
The weak-$H^1$ fermion argument of \Cref{prop:weak-fermion} applies with
$H_h\to H$ in $L^2$ and a uniform $L^4$ bound, as established there.
This passes every fermionic matter and metric variation.  The uniform
$C^1$ test bound and a finite-net argument give convergence on the
$C^{r_0}$ test unit ball.  Consistency and stationarity now imply the
limiting matter Euler identities by \Cref{thm:abstract-closure} restricted
to matter tests.
\end{proof}

\subsection{Positive covariance measures}

\begin{lemma}[Covariance measure of a weak quadratic packet]
\label{lem:positive-packet-defect}
Let $K$ be a compact smooth domain, including its boundary, and let
$Y_h\rightharpoonup Y$ in $L^2(K;E)$ for a fixed finite-rank real bundle with
a smooth positive fiber metric.  After extraction,
\begin{equation}
 Y_h\otimes Y_h\,\dV_0\rightharpoonup^*
 Y\otimes Y\,\dV_0+\mathsf Q_Y,
 \qquad
 |Y_h|^2\dV_0\rightharpoonup^*|Y|^2\dV_0+\mu_Y,
 \label{eq:positive-packet-defect}
\end{equation}
where $\mathsf Q_Y$ is a positive-semidefinite symmetric matrix-valued
Radon measure and $\mu_Y=\operatorname{tr}\mathsf Q_Y\ge0$.
The following conditions are equivalent along this subsequence:
\begin{equation}
 \mu_Y=0
 \quad\Longleftrightarrow\quad \mathsf Q_Y=0
 \quad\Longleftrightarrow\quad Y_h\to Y\ \text{in }L^2(K;E).
 \label{eq:positive-defect-equivalence}
\end{equation}
If $B_h\to B$ uniformly are continuous bilinear coefficient fields, then
\begin{equation}
 B_h(Y_h,Y_h)\,\dV_0\rightharpoonup^*
 B(Y,Y)\,\dV_0+B:\mathsf Q_Y,
 \qquad |B:\mathsf Q_Y|\le C_B\mu_Y.
 \label{eq:defect-contraction}
\end{equation}
Complex packets are treated by realification.
\end{lemma}
\begin{proof}
The quadratic entries are bounded in $L^1$, so extract a common weak-star
measure limit in a finite atlas.  For a continuous real dual section $a$
and a nonnegative continuous function $\varphi$, multiplication by
$\sqrt\varphi\,a$ preserves weak $L^2$ convergence.  Lower semicontinuity
therefore gives
\[
 \int_K\varphi\,a^T\mathsf Q_Ya\ge0.
\]
This proves positivity; taking the trace gives $\mu_Y\ge0$.
Writing the matrix measure relative to its trace, its density is
positive semidefinite with trace one almost everywhere.  In particular
its Hilbert--Schmidt norm is at most one.  This proves the domination in
\eqref{eq:defect-contraction}, and shows that trace zero is equivalent to
matrix-measure zero.  If $\mu_Y=0$, testing the scalar measure convergence
with the constant function one gives convergence of squared norms.
Weak convergence and the Hilbert identity then imply strong convergence.
The converse follows from the product estimate \eqref{eq:quadratic-product}.
Finally, the uniform coefficient error has total variation at most
$\norm{B_h-B}_\infty\sup_h\norm{Y_h}_2^2$, which tends to zero; the
fixed continuous coefficient $B$ passes through weak-star convergence.
\end{proof}

\subsection{Equivariant tests and critical gauge compactness}

\begin{lemma}[Equivariant physical tests]
\label{lem:equivariant-tests}
For every smooth internal gauge transformation $u$,
\begin{equation}
 \mathfrak n_{u\cdot z}(u\cdot v)=\mathfrak n_z(v).
 \label{eq:equivariant-test-invariance}
\end{equation}
On every bounded weak packet for which
$e,e^{-1}\in L^\infty$, $e\in H^1$, $A\in L^4$, $F_A,D_AH\in L^2$,
$H\in H^1$, and $\Psi,\bPsi\in H^1$, the complete minimally coupled first
variation satisfies
\begin{equation}
 \abs{D\Sact_\theta(z)[v]}\le C_K\mathfrak n_z(v).
 \label{eq:equivariant-firstvariation-bound}
\end{equation}
\end{lemma}
\begin{proof}
Internal gauge transformations act unitarily on every matter representation,
and the covariant derivatives of the transformed tests transform
homogeneously.  This proves \eqref{eq:equivariant-test-invariance}.  For the
second assertion, the Yang--Mills row pairs $F_A$ with $D_Aa$; the Higgs
kinetic row pairs $D_AH$ with
$D_A\eta_H+\rho_H(a)H$; and the quartic force is integrable by
$H^1\hookrightarrow L^4$ in four dimensions.  The Dirac terms are controlled
by one $L^2$ spinor derivative and bounded test coefficients.  The metric
variation of the spin connection has the form
$P_1(e)k\,\partial e+P_2(e)\partial k$ and therefore lies in $L^2$; it pairs
with the $L^2$ spinor bilinear.  Finally the first-order Palatini
representative of \Cref{app:palatini} contains only bounded coframe
coefficients, first coframe derivatives in $L^2$, and $k,\partial k$.
Summing the sector estimates proves \eqref{eq:equivariant-firstvariation-bound}.
\end{proof}

\begin{proposition}[Critical Uhlenbeck quotient compactness]
\label{prop:critical-uhlenbeck}
Let $K'\Subset K$ and suppose
\begin{equation}
 \sup_h\norm{F_{A_h}}_{L^2(K)}<\infty
 \label{eq:critical-curvature-bound}
\end{equation}
and \Cref{def:critical-curvature-ui} holds.  Then $K'$ has a finite ball cover
and cutoff-dependent local internal gauges in which
\begin{align}
 d^*A_h&=0,
 &\sup_h\norm{A_h}_{W^{1,2}}&<\infty,
 \label{eq:critical-coulomb-bound}
\end{align}
with the $L^4$ norms uniformly as small as desired after refining the cover.
After extraction,
\begin{align}
 A_h&\rightharpoonup A &&\text{in }W^{1,2},\notag\\
 A_h&\to A &&\text{in }L^q\quad(1\le q<4),\notag\\
 F_{A_h}&\rightharpoonup F_A &&\text{in }L^2.
 \label{eq:critical-coulomb-convergence}
\end{align}
\end{proposition}
\begin{proof}
Uniform integrability gives a radius, independent of the cutoff and the
center, on which every curvature $L^2$ norm lies below the four-dimensional
Uhlenbeck threshold.  The local Coulomb theorem then yields
\[
 \norm{A_h}_{W^{1,2}}+\norm{A_h}_{L^4}
 \le C\norm{F_{A_h}}_{L^2}
\]
on the corresponding smaller balls.  Weak $W^{1,2}$ compactness and Rellich
compactness give the first two limits.  Since $A_h\to A$ strongly in $L^2$,
$A_h\wedge A_h\to A\wedge A$ in $L^1$, while $dA_h\rightharpoonup dA$ in
$L^2$; hence the weak curvature limit is $F_A$.  The smallness of the local
$L^4$ norm follows from the same small-energy estimate.
\end{proof}

\begin{lemma}[Positive covariant graph control on a normalized chart]
\label{lem:critical-positive-graph}
Let $B$ be a fixed bounded smooth four-dimensional chart, with a smooth
positive-metric reference connection $\nabla^{\rm ref}$. Suppose a unitary
internal representation has $\norm{A}_4\le M_A$. For every smooth section $u$,
\begin{align}
 \norm u_4
 &\le C_B\bigl(\norm u_2+
                  \norm{\nabla^{{\rm ref},A}u}_2\bigr),
       \label{eq:critical-Kato-Sobolev}\\
 \norm u_{H^1}
 &\le C_B(1+M_A)\bigl(\norm u_2+
                  \norm{\nabla^{{\rm ref},A}u}_2\bigr).
       \label{eq:critical-vector-graph}
\end{align}
The same estimates hold for the dual representation. The constants do not
depend on the gauge transformation used to produce the displayed $A$.
\end{lemma}
\begin{proof}
Apply the scalar Sobolev inequality to
$(|u|^2+\epsilon^2)^{1/2}$ and then let $\epsilon\downarrow0$.
Metric compatibility and unitarity give the Kato inequality
$|d|u||\le|\nabla^{{\rm ref},A}u|$, proving
\eqref{eq:critical-Kato-Sobolev}. This scalar inequality does not by itself
control the full derivative of $u$. For that derivative use the exact identity
\[
 \nabla^{\rm ref}u=\nabla^{{\rm ref},A}u-\rho(A)u,
 \qquad
 \norm{\nabla^{\rm ref}u}_2
 \le\norm{\nabla^{{\rm ref},A}u}_2+C\norm A_4\norm u_4.
\]
Together with the first estimate and the fixed smooth reference coefficients,
this proves \eqref{eq:critical-vector-graph}. No smallness of $M_A$ is needed.
\end{proof}

\subsection{Full proof of critical gauge-quotient closure}
\label{app:critical-quotient-proof}

\begin{proof}[Complete proof of \Cref{thm:critical-quotient-defect}]
Extract the geometric and parameter subsequence and apply
\Cref{prop:critical-uhlenbeck}.  On each sufficiently small Coulomb ball the
connection has uniformly small $L^4$ norm.  The identities
\[
 dH_h=D_{A_h}H_h-\rho_H(A_h)H_h
\]
and the four-dimensional Sobolev inequality give
\[
 \norm{H_h}_4
 \le C(\norm{D_{A_h}H_h}_2+\norm{H_h}_2)
   +C\norm{A_h}_4\norm{H_h}_4.
\]
The final term is absorbed by the local small-energy choice, giving an
ordinary $H^1$ bound for $H_h$.  Rellich compactness gives strong $L^q$ for
$q<4$ and almost-everywhere convergence.  The weak limit of $D_{A_h}H_h$ is
identified with $D_AH$ distributionally because $A_h\to A$ and
$H_h\to H$ strongly in $L^2$.

Apply \Cref{lem:critical-positive-graph} to
\eqref{eq:critical-spinor-bound} in the normalized Coulomb chart. First Kato
and scalar Sobolev give uniform $L^4$ spinor bounds; then subtracting
$\rho(A_h)\Psi_h$ from the covariant derivative gives the full ordinary
$H^1$ bound. Do the same in the dual representation. After extraction both
spinors converge strongly in every $L^q$, $q<4$, and their bilinears converge
weakly in $L^2$. This uses the positive internal comparison connection, not
an elliptic estimate for the physical Lorentzian Dirac operator.

Fix now a smooth test in a normalized chart and pull it back through the
cutoff-dependent gauge.  Equation \eqref{eq:equivariant-test-invariance}
keeps its test size uniformly bounded, so the finite first variation tends to
zero without differentiating the gauge.  In the Yang--Mills row,
$D_{A_h}a\to D_Aa$ strongly in $L^2$ for fixed smooth $a$ and therefore pairs
with the weak $L^2$ curvature.  In the Higgs row,
$D_{A_h}\eta_H\to D_A\eta_H$ strongly in $L^2$, while the gauge-current term
pairs the weak $L^2$ covariant gradient with the strongly convergent
undifferentiated Higgs field.  The cubic potential force is bounded in
$L^{4/3}$ and converges almost everywhere, hence has the classical weak limit.
For the Dirac rows, derivative terms use weak--strong pairing. For example,
$A_h\bar\Psi_h\Psi_h\to A\bar\Psi\Psi$ and
$H_h\bar\Psi_h\Psi_h\to H\bar\Psi\Psi$ strongly in $L^1$, since all three
factors converge strongly in $L^3$. Terms with a smooth matter test require
no larger exponents. Spin-connection terms pair a strongly convergent $L^2$
coefficient with a weakly convergent $L^2$ bilinear. Consequently the complete
fermionic metric variation passes by the relaxed strong-$L^2$ coefficient
version of \Cref{prop:weak-fermion}.

Only the quadratic Yang--Mills/Higgs metric terms and the critical quartic
potential can retain unresolved products.  Their limits are precisely the
positive covariance and quartic defects of
\Cref{prop:YM-defect,thm:higgs-defect,lem:positive-packet-defect}.  The coframes converge uniformly to their continuous representatives and
strongly in $H^1$. Thus the uniform algebraic metric-coefficient hypotheses
of the stress-defect propositions hold, while strong first-jet convergence
passes the first-order gravitational variation. This proves
\eqref{eq:critical-quotient-einstein}. No conservation of a measure-valued
defect is inferred from this step; such a conclusion would require the
additional connection convergence and Noether hypotheses of
\Cref{lem:defect-conservation-transport}.  Finally, uniform
integrability makes the quadratic covariance measures absolutely continuous;
uniform integrability of $|H_h|^4$ together with convergence in measure gives
strong $L^4$ convergence by Vitali, so the quartic concentration measure
vanishes.
\end{proof}

\subsection{Concentration, oscillation and an alternative Euclidean branch}

\begin{remark}[Zero defect versus the stronger Higgs higher-integrability endpoint]
\label{rem:Higgs-higher-integrability-boundary}
The reduced closure theorem and the exact zero-defect characterization require
strong $L^2$ compactness of $D_{A_h}H_h$, together with the connection
compactness used in \Cref{prop:covariant-higgs-endpoint}; they do not require
a uniform $L^{8/3}$ bound on the Higgs covariant gradient.  The latter is a
strictly stronger reserve.  For the explicit finite Higgs edge packet,
\Cref{prop:positive-Higgs-envelope} identifies it exactly with an
$L^{4/3}$ bound on a positive companion kinetic density.  The same proposition
shows why ordinary quadratic screening cannot be silently upgraded to this
endpoint.
\end{remark}


\begin{example}[A critical bubble]
\label{ex:bubble}
On a four-dimensional chart choose $\varphi\in C_c^\infty(\R^4)$, a point $x_0$, and a unit Higgs vector $v$. Set
\[
 H_\epsilon(x)=\epsilon^{-1}\varphi((x-x_0)/\epsilon)v.
\]
Then $\norm{H_\epsilon}_2=O(\epsilon)$, while $\norm{\dd H_\epsilon}_2$ and $\norm{H_\epsilon}_4$ are independent of $\epsilon$. The fields converge weakly in $H^1$ to zero, but
\[
 \abs{H_\epsilon}^4\dd x\rightharpoonup^*
 \left(\int_{\R^4}\abs\varphi^4\dd x\right)\delta_{x_0}.
\]
These assertions follow directly by $x=x_0+\epsilon y$. This example disproves a compactness inference from bounded critical energy; it does not assert that this bubble is a coupled stationary solution.
\end{example}

\begin{example}[Vanishing matter residual with conserved kinetic defect]
\label{ex:null-wave}
On flat $(0,T)\times\mathbb T^3$, take $v_H=0$, $A=0$, $\Psi=0$, a fixed unit Higgs vector $v$, and
\[
 H_n(t,x)=n^{-1}\sin(n(t+x^1))v,
 \qquad \ell=\dd t+\dd x^1.
\]
Since $\ell$ is null, $\Box H_n=0$. The Higgs Euler residual is the cubic potential force and is $O(n^{-3})$ uniformly. The gauge current vanishes because the amplitude is real and has a fixed internal direction. Although $H_n\to0$ strongly in $L^4$,
\[
 T^H_n
 =2\cos^2(n(t+x^1))\ell\otimes\ell
   -\lambda_Hn^{-4}\sin^4(n(t+x^1))g
 \rightharpoonup\ell\otimes\ell.
\]
Thus $\nu_H=0$ but the kinetic defect is nonzero and conserved. The example concerns the matter system on a fixed background: its nonzero metric residual is exactly what prevents it from being a vacuum Einstein limit. It isolates the error in trying to derive stress convergence from matter stationarity or Bianchi conservation alone.
\end{example}

\begin{proposition}[A sufficient Euclidean defect-removal mechanism]
\label{prop:monotonicity}
Let $V\hookrightarrow L^4$ be a real Hilbert space, and let the complete gauge-fixed Euclidean Higgs operators $\mathcal A_h:V\to V'$ satisfy
\[
 \ip{\mathcal A_h(u)-\mathcal A_h(v)}{u-v}
 \ge c_1\norm{u-v}_V^2+c_2\norm{u-v}_4^4,
 \qquad c_1,c_2>0.
\]
If $H_h\rightharpoonup H$ in $V$, $\mathcal A_h(H_h)=J_h\to J$ in $V'$, and $\mathcal A_h(H)\to J$ in $V'$, then $H_h\to H$ in $V\cap L^4$. If $V$ controls the covariant $H^1$ norm, both Higgs defects vanish.
\end{proposition}
\begin{proof}
Use $(u,v)=(H_h,H)$. The right pairing is
$\ip{J_h-\mathcal A_h(H)}{H_h-H}$, which tends to zero by strong dual convergence and weak boundedness. Each nonnegative term on the left tends to zero.
\end{proof}

This is an alternative sufficient branch, not an assertion about the Lorentzian evolution operator. A positive quartic term alone does not establish the complete monotonicity estimate, particularly with a symmetry-breaking mass term or unfixed gauge directions.

\subsection{Nonminimal Higgs curvature and the assembled improvement}
\label{app:improvement}

The principal reduced theorem concerns the minimally coupled classical action. A retained nonminimal source must be added before taking variations. For the convention
\begin{equation}
 \Sact_\xi=\xi\int f\Scal(g)\,\dV_g,
 \qquad f=H^\dagger H,
 \label{eq:xi-action}
\end{equation}
the compactly supported metric variation is
\begin{equation}
 D_g\Sact_\xi[k]
 =\xi\int\bigl[f\Eins_{\mu\nu}
       +(g_{\mu\nu}\Box_g-\nabla_\mu\nabla_\nu)f\bigr]
       k^{\mu\nu}\,\dV_g.
 \label{eq:improvement-variation}
\end{equation}
Consequently the stress contribution is minus twice the tensor in brackets multiplied by $\xi$. Formula \eqref{eq:improvement-variation} follows by varying $\sqrt{|g|}fR$ and integrating the two derivative terms in $\delta R$ twice. Compact support is essential to omitting the boundary contribution; with a physical boundary the corresponding scalar--GHY term must be included.

\begin{proposition}[An assembled strong--weak improvement criterion]
\label{prop:improvement}
Suppose $g_h\to g$ in $W^{1,\infty}$, $g_h\rightharpoonup g$ in $W^{2,p}$ for $1<p<\infty$, with uniform inverse bounds and fixed signature, and $H_h\to H$ in $L^{2p'}$. Then
\begin{equation}
 f_h\Eins(g_h)+(g_h\Box_{g_h}-\nabla^{h}\nabla^{h})f_h
 \longrightarrow
 f\Eins(g)+(g\Box_g-\nabla\nabla)f
 \label{eq:improvement-limit}
\end{equation}
in distributions, with the volume densities included in the pairing. If the minimal stress also converges and $\xi_h\to\xi$, the full improved metric variation converges.
\end{proposition}
\begin{proof}
Strong $L^{2p'}$ convergence gives $f_h\to f$ in $L^{p'}$. Curvature is linear in second metric derivatives with strongly convergent bounded coefficients, plus quadratic first-derivative terms. Hence $\Eins(g_h)\rightharpoonup\Eins(g)$ in $L^p$. This passes the first product. For the Hessian contribution, test the \emph{assembled} expression and integrate both derivatives onto the smooth test and metric-density coefficients. The resulting coefficient has a weak $L^p$ limit: its second metric derivatives occur linearly, and every lower-order coefficient converges strongly. Pairing with strong $f_h$ in $L^{p'}$ passes the entire expression. This argument does not separately multiply uncontrolled distributions.
\end{proof}

At the critical $p=2$ threshold the Higgs input is strong $L^4$, not merely a bounded $H^1$ sequence. If the assembled improvement has a subsequential distributional limit but not the value in \eqref{eq:improvement-limit}, write its defect as $\mathfrak I$. Then the Higgs stress defect becomes
\begin{equation}
 \mathfrak S_H=\mathfrak K_H-\lambda_Hg\nu_H-2\xi\mathfrak I.
 \label{eq:improved-defect}
\end{equation}
Existence of the assembled limit is an additional hypothesis; a critical Higgs norm alone does not supply it.

For completeness, if the scalar-curvature coefficient is
$F(H)=1/(2\kappa)+\xi f$ and is uniformly positive, the conformal change $\wt g=(F/\chi_*)g$ gives a constant Einstein coefficient $\chi_*>0$. For a real scalar-coordinate kinetic metric $G^0_{AB}$, the transformed kinetic metric is
\[
 \wt G_{AB}=\frac{\chi_*}{F}G^0_{AB}
             +\frac{3\chi_*}{F^2}\partial_AF\,\partial_BF.
\]
This identity follows from the four-dimensional conformal scalar-curvature formula and one integration by parts. With a boundary, the derivative term must be combined with the transformed GHY contribution. The change of frame reorganizes the principal coupling but supplies neither an effective-Planck lower bound nor hyperbolic smoothing. It therefore does not remove the compactness conditions of the main theorem.

\section{The local finite action and comparison reconstructions}
\label{app:local-realization}

We specify the local link action before its critical and quantitative
estimates. The final subsection gives a distinct smooth comparison
reconstruction. Its sampling statements are not substituted for consistency
on arbitrary finite-energy link records.

\subsection{The explicit local finite action and its quantitative estimates}
\label{app:native-action}

This subsection specifies the finite action used in \Cref{sec:native-branch}.  It is the local common-action representative built from the finite gravitational and Standard-Model link data, not a least-squares functional or a discretization chosen after the continuum equations are known.  The consistency and leakage estimates below analyze this fixed action without adding counterterms or filtering its physical variables.  All contractions are based at one lattice vertex.  Let $e^a_\mu$ be an oriented invertible coframe and $q^a_{\lambda\nu}$ an independent first-jet coordinate.  Define
\begin{align}
 G_{\lambda,\mu\nu}(e,q)
 &=\eta_{ab}(q^a_{\lambda\mu}e^b_\nu+e^a_\mu q^b_{\lambda\nu}),\notag\\
 \Gamma^\rho_{\mu\nu}(e,q)
 &=\tfrac12g^{\rho\sigma}(G_{\mu,\nu\sigma}+G_{\nu,\mu\sigma}-G_{\sigma,\mu\nu}),\notag\\
 \Omega_\mu(e,q)^a{}_b
 &=e^a_\rho\Gamma^\rho_{\mu\nu}(e,q)e_b^\nu-q^a_{\mu\nu}e_b^\nu.
 \label{eq:native-LC-reader}
\end{align}
At $q=\partial e$ this is the Levi--Civita connection.  Put
\[
 \omega_{\mu,h}(x)=\Omega_\mu(e(x),\delta_h^+e(x)),
 \qquad L_\mu(x)=\exp(h\omega_{\mu,h}(x)),
 \qquad U_\mu(x)=\exp(hA_\mu(x)).
\]
For $\mu<\nu$ define the logarithmic plaquettes
\begin{align}
 F^h_{\mu\nu}(x)&=h^{-2}\log\!\bigl(U_\mu(x)U_\nu(x+h e_\mu)
 U_\mu(x+h e_\nu)^{-1}U_\nu(x)^{-1}\bigr),\notag\\
 R^h_{\mu\nu}(x)&=h^{-2}\log\!\bigl(L_\mu(x)L_\nu(x+h e_\mu)
 L_\mu(x+h e_\nu)^{-1}L_\nu(x)^{-1}\bigr).
 \label{eq:native-plaquettes}
\end{align}
The logarithm is the analytic branch near the identity.  Let
\begin{equation}
 K^h_\mu=\frac{\rho_H(U_\mu)H(x+h e_\mu)-H(x)}h,
 \qquad
 V_\mu=\exp(h\sigma(\omega_{\mu,h}))\rho_S(U_\mu),
 \label{eq:native-matter-links}
\end{equation}
and
\begin{equation}
 \nabla^h_\mu\Psi=\frac{V_\mu\Psi(x+h e_\mu)-\Psi(x)}h,
 \qquad
 \nabla^{h,\vee}_\mu\bPsi
 =\frac{\bPsi(x+h e_\mu)V_\mu^{-1}-\bPsi(x)}h.
 \label{eq:native-spin-differences}
\end{equation}
With $v(e)=\sqrt{-\det g}$ and the same coefficient normalization as \eqref{eq:total-action}, set
\begin{align}
 \mathcal L_{{\rm g},h}
 &=\frac{v(e)}{2\kappa}e_a^\mu e^{b\nu}(R^h_{\mu\nu})^a{}_b
    -\frac{\Lambda}{\kappa}v(e),\notag\\
 \mathcal L_{{\rm YM},h}
 &=-\tfrac14v(e)g^{\mu\rho}g^{\nu\sigma}
      \ip{F^h_{\mu\nu}}{F^h_{\rho\sigma}}_{\boldsymbol g},\notag\\
 \mathcal L_{{\rm H},h}
 &=-v(e)g^{\mu\nu}\Rea\ip{K^h_\mu}{K^h_\nu}-v(e)V(H),\notag\\
 \mathcal L_{{\rm D},h}
 &=v(e)\Rea\left\{\tfrac{\ii}{2}
 [\bPsi\gamma^\mu(e)\nabla^h_\mu\Psi
 -(\nabla^{h,\vee}_\mu\bPsi)\gamma^\mu(e)\Psi]
 -\bPsi\mathscr M_{\boldsymbol Y}(H)\Psi\right\}.
 \label{eq:native-densities}
\end{align}
The finite local common action is
\begin{equation}
 S_h^{\rm loc}=h^4\sum_{x\in\mathscr Q_h}
 (\mathcal L_{{\rm g},h}+\mathcal L_{{\rm YM},h}
 +\mathcal L_{{\rm H},h}+\mathcal L_{{\rm D},h}).
 \label{eq:native-local-action}
\end{equation}
It has all coframe components and temporal gauge links as variation variables; temporal gauge is imposed only on the evaluated records.  The Cartan and spin links share the same coframe-derived connection.

\subsubsection{A positive companion Higgs kinetic writer}
\label{app:positive-Higgs-envelope}

The Lorentzian Higgs density in \eqref{eq:native-densities} is signed.  For
higher-integrability diagnostics it is useful to retain, on the \emph{same}
covariant Higgs edge packet $K_h^H=(K^h_\mu)_\mu$, a positive companion
writer.  This auxiliary writer does not alter the physical Lorentzian action.
Let $r_h^{\mu\nu}(x)$ and $m_h(x)>0$ be the positive one-form metric and cell
mass obtained by sampling the fixed auxiliary Riemannian comparison data
$(r,\dV_0)$.  On a compact nondegenerate chart assume the usual uniform
comparison
\begin{equation}
 c_r|\xi|^2\le r_h^{\mu\nu}\Rea\ip{\xi_\mu}{\xi_\nu}
       \le C_r|\xi|^2,
 \qquad 0<c_m\le m_h/h^4\le C_m<\infty.
 \label{eq:positive-Higgs-Hodge-comparison}
\end{equation}
Define the cellwise positive density
\begin{equation}
 e_h^{H,+}(x)
 :=r_h^{\mu\nu}(x)\Rea\ip{K^h_\mu(x)}{K^h_\nu(x)},
 \label{eq:positive-Higgs-density}
\end{equation}
and the discrete positive norms
\begin{equation}
 \|f_h\|_{L_h^p}^p:=\sum_xm_h(x)|f_h(x)|^p.
 \label{eq:positive-discrete-Lp}
\end{equation}

\begin{proposition}[Positive kinetic envelope and the $L^{8/3}$ endpoint]
\label{prop:positive-Higgs-envelope}
Under \eqref{eq:positive-Higgs-Hodge-comparison} there are cutoff-independent
constants $0<c_0\le C_0<\infty$ such that
\begin{equation}
 c_0|K_h^H(x)|^2\le e_h^{H,+}(x)\le C_0|K_h^H(x)|^2
 \label{eq:positive-Higgs-envelope}
\end{equation}
for every interior cell.  Consequently
\begin{equation}
 \boxed{
 \|e_h^{H,+}\|_{L_h^{4/3}}^{1/2}
 \asymp
 \|K_h^H\|_{L_h^{8/3}}.}
 \label{eq:positive-Higgs-L43-L83}
\end{equation}
On smooth bounded reconstructed families, the consistency estimate of
\Cref{prop:mesh-consistency} for $K_h^H$ identifies this criterion with the
corresponding local continuum bound for $D_{A_h}H_h$, up to the already
controlled reconstruction error.

If $K_h^H\to K^H$ strongly in $L_h^2$ under a common reconstruction, then
$e_h^{H,+}\to e^{H,+}$ strongly in $L^1_{\rm loc}$ after reconstruction.
This conclusion alone does not imply a cutoff-uniform $L^{4/3}$ bound for
$e_h^{H,+}$ or an $L^{8/3}$ bound for $K_h^H$.
\end{proposition}
\begin{proof}
The pointwise comparison \eqref{eq:positive-Higgs-envelope} is exactly the
uniform positive definiteness of the sampled comparison Hodge form.  Raising
it to the power $4/3$, multiplying by $m_h$, and summing gives
\[
 \|e_h^{H,+}\|_{L_h^{4/3}}^{4/3}
 \asymp
 \sum_xm_h|K_h^H|^{8/3}
 =\|K_h^H\|_{L_h^{8/3}}^{8/3},
\]
which is equivalent to \eqref{eq:positive-Higgs-L43-L83}.  The smooth-family
continuum statement follows from $K_h^H=D_AH+o(1)$ in the corresponding
local norm and the uniform equivalence of the positive cell masses.

For the $L^1$ statement, use
\[
 \bigl||u|^2-|v|^2\bigr|
 \le |u-v|(|u|+|v|)
\]
and Cauchy--Schwarz, together with the uniform coefficient convergence of
$r_h$.  The failure of an automatic $L^{4/3}$ upgrade is genuine.  In a
four-dimensional coordinate ball choose a nonzero
$\varphi\in C_c^\infty(B_1)$ and set
\[
 Y_n(x)=n^{7/4}\varphi(n(x-x_0))v
\]
for a fixed unit fiber vector $v$.  Then
$\|Y_n\|_2=n^{-1/4}\|\varphi\|_2\to0$, while
$\|Y_n\|_{8/3}=n^{1/4}\|\varphi\|_{8/3}\to\infty$.
Thus $|Y_n|^2\to0$ strongly in $L^1$ although
$\||Y_n|^2\|_{4/3}\to\infty$.  Sampling a smooth version of this
concentrating family on resolving meshes gives the same discrete conclusion.
\end{proof}

\begin{corollary}[A positive amplitude--energy compiler]
\label{cor:positive-Higgs-amplitude-criterion}
If, on every compact interior chart,
\begin{equation}
 \sup_h\|e_h^{H,+}\|_{L^1_h}<\infty,
 \qquad
 \sup_h\|e_h^{H,+}\|_{L^\infty_h}<\infty,
 \label{eq:positive-Higgs-amplitude-energy}
\end{equation}
then
\begin{equation}
 \sup_h\|e_h^{H,+}\|_{L_h^{4/3}}<\infty,
 \qquad
 \sup_h\|K_h^H\|_{L_h^{8/3}}<\infty.
 \label{eq:positive-Higgs-compiled-endpoint}
\end{equation}
Any finite microscopic mechanism used to verify
\eqref{eq:positive-Higgs-amplitude-energy} must control this same Higgs kinetic
writer; a bound for an unrelated positive observable is not a substitute.
\end{corollary}
\begin{proof}
Since $e_h^{H,+}\ge0$,
\[
 \int (e_h^{H,+})^{4/3}
 \le \|e_h^{H,+}\|_\infty^{1/3}\|e_h^{H,+}\|_1.
\]
Apply \Cref{prop:positive-Higgs-envelope}.
\end{proof}

\begin{remark}[Role in the present closure theorem]
\label{rem:positive-Higgs-endpoint-scope}
The criterion \eqref{eq:positive-Higgs-compiled-endpoint} is a stronger
higher-integrability card, not a hidden hypothesis of
\Cref{thm:reduced-closure,thm:zero-defect-characterization}.  Those results
close the minimally coupled Higgs stress from strong $L^2$ compactness of the
covariant gradient and the associated $H^1\cap L^4$ convergence of the Higgs
field.  The $L^{8/3}$ reserve becomes relevant only in stronger weak-limit or
compensated-compactness arguments that require additional integrability of
the matter source.  Establishing such a reserve from a particular microscopic
source law is a separate selection problem.
\end{remark}

\subsection{Structural realization and representation packet}
\label{app:structural-packet}

\label{subsec:SM-packet}
\begin{proposition}[Renewal realization of the finite interface]
\label{prop:renewal-interface}
The Renewal Geometry constructions \cite{PelissierRenewalSpacetime2026,PelissierSpectral2026,PelissierDuality2026} realize \Cref{def:finite-interface} on a selected Standard-Model branch.  In particular, they provide a nonempty family of finite interfaces with the structural bundle content and common-action normalization required here.  This proposition is a realization statement only; none of the compactness, first-variation consistency, stationarity, or PDE hypotheses of \Cref{sec:limit} is inferred from those constructions.
\end{proposition}
\begin{proof}
The three cited constructions provide respectively the Lorentzian carrier
and gravitational normalization, the differential/spin realization, and the
Standard-Model bundle, generation factor, Dirac/Yukawa incidence, positive
source quotient and common action. These are precisely
\textup{(I1)}--\textup{(I5)}; the analytic hypotheses are imposed separately.
\end{proof}

\begin{remark}[Structural provenance versus analytic input]
\label{rem:structural-provenance}
In the structural gravity--matter construction \cite{PelissierDuality2026}, items \textup{(I2)}--\textup{(I5)} are organized in part by the double-centralizer relation
\begin{equation}
 \A_{\mathrm{act},h}'=\M_{\mathrm{type},h},
 \qquad \M_{\mathrm{type},h}'=\A_{\mathrm{act},h},
 \label{eq:imported-duality}
\end{equation}
and from the common-source Schur quotient
\begin{equation}
 G_{\mathrm{i}\mid\mathrm g,h}
 =G_{\mathrm i,h}-C_hG_{\mathrm g,h}^{-1}C_h^*
 =S_{\mathrm i,h}^*(I-P_{\mathrm g,h})S_{\mathrm i,h}.
 \label{eq:source-schur}
\end{equation}
These identities explain why the finite matter packet and the physical source normalization are not appended by hand. They are not used to prove spatial or temporal compactness, convergence of $F_{A_h}$ or $D_{A_h}H_h$, spinor $H^1$ convergence, or any nonlinear stress limit. Those are analytic questions addressed independently below.
\end{remark}

To specify the matter carrier independently of its structural origin, use
$Q_{\rm em}=T_3+Y$ and the representations in \Cref{tab:SM-representations}.
They give the usual quark, lepton and Higgs assignments
\cite{ErlerFreitas2024}.  The integer weight of the covering $\U(1)$ factor
is $6Y$; the fractions in the table therefore do not denote multivalued
characters of $\U(1)$.

\begin{table}[htbp]
\centering
\caption{The fixed Standard-Model representation packet.  Every fermion row
carries a generation factor $\C^3$; the Higgs carries no generation factor.
The neutral row is included only on the declared neutrino extension.}
\label{tab:SM-representations}
\begin{tabular}{@{}lllll@{}}
\toprule
Field & $\SU(3)$ & $\SU(2)$ & $Y$ & Spacetime chirality\\
\midrule
$Q_L$ & $\boldsymbol 3$ & $\boldsymbol 2$ & $1/6$ & left\\
$u_R$ & $\boldsymbol 3$ & $\boldsymbol 1$ & $2/3$ & right\\
$d_R$ & $\boldsymbol 3$ & $\boldsymbol 1$ & $-1/3$ & right\\
$L_L$ & $\boldsymbol 1$ & $\boldsymbol 2$ & $-1/2$ & left\\
$e_R$ & $\boldsymbol 1$ & $\boldsymbol 1$ & $-1$ & right\\
$\nu_R$ (optional) & $\boldsymbol 1$ & $\boldsymbol 1$ & $0$ & right\\
$H$ & $\boldsymbol 1$ & $\boldsymbol 2$ & $1/2$ & scalar\\
\bottomrule
\end{tabular}
\end{table}

\begin{lemma}[Descent and covariant Yukawa blocks]
\label{lem:SM-descent}
The representations in \Cref{tab:SM-representations} descend to
$G_{\rm SM}$ in \eqref{eq:gauge-group}.  With
$\widetilde H=\ii\sigma_2\overline H$ and generation matrices
$Y_u,Y_d,Y_e\in M_3(\C)$, the mass bilinear is assembled from
\begin{equation}
 \overline Q_L\widetilde H Y_u u_R+
 \overline Q_L H Y_d d_R+\overline L_L H Y_e e_R
 +\text{Hermitian conjugates}.
 \label{eq:explicit-Yukawa-blocks}
\end{equation}
On the neutral extension one adds
$\overline L_L\widetilde H Y_\nu\nu_R$ and a constant symmetric Majorana
block on the charge-conjugate doubled carrier.  The resulting map
$\mathscr M_{\boldsymbol Y}(H)$ is affine in real Higgs coordinates and
satisfies \eqref{eq:yukawa-bound} on every bounded coefficient bank.
\end{lemma}
\begin{proof}
The covering homomorphism is
\[
 (g_3,g_2,z)\longmapsto(z^{-2}g_3,z^3g_2).
\]
Its kernel is generated by
$(e^{2\pi\ii/3}I_3,-I_2,e^{\pi\ii/3})$.
In each row the product of the color-center character, the weak-center
character and $z^{6Y}$ is one, proving descent.  Each monomial in
\eqref{eq:explicit-Yukawa-blocks} has zero total hypercharge, and its color
and weak indices are contracted invariantly.  The same check applies to
the optional neutrino coupling; the Majorana block is gauge neutral.
The entries are linear in $H$ or $\overline H$, apart from that constant
block.  Finite-dimensional operator bounds give \eqref{eq:yukawa-bound}.
\end{proof}

No value of a generation matrix, gauge coupling or gravitational coefficient
is inferred from this representation statement.  Omitting the neutral row
and its blocks gives the minimal packet; retaining them gives the declared
neutral extension throughout the same analytic argument.

\subsection{Concrete reconstruction and smooth first-variation consistency}
\label{app:reconstruction}

This smooth comparison regulator verifies the reconstruction and
first-variation consistency hypotheses of \Cref{def:regulator}, with
$c_h(K)=O(h)$ on bounded smooth families. It is distinct from the rough-link
consistency analysis and does not rederive the structural interface.

\subsubsection{Whitney maps and a signed Hodge realization}

On a contractible coordinate region choose a quasiuniform shape-regular simplicial mesh of size $h$. Write $C_h^k$ for oriented $k$-cochains. For an oriented simplex $\sigma=[i_0,\ldots,i_k]$ with barycentric coordinates $\lambda_i$, define the Whitney form
\begin{equation}
 w_\sigma=k!\sum_{j=0}^k(-1)^j\lambda_{i_j}
  \dd\lambda_{i_0}\wedge\cdots\wedge
  \widehat{\dd\lambda_{i_j}}\wedge\cdots\wedge\dd\lambda_{i_k},
 \qquad W_hc=\sum_\sigma c_\sigma w_\sigma.
 \label{eq:Whitney}
\end{equation}
The de Rham map $R_h\alpha=(\int_\sigma\alpha)_\sigma$ is used for smooth test forms. It is not asserted to be a bounded point- or trace-evaluation map on arbitrary $L^2$ forms.

\begin{lemma}[Commuting reconstruction and positive comparison norm]
\label{lem:Whitney}
The maps satisfy $R_hW_h=I$ and $\dd W_h=W_h\dd_h$. If
\[
 \norm c_h^2:=\int_K\abs{W_hc}_r^2\,\dV_0,
\]
then $W_h$ is an isometry onto its finite-element range. The reference mass matrix is positive definite. Shape regularity gives the usual cellwise scaling constants uniformly in $h$. Orthogonal projections $P_h^k$ onto these ranges converge strongly to the identity on $L^2$ $k$-forms.
\end{lemma}
\begin{proof}
Integration of \eqref{eq:Whitney} over oriented simplices gives the Kronecker pairing with the cochain basis. Differentiating the barycentric formula and using $\sum_i\lambda_i=1$ gives the coboundary identity; equivalently, integrate it and use Stokes on each simplex. The norm statement is its definition, and injectivity follows from $R_hW_h=I$. Scaling to a reference simplex and finite-dimensional norm equivalence give the uniform local mass bounds. Smooth forms are approximated by their Whitney interpolants in $L^2$; this follows by Taylor expansion on each cell and shape-regular scaling. Density of smooth forms and $\norm{P_h^k}\le1$ give strong convergence of the projections.
\end{proof}

This is the standard untwisted de Rham structure \cite{ArnoldFalkWinther2006,ArnoldFalkWinther2010}. Gauge curvature is treated separately: for a nonflat connection $D_A^2$ is curvature, and the ordinary commuting-complex identity is not a claim that $D_A^2=0$.

Let $g_h\to g$ uniformly with fixed signature and uniform inverse bounds. Using the fixed positive reference projections, define a discrete signed Hodge map by
\begin{equation}
 W_h^{4-k}\star_h^k
 =P_h^{4-k}\star_{g_h}W_h^k.
 \label{eq:Hodge-definition}
\end{equation}
For $W_h^kc_h\to\alpha$ in $L^2$,
\begin{equation}
 W_h^{4-k}\star_h^kc_h\longrightarrow\star_g\alpha\quad\text{in }L^2.
 \label{eq:Hodge-convergence}
\end{equation}
Indeed, subtract $\star_g\alpha$, use uniform boundedness of $\star_{g_h}$, uniform convergence of its algebraic coefficients, and strong convergence of $P_h^{4-k}$. The same proof applies to a fixed smooth metric variation because $\dot\star_g[k]$ is algebraic in $g,g^{-1},k$. The projections here are fixed by the positive reference metric, so no omitted derivative of a metric-dependent projection occurs.

This is a concrete signed-Hodge realization. It neither guarantees strong convergence of $W_hc_h$ nor excludes concentration of its quadratic products. Those are independent compactness requirements on the low-regularity branch; the regular a-priori branch of \Cref{thm:regular-branch} derives them from stronger Sobolev control. Exact mass integration can be replaced by quadrature only after its action and first-variation errors have been estimated.

\subsubsection{A finite smooth comparison reconstruction}

The following construction suffices to populate the consistency statement; it is not a claim of a canonical interpolation for every renewal graph. On a regular coordinate mesh let $I_hf$ be the continuous piecewise-affine interpolant of nodal samples. Choose a fixed smooth mollifier $\rho$ supported in the unit ball and put
\begin{equation}
 \mathcal J_hf=\rho_h*I_hf,
 \qquad \rho_h(x)=h^{-4}\rho(x/h).
 \label{eq:smooth-reconstruction}
\end{equation}
Work on a slightly larger region, so this local convolution creates no physical boundary condition. A fixed partition of unity handles a finite chart cover. The reconstruction is finite dimensional because $I_hf$ is determined by finitely many samples.

For $f$ in a bounded $C^3$ set, reference-cell interpolation and convolution give
\begin{equation}
 \norm{\mathcal J_hf-f}_{W^{1,\infty}}\le Ch\norm f_{C^2},
 \qquad \norm{\mathcal J_hf}_{W^{2,\infty}}\le C\norm f_{C^2}.
 \label{eq:smooth-reconstruction-estimate}
\end{equation}
For the second bound, differentiate the convolution of $\partial I_hf-\partial f$, whose $L^\infty$ norm is $O(h)$; the derivative of $\rho_h$ has $L^1$ norm $O(h^{-1})$. The remaining term is the convolution of $\partial^2f$. The same estimates hold for a smooth parameter derivative $\dot f$. Nondegeneracy of a smooth reference coframe is preserved for small $h$ by uniform approximation.

Use this map for coframe and matter comparison fields, with fixed local spin/gauge frames. For a smooth gauge connection let $U_\mu(x)$ denote backward parallel transport from $x+h e_\mu$ to $x$. Define
\begin{align}
 U_{\mu\nu}(x)
 &=U_\mu(x)U_\nu(x+h e_\mu)
          U_\mu(x+h e_\nu)^{-1}U_\nu(x)^{-1},\notag\\
 F^h_{\mu\nu}(x)&=h^{-2}\log U_{\mu\nu}(x),\notag\\
 K^h_\mu(x)&=h^{-1}\bigl(U_\mu(x)H(x+h e_\mu)-H(x)\bigr).
 \label{eq:lattice-comparison}
\end{align}
The logarithm is evaluated on the identity chart, valid uniformly for sufficiently small $h$ on the smooth bounded family. In \eqref{eq:lattice-comparison} the fields may be the reconstructed smooth fields $\mathcal J_hf$.

For spinors use backward spin--gauge transport $\mathcal U_\mu$ and the centered covariant difference
\begin{equation}
 \nabla^h_\mu\Psi(x)=\frac{
 \mathcal U_\mu(x)\Psi(x+h e_\mu)
 -\mathcal U_\mu(x-h e_\mu)^{-1}\Psi(x-h e_\mu)}{2h}.
 \label{eq:spin-difference}
\end{equation}
The dual-spinor difference uses the dual transport. Insert these quantities into the symmetric density \eqref{eq:dirac-density}, the Higgs density, and the Yang--Mills density, and sum with the cell-volume/Hodge coefficients computed from the same coframe. Define the gravitational comparison action by the integrated first-derivative torsion-free density \eqref{eq:palatini-first-order} on the same finite reconstruction. This is a finite-dimensional differentiable action; numerical integration, if used, must satisfy the stated derivative consistency.

\begin{proposition}[Smooth local action consistency]
\label{prop:mesh-consistency}
On a bounded smooth coframe/gauge/matter family with a common inverse-coframe bound, the preceding comparison action and its physical variations satisfy, uniformly for $\norm v_{C^{r_0}}\le1$,
\begin{equation}
 \abs{D\Sact_{b,h}[\mathcal I_hv]
       -D\Sact_{b,\theta_h}(z_h)[v]}\le C_Kh,
 \qquad b\in\{\mathrm g,\mathrm{SM}\}.
 \label{eq:mesh-C1}
\end{equation}
The reconstruction converges in the strong packet. The estimates are local smooth consistency estimates; they do not assert stability on arbitrary finite-energy link configurations.
\end{proposition}
\begin{proof}
The integral equation for parallel transport and Taylor expansion along an edge give
\[
 U_\mu(x)=I+hA_\mu(x)+O(h^2).
\]
Multiplying the four edge expansions cancels the terms linear in $h$ and leaves
\[
 U_{\mu\nu}(x)=I+h^2\Fcurv_{A,\mu\nu}(x)+O(h^3).
\]
The second-order contribution is $\partial_\mu A_\nu-\partial_\nu A_\mu+[A_\mu,A_\nu]$. Smoothness and the uniform identity chart therefore give $F^h=\Fcurv_A+O(h)$. Covariant Taylor expansion similarly gives $K^h=D_AH+O(h)$ and $\nabla^h\Psi=\nabla^{e,A}\Psi+O(h)$.

These expansions also hold after one field or metric variation. To see this without interchanging an uncontrolled limit and derivative, differentiate the transport equation itself. If $U(t,s)$ is transport along a path and $a$ is the connection variation, its derivative is a single integral of $U(t,u)a(u)U(u,s)$ with the fixed orientation sign. A second Taylor expansion gives the derivative of each displayed remainder, bounded by $Ch$ after the same normalizations. For a metric variation the spin-connection derivative contains only the first jet of the coframe variation; \eqref{eq:smooth-reconstruction-estimate} supplies the needed uniform bounds. The reconstruction of \eqref{eq:metric-lift} differs from the same lift evaluated at the reconstructed metric by $O(h)$ in the norms entering the density.

The pointwise densities and their first variations are finite sums of smooth coefficient functions and products of these quantities. Their difference is $O(h)\norm v_{C^{r_0}}$. A Riemann sum over cells of volume $O(h^4)$ has $O(h^{-4})$ terms on a fixed region, so the integrated error is $O(h)$, not an extensive error. The first-derivative gravitational expression obeys the same estimate directly from \eqref{eq:smooth-reconstruction-estimate}; exact integration eliminates its quadrature error. Differentiated consistent quadrature adds a further $O(h)$ term. This proves \eqref{eq:mesh-C1}.

The reconstructed fields converge in $W^{1,\infty}$ on the smooth family, and their curvature and covariant derivatives converge by the same expansions. This implies every norm in the strong packet on a compact region.
\end{proof}

The link-variable action has its usual finite gauge covariance. The comparison interpolation is only a means of comparing smooth fields and physical variations; arbitrary node-wise gauge transformations need not preserve that finite-dimensional interpolation chart. No exact lattice gauge identity is inferred merely from restriction to a non-invariant interpolation subspace. Equivariant reconstruction and nonlinear stability for arbitrary rough links are additional questions, distinct from the proved smooth estimate. Discretely gauge-invariant simplicial formulations provide another route, subject to their own consistency estimates \cite{ChristiansenHalvorsen2012}.

\begin{remark}[Intrinsic finite action and analytical comparison devices]
The local action combines the Cartan/gauge plaquettes, Higgs links,
spin--gauge transport and calibrated relative normalization. Its structural
provenance is \Cref{prop:renewal-interface}; the explicit representative is
defined here. The auxiliary periodic box, trigonometric reconstruction and
$P_{\le K}$ in the leakage proof are comparison devices. They neither alter
the finite Euler covector nor replace the physical record.
\end{remark}

\section{Critical finite graphs and first-variation consistency}
\label{app:native-critical-closure}

This appendix gives a first-jet proof of consistency for the unchanged local
action of \Cref{app:native-action}.  Its purpose is qualitative closure at the
critical four-dimensional endpoint.  It is independent of the growing-band
estimates used later for quantitative rates.  We work on the periodic
comparison box fixed in \Cref{subsec:local-action-conventions}; all conclusions may be
restricted to an interior physical cylinder.  Bundle components are expressed
in one fixed compatible frame during the compactness argument.  Internal site
gauges are treated intrinsically below.

For a nodal array $u_h$ set
\begin{equation}
 \norm{u_h}_{p,h}^p=h^4\sum_x|u_h(x)|^p,
 \qquad
 D^+_{\mu,h}u_h(x)=\frac{u_h(x+h e_\mu)-u_h(x)}h,
 \label{eq:native-critical-norms}
\end{equation}
with the evident modification for $p=\infty$, and let $\mathsf R_h^0u_h$
denote the piecewise-constant reconstruction.  Write
$T_{\mu,m}u_h(x)=u_h(x+mhe_\mu)$.  These are positive comparison norms;
Lorentzian signs occur only in the physical action.

\begin{lemma}[Discrete translation compactness]
\label{lem:native-discrete-KR}
Let $u_h$ be uniformly bounded in $L_h^2$.  If
\begin{equation}
 \lim_{\rho\downarrow0}\limsup_{h\downarrow0}
 \max_\mu\max_{0<|m|h\le\rho}
 \norm{T_{\mu,m}u_h-u_h}_{2,h}=0,
 \label{eq:native-discrete-KR}
\end{equation}
then $\mathsf R_h^0u_h$ is precompact in $L^2$.
\end{lemma}
\begin{proof}
For $0\le\vartheta\le1$, direct integration on the two portions of each
cell gives
\[
 \norm{\tau_{\vartheta h e_\mu}\mathsf R_h^0u_h-
       \mathsf R_h^0u_h}_2^2
 =\vartheta\norm{T_{\mu,1}u_h-u_h}_{2,h}^2.
\]
An arbitrary displacement is an integer grid displacement plus a subcell
remainder.  Decomposing a four-dimensional displacement into coordinate
parts therefore transfers \eqref{eq:native-discrete-KR} to a uniform
continuum translation modulus for the piecewise-constant reconstructions.
The Fr\'echet--Kolmogorov criterion yields the result.  Equivalently, one may
first mollify at fixed scale, approximate the mollified family by finitely
many Fourier modes, and then remove the mollification using the translation
estimate.
\end{proof}

\subsection{Critical finite graphs for the matter packets}

Let $U_{\mu,h}=e^{hA_{\mu,h}}$ be the actual internal links.  For any retained
unitary representation $\rho$ define
\begin{equation}
 \mathcal D^U_{\mu,h}u_h
 =\frac{\rho(U_{\mu,h})T_{\mu,1}u_h-u_h}{h},
 \qquad
 B^\rho_{\mu,h}=\frac{\rho(U_{\mu,h})-I}{h}.
 \label{eq:native-internal-graph}
\end{equation}
For spinors the positive graph uses only the compact internal representation;
the coframe-derived Lorentz spin transport remains in the physical Dirac
term and is treated separately.

\begin{lemma}[Critical grid bounds and link coefficients]
\label{lem:native-critical-grid}
There is a cutoff-independent constant such that
\begin{equation}
 \norm{u_h}_{4,h}\le C\left(\norm{u_h}_{2,h}
           +\sum_\mu\norm{D^+_{\mu,h}u_h}_{2,h}\right),
 \label{eq:native-grid-Sobolev-a}
\end{equation}
and, for a unitary internal graph,
\begin{equation}
 \norm{u_h}_{4,h}\le C\left(\norm{u_h}_{2,h}
           +\sum_\mu\norm{\mathcal D^U_{\mu,h}u_h}_{2,h}\right).
 \label{eq:native-grid-Sobolev-b}
\end{equation}
A family bounded in the discrete $H^1$ norm has a subsequence for which
$\mathsf R_h^0u_h\to u$ strongly in $L^2$, the forward differences converge
weakly to $\partial u$ in $L^2$, and the trigonometric reconstructions are
weakly compact in $H^1$.  If all forward differences converge strongly,
then the trigonometric reconstructions converge strongly in $H^1$ and
$\mathsf R_h^0u_h\to u$ strongly in $L^4$.

If $\mathsf R_h^0A_h\to A$ strongly in $L^4$, then
\begin{equation}
 h\norm{A_h}_{\infty,h}\to0,
 \qquad
 \mathsf R_h^0B^\rho_h\to\rho(A)
       \quad\text{strongly in }L^4.
 \label{eq:native-link-coeff}
\end{equation}
Consequently a bounded positive graph norm in
\eqref{eq:native-internal-graph}, together with a bounded $L_h^4$ connection,
controls the ordinary discrete $H^1$ norm.
\end{lemma}
\begin{proof}
On a reference four-cube, integration of the fourth power of the multilinear
interpolant is a norm on the finite vector of vertex values.  Scaling and
summation reduce \eqref{eq:native-grid-Sobolev-a} to the continuum embedding
$H^1\hookrightarrow L^4$ in four dimensions.  Pointwise unitarity gives the
discrete Kato inequality
$|D^+_{\mu,h}|u_h||\le|\mathcal D^U_{\mu,h}u_h|$, proving
\eqref{eq:native-grid-Sobolev-b}.  The compactness statements follow from
Rellich compactness, discrete summation by parts and the sine-symbol bound
for trigonometric interpolation.  Strong derivative convergence and a smooth
approximation of the limit give the endpoint $L^4$ convergence.

Strong $L^4$ convergence of the piecewise-constant connections makes their
fourth powers uniformly integrable.  A single cell has volume $h^4$, hence
$h\norm{A_h}_{\infty,h}\to0$.  The exact identity
\[
 B^\rho_h=\int_0^1e^{th\rho(A_h)}\rho(A_h)\,\dd t
\]
and bounded-coefficient multiplication prove the second assertion in
\eqref{eq:native-link-coeff}.  Finally
$D_h^+u_h=\mathcal D_h^Uu_h-B_h^\rho T u_h$ and H\"older give the ordinary
first-difference bound.
\end{proof}


\begin{lemma}[Identification of raw and trigonometric first jets]
\label{lem:native-reconstruction-identification}
On the odd periodic grids of side $L=2\pi$, define
$\norm{u_h}_{1,h}^2=\norm{u_h}_{2,h}^2+
\sum_\mu\norm{D^+_{\mu,h}u_h}_{2,h}^2$.
Then
\begin{equation}
 \norm{\mathcal I_h^{\rm trig}u_h-\mathsf R_h^0u_h}_2
       \le Ch\norm{D_h^+u_h}_{2,h},\qquad
 \norm{\mathcal I_h^{\rm trig}u_h}_{H^1}\le C\norm{u_h}_{1,h}.
 \label{eq:native-reconstruction-comparison}
\end{equation}
If $\mathsf R_h^0u_h\to u$ and
$\mathsf R_h^0D^+_{\mu,h}u_h\to v_\mu$ strongly in $L^2$ for every $\mu$,
then $v_\mu=\partial_\mu u$,
$\mathcal I_h^{\rm trig}u_h\to u$ strongly in $H^1$, and
$\mathsf R_h^0u_h\to u$ strongly in $L^4$.
\end{lemma}
\begin{proof}
Use the normalized discrete Fourier coefficients with physical frequencies
$\xi_\ell=2\pi\ell/L$.  For a fixed subcell displacement $y\in[0,h)^4$,
discrete orthogonality gives
\[
 h^4\sum_x|\mathcal I_h^{\rm trig}u_h(x+y)-u_h(x)|^2
 =L^4\sum_\ell|\widehat u_h(\ell)|^2
                         |e^{\ii\xi_\ell\cdot y}-1|^2.
\]
Average over $y$.  On the represented frequency cube,
$|\xi_\ell|\le C(\sum_\mu
|(e^{\ii h\xi_{\ell,\mu}}-1)/h|^2)^{1/2}$ by the sine bound.
This proves both estimates in \eqref{eq:native-reconstruction-comparison}.
Summation by parts against a fixed sampled smooth test identifies $v_\mu$
as $\partial_\mu u$.

For any strongly convergent raw $L^2$ array, its fixed discrete Fourier
coefficients converge to those of the limit: replace the cellwise constant
Fourier test by its continuum character, whose error tends uniformly to
zero.  Parseval and convergence of the total squared norm then give a
vanishing high-frequency tail after first taking $h\downarrow0$ and then
the frequency threshold to infinity.  Apply this observation to all first
differences.  The ratio of the trigonometric derivative symbol
$\ii\xi_{\ell,\mu}$ to the forward symbol is uniformly bounded and tends
to one at each fixed frequency.  The finite-mode/tail argument therefore
gives strong convergence of the trigonometric derivatives.

Finally choose smooth periodic $u^j\to u$ in $H^1$ and compare $u_h$ with
its sampled values $\mathsf S_hu^j$.  Discrete Sobolev bounds the raw
$L^4$ difference by their discrete $H^1$ difference.  For fixed $j$ its
limsup is at most $C\norm{u-u^j}_{H^1}$.  Let $j\to\infty$ and use
continuous $H^1\hookrightarrow L^4$ for the limit.  This proves the raw
endpoint statement without identifying raw arrays with trigonometric
fields at a finite cutoff.
\end{proof}

\begin{proposition}[Native Higgs graph compactness]
\label{prop:native-Higgs-compactness}
Suppose
$\mathsf R_h^0A_h\to A$ strongly in $L^4$,
$\sup_h\norm{H_h}_{2,h}<\infty$, and the literal Higgs-link packet
$K_{\mu,h}=\mathcal D^U_{\mu,h}H_h$ satisfies
$\mathsf R_h^0K_h\to K$ strongly in $L^2$.  After extraction,
\begin{align}
 \mathsf R_h^0H_h&\to H &&\text{strongly in }L^4,\notag\\
 \mathsf R_h^0D_h^+H_h&\to\dd H &&\text{strongly in }L^2,\notag\\
 \mathcal I_h^{\rm trig}H_h&\to H &&\text{strongly in }H^1,
 \qquad K=D_AH.
 \label{eq:native-Higgs-compactness}
\end{align}
Thus the literal covariant-gradient packet closes the quartic Higgs endpoint
without a separate amplitude compactness or high-derivative hypothesis.
\end{proposition}
\begin{proof}
\Cref{lem:native-critical-grid} gives bounded ordinary discrete $H^1$ norms,
strong $L^2$ compactness of the values, and bounded $L_h^4$ norms.  Since
one-step shifts have the same strong $L^2$ limit, the product
$B_h^{\rho_H}T H_h$ converges strongly in $L^2$ to $\rho_H(A)H$; at the
critical exponent this follows by first approximating $A$ in $L^4$ by a
bounded smooth coefficient.  The exact graph identity
$D_h^+H_h=K_h-B_h^{\rho_H}T H_h$ then gives strong ordinary first-difference
convergence.  Discrete summation by parts identifies the limit derivative,
and the final assertions of \Cref{lem:native-critical-grid} give strong
$H^1$ and $L^4$ convergence.
\end{proof}

For the literal logarithmic plaquette write
\begin{equation}
 F^h_{\mu\nu}(A_h)
 =(D^+_{\mu,h}A_{\nu,h}-D^+_{\nu,h}A_{\mu,h})
       +\mathcal N_{\mu\nu,h}(A_h).
 \label{eq:native-YM-split}
\end{equation}
For the four slots
$W_{\mu\nu,h}(A)=(A_\mu,T_{\mu,1}A_\nu,T_{\nu,1}A_\mu,A_\nu)$,
the analytic logarithm has linear part
$h(A_\mu+T_{\mu,1}A_\nu-T_{\nu,1}A_\mu-A_\nu)$.
Subtracting that part and applying its integral second-order remainder
gives, on a fixed small scaled logarithm chart,
\begin{align}
 |\mathcal N_{\mu\nu,h}(A)(x)|&\le C|W_{\mu\nu,h}(A)(x)|^2,\notag\\
 \norm{\mathcal N_h(A)-\mathcal N_h(B)}_{2,h}
 &\le C(\norm A_{4,h}+\norm B_{4,h})\norm{A-B}_{4,h}.
 \label{eq:native-YM-envelope}
\end{align}
The same analytic expansion holds after one physical variation.

\begin{proposition}[Identification and variation of the literal curvature]
\label{prop:native-YM-identification}
Suppose $\mathsf R_h^0A_h\to A$ strongly in $L^4$ and
$\mathsf R_h^0F_h\to F$ strongly in $L^2$.  Then $F=F_A$.  For each fixed
smooth gauge test $a$,
\begin{equation}
 \mathsf R_h^0\left(D_AF_h[\mathsf S_ha]\right)
 \longrightarrow \dd a+[A,a]
 \quad\text{strongly in }L^2.
 \label{eq:native-YM-variation}
\end{equation}
If the coframes converge strongly in $H^1$ on one compact nondegenerate
chart, the complete finite Yang--Mills first variation, including its metric
variation, converges to the continuum Yang--Mills first variation uniformly
on bounded $C^r$ test sets, $r\ge2$.
\end{proposition}
\begin{proof}
By \eqref{eq:native-link-coeff}, the scaled slots eventually lie in the
analytic plaquette domain.  One-step shifts of $A_h$ have the same strong
$L^4$ limit.  The quadratic plaquette envelope therefore gives
$\mathsf R_h^0\mathcal N_h(A_h)\to A\wedge A$ strongly in $L^2$.
Testing \eqref{eq:native-YM-split} against a sampled smooth function and
summing by parts identifies $F=\dd A+A\wedge A$.  Differentiating the exact
split in the nodal direction $\mathsf S_ha$ yields
\eqref{eq:native-YM-variation}.  The gauge row is a strong $L^2$--strong
$L^2$ pairing.  The metric row contains bounded coframe coefficients
converging in measure multiplying the strongly convergent $L^1$ packet
$F_h\otimes F_h$; \Cref{lem:products} passes it.  Uniform $C^1$ bounds and a
finite-net argument give the stated $C^r$-dual convergence.
\end{proof}

For the spinor and dual-spinor blocks define the positive internal-link graph
bound
\begin{equation}
 \sup_h\left(\norm{\Psi_h}_{2,h}+\norm{\mathcal D_h^U\Psi_h}_{2,h}
 +\norm{\bPsi_h}_{2,h}+\norm{\mathcal D_h^{U,\vee}\bPsi_h}_{2,h}\right)<\infty.
 \label{eq:native-spinor-graph}
\end{equation}

\begin{proposition}[Weak native spinor graphs and the complete Dirac variation]
\label{prop:native-spinor-variation}
Assume \eqref{eq:native-spinor-graph}, strong $L^4$ connection convergence,
the Higgs conclusions of \Cref{prop:native-Higgs-compactness}, strong
$H^1$ coframe convergence on one compact nondegenerate chart, and the scaled
spin-link margin \eqref{eq:native-gravity-log-margin}. After
extraction the spinors converge weakly in the reconstructed $H^1$ class, and
the complete finite Dirac--Yukawa first variation converges to its continuum
counterpart in $(\Vtest_K^r)^*$ for every $r\ge2$.  This includes the full
metric variation of the coframe-derived spin connection.  Strong convergence
of the spinor first differences is not required.
\end{proposition}
\begin{proof}
The graph identity in \Cref{lem:native-critical-grid} gives ordinary
discrete $H^1$ bounds.  Hence the values converge strongly in $L^2$ after
extraction, their first differences converge weakly in $L^2$, and the
undifferentiated spinor bilinears converge strongly in $L^1$ while remaining
bounded in $L^2$.  The physical spin--gauge link can be written exactly as
\[
 \nabla^h_\mu\Psi_h=D^+_{\mu,h}\Psi_h+B_{\mu,h}T_{\mu,1}\Psi_h,
\]
with the analogous dual identity.  The coefficient $B_{\mu,h}$ splits into
the coframe-derived spin coefficient and the compact internal-link
coefficient.  Strong first-jet coframe convergence gives the former and its
metric variation strongly in $L^2$ by \Cref{lem:products}; strong $L^4$
connection convergence gives the latter by \eqref{eq:native-link-coeff}.
Derivative terms therefore use weak--strong $L^2$ pairing, while link and
Yukawa terms pair strongly convergent $L^2$ coefficients with weakly
convergent $L^2$ bilinears.  The same structures occur after metric, gauge,
Higgs, spinor and dual-spinor variation.  Uniform $C^1$ bounds and a finite
$C^1$ net inside the $C^r$ unit ball complete the proof.
\end{proof}

\subsection{The literal gravitational row at first-jet regularity}

For the nodal coframes of \Cref{app:native-action}, write
$\omega_{\mu,h}=\Omega_\mu(e_h,D_h^+e_h)$.  We say that the common
logarithm-chart condition holds when the nodal coframes remain in one fixed
compact nondegenerate chart and
\begin{equation}
 \max_{x,\mu}h|\omega_{\mu,h}(x)|\le c_*
 \label{eq:native-gravity-log-margin}
\end{equation}
for a fixed $c_*$ on which the four-exponential plaquette product and its
logarithm are analytic along the relevant segments.  This is a branch
condition needed to define the literal logarithmic action; it is not a
second-derivative bound.

\begin{proposition}[Literal gravitational consistency from strong first jets]
\label{prop:native-gravity-firstjet}
Suppose the actual nodal coframes satisfy
\begin{equation}
 \mathsf R_h^0e_h\to e\quad\text{in }L^2,
 \qquad
 \mathsf R_h^0D^+_{\mu,h}e_h\to\partial_\mu e
       \quad\text{strongly in }L^2,
 \label{eq:native-coframe-firstjet}
\end{equation}
with values in one fixed compact nondegenerate coframe chart, and assume
\eqref{eq:native-gravity-log-margin}.  Let an inverse-metric test $k$ be
lifted nodally by the symmetric lift \eqref{eq:metric-lift}.  Then for every
$r\ge2$,
\begin{equation}
 \sup_{\norm{k}_{C^r}\le1}
 \left|D S_{{\rm g},h}^{\rm loc}(e_h)[\mathcal I_hk]
       -D\Sact_{{\rm g},\theta}(e)[k]\right|\to0.
 \label{eq:native-gravity-firstvariation}
\end{equation}
If the complete unfiltered reconstructed coframes converge strongly in
$H^1$ to the same $e$ and remain in the same compact chart, then the
sectorwise consistency error relative to those reconstructed coframes tends
to zero.  No growing-band, $C^5$, or strong curvature-discrepancy hypothesis
is required.
\end{proposition}
\begin{proof}
Let
\[
 \Phi(Z_1,Z_2,Z_3,Z_4)
 =\log(e^{Z_1}e^{Z_2}e^{-Z_3}e^{-Z_4}),
 \qquad \mathsf LZ=Z_1+Z_2-Z_3-Z_4,
\]
and define the exact quadratic remainder
\[
 \mathcal Q_h(W)=h^{-2}(\Phi(hW)-h\mathsf LW)
 =\int_0^1(1-t)D^2\Phi(thW)[W,W]\,\dd t.
\]
On the fixed scaled logarithm chart, $|\mathcal Q_h(W)|\le C|W|^2$.
If $W_h\to W$ and $Z_h\to Z$ strongly in $L^2$, bounded-coefficient
multiplication gives
\[
 \mathcal Q_h(W_h)\to\tfrac12D^2\Phi(0)[W,W],\qquad
 D\mathcal Q_h(W_h)[Z_h]\to D^2\Phi(0)[W,Z]
 \quad\text{in }L^1.
\]
For the actual slots
$W_{\mu\nu,h}=(\omega_{\mu,h},T_{\mu,1}\omega_{\nu,h},
T_{\nu,1}\omega_{\mu,h},\omega_{\nu,h})$, the literal Cartan plaquette has
the exact split
\[
 R^h_{\mu\nu}
 =D^+_{\mu,h}\omega_{\nu,h}-D^+_{\nu,h}\omega_{\mu,h}
   +\mathcal Q_h(W_{\mu\nu,h}).
\]
Discrete summation by parts moves the two difference terms onto the smooth
Palatini coframe coefficient.  Because $\Omega(e,p)$ is smooth in $e$ and
linear in $p$, \eqref{eq:native-coframe-firstjet} and
\Cref{lem:products} give strong $L^2$ convergence of the literal connection,
its finite shifts, the differentiated coframe coefficient, and all of their
first physical variations.  Differentiating the exact summed-by-parts action
therefore leaves only products of strongly convergent $L^2$ arrays and the
quadratic remainder terms above.  All pass in $L^1$ and give precisely the
first-derivative Palatini variation of \Cref{app:palatini}.  The same
expressions are bounded by $C\norm{k}_{C^1}$; a finite-net argument gives
\eqref{eq:native-gravity-firstvariation}.  The final consistency statement
follows by subtracting the continuum variation at the reconstructed coframe
and applying \Cref{prop:reduced-continuity} to the geometric sector.
\end{proof}

\begin{theorem}[All-sector native consistency without a growing band]
\label{thm:native-firstvariation-no-band}
Let $z_h^{\rm d}=(e_h,A_h,H_h,\Psi_h,\bPsi_h)$ be actual nodal records of
the unchanged local action \eqref{eq:native-local-action}.  Assume:
\begin{enumerate}[label=\textup{(N\arabic*)},leftmargin=3.2em]
\item the coframes satisfy \eqref{eq:native-coframe-firstjet}, the common
      compact chart and \eqref{eq:native-gravity-log-margin};
\item $\mathsf R_h^0A_h\to A$ strongly in $L^4$ and the literal logarithmic
      curvature satisfies $\mathsf R_h^0F_h\to F$ strongly in $L^2$;
\item $\sup_h\norm{H_h}_{2,h}<\infty$ and
      $\mathsf R_h^0K_h\to K$ strongly in $L^2$ for the literal Higgs-link
      packet;
\item the spinors satisfy \eqref{eq:native-spinor-graph};
\item the coefficient banks converge in a compact physical parameter set.
\end{enumerate}
Then, after extraction, $F=F_A$, $K=D_AH$, the Higgs fields converge as in
\eqref{eq:native-Higgs-compactness}, and the spinors converge weakly in the
reconstructed $H^1$ class.  For every $r\ge2$,
\begin{equation}
 \sup_{\norm v_{C^r}\le1}
 \left|D S_h^{\rm loc}(z_h^{\rm d})[\mathcal I_hv]
       -D\Sact_\theta(z)[v]\right|\longrightarrow0.
 \label{eq:native-all-sector-limit}
\end{equation}
The convergence holds sectorwise.  If, in addition, the complete physical
finite variation tends to zero on the corresponding unit test family, then
$z$ satisfies all minimally coupled torsion-free Einstein--Standard-Model
Euler equations in distributions.

If the complete unfiltered reconstructed coframes stay in the same compact
nondegenerate chart and the ordinary connection first differences converge
strongly to $\partial A$, then the sectorwise consistency errors relative to
the complete unfiltered reconstructed fields tend to zero.  No independent
matter-consistency or growing-band premise remains in this branch.
\end{theorem}
\begin{proof}
The gravitational sector is \Cref{prop:native-gravity-firstjet}.  The
Yang--Mills identification and variation are
\Cref{prop:native-YM-identification}; the Higgs sector is
\Cref{prop:native-Higgs-compactness}; and the complete fermionic sector is
\Cref{prop:native-spinor-variation}.  Each sector has a cutoff-independent
$C^1$ test bound, so their sum converges uniformly on the $C^r$ test unit
ball by a finite-net argument.  Vanishing finite stationarity gives zero
limiting first variation.  For the final assertion, strong ordinary
connection first-jet convergence identifies the complete reconstructed
curvature and covariant Higgs gradient, while the reconstructed coframe
convergence is covered by the reduced continuity theorem.  Subtracting the
native and reconstructed limits yields the sectorwise consistency errors.
\end{proof}

The next appendix supplies the Wilson translation argument and the actual
finite Coulomb normalization that discharge the packet hypotheses in
\Cref{thm:finite-Wilson-zero-defect}.

\section{Wilson compactness and same-record Coulomb normalization}
\label{app:finite-gauge-normalization}

The first result obtains scalar compactness before any gauge is chosen.
A controlled representative then converts Wilson shifts to ordinary
translations. The remaining estimates construct that representative from
a small logarithmic seed and an exact site-gauge path.

\subsection{Wilson compactness directly on the finite packets}

\begin{theorem}[Direct native Wilson compactness]
\label{thm:native-Wilson-compactness}
Let $Y_h$ be bounded in $L_h^2$ on the entire periodic comparison box and
let its internal representation be unitary.  Suppose
\begin{equation}
 \lim_{\rho\downarrow0}\limsup_{h\downarrow0}
       \Omega_h(Y_h;\rho)=0.
 \label{eq:native-Wilson-compactness}
\end{equation}
Then $|\mathsf R_h^0Y_h|$ is precompact in $L^2$ and its square is uniformly
integrable.  This assertion requires no connection bound and holds in any
site gauge.  If, in addition, the links have coordinates $U_h=e^{hA_h}$
with $\sup_h\norm{A_h}_{2,h}=M_A<\infty$, then
$\mathsf R_h^0Y_h$ itself is precompact in $L^2$.  No smooth link
reconstruction or growing band is required.
\end{theorem}
\begin{proof}
Unitarity gives the pointwise inequality
$|T_{\mu,m}|Y_h|-|Y_h||\le|\mathsf W^U_{\mu,m}Y_h-Y_h|$.
\Cref{lem:native-discrete-KR} therefore makes the scalar magnitudes
precompact in $L^2$.  The squares are uniformly integrable by the finite-net
argument of \Cref{lem:compact-magnitude-ui}.  With
\[
 q_h(R)=\norm{Y_h\mathbf1_{|Y_h|>R}}_{2,h},
 \qquad \lim_{R\to\infty}\limsup_{h\downarrow0}q_h(R)=0,
\]
the scalar assertion is complete before a logarithmic representative is used.

In a representative with the stated connection bound, let $P_{\mu,m,h}$
be the represented open-link product.  Telescoping and unitarity imply
\[
 \norm{P_{\mu,m,h}-I}_{2,h}
 \le\sum_{j=0}^{|m|-1}\norm{\rho(U_\mu)-I}_{2,h}
 \le C|m|hM_A.
\]
Here translated norms agree by periodicity; negative displacements use the
inverse product.  Split the translated packet at amplitude $R$.  Since
a unitary matrix minus the identity has operator norm at most two,
\begin{equation}
 \norm{T_{\mu,m}Y_h-Y_h}_{2,h}
 \le\Omega_h(Y_h;\rho)+CR\rho M_A+2q_h(R)
 \quad (|m|h\le\rho).
 \label{eq:native-Wilson-ordinary-modulus}
\end{equation}
For fixed $R$, first take the limsup in $h$, then let $\rho\downarrow0$,
and only then let $R\to\infty$.  The right-hand side tends to zero.
Thus the ordinary integer-shift modulus vanishes.  The exact subcell
identity in \Cref{lem:native-discrete-KR} supplies arbitrary translations,
proving vector $L^2$ precompactness.  On buffered charts the same argument
is local, with all translated endpoints kept in the buffer; the global
periodic Hodge conclusions below use full-box screens.
\end{proof}

\subsection{Discrete Coulomb absorption and a finite same-record gauge}

For a nodal Lie-algebra one-form put
\begin{equation}
 \delta_hA_h=-\sum_{\mu=0}^3D^-_{\mu,h}A_{\mu,h},
 \qquad
 D^-_{\mu,h}=\frac{I-T_{\mu,-1}}h.
 \label{eq:native-discrete-divergence}
\end{equation}

\begin{lemma}[Exact periodic discrete Hodge identity]
\label{lem:native-discrete-Hodge}
For every nodal one-form,
\begin{equation}
 \sum_{\mu,\nu}\norm{D^+_{\mu,h}A_{\nu,h}}_{2,h}^2
 =\norm{\dd_h^+A_h}_{2,h}^2+\norm{\delta_hA_h}_{2,h}^2.
 \label{eq:native-discrete-Hodge}
\end{equation}
Consequently
$\norm{A_h}_{1,h}\le C(\norm{\dd_h^+A_h}_{2,h}
+\norm{\delta_hA_h}_{2,h}+\norm{A_h}_{2,h})$ uniformly in $h$.
\end{lemma}
\begin{proof}
At one Fourier frequency let
$d_\mu=(e^{\ii\ell_\mu h}-1)/h$ on the side-$2\pi$ box.  Exterior multiplication by $d$ and
contraction by $\bar d$ satisfy
\[
 \sum_{\mu<\nu}|d_\mu a_\nu-d_\nu a_\mu|^2
 +\left|\sum_\mu\bar d_\mu a_\mu\right|^2
 =\left(\sum_\mu|d_\mu|^2\right)\left(\sum_\nu|a_\nu|^2\right).
\]
Parseval gives \eqref{eq:native-discrete-Hodge}; adding the zero-mode
$L^2$ norm gives the estimate.
\end{proof}

\begin{theorem}[Critical discrete Coulomb compactness]
\label{thm:native-discrete-Coulomb}
There is $\epsilon_*>0$ such that, if
\begin{equation}
 \delta_hA_h=0,
 \qquad \norm{A_h}_{4,h}\le\epsilon_*,
 \qquad \mathsf R_h^0F_h(A_h)\to F\quad\text{strongly in }L^2,
 \label{eq:native-small-Coulomb}
\end{equation}
then, after extraction,
\begin{align}
 \mathsf R_h^0A_h&\to A &&\text{strongly in }L^4,\notag\\
 \mathsf R_h^0D_h^+A_h&\to\partial A &&\text{strongly in }L^2,\notag\\
 \mathcal I_h^{\rm trig}A_h&\to A &&\text{strongly in }H^1,
 \qquad F=F_A,
 \label{eq:native-Coulomb-strong}
\end{align}
with $\delta A=0$.  The result is purely discrete; no continuum Coulomb
theorem or high-derivative reserve is used.
\end{theorem}
\begin{proof}
The exact curvature split, its shifted-slot quadratic bound, and the Hodge
identity give
\[
 \norm{A_h}_{1,h}\le C\bigl(\norm{F_h(A_h)}_{2,h}
                   +\norm{A_h}_{4,h}^2+\norm{A_h}_{2,h}\bigr).
\]
Hence $A_h$ is bounded in the discrete $H^1$ norm.  Extract raw strong
$L^2$ convergence to $A$ and weak convergence of the first differences.
The one-step shifts have the same strong $L^2$ limit.
The scaled slots $hA_h$ are uniformly small because
$h\norm{A_h}_{\infty,h}\le\norm{A_h}_{4,h}$, and tend to zero in measure.
The integral quadratic expansion therefore gives
$\mathsf R_h^0\mathcal N_h(A_h)\to A\wedge A$ in $L^1$.
Indeed its slot products converge strongly in $L^1$, while its bounded
analytic coefficients converge in measure.  Discrete summation by parts
now identifies $F=F_A$ and $\delta A=0$.

Choose periodic mollifications $A^j$ of $A$ preserving $\delta A^j=0$,
with $A^j\to A$ in $H^1$ and
$\norm{A^j}_4\le\norm A_4\le\epsilon_*$.  Put
$B_h^j=\mathsf S_hA^j$.  For each fixed $j$, and all sufficiently small
$h$, $\norm{B_h^j}_{4,h}\le2\epsilon_*$.  Apply the discrete Hodge
estimate and \eqref{eq:native-YM-envelope} to $A_h-B_h^j$:
\begin{align}
 \norm{A_h-B_h^j}_{1,h}
 \le C\bigl(&\norm{F_h(A_h)-F_h(B_h^j)}_{2,h}
       +\norm{\delta_hB_h^j}_{2,h}
       +\norm{A_h-B_h^j}_{2,h}\bigr)\notag\\
       &+3C_NC_S\epsilon_*\norm{A_h-B_h^j}_{1,h}.
 \label{eq:native-Coulomb-explicit-absorption}
\end{align}
Here $C_N,C_S$ depend on the fixed logarithm chart and the periodic
Sobolev/Hodge constants, not on $h$ or $j$.  Choose the threshold so that
$3C_NC_S\epsilon_*<1/2$ and absorb the last term.  Fixed-smooth-field
consistency, strong raw curvature convergence, and
$\delta A^j=0$ then imply
\[
 \limsup_{h\downarrow0}\norm{A_h-B_h^j}_{1,h}
 \le C\bigl(\norm{F_A-F_{A^j}}_2+\norm{A-A^j}_2\bigr).
\]
Only after this limit is taken do we let $j\to\infty$.
Strong $H^1$ convergence of $A^j$ closes its quadratic curvature product,
so the right-hand side tends to zero.  Comparison with these fixed smooth
arrays proves strong convergence of the raw first differences and of the
raw $L^4$ values.  \Cref{lem:native-reconstruction-identification} supplies
the trigonometric $H^1$ statement.  No cutoff-dependent smooth comparison
or estimate for its high derivatives is used.
\end{proof}

For a small operator $Z$ define
\begin{equation}
 \mathcal J(Z)=\int_0^1e^{tZ}\,\dd t,
 \qquad
 \mathcal K(Z)=\frac Z2\coth\!\left(\frac Z2\right),
 \qquad \mathcal K(0)=I.
 \label{eq:native-JK}
\end{equation}
\begin{lemma}[Exact logarithmic site-gauge differential]
\label{lem:native-log-gauge-differential}
Use $U^q_\mu(x)=q_xU_\mu(x)q_{x+h e_\mu}^{-1}$ and a fixed
conjugation-invariant logarithm chart. Put
$Z_\mu=h\operatorname{ad}_{A_\mu}$ and
$m_\mu\xi=\tfrac12(\xi+T_{\mu,1}\xi)$. For $q_x(t)=e^{t\xi_x}$,
\begin{equation}
 \left.\frac{\dd}{\dd t}A_\mu^{q(t)}\right|_{t=0}
 =-\mathcal K(Z_\mu)D_\mu^+\xi-[A_\mu,m_\mu\xi].
 \label{eq:native-gauge-differential}
\end{equation}
In particular
\begin{equation}
 \left.\frac{\dd}{\dd t}\frac12\norm{A^{q(t)}}_{2,h}^2\right|_{t=0}
 =-\langle\delta_hA,\xi\rangle_{2,h}.
 \label{eq:native-gauge-energy-differential}
\end{equation}
These identities concern logarithmic link coordinates, not an additive
approximation to a finite gauge transformation.
\end{lemma}
\begin{proof}
The differential of the exponential in left-trivialized coordinates is
\[
 \dot U_\mu U_\mu^{-1}
   =h\mathcal J(Z_\mu)\dot A_\mu.
\]
For a site-gauge variation the same quantity is
$\xi_x-e^{Z_\mu}\xi_{x+h e_\mu}$. Using
$\xi_x=m_\mu\xi-\tfrac h2D_\mu^+\xi$ and
$\xi_{x+h e_\mu}=m_\mu\xi+\tfrac h2D_\mu^+\xi$ gives
\[
 \mathcal J(Z)^{-1}(I-e^Z)=-Z,\qquad
 \tfrac12\mathcal J(Z)^{-1}(I+e^Z)=\mathcal K(Z),
\]
which proves \eqref{eq:native-gauge-differential}. On the real Lie algebra
with its positive invariant metric, $\operatorname{ad}_A$ is skew-adjoint;
$\mathcal K$ is an even analytic function and
$\mathcal K(h\operatorname{ad}_A)^*A=A$. Moreover
$\langle A,[A,\eta]\rangle=0$. Pairing the differential with $A$ and using
$\delta_h=(D_h^+)^*$ therefore gives
\eqref{eq:native-gauge-energy-differential}.
\end{proof}

Define the mean-zero site space
$\mathcal Z_h=\{\xi:h^4\sum_x\xi_x=0\}$, with norm
$\norm\xi_{\dot H_h^1}=\norm{D_h^+\xi}_{2,h}$ and dual norm
\[
 \norm f_{\dot H_h^{-1}}
 =\sup_{\xi\in\mathcal Z_h,\,\norm{D_h^+\xi}_{2,h}=1}
       |\langle f,\xi\rangle_{2,h}|.
\]
Periodic Poincar\'e and Sobolev give
$\norm\xi_{2,h}+\norm\xi_{4,h}\le C\norm{D_h^+\xi}_{2,h}$ on this space,
with $C$ independent of $h$. The logarithmic discrete Faddeev--Popov operator is
\begin{equation}
 \mathcal M_A\xi
 =\delta_h\left(\mathcal K(h\operatorname{ad}_{A_\mu})D_\mu^+\xi
       +[A_\mu,m_\mu\xi]\right)_\mu.
 \label{eq:native-FP}
\end{equation}
Its range has zero mean because it is a discrete divergence.

\begin{lemma}[Uniform transverse inverse]
\label{lem:native-FP-inverse}
There is $r_{\rm FP}>0$, independent of $h$, such that
$\norm A_{4,h}<r_{\rm FP}$ implies
\begin{equation}
 \langle\mathcal M_A\xi,\xi\rangle_{2,h}
 \ge\tfrac12\norm{D_h^+\xi}_{2,h}^2\qquad(\xi\in\mathcal Z_h).
 \label{eq:native-FP-coercivity}
\end{equation}
Thus $\mathcal M_A:\mathcal Z_h\to\mathcal Z_h$ is invertible and
\begin{equation}
 \norm{D_h^+\mathcal M_A^{-1}f}_{2,h}
       \le2\norm f_{\dot H_h^{-1}}.
 \label{eq:native-FP-dual-inverse}
\end{equation}
In particular, for any link array $R$,
\begin{equation}
 \norm{\mathcal M_A^{-1}\delta_hR}_{2,h}
 +\norm{\mathcal M_A^{-1}\delta_hR}_{4,h}
 +\norm{D_h^+\mathcal M_A^{-1}\delta_hR}_{2,h}
 \le C\norm R_{2,h}.
 \label{eq:native-FP-divergence-inverse}
\end{equation}
The inverse depends smoothly on $A$ on this finite-dimensional open set.
Self-adjointness of $\mathcal M_A$ away from Coulomb gauge is not required.
\end{lemma}
\begin{proof}
On a fixed sufficiently small scaled chart,
$\norm{\mathcal K(h\operatorname{ad}_A)-I}\le Ch^2|A|^2$.
The four-dimensional cell mass gives
$h\norm A_{\infty,h}\le\norm A_{4,h}$. Consequently
\[
 |\langle(\mathcal K-I)D^+\xi,D^+\xi\rangle_{2,h}|
 \le C\norm A_{4,h}^2\norm{D^+\xi}_{2,h}^2.
\]
For the other perturbation, discrete H\"older, translation invariance and
the mean-zero Sobolev estimate give
\[
 |\langle[A,m\xi],D^+\xi\rangle_{2,h}|
 \le C\norm A_{4,h}\norm{m\xi}_{4,h}\norm{D^+\xi}_{2,h}
 \le C\norm A_{4,h}\norm{D^+\xi}_{2,h}^2.
\]
Choose $r_{\rm FP}$ so that $C(r_{\rm FP}+r_{\rm FP}^2)\le1/2$ and so
that the scaled analytic chart is retained. This proves coercivity.
The restriction to $\mathcal Z_h$ is injective, hence bijective because its
domain and codomain have equal finite dimension. Pairing
$\mathcal M_A\xi=f$ with $\xi$ gives
\eqref{eq:native-FP-dual-inverse}. Also
$\norm{\delta_hR}_{\dot H_h^{-1}}\le\norm R_{2,h}$, proving
\eqref{eq:native-FP-divergence-inverse}. Smooth dependence follows from the
finite matrix inverse on its invertibility set. All displayed norm constants
are independent of the number of sites.
\end{proof}

\begin{theorem}[Finite critical Coulomb normalization]
\label{thm:finite-Coulomb-normalization}
There are $\epsilon_c,C_c>0$, independent of $h$, such that any logarithmic
seed $B_h$ satisfying
\begin{equation}
 \norm{B_h}_{4,h}+\norm{F_h(B_h)}_{2,h}\le\epsilon_c
 \label{eq:native-critical-seed}
\end{equation}
admits a finite site gauge $q_h$ for which
\begin{equation}
 A_{\mu,h}=h^{-1}\operatorname{Log}
 \bigl(q_xe^{hB_{\mu,h}(x)}q_{x+h e_\mu}^{-1}\bigr),
 \qquad \delta_hA_h=0,
 \label{eq:native-constructed-Coulomb}
\end{equation}
and
\begin{equation}
 \norm{A_h}_{1,h}+\norm{A_h}_{4,h}
 \le C_c\bigl(\norm{B_h}_{2,h}+\norm{F_h(B_h)}_{2,h}
                 +\norm{B_h}_{4,h}^2\bigr).
 \label{eq:native-Coulomb-estimate}
\end{equation}
The threshold may be chosen so that $\norm{A_h}_{4,h}\le\epsilon_*$ in
\Cref{thm:native-discrete-Coulomb}.  No bound on $\delta_hB_h$ or on the
complete first differences of $B_h$ is assumed.
\end{theorem}
\begin{proof}
All constants below refer to the fixed periodic box, invariant Lie-algebra
metric and logarithm chart. They are independent of $h$. Choose $r>0$ below
the transverse threshold of \Cref{lem:native-FP-inverse} and below the scaled
four-exponential logarithm threshold. Shrink it further in the absorption
step below. For a fixed grid consider the auxiliary links
$U^s_\mu=e^{shB_\mu}$, $0\le s\le1$.

\emph{Local gauge path.}
On the open set of pairs $(s,q)$ for which
\[
 A_{s,\mu}=h^{-1}\operatorname{Log}
       (q_xe^{shB_\mu}q_{x+h e_\mu}^{-1}),\qquad\norm{A_s}_{4,h}<r,
\]
define
\begin{align}
 R_{s,\mu}
  &=\mathcal J(h\operatorname{ad}_{A_{s,\mu}})^{-1}
            \operatorname{Ad}_{q_x}B_\mu,\notag\\
 \xi_s&=\mathcal M_{A_s}^{-1}\delta_hR_s\in\mathcal Z_h,
 &\dot q_x&=\xi_s(x)q_x,
 &q_x(0)&=I.
 \label{eq:native-Coulomb-homotopy-ODE}
\end{align}
The right-hand side is a smooth tangent field on a finite product of compact
Lie groups. The logarithms are defined at $(s,q)=(0,I)$, where $A_0=0$;
finite-dimensional ODE existence gives a local solution. The exact chain
rule of \Cref{lem:native-log-gauge-differential} gives
\[
 \dot A_s=R_s-
   \bigl(\mathcal K(h\operatorname{ad}_{A_{s,\mu}})D_\mu^+\xi_s
                +[A_{s,\mu},m_\mu\xi_s]\bigr)_\mu,
 \qquad\delta_h\dot A_s=0.
\]
Hence $\delta_hA_s=0$ for as long as the path remains in the open set.

\emph{The physical $L_h^2$ estimate.}
The gauge tangent contributes zero when paired with $A_s$, by
\eqref{eq:native-gauge-energy-differential}. The adjoint of
$\mathcal J(h\operatorname{ad}_{A_{s,\mu}})^{-1}$ fixes $A_{s,\mu}$.
Therefore
\[
 \frac12\frac{\dd}{\dd s}\norm{A_s}_{2,h}^2
 =\sum_\mu\langle A_{s,\mu},\operatorname{Ad}_{q_x}B_\mu\rangle_{2,h}
 \le\norm{A_s}_{2,h}\norm B_{2,h}.
\]
Apply this inequality first to
$(\norm{A_s}_{2,h}^2+\epsilon^2)^{1/2}$ and then let
$\epsilon\downarrow0$. It follows that
$\norm{A_s}_{2,h}\le s\norm B_{2,h}$.

\emph{Curvature and the cutoff-independent first-exit bound.}
Exact conjugation covariance of the plaquette logarithm gives
$\norm{F_h(A_s)}_{2,h}=\norm{F_h(sB)}_{2,h}$ on the declared chart.
The exact split \eqref{eq:native-YM-split} yields
\[
 F_h(sB)=sF_h(B)-s\mathcal N_h(B)+\mathcal N_h(sB),
 \qquad
 \norm{F_h(sB)}_{2,h}
 \le\norm{F_h(B)}_{2,h}+C\norm B_{4,h}^2.
\]
Put $b_h=\norm B_{2,h}+\norm{F_h(B)}_{2,h}+\norm B_{4,h}^2$.
The discrete Hodge identity, Coulomb condition and quadratic plaquette bound
then give, throughout this path,
\begin{equation}
 \norm{A_s}_{1,h}\le C_H(b_h+\norm{A_s}_{4,h}^2),\qquad
 a_s:=\norm{A_s}_{4,h}\le C_0b_h+C_1a_s^2.
 \label{eq:native-Coulomb-first-exit}
\end{equation}
Choose $r$ so that $C_1r\le1/4$. While $a_s<r$, the last inequality gives
$a_s\le(4/3)C_0b_h$. On the fixed-volume box
$\norm B_{2,h}\le C\norm B_{4,h}$, so the seed assumption makes
$b_h\le C(\epsilon_c+\epsilon_c^2)$. Choose $\epsilon_c$ so that
$(4/3)C_0b_h\le r/2$. Thus the path cannot reach the boundary $a_s=r$.
This choice also makes the right side smaller than the desired
$\epsilon_*$, after shrinking $r$ if necessary.

\emph{Continuation without leaving the logarithm chart.}
The improved bound gives $h\norm{A_s}_{\infty,h}\le r/2$ and therefore
places every normalized link strictly inside a common logarithm domain;
all four-link products stay inside its corresponding plaquette domain.
The seed links $e^{shB}$ stay in their analytic domain by the same cell-mass
bound and the small seed threshold. Moreover
$\norm{R_s}_{2,h}\le C\norm B_{2,h}$, since $\mathcal J^{-1}$ is uniformly
bounded on this domain. By \eqref{eq:native-FP-divergence-inverse},
\begin{equation}
 \norm{\xi_s}_{2,h}+\norm{\xi_s}_{4,h}
       +\norm{D_h^+\xi_s}_{2,h}\le C\norm B_{2,h}.
 \label{eq:native-Coulomb-generator-control}
\end{equation}
For each fixed $h$ this also bounds all site velocities in the finite ODE.
Thus if the maximal path ended at some $s_0\le1$, the group variables would
have limits there. The strict logarithm and transverse-inverse margins would
put the limiting point inside the smooth ODE domain, allowing continuation.
This contradicts maximality unless the path reaches $s=1$. No bound uniform
in $h$ on derivatives of the gauge itself is needed for this continuation.

At $s=1$ the path gives the actual site gauge in
\eqref{eq:native-constructed-Coulomb}. From $a_s\le Cb_h$ and the first
inequality in \eqref{eq:native-Coulomb-first-exit}, with $b_h$ small, we obtain
\eqref{eq:native-Coulomb-estimate}. The conclusion is existence of a small
Coulomb representative on this seed branch; it is not uniqueness among all
site gauges and makes no assertion outside the logarithm chart. The auxiliary
homotopy is a proof construction, not a modification of the observed record
or of its physical dynamics.
\end{proof}

\begin{proposition}[Rooted same-record gauge certificate]
\label{prop:rooted-gauge-certificate}
Let $\mathfrak T_h$ be the fixed rooted spanning tree and define
$B_h^{\mathfrak T}$ by \eqref{eq:rooted-Wilson-seed}.  On its admissible
logarithm chart the quantity \eqref{eq:rooted-Wilson-certificate} is invariant
under finite site gauges.  If
$\mathfrak c_h^{\mathfrak T}(U_h)\le\epsilon_c$, then the exact tree gauge
followed by \Cref{thm:finite-Coulomb-normalization} puts the original links
in exact small discrete Coulomb gauge.  Under this normalization:
\begin{enumerate}[label=\textup{(a)},leftmargin=2.8em]
\item every sector of the finite common action is unchanged;
\item the literal curvature and Higgs-gradient packets, the positive
      internal-link spinor graphs, and the Wilson moduli are unchanged in
      norm;
\item the norm of the complete action covector is unchanged in the positive
      physical mass metric whose link tangent is
      $\alpha_\mu=h^{-1}\dot U_\mu U_\mu^{-1}$.
\end{enumerate}
No derivative of the normalizing site gauge appears in these norm identities.
\end{proposition}
\begin{proof}
Endpoint cancellation along the tree gives
$C_x(U^k)=k_oC_x(U)k_x^{-1}$ and hence
$W_\mu^{\mathfrak T}(U^k)=k_oW_\mu^{\mathfrak T}(U)k_o^{-1}$.
The logarithm is equivariant on the declared chart and the invariant inner
product makes the certificate gauge invariant.  The tree representative is
therefore an exact gauge transform of the original links, not an
approximation, and \Cref{thm:finite-Coulomb-normalization} applies to it.

For the physical statements, plaquette and open-link endpoint cancellation
give conjugation of the curvature packet at its base vertex.  Direct
substitution gives the corresponding unitary transformation of the
Higgs-gradient and internal-link spinor packets.  The coframe-derived spin
transport commutes with the internal group, and the covariant Yukawa blocks
transform in their prescribed representations; hence the full finite action
is unchanged sector by sector.  The physical link tangent transforms by
$\alpha_\mu\mapsto\operatorname{Ad}_{k_x}\alpha_\mu$ and the matter tangents
unitarily.  Thus the positive tangent map is an isometry, and differentiating
the exact action invariance proves invariance of the dual covector norm.
\end{proof}

\section{Determinant normalization and covariant Wilson transfer}
\label{app:determinant-transfer}

\subsection{The determinant sector: flux, harmonic coordinates, and exact assembly}
\label{app:determinant-sector}

We work on the periodic comparison box of side $L=2\pi$, with $h=L/n$
and positive cell masses $h^4$.  Keeping $L$ in the formulas distinguishes
physical differences from lattice-unit differences.  The Abelian
flux and flat-holonomy decomposition is classical
\cite[Section~7]{Luscher1999}.  The proofs below record its physical
normalization, its critical compactness estimates, and its application to
the same Standard-Model links and physical variations.  No change to the
Lorentzian common action is involved.

\subsubsection{The determinant screen and the integer flux}

For $u_h=\chi(U_h)$ let $p_{\mu\nu,h}$ be its oriented plaquette and put
\begin{equation}
 \phi_{\mu\nu,h}=\operatorname{Arg}p_{\mu\nu,h}\in(-\pi,\pi],\qquad
 f_{\mu\nu,h}=h^{-2}\phi_{\mu\nu,h}.
 \label{eq:determinant-real-curvature}
\end{equation}
Let $\dd_h$ denote the forward exterior derivative on the periodic cubical
complex, $\delta_h=\dd_h^*$ its positive-mass adjoint, and
$\Delta_h=\dd_h\delta_h+\delta_h\dd_h$ the nonnegative Hodge Laplacian.
On one-forms $\delta_h$ is \eqref{eq:native-discrete-divergence}.
For an anti-Hermitian $X=(X_3,X_2)\in\mathfrak g_{\rm SM}$ define
\begin{equation}
 \ell_c(X)=\frac1{\ii}\Tr X_2,\qquad
 \Pi_{\rm ss}X=X-Z_c\ell_c(X).
 \label{eq:determinant-projections}
\end{equation}
Both maps are bounded and adjoint equivariant, with $\ell_c$ adjoint
invariant and $\ell_c(Z_c)=1$.

\begin{lemma}[Determinant compactness, admissibility, and flux]
\label{lem:determinant-screen-flux}
Suppose the literal full curvature has bounded $L_h^2$ norm and vanishing
adjoint Wilson modulus.  Then, for all sufficiently small $h$,
\begin{equation}
 f_h=\ell_c(\mathbb F_h),\qquad
 \omega_h(f;\rho)\le C\Omega_h(\mathbb F_h;\rho),\qquad
 \max_{x,\mu<\nu}|\phi_{\mu\nu,h}(x)|\longrightarrow0,
 \label{eq:determinant-inherited-screen}
\end{equation}
where $\omega_h$ is the ordinary discrete translation modulus.  In
particular, $\mathsf R_h^0f_h$ is precompact in $L^2$ and
$\dd_h f_h=0$ exactly for sufficiently small $h$.  Its coordinate-plane
fluxes are transverse-base-point-independent integers:
\begin{equation}
 h^2\sum_{x\in\mathbb T^2_{\mu\nu,h}} f_{\mu\nu,h}(x)
       =2\pi m_{\mu\nu,h},\qquad
 \overline f_{\mu\nu,h}
       :=L^{-4}h^4\sum_x f_{\mu\nu,h}(x)
       =\frac{2\pi}{L^2}m_{\mu\nu,h}.
 \label{eq:determinant-flux}
\end{equation}
The integer vector obeys $2\pi|m_h|\le\norm{f_h}_{2,h}$, and hence is
constant along a subsequence.  None of these conclusions uses a bound on
the logarithms of the original links.
\end{lemma}
\begin{proof}
The scalar-magnitude part of \Cref{thm:native-Wilson-compactness} gives
uniform integrability of $|\mathsf R_h^0\mathbb F_h|^2$ without a connection
bound.  A cell has mass $h^4$, so
\[
 \max_x h^4|\mathbb F_h(x)|^2\longrightarrow0.
\]
Thus the full scaled curvature $h^2\mathbb F_h$ tends uniformly to zero.
Since $\chi(e^X)=e^{\ii\ell_c(X)}$, the principal determinant argument is
$h^2\ell_c(\mathbb F_h)$ eventually.  Applying the invariant functional
$\ell_c$ to an adjoint Wilson difference gives exactly the corresponding
ordinary difference of $f_h$.  \Cref{lem:native-discrete-KR} therefore
proves its compactness and the first two assertions.

The product of the six oriented Abelian plaquettes around a cube is one.
Their signed arguments sum to $2\pi$ times an integer; when
$\max|\phi_h|<\pi/3$, the absolute value of that sum is less than $2\pi$.
The integer is zero, giving $\dd_hf_h=0$.  The product over a periodic
coordinate plane is also one, which gives its integer flux.  Summing the
real cube identity over a plane shows that its flux is unchanged by a
transverse shift.  Averaging all parallel planes gives the second identity
in \eqref{eq:determinant-flux}.  Finally
$L^2|\overline f_h|\le\norm{f_h}_{2,h}$ proves the integer bound.  The
finite set of permitted integer vectors admits a constant subsequence.
\end{proof}

\subsubsection{A uniformly controlled transverse potential}

\begin{lemma}[Strong stability of the periodic discrete Hodge inverse]
\label{lem:determinant-Hodge-stability}
For a closed, zero-mean real two-form $f_h$, set
\begin{equation}
 a_h^0=\delta_h\Delta_h^\dagger f_h,
 \label{eq:determinant-Hodge-primitive}
\end{equation}
where $\Delta_h^\dagger$ vanishes on the constant modes.  Then
\begin{gather}
 \dd_h a_h^0=f_h,\qquad \delta_h a_h^0=0,\qquad
 \overline{a_h^0}=0,\notag\\
 \norm{D_h^+a_h^0}_{2,h}=\norm{f_h}_{2,h},\qquad
 L^{-1}\norm{a_h^0}_{2,h}+\norm{a_h^0}_{4,h}
       \le C\norm{f_h}_{2,h}.
 \label{eq:determinant-Hodge-bounds}
\end{gather}
If $\mathsf R_h^0f_h\to f$ strongly in $L^2$, then
$\mathsf R_h^0a_h^0\to a^0$ strongly in $L^4$ and
$\mathsf R_h^0D_h^+a_h^0\to\partial a^0$ strongly in $L^2$, where
$a^0=\delta\Delta^\dagger f$.  On the odd grids used for trigonometric
reconstruction, $\mathcal I_h^{\rm trig}a_h^0\to a^0$ strongly in $H^1$.
\end{lemma}
\begin{proof}
The Hodge operators commute with $\Delta_h^\dagger$.  Since $f_h$ is
closed and has no harmonic part,
$\dd_h\delta_h\Delta_h^\dagger f_h=f_h$ and
$\delta_h^2\Delta_h^\dagger f_h=0$.  The discrete Hodge identity gives
the exact derivative norm.  On a Fourier mode with
$\xi_\ell=2\pi\ell/L$ the symbols are
$d_{h,\mu}=(e^{\ii h\xi_{\ell,\mu}}-1)/h$ and
$\lambda_h=\sum_\mu|d_{h,\mu}|^2$.
The smallest nonzero $\lambda_h$ is bounded below by $cL^{-2}$,
independently of $n$.  Poincar\'e and the scaled discrete Sobolev inequality
therefore give \eqref{eq:determinant-Hodge-bounds}.

The multiplier of $\delta_h\Delta_h^\dagger$ is bounded by
$C/|\xi_\ell|$, and that of each of its full first differences is uniformly
bounded.  Both converge at each fixed nonzero frequency.  Strong raw
$L^2$ convergence gives convergence of every fixed discrete Fourier
coefficient and of the total Parseval mass.  Consequently
\[
 \lim_{R\to\infty}\limsup_{h\downarrow0}
    L^4\sum_{|\xi_\ell|>R}|\widehat f_h(\ell)|^2=0.
\]
Apply the multiplier bounds to this tail and pass the finitely many low
modes first.  This proves strong $L^2$ convergence of $a_h^0$ and of all
its first differences.  \Cref{lem:native-reconstruction-identification}
then gives the $L^4$ and reconstructed $H^1$ conclusions.  This proof uses
no rate of decay of the tail and no growing frequency band.
\end{proof}

\begin{proposition}[Zero-flux Abelian normalization with unrestricted flat holonomy]
\label{prop:periodic-determinant-normalization}
Let periodic $U(1)$ links have $\max|\phi_h|<\pi/3$ and zero flux.  The
potential \eqref{eq:determinant-Hodge-primitive}, together with constants
$c_{\mu,h}\in[-\pi/L,\pi/L]$, determines a periodic site gauge $s_h$ such
that
\begin{equation}
 u_{\mu,h}^{s_h}=e^{\ii h a_{\mu,h}},\qquad
 a_h=a_h^0+c_h,\qquad \delta_h a_h=0.
 \label{eq:determinant-Abelian-normalization}
\end{equation}
The constants retain the four cycle holonomies of the flat residual.
For a screened bounded-curvature sequence, these coordinates are strongly
precompact in raw $L^4$, their first differences are strongly precompact
in $L^2$, and their trigonometric reconstructions are strongly precompact
in $H^1$.  Moreover $h\norm{a_h}_{\infty,h}\to0$; thus the coordinates are
eventually principal.  No smallness of $\norm{f_h}_{2,h}$,
$\norm{a_h}_{4,h}$, or the cycle phases is required.
\end{proposition}
\begin{proof}
Zero flux gives $\overline f_h=0$ by \eqref{eq:determinant-flux}.  The
plaquette of $e^{\ii h a_h^0}$ is exactly
$e^{\ii h^2\dd_h a_h^0}=e^{\ii\phi_h}$.  Hence
$r_{\mu,h}=u_{\mu,h}e^{-\ii h a^0_{\mu,h}}$ is flat.  Its four cycle
products $w_{\mu,h}$ are independent of the base point, since moving a
cycle across a strip multiplies it by a product of trivial plaquettes.
Choose $Lc_{\mu,h}=\operatorname{Arg}w_{\mu,h}$ using one fixed endpoint
convention.  The ratio $r_he^{-\ii h c_h}$ is flat with trivial cycle
products.  Its ordered product from a fixed root to a vertex defines a
periodic, path-independent site phase $s_h$ with
$s_h(x+h e_\mu)=s_h(x)r_{\mu,h}(x)e^{-\ii h c_{\mu,h}}$.
This proves \eqref{eq:determinant-Abelian-normalization} exactly.

For sequences, \Cref{lem:determinant-screen-flux} gives strong curvature
compactness and \Cref{lem:determinant-Hodge-stability} gives strong
potential compactness.  Extract also the bounded four-vector $c_h$.
Uniform integrability of the fourth powers of the raw potentials implies
$\max_x h^4|a_h(x)|^4\to0$, proving the resolved-link assertion.  The
compactness is a statement about the entire family as $h\downarrow0$:
otherwise a violating sequence would have the convergent subsequence just
constructed.
\end{proof}

\begin{remark}[Nonzero flux is not removed by screening]
\label{rem:determinant-nonzero-flux}
For fixed flux $m$, the fluctuation potential
$\delta_h\Delta_h^\dagger(f_h-2\pi m/L^2)$ still satisfies the preceding
bounds and compactness statement.  The remaining flux must be kept in a
background or transition description.  For example, for one pair
$\mu<\nu$, lattice indices $k\in\{0,\ldots,n-1\}^4$ and integer $m$, the
periodic links
\[
 v_\nu(k)=e^{2\pi\ii m k_\mu/n^2},\qquad
 v_\mu(k)=
 \begin{cases}
 e^{-2\pi\ii m k_\nu/n},&k_\mu=n-1,\\
 1,&k_\mu\ne n-1,
 \end{cases}
\]
with all other components one have constant physical curvature
$f_{\mu\nu}=2\pi m/L^2$ on every plaquette, including the seam, for large
$n$.  Their curvature translation modulus is zero.  Conversely, links
$e^{\ii h a_h}$ with one periodic real potential and
$h\norm{a_h}_\infty\to0$ have zero flux eventually: the principal
plaquette argument equals the real curl of $ha_h$, and summing that curl
over a coordinate plane telescopes to zero.  Thus the global hypothesis
in \Cref{cor:determinant-resolved-closure} cannot be dropped on regularity
grounds.  No global trivialization of a nonzero-flux sector is claimed.
\end{remark}

\subsubsection{Assembly in the actual Standard-Model group}

\begin{proposition}[Exact determinant normalization and semisimple remainder]
\label{prop:native-determinant-split}
Let the full curvature screen hold and the determinant flux vanish.
Construct $s_h,a_h$ by \Cref{prop:periodic-determinant-normalization}.
At each site choose $t_h(x)\in\R$ with $e^{\ii t_h(x)}=s_h(x)$ and define
\begin{equation}
 g_h(x)=e^{t_h(x)Z_c},\qquad
 U'_h=U_h^{g_h},\qquad
 V_{\mu,h}=e^{-h a_{\mu,h}Z_c}U'_{\mu,h}.
 \label{eq:determinant-exact-split}
\end{equation}
Then $V_h\in G_{\rm ss}$ exactly and, for sufficiently small $h$,
\begin{equation}
 F_h(V_h)=\Pi_{\rm ss}F_h(U'_h).
 \label{eq:determinant-semisimple-curvature}
\end{equation}
The adjoint Wilson transport of $U'_h$ equals that of $V_h$.  Its
semisimple curvature screen is therefore bounded by the original full
screen, up to a fixed norm constant.  A change of the site lifts $t_h$
changes $V_h$ only by a central $G_{\rm ss}$ site gauge, and does not change
its rooted certificate.  All transformations applied to matter and to the
common action are transformations of the combined $G_{\rm SM}$ record.
\end{proposition}
\begin{proof}
By \eqref{eq:determinant-character},
$\chi(e^{tZ_c})=e^{\ii t}$, so
$\chi(V_{\mu,h})=e^{-\ii h a_{\mu,h}}\chi(U'_{\mu,h})=1$.
Centrality gives the exact plaquette relation
\[
 P_{U',\mu\nu}=e^{h^2f_{\mu\nu,h}Z_c}P_{V,\mu\nu}.
\]
All scaled full-curvature arguments tend uniformly to zero by
\Cref{lem:determinant-screen-flux}.  The commuting central factor and the
unique local logarithm then give
$F_h(V_h)=F_h(U'_h)-Z_cf_h$, which is
\eqref{eq:determinant-semisimple-curvature}.  Central factors act trivially
in the adjoint representation, so the open-link transports and their
projected curvature differences have the asserted relation.

The central one-parameter subgroup has period $12\pi$, not $2\pi$:
\begin{equation}
 z_6=e^{2\pi Z_c}=(e^{-2\pi\ii/3}I_3,-I_2)\in G_{\rm ss},\qquad
 z_6^6=I.
 \label{eq:determinant-central-six}
\end{equation}
Changing $t_h(x)$ to $t_h(x)+2\pi k_x$ thus multiplies $g_h(x)$ by the
central site element $z_6^{k_x}$.  The corresponding change of $V_h$ is an
actual $G_{\rm ss}$ gauge transformation; apply
\Cref{prop:rooted-gauge-certificate} for certificate invariance.  A
semisimple normalization $q_h$ of $V_h$ acts on the full record by
$q_hg_h$, with its central factor retained.  Exact covariance preserves
all sector actions, the full physical covector norm, and the full matter
Wilson and graph norms.  In particular no fractional hypercharge is
regarded as a single-valued character of the determinant circle, and $V_h$
alone is never substituted for the matter connection.
\end{proof}

\begin{lemma}[Flat semisimple normalization without a small-holonomy assumption]
\label{lem:flat-semisimple-normalization}
If $V_h$ is a flat periodic $\SU(3)\times\SU(2)$ link field, it is gauge
equivalent to constant links $e^{h b_{\mu,h}}$ with commuting, uniformly
bounded $b_{\mu,h}\in\mathfrak g_{\rm ss}$.  After extraction these
constant coordinates converge, and their first differences and curvatures
vanish identically.  Their cycle holonomies need not be close to the
identity.
\end{lemma}
\begin{proof}
Flatness implies that the based coordinate-cycle holonomies commute.
In each $\SU(N)$ factor, simultaneously diagonalize them and choose real
eigenangles for each holonomy.  Their sum is an integer multiple $2\pi m$; subtracting $2\pi m$ from
one eigenangle preserves its exponential and gives a traceless
anti-Hermitian logarithm.  All angles remain bounded by a constant
depending only on $N$.  Dividing the four logarithms by $L$ gives
commuting $b_{\mu,h}$ with the prescribed cycle holonomies.  Comparison of
path transports for the two flat link fields gives a site gauge: plaquette
moves preserve the comparison, and their equal based cycle holonomies
make it periodic.  This argument compares two flat connections; it does
not assume that the pointwise ratio of arbitrary non-Abelian connections
is flat.  The bounded finite-dimensional tuple $(b_{\mu,h})_\mu$ has a
convergent subsequence.  Constant conjugations needed for diagonalization
are also actual internal gauges.
\end{proof}

\begin{example}[Non-small determinant curvature with compact gauge data]
\label{ex:determinant-large-field}
Let $M_0,c$ be arbitrary fixed real numbers, and on the side-$L$ grid set
\[
 a_{1,h}(x)=M_0\sin(2\pi x^0/L)+c,\qquad
 a_{\mu,h}=0\ (\mu\ne1),\qquad U_{\mu,h}=e^{hZ_ca_{\mu,h}}.
\]
The divergence is zero and the determinant flux vanishes.  Its only
nonzero determinant curvature component is
$M_0D^+_{0,h}\sin(2\pi x^0/L)$, which converges strongly in $L^2$ and
has a vanishing translation modulus.  Its norm need not be small, and the
flat cycle phase contributed by $c$ is unrestricted.  The remainder of
\eqref{eq:determinant-exact-split} is flat.  Thus the gauge and compactness
conditions of \Cref{cor:determinant-resolved-closure} include bounded
families excluded by a small full-curvature certificate.  Arbitrary site
gauges do not change these conclusions.  This example verifies the range
of the gauge criterion only; it does not assert coupled finite
Einstein--Standard-Model stationarity.
\end{example}

\subsection{Determinant-resolved closure}
\label{subsec:determinant-closure}

The smallness condition in \Cref{thm:finite-Wilson-zero-defect} can be
restricted to the semisimple factor.  For the group in
\eqref{eq:gauge-group}, put
\begin{equation}
 \chi(U_3,U_2)=\det U_2=(\det U_3)^{-1},\qquad
 G_{\rm ss}=\ker\chi=\SU(3)\times\SU(2).
 \label{eq:determinant-character}
\end{equation}
The Abelian curvature of $\chi(U_h)$ is its principal plaquette argument
divided by $h^2$.  Its six integer coordinate-plane fluxes are defined in
\eqref{eq:determinant-flux}.  The full curvature Wilson screen supplies
both compactness of this scalar curvature and eventual admissibility of
its arguments; no independent determinant screen is imposed.  On zero flux,
\Cref{prop:periodic-determinant-normalization,prop:native-determinant-split}
construct from the links a real co-closed determinant potential $a_h$ and
an exactly defined semisimple remainder $V_h\in G_{\rm ss}$.  The potential
retains its four harmonic coordinates.  All matter fields continue to use
the combined $G_{\rm SM}$ links.

\begin{corollary}[Determinant-resolved finite Wilson closure]
\label{cor:determinant-resolved-closure}
Assume \textup{(F1)}, \textup{(F3)}, \textup{(F4)} and \textup{(F5)} of
\Cref{thm:finite-Wilson-zero-defect}.  In place of \textup{(F2)}, assume
that the determinant flux is zero for all sufficiently small $h$ and that
the semisimple remainder constructed above satisfies
\begin{equation}
 \norm{B_h^{\mathfrak T}(V_h)}_{4,h}
       +\norm{F_h(V_h)}_{2,h}\le\epsilon_c
 \label{eq:semisimple-only-certificate}
\end{equation}
for a predeclared rooted tree and admissible logarithms.  Here $\epsilon_c$
is the threshold of \Cref{thm:finite-Coulomb-normalization} for
$G_{\rm ss}$.  Then all the convergence and variational conclusions of
\Cref{thm:finite-Wilson-zero-defect} hold after actual $G_{\rm SM}$ site
gauges.  The full logarithmic connection is co-closed and has the form
\begin{equation}
 A_h=b_h+Z_c a_h,\qquad
 Z_c=\left(-\frac{\ii}{3}I_3,\frac{\ii}{2}I_2\right),\qquad
 \delta_h b_h=\delta_h a_h=0.
 \label{eq:determinant-combined-connection}
\end{equation}
Only $b_h$ needs to lie in the small critical ball.  Neither the determinant
curvature norm, the determinant potential $L_h^4$ norm, nor its flat cycle
holonomies is required to be small.

Alternatively, \eqref{eq:semisimple-only-certificate} may be replaced by
exact flatness of $V_h$.  In that case its semisimple cycle holonomies also
need not be small.  Zero determinant flux is a topological restriction of
the single periodic-coordinate conclusion, not a consequence of the
curvature screen.
\end{corollary}
\begin{proof}
\Cref{lem:determinant-screen-flux} first extracts the determinant curvature
without selecting a gauge.  On zero flux,
\Cref{prop:periodic-determinant-normalization} constructs $a_h$ with strong
$L^4$ and first-difference $L^2$ compactness.  The exact group split of
\Cref{prop:native-determinant-split} yields $V_h$ with its semisimple
curvature Wilson screen inherited from the original full record.  Under
\eqref{eq:semisimple-only-certificate},
\Cref{prop:rooted-gauge-certificate} supplies $b_h$ in small discrete
Coulomb gauge.  \Cref{thm:native-Wilson-compactness} gives strong
semisimple curvature compactness, and
\Cref{thm:native-discrete-Coulomb} gives strong $L^4$ and full
first-difference $L^2$ convergence of $b_h$.  When $V_h$ is flat, the
bounded constant representatives of \Cref{lem:flat-semisimple-normalization}
give these conclusions directly, without smallness.

Since $Z_c$ is central, adding $Z_ca_h$ preserves the Coulomb condition and
introduces no mixed curvature commutator.  The normalized actual links are
exactly $e^{h b_h}e^{hZ_ca_h}=e^{hA_h}$.  Their logarithms are resolved:
$h\norm{a_h}_\infty\to0$, and the semisimple part is either uniformly in
the small logarithm chart or bounded and constant.  The full $A_h$ is thus
bounded in $L_h^4$ and converges strongly there, although its norm need not
be small.  The full literal curvature is the sum of the semisimple
curvature and $Z_cf_h$.

The unchanged Higgs Wilson screen now gives strong Higgs-gradient
compactness by \Cref{thm:native-Wilson-compactness}, which requires only a
bounded full connection $L_h^2$ norm.  Apply
\Cref{prop:native-Higgs-compactness,prop:native-spinor-variation} and then
\Cref{thm:native-firstvariation-no-band}.  Both applied normalizations are
site transformations of $G_{\rm SM}$, so the complete physical covector norm
is unchanged.  Assumption \textup{(F5)} therefore gives all limiting Euler
identities.  The reconstructed consistency and stress conclusions follow
from the same strong-packet products as in the preceding theorem; no
small-Coulomb theorem is applied to the full, possibly large, $A_h$.
\end{proof}

\begin{remark}[Logical scope of the finite certificates]
\label{rem:native-dependency-scope}
The order of the argument is gauge normalization, compactness of the
literal packets, identification of the limiting connection and covariant
derivatives, and convergence of the complete first variation.  Physical
stationarity is used only in the last step.  Scalar-magnitude compactness,
which precedes gauge normalization in the determinant argument, follows
from unitarity alone.  The semisimple certificate is still an independent
finite hypothesis: the result is not a non-Abelian curvature-only gauge
construction.  The geometric and spinor conditions are unchanged.
Neither this corollary nor \Cref{thm:finite-Wilson-zero-defect} uses a
verified Cauchy tube, a source-drift law, or the high-order reader bounds
needed for the separate quantitative shadow theorem.
\end{remark}

\subsection{Gauge-intrinsic test budgets and Wilson-screen transfer}
\label{app:gauge-intrinsic-screen}

This subsection gives a regular reconstruction route to the equivariant
budgets of \Cref{thm:critical-quotient-defect} and compares finite Wilson
shifts with continuum parallel transport.  Its regularity-dependent
transfer estimates are not premises of the native closure proof in
\Cref{app:native-critical-closure,app:determinant-sector}.  In the
equivariant-budget argument a normalized continuum test is pulled back to
the original representation, with exact covariance expressed by
\eqref{eq:equivariant-test-size}.  This test-based argument is distinct
from the actual finite site-gauge normalizations used in those appendices.

\subsubsection{Equivariant consistency of the explicit local action}

For a metric bundle connection $\mathbb D$ on a finite direct sum of the
coframe, adjoint, Higgs and spinor test bundles, the ordinary cellwise
quadrature estimate has the following covariant form.

\begin{lemma}[Covariant cell-average quadrature]
\label{lem:covariant-quadrature}
Let $\mathbb D$ be a smooth metric bundle connection on the buffered chart.
For a cell $C_x=x+[0,h)^4$, use parallel transport along a fixed family of
paths inside the cell, each of length at most $Ch$, and define
\begin{equation}
 (\mathsf Q_h^{\mathbb D}v)(x)
    =h^{-4}\int_{C_x}P^{\mathbb D}_{x\leftarrow y}v(y)\,\dd y.
 \label{eq:covariant-test-average}
\end{equation}
All integrals and norms in this lemma use coordinate Lebesgue measure;
a fixed smooth comparison density can be incorporated into the Euler
density $X$.
For a smooth $X$ with bounded covariant derivative,
\begin{equation}
 \left|h^4\sum_x\ip{X(x)}{\mathsf Q_h^{\mathbb D}v(x)}
       -\int\ip Xv\,\dd x\right|
 \le Ch\norm{\mathbb D X}_{\infty}\norm v_1.
 \label{eq:covariant-quadrature}
\end{equation}
Moreover
$\norm{\mathsf Q_h^{\mathbb D}v}_{\infty,h}\le\norm v_\infty$ and
$\norm{\mathsf Q_h^{\mathbb D}v}_{2,h}\le\norm v_2$.
The formulas are equivariant under unitary bundle changes of frame.
\end{lemma}
\begin{proof}
Metric compatibility rewrites the difference as
\[
 \sum_x\int_{C_x}
 \ip{P^{\mathbb D}_{y\leftarrow x}X(x)-X(y)}{v(y)}\,\dd y.
\]
The covariant fundamental theorem of calculus along the chosen path bounds
the first factor by $Ch\norm{\mathbb D X}_\infty$.
Integrating proves \eqref{eq:covariant-quadrature}.
The two norm estimates follow from unitarity and Jensen's inequality.
Endpoint transformation of parallel transport proves equivariance.
\end{proof}

The averaging in \eqref{eq:covariant-test-average} is applied only to
physical tests, not to the observed fields or the finite action.
It avoids point evaluation on a four-dimensional $W^{1,2}$ test class;
a nodal quadrature estimate cannot be inferred from that norm alone.
The native $C^r$ consistency theorem uses the separately specified smooth
nodal test lift and does not require this averaging construction.

Let $B\ge1$ be a derivative scale and suppose a reconstructed field tuple
satisfies
\begin{equation}
 \norm{\partial^\alpha z}_{L^\infty(Q')}
 \le C_qB^{|\alpha|},
 \qquad |\alpha|\le q,\quad q\ge5,
 \qquad hB\le c_{\rm res}.
 \label{eq:eq-growing-band}
\end{equation}
The complete continuum Euler density then has differential degree at most two
in the bosonic rows and one in the Dirac rows, so on the fixed chart
\begin{equation}
 \norm{\mathcal E_0(z)}_{L^2}\le CB^2,
 \qquad
 \norm{\mathbb D_z\mathcal E_0(z)}_{L^\infty}\le CB^3.
 \label{eq:eq-euler-jets}
\end{equation}
For the physical coframe lift of a test $v$,
\begin{equation}
 \norm{\widetilde v_z}_{L^\infty}
 \le C\mathfrak n_z(v),
 \qquad
 \norm{\mathbb D_z\widetilde v_z}_{L^2}
 \le CB\mathfrak n_z(v).
 \label{eq:eq-test-jets}
\end{equation}
The second estimate follows because differentiating the symmetric coframe lift
introduces at most one first coframe derivative, while the matter components
are controlled directly by \eqref{eq:equivariant-test-size}.

For the gauge links it is convenient to use physical left-link tangents.  If
$U_\mu=e^{hA_\mu}$ and
$\dot U_\mu U_\mu^{-1}=h\alpha_\mu$, let
\begin{equation}
 \norm{V_h}_{0,h,{\rm link}}^2
 :=h^4\sum_x\left(
 |\dot e|^2+|\alpha|^2+|\eta_H|^2
 +|\eta_\Psi|^2+|\eta_{\bPsi}|^2\right)
 \label{eq:link-tangent-mass}
\end{equation}
with the fixed positive coefficient metrics, and let
$\sigma_h^{\rm link}$ be the dual norm of $D S_h^{\rm loc}$ in this metric.
On every compact logarithm chart the differential of the exponential map is
uniformly invertible, so $\sigma_h^{\rm link}$ is uniformly equivalent to
the logarithmic-coordinate row norm used in
\eqref{eq:native-sigma}.

\begin{proposition}[Equivariant native first-variation budgets]
\label{prop:equivariant-native-budgets}
Under \eqref{eq:eq-growing-band}, the unchanged local action of
\Cref{app:native-action} and the covariant cell-average physical test lift
$\mathcal I_h^{\rm cov}$ defined below satisfy
\begin{align}
 \abs{D S_h^{\rm loc}[\mathcal I_h^{\rm cov}v]
      -D\Sact_\theta(z)[v]}
 &\le ChB^3\mathfrak n_z(v),
 \label{eq:eq-native-consistency}\\
 \abs{D S_h^{\rm loc}[\mathcal I_h^{\rm cov}v]}
 &\le C\sigma_h^{\rm link}\mathfrak n_z(v).
 \label{eq:eq-native-stationarity}
\end{align}
Thus the budgets \eqref{eq:equivariant-consistency}--
\eqref{eq:equivariant-stationarity} may be chosen with
\begin{equation}
 c_h^{\rm eq}\le ChB^3,
 \qquad
 \varepsilon_h^{\rm eq}\le C\sigma_h^{\rm link}.
 \label{eq:eq-native-budgets}
\end{equation}
The link-row norm is exactly invariant under finite site gauge transformations.
\end{proposition}
\begin{proof}
Express the complete raw Euler covector in the physical coframe,
logarithmic-link and matter coordinates of \Cref{prop:native-consistency}.
Let $\widetilde v_z$ be the corresponding continuum physical-coordinate
test, including its symmetric coframe lift.  Choose the positive comparison
connection $\mathbb D_z$ acting on these coordinate bundles and define the
finite coordinate test by $\mathsf Q_h^{\mathbb D_z}\widetilde v_z$.
Its gauge-coordinate component $a_h$ is converted to the actual link tangent
by $h^{-1}\dot U U^{-1}=\mathcal J(h\operatorname{ad}_A)a_h$.
This specifies $\mathcal I_h^{\rm cov}$ as a linear physical test lift;
no finite field is replaced by a cell average.

Pair the $O(hB^3)$ pointwise Euler-row consistency estimate with this
averaged test.  Jensen and \eqref{eq:eq-test-jets} bound the contribution
by $ChB^3\mathfrak n_z(v)$.  For the remaining comparison apply
\Cref{lem:covariant-quadrature} to $X=\mathcal E_0(z)$ and
$\widetilde v_z$.  Equation \eqref{eq:eq-euler-jets} gives
$\norm{\mathbb D_z X}_\infty\le CB^3$, and on the fixed box
$\norm{\widetilde v_z}_1\le C\mathfrak n_z(v)$.  Thus this term has the
same bound, proving \eqref{eq:eq-native-consistency} uniformly even when
the ordinary derivatives of the pulled-back test are large.

The averaged coordinate lift has mass norm at most $C\mathfrak n_z(v)$.
The logarithmic-to-left-link Jacobian is uniformly bounded on the fixed
scaled logarithm chart, so the same bound holds in the physical link/matter
mass metric.  Cauchy--Schwarz proves \eqref{eq:eq-native-stationarity}.
Exact site-gauge invariance of the finite action and the isometry of its
physical tangent map preserve the complete dual link-row norm.
Uniform invertibility of the same Jacobian proves the claimed equivalence
with the logarithmic-coordinate row norm.  This last invariance is a
statement about the finite covector, independent of the auxiliary test
averaging.
\end{proof}

\subsubsection{Finite Wilson packets and continuum covariant translations}

For a continuum connection $A$ and a coordinate displacement $s e_\mu$, let
$P^A_{\mu,s}(x)$ denote parallel transport from $x+s e_\mu$ back to $x$ and
put
\begin{equation}
 \mathscr T^A_{\mu,s}Y(x)
 :=\rho(P^A_{\mu,s}(x))Y(x+s e_\mu).
 \label{eq:continuum-covariant-translation}
\end{equation}
This operation is exactly internal-gauge covariant.

\begin{proposition}[Finite packet and open-link transfer]
\label{prop:Wilson-packet-transfer}
Assume \eqref{eq:eq-growing-band} and $hB_h^2\to0$.  For the literal
packets \eqref{eq:literal-finite-bosonic-packets},
\begin{align}
 \max_x\abs{\mathbb F_h(x)-F_{A_h}(x)}
 +\max_x\abs{\mathbb K_h(x)-D_{A_h}H_h(x)}
 &\le ChB_h^2,
 \label{eq:Wilson-packet-consistency}
\end{align}
and for $|m|h\le\varrho_0$,
\begin{equation}
 \left\|\rho(\mathcal U_{\mu,m}(x))
       -\rho(P^{A_h}_{\mu,mh}(x))\right\|
 \le C\varrho_0 hB_h.
 \label{eq:Wilson-transport-consistency}
\end{equation}
Consequently, if a literal packet $Y_h^{\rm d}$ obeys
\begin{equation}
 \lim_{\varrho\downarrow0}\limsup_{h\downarrow0}
 \Omega_h(Y_h^{\rm d};\varrho)=0,
 \label{eq:Wilson-finite-screen}
\end{equation}
then its continuum packet $Y_h$ satisfies, on every smaller chart,
\begin{equation}
 \lim_{\varrho\downarrow0}\limsup_{h\downarrow0}
 \max_\mu\sup_{|s|\le\varrho}
 \norm{\mathscr T^{A_h}_{\mu,s}Y_h-Y_h}_{L^2}=0.
 \label{eq:Wilson-continuum-screen}
\end{equation}
\end{proposition}
\begin{proof}
The plaquette expansion of \Cref{lem:native-plaquette} and the
amplitude-preserving scaling of \Cref{lem:native-scaling} give
\[
 U_{\alpha\beta}(x)
 =I+h^2F_{A_h,\alpha\beta}(x)+O(h^3B_h^2),
\]
which yields the curvature part of
\eqref{eq:Wilson-packet-consistency} after applying the analytic logarithm.
Covariant Taylor expansion of the Higgs link gives the second part.

For one edge, Duhamel's formula compares the exact parallel-transport ODE with
the frozen coefficient entering $U_\mu(x)$; the error is $O(h^2B_h)$.
Unitary factors do not amplify operator norm, so telescoping over $|m|$ edges
gives \eqref{eq:Wilson-transport-consistency}.

For a lattice displacement $s=mh$, replace both endpoint continuum packets by
the finite readers and continuum transport by the open finite link.  The
resulting errors are $O(hB_h^2)$ and
$O(\varrho hB_h)\norm{Y_h^{\rm d}}_2$; the remaining term is exactly the
Wilson difference in \eqref{eq:finite-Wilson-modulus}.  For arbitrary $s$
choose $m$ with $|s-mh|\le h$.  The growing-band reserve controls the final
sub-cell displacement by $O(hB_h^2)$.  Taking first $h\downarrow0$ and then
$\varrho\downarrow0$ proves \eqref{eq:Wilson-continuum-screen}.
\end{proof}

\begin{lemma}[Translation Kato inequality]
\label{lem:translation-Kato}
For a unitary bundle representation,
\begin{equation}
 \bigl|\,|Y(x+s e_\mu)|-|Y(x)|\,\bigr|
 \le\abs{\mathscr T^A_{\mu,s}Y(x)-Y(x)}.
 \label{eq:translation-Kato}
\end{equation}
Thus a vanishing covariant $L^2$ translation modulus controls the ordinary
translation modulus of the scalar magnitude $|Y|$.
\end{lemma}
\begin{proof}
Parallel transport preserves the fiber norm, and the result is the reverse
triangle inequality.
\end{proof}

\begin{lemma}[Compact scalar magnitude implies energy uniform integrability]
\label{lem:compact-magnitude-ui}
If a family is relatively compact in $L^2(K)$ on a finite-measure region,
then the squares of its elements are uniformly integrable:
\begin{equation}
 \lim_{|E|\downarrow0}\sup_f\int_E|f|^2=0.
 \label{eq:compact-magnitude-ui}
\end{equation}
\end{lemma}
\begin{proof}
Choose a finite $L^2$ net.  Absolute continuity of the integral for the
finitely many net elements and
$\int_E|f|^2\le2\norm{f-f_j}_2^2+2\int_E|f_j|^2$ give the claim.
\end{proof}

\begin{proposition}[Covariant screen compactness at the critical endpoint]
\label{prop:Wilson-screen-compactness}
Let $Y_h$ be bounded in $L^2(B)$ and satisfy
\eqref{eq:Wilson-continuum-screen}.  Then $|Y_h|$ is relatively compact in
$L^2_{\rm loc}$ and $|Y_h|^2$ is uniformly integrable locally.  If, in
addition, after local gauge normalization
\begin{equation}
 \sup_h\norm{A_h}_{L^4(B)}\le M_A,
 \label{eq:critical-A4-bound}
\end{equation}
then $Y_h$ itself is precompact in $L^2_{\rm loc}(B)$.
\end{proposition}
\begin{proof}
By \Cref{lem:translation-Kato}, the scalar magnitudes have a vanishing
ordinary translation modulus.  Their $L^2$ boundedness and the boundedness of
the chart allow the Fr\'echet--Kolmogorov criterion, proving relative
$L^2$ compactness.  Apply \Cref{lem:compact-magnitude-ui} to obtain uniform
integrability of $|Y_h|^2$.

It remains to remove the parallel transport.  The transport ODE gives
\begin{equation}
 \norm{\rho(P^{A_h}_{\mu,s}(x))-I}
 \le C\min\{1,L_{h,s}(x)\},
 \qquad
 L_{h,s}(x)=|s|\int_0^1
 |A_h(x+\theta s e_\mu)|\,d\theta.
 \label{eq:short-line-integral}
\end{equation}
Minkowski gives $\norm{L_{h,s}}_4\le |s|M_A$.  For
$E_{h,s,\delta}=\{L_{h,s}>\delta\}$, Chebyshev gives
$|E_{h,s,\delta}|\le M_A^4|s|^4\delta^{-4}$.  On the complement, the square
of the transport error is bounded by $C\delta^2|Y_h|^2$; on the exceptional
set use the uniform bound on the unitary transport and the uniform
integrability just proved.  Hence
\begin{equation}
 \sup_h\norm{(\rho(P^{A_h}_{\mu,s})-I)
 Y_h(\cdot+s e_\mu)}_2\longrightarrow0
 \quad(s\to0).
 \label{eq:short-transport-harmless}
\end{equation}
Combining this with the covariant screen gives an ordinary vector translation
modulus for $Y_h$.  Fr\'echet--Kolmogorov now yields local strong $L^2$
precompactness.
\end{proof}

\begin{lemma}[Critical Coulomb Hodge absorption]
\label{lem:critical-Coulomb-Hodge}
On a fixed small-energy Coulomb ball $B$ suppose $A_h,A$ are co-closed,
$A_h\to A$ in $L^2(B)$, their $L^4(B)$ norms are uniformly sufficiently
small, and $F_{A_h}\to F_A$ in $L^2(B)$.  Then, for every $B'\Subset B$,
\begin{equation}
 A_h\to A\qquad\text{strongly in }L^4(B').
 \label{eq:critical-A4-strong}
\end{equation}
\end{lemma}
\begin{proof}
Choose $\chi\in C_c^\infty(B)$ with $\chi=1$ on $B'$ and put
$a_h=A_h-A$.  Apply the compactly supported Hodge and Sobolev estimates to
$\chi a_h$, extended by zero.  Since $d^*a_h=0$, differentiation of the
cutoff produces only terms bounded by $C_\chi\norm{a_h}_2$.
The curvature identity therefore gives
\[
 \norm{\chi a_h}_4
 \le C\norm{F_{A_h}-F_A}_2+C_\chi\norm{a_h}_2
      +C(\norm{A_h}_4+\norm A_4)\norm{\chi a_h}_4.
\]
Choose the small-energy threshold to make the last coefficient less than
one half.  Absorb that term and pass to the limit.  No unrestricted Hodge
estimate up to the boundary of $B$ is being used.
\end{proof}

\section{Quantitative consistency and physical forcing estimates}
\label{app:quantitative-residuals}

This appendix proves the two finite-action estimates used in
\Cref{sec:native-branch}. Their assumptions concern the full reconstructed
record and the comparison head on the same buffered chart. The qualitative
Wilson theorem does not use this growing-band route.

\subsection{Consistency and forcing statements}

\begin{proposition}[Growing-band consistency of the unchanged finite action]
\label{prop:native-consistency}
Let $y$ be a smooth physical field tuple on $Q'$ with uniformly bounded head amplitudes and inverse metric, and assume
\begin{equation}
 \norm{\partial^\alpha y}_\infty\le C_qK^{|\alpha|},
 \qquad |\alpha|\le q,\quad q\ge5,
 \qquad hK\le c_{\rm res}.
 \label{eq:growing-derivatives}
\end{equation}
Then the local finite action of \Cref{app:native-action} satisfies
\begin{equation}
 \max_{x\in Q'_h}\left|E_h^{\rm raw}(\mathsf S_hy)(x)
 -\mathsf S_h\mathcal E_0(y)(x)\right|\le ChK^3,
 \label{eq:native-consistency}
\end{equation}
where $\mathcal E_0$ is the complete minimally coupled torsion-free Einstein--Standard-Model Euler density in the same physical coordinates.  The literal Cartan plaquette curvature also obeys
\begin{equation}
 \norm{I_h^0R_h^{\rm raw}(\mathsf S_hy)-\Riem(g)}_{L^2(Q),+}
 \le ChK^3.
 \label{eq:native-Cartan}
\end{equation}
The constants depend on the bounded chart but not on $h$ or $K$.
\end{proposition}
\begin{proof}
The proof is given in \Cref{app:native-scaling}.  Its central point is an amplitude-preserving rescaling $\xi=Kx$, $\rho=hK$, $\nu=K^{-1}$.  The exact local density splits as $K^2$ times a uniformly regular bosonic stencil plus $K$ times a uniformly regular Dirac stencil.  First-variation consistency at normalized mesh $\rho$ is $O(\rho)$, giving $O(hK^3)$ after restoring the physical derivative factors.  No field amplitude is divided by $K$ and the metric--spinor chain is included.
\end{proof}

\begin{theorem}[Leakage-tolerant physical forcing estimate]
\label{thm:native-source}
Assume the complete record and its low-frequency comparison head lie in the fixed calibrated field-value chart on the buffered box, with uniform amplitude and inverse-metric margins.  Assume $hK\le c_{\rm res}$, $\tau_{h,m}(K)\le\tau_*$, and $\varepsilon_{c,h}\le1$, with $\tau_*$ sufficiently small for these margins.  For the complete, unfiltered reconstruction $z_h=\mathcal I_h^{\rm trig}u_h$,
\begin{align}
 \epsilon^{\rm phys}(z_h)
 &:={\norm{\mathcal R_{\rm B}(z_h)}_{L^2(Q)}}
   +{\norm{\mathcal R_{\rm D}(z_h)}_{L^2(Q)}}
 \le C\varepsilon_{0,h},\label{eq:native-zero-source}\\
 \mathcal Y_k(z_h)
 &:={\norm{\mathcal R_{\rm B}(z_h)}_{L^2_tH_x^k(Q)}}
   +{\norm{\mathcal R_{\rm D}(z_h)}_{L^2_tH_x^{k+1}(Q)}}
 \le C F_{h,k},\label{eq:native-strong-source}\\
 \norm{I_h^0R_h^{\rm raw}(u_h)-\Riem(g_h)}_{L^2(Q),+}
 &\le ChK^3.\label{eq:native-full-Cartan}
\end{align}
Here $\mathcal R_{\rm B}$ collects the Einstein, Yang--Mills and Higgs physical Euler rows, while $\mathcal R_{\rm D}$ contains the Dirac and dual-Dirac rows.  No filtering is applied to the field entering these residuals or to the finite action.
\end{theorem}
\begin{proof}
\Cref{app:native-tails} proves the finite-order tail transfer.  The low-frequency proof field has the same chart margins and a complex spatial strip of width $a/K$; \Cref{lem:log-source-upgrade} upgrades its small $L^2$ source to the differentiated source norm with only the logarithmic factor in \eqref{eq:native-forcing-budget}.  The full measured tail is then restored at cost $K^{k+2}\tau_{h,m}$.  The zero-order estimate follows from \Cref{prop:native-consistency,lem:nodal-source-transfer} applied to the full reconstructed field, whose first five derivatives satisfy \eqref{eq:growing-derivatives} by the tail bounds.  The same cellwise comparison gives the Cartan estimate.
\end{proof}

\subsection{Regular plaquette expansion}

\begin{lemma}[Regularity of the plaquette map]
\label{lem:native-plaquette}
For bounded matrices $X,Y,a,b$ the map
\[
 \Phi(h;X,Y,a,b)=h^{-2}\log(e^{hX}e^{h(Y+ha)}e^{-h(X+hb)}e^{-hY})
\]
extends analytically through $h=0$ with $\Phi(0)=a-b+[X,Y]$.  All fixed-order derivatives are uniformly bounded on compact argument sets for sufficiently small $h$.
\end{lemma}
\begin{proof}
The product inside the logarithm is $I+h^2(a-b+[X,Y])+O(h^3)$.  Since the logarithm is analytic near the identity and the linear term vanishes, the removable quotient may be written as an integral of the second $h$ derivative.  This also proves uniform bounds for parameter derivatives.
\end{proof}

\subsection{Amplitude-preserving scaling}
\label{app:native-scaling}

\begin{lemma}[Differential-order scaling]
\label{lem:native-scaling}
Under \eqref{eq:growing-derivatives}, put $\rho=hK$, $\nu=K^{-1}$ and $\widetilde y(\xi)=y(\xi/K)$, without rescaling the field amplitudes.  The exact density \eqref{eq:native-densities} has the form
\begin{equation}
 \mathcal L_h(y)=K^2\widetilde{\mathcal L}^{\rm B}_{\rho,\nu}(\widetilde y)
 +K\widetilde{\mathcal L}^{\rm D}_{\rho,\nu}(\widetilde y),
 \label{eq:native-density-scaling}
\end{equation}
where both normalized stencil densities and their required finite jet derivatives are uniformly regular for small $\rho$ and $0\le\nu\le1$.  The normalized Euler rows carry the same factors $K^2$ and $K$.
\end{lemma}
\begin{proof}
The algebraic Cartan reader is linear in a first difference, so $\omega_h=K\widetilde\omega_\rho$ and $e^{h\omega_h}=e^{\rho\widetilde\omega_\rho}$.  The Cartan plaquette therefore has factor $K^2$.  For the gauge plaquette use \Cref{lem:native-plaquette} with arguments $\nu A$; the removable quotient shows $F^h=K\widetilde F_{\rho,\nu}$, so the Yang--Mills density has factor $K^2$.  The Higgs link is
\[
 K^h_\mu=K\frac{e^{\rho\nu\rho_H(A_\mu)}H(\xi+\rho e_\mu)-H(\xi)}\rho,
\]
which is a regular first-jet function times $K$.  The spin link is
$e^{\rho(\sigma(\widetilde\omega_\rho)+\nu\rho_S(A))}$, so each Dirac difference has factor $K$ and the bilinear density has factor $K$.  Potential, mass and cosmological coefficients appear with nonnegative powers of $\nu$ inside the normalized densities.  Since the varied head amplitudes are not rescaled, differentiation of the finite sum leaves exactly the factors in \eqref{eq:native-density-scaling}.
\end{proof}

\begin{proof}[Proof of \Cref{prop:native-consistency}]
At normalized mesh $\rho$, Taylor expansion of the regular shifted-jet densities and their first variations gives an $O(\rho)$ Euler error with constants uniform in $\nu$.  Restoring the factors of \Cref{lem:native-scaling} gives
$C(K^2+K)\rho\le ChK^3$.  The same normalized plaquette expansion gives $O(\rho)$ for Cartan curvature, which scales by $K^2$, proving \eqref{eq:native-Cartan}.  The metric row differentiates the Cartan reader and the spin link throughout; no connection is frozen in this argument.
\end{proof}

\begin{lemma}[From nodal Euler bounds to a physical source norm]
\label{lem:nodal-source-transfer}
Let $f$ be a $C^1$ vector-valued function on the buffered slab $Q'$ and let
$I_h^0\mathsf S_hf$ be its piecewise-constant reconstruction on cells
meeting $Q\Subset Q'$.  For all sufficiently small $h$,
\begin{equation}
 \norm{f}_{L^2(Q)}
 \le C\norm{\mathsf S_hf}_{0,h;Q'}
       +Ch\norm{\partial f}_{L^\infty(Q')}.
 \label{eq:nodal-source-transfer}
\end{equation}
For the complete physical Euler density of a field satisfying
\eqref{eq:growing-derivatives}, the last term is $O(hK^3)$.
Thus a nodal consistency estimate and a finite mass-norm Euler bound give
a continuum $L^2$ source estimate without an assumption that the nonlinear
Euler density is band limited.
\end{lemma}
\begin{proof}
On a cell based at $x$, the fundamental theorem of calculus gives
$|f(y)-f(x)|\le Ch\norm{\partial f}_\infty$.
Square, integrate over the fixed-volume union of cells, and use the exact
mass normalization of the piecewise-constant reconstruction.
The bosonic Euler density has differential degree at most two and the
Dirac density at most one; differentiating once and applying
\eqref{eq:growing-derivatives} bounds its first derivative by $CK^3$ on
the fixed coefficient chart.  Bounded invertible algebraic changes from
Euler densities to physical rows preserve the asserted $L^2$ estimate.
\end{proof}


\subsection{Measured spectral leakage}
\label{app:native-tails}

For the proof only, write $u=P_{\le K}u+r$ and $z_h^{\rm lo}=\mathcal I_h^{\rm trig}P_{\le K}u_h$.  The complete record remains $u_h$.

\begin{lemma}[Finite-order tail bounds]
\label{lem:native-tail-transfer}
For $0\le j\le m=k+2$,
\begin{equation}
 \max_{|\alpha|\le j}K^{-|\alpha|}\norm{\partial^\alpha(z_h-z_h^{\rm lo})}_\infty
 \le C_j\tau_{h,j}(K).
 \label{eq:native-tail-derivatives}
\end{equation}
For sufficiently small $hK$ and $\tau_{h,m}$ the full and low-frequency fields stay in the same nondegenerate chart.  Moreover
\begin{align}
 \norm{E_h^{\rm raw}(u_h)-E_h^{\rm raw}(\mathsf S_hz_h^{\rm lo})}_{0,h;Q'}
 &\le CK^2\tau_{h,2},\label{eq:native-Euler-tail}\\
 \norm{\mathcal R_{\rm B}(z_h)-\mathcal R_{\rm B}(z_h^{\rm lo})}_{L^2_tH_x^k}
 +\norm{\mathcal R_{\rm D}(z_h)-\mathcal R_{\rm D}(z_h^{\rm lo})}_{L^2_tH_x^{k+1}}
 &\le CK^{k+2}\tau_{h,m}.\label{eq:native-source-tail}
\end{align}
The corresponding zero-order physical source difference is bounded by $CK^2\tau_{h,2}$.
\end{lemma}
\begin{proof}
Differentiate the finite Fourier tail and use
$|\ell^\alpha|\le K^{|\alpha|}(1+|\ell|_1/K)^j$.  Smooth algebraic dependence of the metric, inverse metric, frame and finite representation coefficients on the field heads gives \eqref{eq:native-tail-derivatives}.  The finite Euler row is, after the exact shifted-first-jet reduction in \Cref{lem:native-scaling}, a smooth function of finitely many shifted jets through second order with overall bosonic factor at most $K^2$; the mean-value theorem gives \eqref{eq:native-Euler-tail}.  The continuum bosonic Euler rows have differential degree at most two and the Dirac rows degree one.  Differentiate their difference $k$ and $k+1$ spatial times, respectively; every product contains at least one tail factor and at most $k+2$ total field derivatives.  Equation \eqref{eq:native-tail-derivatives} gives \eqref{eq:native-source-tail}.
\end{proof}

\begin{lemma}[Logarithmic derivative upgrade for a band-limited source]
\label{lem:log-source-upgrade}
Let $f(t,x)$ extend holomorphically in the three spatial variables to width $a/K$, be bounded there by $A$, and satisfy $\norm f_{L^2(Q)}\le\epsilon\le1$.  Then for every fixed integer $j\ge1$,
\begin{equation}
 \norm f_{L^2_tH_x^j(Q)}
 \le C_j\epsilon K^j
 \left[1+\log\left(2+\frac{C_jAK^{j+3}}\epsilon\right)\right]^j.
 \label{eq:log-source-upgrade}
\end{equation}
When $\epsilon=0$ the right side is interpreted as zero.
\end{lemma}
\begin{proof}
Spatial contour shifting gives $|\widehat f(t,\ell)|\le Ae^{-a|\ell|_1/K}$.  For a frequency threshold $B$, Plancherel bounds the low frequencies by $C_j(1+B)^j\epsilon$, while the exponentially small tail is bounded by $C_jAK^{j+3}e^{-aB/(2K)}$.  Choosing $B$ proportional to $K[1+\log(2+C_jAK^{j+3}/\epsilon)]$ proves the claim.
\end{proof}

\begin{proof}[Proof of \Cref{thm:native-source}]
The full field satisfies the growing $C^5$ reserve by \Cref{lem:native-tail-transfer}; hence \Cref{prop:native-consistency} and \eqref{eq:native-sigma}, followed by the nodal-to-continuum estimate \Cref{lem:nodal-source-transfer}, give \eqref{eq:native-zero-source} for the unfiltered field.  For the differentiated source, \eqref{eq:native-Euler-tail} first transfers the small finite row to the low-frequency field at cost $K^2\tau_{h,2}$.  By Bernstein estimates the low-frequency field changes by at most a fixed constant times $K|\operatorname{Im}x|$ in a complex spatial strip.  Taking its width $a/K$ with $a$ small relative to the declared field-value and inverse-metric margins preserves the analytic branches of every algebraic coefficient.  Complexification is in the real field coordinates, including real and imaginary spinor and Higgs components.  Thus the bosonic and Dirac source amplitudes in this strip are bounded by $CK^2$ and $CK$.  \Cref{lem:nodal-source-transfer} supplies the needed continuum $L^2$ bound for these nonlinear sources before the logarithmic upgrade.  Apply \Cref{lem:log-source-upgrade} with $j=k$ and $j=k+1$.  Finally restore the measured full tail by \eqref{eq:native-source-tail}.  This is exactly the budget \eqref{eq:native-forcing-budget}.  The Cartan assertion is the second estimate of \Cref{prop:native-consistency} applied to the full field.
\end{proof}

\begin{proposition}[A sufficient nodal-accuracy condition]
\label{prop:native-rounding}
Let $u_h^b$ have Fourier support in $|\ell|_\infty\le K$ and let $u_h=u_h^b+e_h$ with $\max_x|e_h(x)|\le e$, preserving the affine field-chart constraints.  Then
\begin{equation}
 \tau_{h,j}(K)\le C_j e\,h^{-2}(hK)^{-j}.
 \label{eq:native-rounding-tail}
\end{equation}
In particular, with $m=k+2$, the sufficient tolerance
\begin{equation}
 e\le c\,h^{m+3}K^{m+2}
 \label{eq:native-rounding-tolerance}
\end{equation}
implies $\tau_{h,m}=O(hK^2)$ and, on the rate families of \Cref{cor:native-rate}, $\tau_{h,2}=O(hK)$.
\end{proposition}
\begin{proof}
Parseval for the normalized four-dimensional discrete Fourier transform gives $\sum_\ell|\widehat e_h(\ell)|^2\le e^2$.  There are $O(h^{-4})$ Fourier modes, so Cauchy--Schwarz bounds the absolute coefficient sum by $Ch^{-2}e$.  On the grid spectrum $(1+|\ell|_1/K)^j\le C_j(hK)^{-j}$, which proves \eqref{eq:native-rounding-tail}.  Substitute \eqref{eq:native-rounding-tolerance} for the final assertion.
\end{proof}

\subsection{An explicit cofinal rate}

\begin{corollary}[An explicit cofinal rate]
\label{cor:native-rate}
Suppose $\sigma_h=O(h)$, $K_h\asymp h^{-\beta}$ with $0<\beta<1/(k+4)$, and
\begin{equation}
 \tau_{h,2}=O(hK_h),\qquad
 \tau_{h,k+2}=O(hK_h^2).
 \label{eq:native-tail-rate}
\end{equation}
Then
\[
 F_{h,k}=O\!\left(h^{1-\beta(k+4)}(1+|\log h|)^{k+1}\right),
 \qquad
 \epsilon^{\rm phys}(z_h)=O(h^{1-3\beta}).
\]
For example,
\begin{equation}
 k=4,\qquad K_h\asymp h^{-1/18},\qquad
 i_{h,4}+\gamma_{h,4}=O(h^{1/2})
 \label{eq:native-example-rate}
\end{equation}
with \eqref{eq:native-tail-rate} gives $O(h^{1/2})$ state, curvature and Standard-Model stress convergence; the strong forcing is $O(h^{5/9}(1+|\log h|)^5)$ and the zero-order physical and Cartan residuals are $O(h^{5/6})$.  Along a dyadic sequence the adjacent state, curvature and stress differences are summable, so cofinal transport is a conclusion on this branch.
\end{corollary}
\begin{proof}
The hypotheses give $\varepsilon_{c,h}=O(hK_h^3)$ and
$L_{h,k}=O(1+|\log h|)$. Insert these bounds in
\eqref{eq:native-forcing-budget} and apply \Cref{thm:native-closure}.
For $k=4$, $\beta=1/18$, the exponents are $5/9$ and $5/6$ as stated;
comparison of adjacent records through the same reference gives dyadic
summability.
\end{proof}

\section{Initial-constraint retraction and shadow refinements}
\label{app:initial-shadow-refinements}

The exact constraint reduction and verified Cauchy tube are defined in
\Cref{subsec:aposteriori-shadow}. We first solve the retained exact
constraint map, then distinguish evaluation and complementary-correction
errors, and finally derive the full-record forcing interpolation.

\subsection{Exact and validated initial correction}

\begin{proposition}[Uniform initial-constraint retraction]
\label{prop:initial-constraint-retraction}
Assume \Cref{def:exact-initial-reduction}. Suppose
\begin{equation}
 A_h:=D_\lambda\Phi_h(x_h,0),\qquad
 \sigma_{\min}(A_h)\ge\gamma_*>0,
 \label{eq:constraint-linear-margin}
\end{equation}
and, on $|\lambda|\le r_*$,
\begin{equation}
 \norm{D_\lambda\Phi_h(x_h,\lambda)-A_h}\le\frac{\gamma_*}{2}.
 \label{eq:constraint-derivative-margin}
\end{equation}
If
\begin{equation}
 c_h:=|\Phi_h(x_h,0)|\le\beta_h,\qquad
 2\gamma_*^{-1}\beta_h\le r_*,
 \label{eq:constraint-small-residual}
\end{equation}
there is a unique zero $\lambda_h$ of the reduced exact map in the closed
ball of radius $2\gamma_*^{-1}\beta_h$. The corrected datum
$d_h=\iota_h(x_h,\lambda_h)$ belongs to $\mathcal C_k$ and satisfies
\begin{equation}
 \norm{d_h-\mathcal U(\widehat z_h)(0)}_{H^k}
       \le2B_*\gamma_*^{-1}\beta_h.
 \label{eq:initial-data-retraction}
\end{equation}
Only the initial comparison datum is corrected; the observed spacetime
record is not projected or filtered. Membership in the regular Cauchy tube
is a separate condition in the shadow theorem.
\end{proposition}
\begin{proof}
Write
$\Phi_h(x_h,\lambda)=\Phi_h(x_h,0)+A_h\lambda+N_h(\lambda)$.
Then $N_h(0)=0$ and
$\norm{N_h(\lambda)-N_h(\mu)}\le(\gamma_*/2)|\lambda-\mu|$ on the ball.
The frozen-inverse map
\[
 \mathcal T_h(\lambda)
 =-A_h^{-1}\Phi_h(x_h,0)-A_h^{-1}N_h(\lambda)
\]
is $1/2$-Lipschitz. On $|\lambda|\le2\gamma_*^{-1}\beta_h$ it satisfies
\[
 |\mathcal T_h(\lambda)|
 \le\gamma_*^{-1}\beta_h+\tfrac12|\lambda|
 \le2\gamma_*^{-1}\beta_h.
\]
Banach contraction gives a unique fixed point in this ball, and invertibility
of $A_h$ identifies it with a zero of the exact reduced map. The exact
factorization, not the contraction alone, now gives
$\mathfrak C_k(d_h)=\jmath_h\Phi_h(x_h,\lambda_h)=0$.
Finally \eqref{eq:exact-initial-basepoint} and
\eqref{eq:exact-initial-reader-Lipschitz} give
\eqref{eq:initial-data-retraction}. All constants are independent of $m_h$
when the displayed physical norm margins are uniform.
\end{proof}

\begin{corollary}[Validated constraint evaluation with an explicit error budget]
\label{cor:validated-initial-retraction}
Assume the exact family and factorization of
\Cref{def:exact-initial-reduction}. Let $\widetilde\Phi_h$ be an evaluated
$C^1$ reduced map and put
$\widetilde A_h=D_\lambda\widetilde\Phi_h(x_h,0)$.
Suppose, on the correction ball,
\begin{align}
 \sigma_{\min}(\widetilde A_h)&\ge\gamma_*,\notag\\
 \sup_\lambda\norm{D_\lambda\widetilde\Phi_h(x_h,\lambda)
                          -\widetilde A_h}&\le\gamma_*/4,\notag\\
 \sup_\lambda\norm{D_\lambda\Phi_h(x_h,\lambda)
                 -D_\lambda\widetilde\Phi_h(x_h,\lambda)}
                         &\le\gamma_*/4,\notag\\
 |\Phi_h(x_h,0)-\widetilde\Phi_h(x_h,0)|&\le\zeta_{0,h}.
 \label{eq:validated-initial-margins}
\end{align}
Then the retraction and displacement conclusions of
\Cref{prop:initial-constraint-retraction} hold with the effective
residual bound
\begin{equation}
 \beta_h^{\rm eff}=|\widetilde\Phi_h(x_h,0)|+\zeta_{0,h},
 \qquad 2\gamma_*^{-1}\beta_h^{\rm eff}\le r_*.
 \label{eq:validated-initial-budget}
\end{equation}
In every shadow estimate using this evaluated map, $\beta_h$ is replaced
by $\beta_h^{\rm eff}$. In particular an exact zero of
$\widetilde\Phi_h$ does not remove $\zeta_{0,h}$.
\end{corollary}
\begin{proof}
The two derivative bounds give
$\sup_\lambda\norm{D_\lambda\Phi_h(x_h,\lambda)-\widetilde A_h}
\le\gamma_*/2$. Use the frozen matrix $\widetilde A_h$ in the preceding
contraction proof, applied to the exact map $\Phi_h$. The initial residual
has norm at most $\beta_h^{\rm eff}$. The fixed point therefore solves the
exact constraints, with displacement at most
$2B_*\gamma_*^{-1}\beta_h^{\rm eff}$.
\end{proof}

\begin{remark}[Initial normalization and complementary corrections]
\label{rem:initial-correction-bookkeeping}
The base point in \eqref{eq:exact-initial-basepoint} is the actual initial
jet of the normalized record. Fixing the initial hypersurface and the first
differential of a coordinate transformation is not used to identify all
metric or spinor first jets before and after that transformation.
If a preliminary complementary solve instead supplies a base point with
\[
 \norm{\iota_h(x_h,0)-\mathcal U(\widehat z_h)(0)}_{H^k}
       \le\zeta_h^{\rm comp},
\]
the triangle inequality replaces the initial bound by
$\zeta_h^{\rm comp}+2B_*\gamma_*^{-1}\beta_h^{\rm eff}$.
That full quantity must enter the shadow budget. The formulas stated below
use the exact-basepoint case \eqref{eq:exact-initial-basepoint}; any
nonzero complementary correction requires this explicit additional charge.
\end{remark}

\subsection{Residual interpolation}

\begin{proposition}[Full-record residual interpolation]
\label{prop:shadow-residual-upgrade}
Assume \eqref{eq:shadow-regularity-indices}--
\eqref{eq:shadow-reader-reserve} and let the complete finite Euler-row norm
satisfy $\sigma_h\le s_h$.  Then on a common smaller cylinder
\begin{equation}
 \varepsilon_h
 :=\norm{\mathcal R_{\rm B}(z_h)}_{L^2}
   +\norm{(r_D,\bar r_D)(z_h)}_{L^2}
 \le C_R(s_h+h).
 \label{eq:shadow-zero-residual}
\end{equation}
Put
\begin{equation}
 \eta_{h,k}
 :=C_R\left(
 \varepsilon_h^{1-k/s}
 +\varepsilon_h^{1-(k+1)/s}
 \right).
 \label{eq:shadow-eta}
\end{equation}
Then
\begin{equation}
 \norm{\mathcal R_{\rm B}(z_h)}_{H^k(Q)}
 +\norm{(r_D,\bar r_D)(z_h)}_{H^{k+1}(Q)}
 \le\eta_{h,k}.
 \label{eq:shadow-strong-residual}
\end{equation}
For $\varepsilon_h\le1$,
$\eta_{h,k}\le C_R\varepsilon_h^{1-(k+1)/s}$.
\end{proposition}
\begin{proof}
The uniform reader reserve permits \Cref{prop:native-consistency} and
\Cref{lem:nodal-source-transfer} to be applied with a fixed derivative scale,
which gives \eqref{eq:shadow-zero-residual}.  The bosonic Euler rows have
differential order at most two and the Dirac rows order at most one.
Equation \eqref{eq:shadow-reader-reserve} therefore gives a cutoff-independent
$H^s$ bound for the differentiated residuals.  Sobolev/Fourier interpolation
between this bound and \eqref{eq:shadow-zero-residual} gives exponents
$1-k/s$ for the bosonic $H^k$ residual and $1-(k+1)/s$ for the Dirac
$H^{k+1}$ residual.  The second exponent is smaller and dominates when
$\varepsilon_h\le1$.
\end{proof}

\subsection{Harmonic defects and an illustrative rate}

\begin{corollary}[Exact shadows with a monitored harmonic defect]
\label{cor:shadow-harmonic-defect}
Retain the exact-reduction, regular Cauchy-tube and physical-forcing
hypotheses of \Cref{thm:aposteriori-physical-shadow}, but allow the normalized
record $\widehat z_h$ to be only approximately harmonic. Keep exact temporal
internal gauge and the common foliation margins. Put
\begin{equation}
 \begin{aligned}
 \widehat\gamma_{h,k}
 &=\norm{C(\widehat g_h)}_{L_t^2H_x^{k+1}}
       +\norm{\partial_tC(\widehat g_h)}_{L_t^2H_x^k},\\
 \mathfrak d_h^C
 &=2B_*\gamma_*^{-1}\beta_h
       +C_R\sqrt T(\eta_{h,k}+\widehat\gamma_{h,k}).
 \end{aligned}
 \label{eq:shadow-harmonic-budget}
\end{equation}
If the enlarged budget lies within the same stability radius, the exact
solution with corrected initial datum $d_h$ satisfies
\eqref{eq:aposteriori-physical-shadow} with $\mathfrak d_h^C$ in place of
$\mathfrak d_h$. A validated constraint-evaluation error or a nonzero
complementary initial correction is included as in
\Cref{cor:validated-initial-retraction,rem:initial-correction-bookkeeping}.
\end{corollary}
\begin{proof}
Generate the exact solution from the same corrected compatible datum.
Use \eqref{eq:actual-jet-integrated-stability} with the additional forcing
in \eqref{eq:actual-jet-gauge-forcing}. Cauchy--Schwarz in time charges it
by $C_R\sqrt T\widehat\gamma_{h,k}$. The first-exit and curvature/stress
arguments are unchanged. No projection of the approximate metric onto
exact harmonic gauge is used.
\end{proof}

\begin{corollary}[A transparent data-centred rate]
\label{cor:aposteriori-shadow-rate}
Take $k=4$, $s=10$, $q\ge15$ and fix one certified reader-energy radius $R$.
If
\begin{equation}
 s_h=O(h),
 \qquad
 \beta_h=O(h^{1/2}),
 \label{eq:aposteriori-rate-inputs}
\end{equation}
then
\begin{equation}
 \varepsilon_h=O_R(h),
 \qquad
 \eta_{h,4}=O_R(h^{1/2}),
 \qquad
 \mathfrak d_h=O_R(h^{1/2}),
 \label{eq:aposteriori-rate-output}
\end{equation}
and the exact physical shadow differs from the normalized finite record by
$O_R(h^{1/2})$ in the norms of
\eqref{eq:aposteriori-physical-shadow}.
\end{corollary}
\begin{proof}
For $k=4$ and $s=10$, the bosonic interpolation exponent is $3/5$ and the
Dirac exponent is $1/2$.  Substitute
\eqref{eq:aposteriori-rate-inputs} into
\eqref{eq:shadow-zero-residual}--\eqref{eq:aposteriori-shadow-budget}.
\end{proof}

\section{Coupled stability and constraint propagation}
\label{app:stability-package}

\subsection{Coupled Einstein--matter actual-jet stability}
\label{app:coupled-stability}

This subsection proves the geometric-analysis stability estimate used in \Cref{thm:native-closure}.  Its variables are actual derivatives and curvatures of the reconstructed fields, so the argument is a coupled Einstein--matter estimate rather than an evolution of auxiliary jets followed by an identification step.

Use one common spacelike foliation and its adapted orthonormal frame
\begin{equation}
 e_0=N^{-1}(\partial_t-\beta^j\partial_j),\qquad
 e_a=e_a{}^j\partial_j\quad(a=1,2,3),\qquad 0<N_-\le N\le N_+.
 \label{eq:actual-jet-adapted-frame}
\end{equation}
Here the spatial frame is a fixed smooth algebraic choice on the positive
spatial-metric chart. In particular the $e_a$, $a=1,2,3$, are tangent to the
Cauchy slices. This choice is needed when counting derivatives of the Dirac
residual. The lapse, shift and frame coefficients have the uniform margins
of the regular chart; harmonic coordinates are not imposed on the approximate
metric. Choose Clifford multiplication $c_a$ satisfying
$c_ac_b+c_bc_a=-2\eta_{ab}I$. Let $\nabla$ be the torsion-free spin connection
twisted by the retained gauge representation and write
\[
 \mathcal R^\nabla_{ab}=\mathcal R^S_{ab}+\rho(F_{ab}),\qquad
 \mathcal R^S_{ab}=\frac14\sum_{c,d}\epsilon_c\epsilon_dR_{ab cd}c_cc_d.
\]

\begin{lemma}[Twisted half-Ricci contraction]
\label{lem:half-ricci}
With the preceding conventions,
\begin{equation}
 \sum_b\epsilon_b c_b\mathcal R^\nabla_{ba}
 =\frac12\sum_b\epsilon_b\Ric_{ab}c_b
 +\sum_b\epsilon_bc_b\rho(F_{ba}).
 \label{eq:half-ricci}
\end{equation}
\end{lemma}
\begin{proof}
At one point choose a normal orthonormal frame.  Insert
$\mathcal R^S_{ba}=\frac14R_{ba cd}c^cc^d$ and use
$c_bc_cc_d=c_{[b}c_cc_{d]}-\eta_{bc}c_d+\eta_{bd}c_c-\eta_{cd}c_b$.
The fully alternating term vanishes by the first Bianchi identity, the last contraction vanishes by antisymmetry, and the remaining two contractions give one half of the Ricci contraction.  The gauge curvature acts on the twisting factor and contributes the last term directly.
\end{proof}

Put
\[
 r_D=\mathscr D_{g,A}\Psi-\mathscr M_{\boldsymbol Y}(H)\Psi,
 \qquad X_a=\nabla_{e_a}\Psi.
\]
\begin{proposition}[Residual-sensitive tangential spinor prolongation]
\label{prop:spinor-prolongation}
For each frame index $a$,
\begin{equation}
 (\mathscr D^{T^*}X)_a
 =\mathscr M X_a+\mathscr M_H[D_aH]\Psi+\nabla_a r_D
 +\frac12\Ric_{ab}c^b\Psi+c^b\rho(F_{ba})\Psi.
 \label{eq:spinor-prolongation}
\end{equation}
Moreover the normal jet is algebraic,
\begin{equation}
 X_0=c_0c_iX_i-c_0\mathscr M\Psi-c_0r_D,
 \label{eq:normal-spinor-jet}
\end{equation}
so the tangential prolonged equations require spatial derivatives of $r_D$ but no time derivative of $r_D$.
\end{proposition}
\begin{proof}
Commute $\nabla_a$ with the twisted Dirac operator and use \Cref{lem:half-ricci}.  The covariant product rule differentiates the affine mass map to $\mathscr M_H[D_aH]$.  The Dirac equation itself gives \eqref{eq:normal-spinor-jet}.  In a moving frame, the normal covector component enters the tangential covector connection only algebraically, so substituting \eqref{eq:normal-spinor-jet} introduces no normal derivative of the residual.
\end{proof}

Let
\[
 \mathcal E_{ab}=G_{ab}+\Lambda g_{ab}-\kappa T^{\rm SM}_{ab},
 \qquad
 \mathcal E^{\rm tr}_{ab}=\mathcal E_{ab}-\tfrac12g_{ab}\operatorname{tr}_g\mathcal E.
\]
Then
\begin{equation}
 \Ric_{ab}=\kappa\left(T^{\rm SM}_{ab}-\tfrac12g_{ab}\operatorname{tr}_gT^{\rm SM}\right)
 +\Lambda g_{ab}+\mathcal E^{\rm tr}_{ab}.
 \label{eq:trace-reversal-residual}
\end{equation}
The complete torsion-free Standard-Model stress is algebraic in the actual first jets
\[
 (g,F,H,DH,\Psi,\bPsi,X_a,\bar X_a)
\]
and contains no derivative of $F$, $DH$ or $X$.  After \eqref{eq:normal-spinor-jet} it is a smooth algebraic function of the tangential actual-jet state and is at most linear in the Dirac residual.

For an arbitrary metric in this foliated coordinate chart, in exact temporal
internal gauge $A_0=0$, define for every coordinate metric component
$u=g_{\mu\nu}$,
$p_u=e_0u$, $q_{u,a}=e_au$, and
\[
 E_a=F(e_0,e_a),\quad B_a=-\tfrac12\varepsilon_{abc}F(e_b,e_c),
 \quad \Pi=D_{e_0}H,\quad Q_a=D_{e_a}H.
\]
The realified actual-jet state is
\begin{equation}
 \mathcal U=\bigl(g_{\mu\nu},p_u,q_{u,a},A_i,E_a,B_a,H,\Pi,Q_a,
                 \Psi,(X_a)_{a=1}^3,\bPsi,(\bar X_a)_{a=1}^3\bigr).
 \label{eq:actual-jet-state}
\end{equation}
The dual spinor and its tangential covariant jets are included explicitly.
On the physical duality branch these are redundant blocks fixed by the
adjoint identification; otherwise the dual equation is retained as an
additional symmetric-hyperbolic block. All state norms include the retained
dual variables.  Let $r^A$, $r_H$ and $r_D$ denote the physical Yang--Mills, Higgs and Dirac residuals.

\begin{lemma}[Exact harmonic-defect forcing identity]
\label{lem:harmonic-defect-forcing}
For a smooth metric in the common coordinate chart put
$C^\mu=g^{\alpha\beta}\Gamma^\mu_{\alpha\beta}$ and
$C_\nu=g_{\nu\mu}C^\mu$. Define
\begin{equation}
 \mathscr H_{\mu\nu}(g,C)
   =\nabla_{(\mu}C_{\nu)}
       -\tfrac12g_{\mu\nu}\nabla_\alpha C^\alpha.
 \label{eq:harmonic-defect-tensor}
\end{equation}
Then the reduced tensor of \eqref{eq:reduced-Einstein} satisfies
$\widehat\Eins=\Eins-\mathscr H$. In four dimensions its residual obeys
\begin{equation}
 (\widehat\Eins+\Lambda g-\kappa T)^{\rm tr}_{\mu\nu}
     =\mathcal E^{\rm tr}_{\mu\nu}-\nabla_{(\mu}C_{\nu)}.
 \label{eq:harmonic-residual-trace-reversal}
\end{equation}
There is a smooth expression $\mathscr Q_{\mu\nu}(g,\partial g)$, quadratic
in first metric derivatives, such that every smooth field tuple satisfies
\begin{equation}
 -\tfrac12g^{\alpha\beta}\partial_\alpha\partial_\beta g_{\mu\nu}
 +\mathscr Q_{\mu\nu}(g,\partial g)
 =\kappa\bigl(T_{\mu\nu}-\tfrac12g_{\mu\nu}\operatorname{tr}_gT\bigr)
   +\Lambda g_{\mu\nu}+\mathcal E^{\rm tr}_{\mu\nu}
   -\nabla_{(\mu}C_{\nu)}.
 \label{eq:off-harmonic-metric-wave}
\end{equation}
Thus failure of harmonic gauge contributes only $C$ and its first derivatives
to the metric wave source.
\end{lemma}
\begin{proof}
The first identity is the definition of the reduced tensor. In dimension four,
$\operatorname{tr}_g\mathscr H=-\nabla_\alpha C^\alpha$, so trace reversal
of $\mathscr H$ is exactly $\nabla_{(\mu}C_{\nu)}$. This proves
\eqref{eq:harmonic-residual-trace-reversal}.
To verify the derivative count, expand
\[
 \Ric_{\mu\nu}
 =\partial_\alpha\Gamma^\alpha_{\mu\nu}
  -\partial_\nu\Gamma^\alpha_{\mu\alpha}
  +\Gamma^\alpha_{\alpha\lambda}\Gamma^\lambda_{\mu\nu}
  -\Gamma^\alpha_{\nu\lambda}\Gamma^\lambda_{\mu\alpha}.
\]
After subtracting $\nabla_{(\mu}C_{\nu)}$, the second-derivative part is
$-\tfrac12g^{\alpha\beta}\partial_\alpha\partial_\beta g_{\mu\nu}$.
Every remaining term comes from a differentiated inverse metric or a product
of Christoffel symbols, hence is smooth in $g,g^{-1}$ and quadratic in
$\partial g$. This defines $\mathscr Q$. Combine the resulting reduced
Ricci identity with \eqref{eq:trace-reversal-residual}. This is the usual
harmonic reduction \cite{FriedrichRendall2000}, here retained as an identity
with its nonzero gauge defect rather than restricted to $C=0$.
\end{proof}

\begin{proposition}[Gauge-defect-aware block-symmetric actual-jet equation]
\label{prop:actual-jet-writer}
On a compact nondegenerate foliated chart, every smooth actual field tuple
in exact temporal internal gauge, with both spinor rows retained, satisfies
\begin{align}
 \partial_t\mathcal U+\sum_{j=1}^3\mathcal A^j(g)\partial_j\mathcal U
 &=\mathcal F(\mathcal U)
   +\mathcal B(\mathcal U)(\mathcal E,r^A,r_H,r_D,\bar r_D)\notag\\
 &\quad+\sum_{j=1}^3\mathcal Q^j(g)\partial_j(r_D,\bar r_D)
        +\mathfrak G(\mathcal U;C,\partial C),
 \label{eq:actual-jet-writer}\\
 \mathfrak G(\mathcal U;C,\partial C)
 &=\mathcal J_C(\mathcal U)C+
       \mathcal J_C^0(g)\partial_tC+
       \sum_{j=1}^3\mathcal J_C^j(g)\partial_jC.
 \label{eq:actual-jet-gauge-forcing}
\end{align}
Here $C=(C^\mu(g))$ is the actual harmonic defect, not an independent state
variable. The matrices $\mathcal A^j$ are symmetric in the fixed positive
component norm after realification; the other coefficient maps are smooth
on bounded actual-jet sets. No derivative of $\mathcal E$, $r^A$ or $r_H$
and no time derivative of $(r_D,\bar r_D)$ occurs. If $C=0$, the gauge-forcing
term vanishes and the exactly harmonic system is recovered.
\end{proposition}
\begin{proof}
Use \eqref{eq:off-harmonic-metric-wave} for each metric component. First-order
reduction in the adapted frame gives symmetric wave blocks for $(p_u,q_u)$.
The additional source $-\nabla_{(\mu}C_{\nu)}$ is a linear combination of
$\partial_\alpha C^\lambda$ with coefficients algebraic in $g$, and of
$C^\lambda$ with coefficients algebraic in $g,\partial g$. Multiplication by
the lapse and frame matrices of the wave reduction therefore gives exactly
\eqref{eq:actual-jet-gauge-forcing}. It changes no principal matrix and
introduces no second derivative of $C$.

The Yang--Mills equation and Bianchi identity give the symmetric Maxwell
blocks for $(E,B)$. The Higgs equation and covariant commutator identity give
symmetric wave blocks for $(\Pi,Q)$. Temporal gauge gives
$\partial_tA_i=F_{0i}$, an algebraic head equation in this state. Metric and
Higgs head equations follow from their defining first jets. The Dirac head
has Hermitian normal-frame principal matrices $-c_0c_a$; realification makes
them real symmetric. Transposing the dual row gives the transposed principal
matrices, symmetric in the dual positive component norm, so it has the same
energy estimate. Passing from the adapted frame to coordinates multiplies
each wave, Maxwell or Dirac block by the common real lapse/frame coefficients
and adds the scalar shift contribution $-\beta^jI$, preserving symmetry.

For the tangential spinor jets use \Cref{prop:spinor-prolongation}.
The contracted curvature there is the physical Ricci tensor, so
\eqref{eq:trace-reversal-residual} eliminates it without assuming harmonic
gauge and without differentiating stress. Substitute the algebraic normal
jet \eqref{eq:normal-spinor-jet} throughout. Since $e_a=e_a{}^j\partial_j$
for $a=1,2,3$, its residual derivative is
\[
 \nabla_a r_D=e_a{}^j\partial_j r_D+\text{a connection coefficient times }r_D,
\]
and similarly for the dual row. The normal covector component enters the
tangential prolonged connection only algebraically. Hence no
$\partial_t r_D$ is introduced. Frame commutators and spin-connection
coefficients depend only on the actual first metric jets. The off-shell
stress is algebraic in those jets and at most linear in the normal Dirac
residual. Consequently all residual terms have the form displayed in
\eqref{eq:actual-jet-writer}, with no differentiated bosonic residual.
This constructs the same symmetric principal blocks for exactly and
approximately harmonic records.
\end{proof}

\begin{lemma}[Derivative-counted complete forcing bound]
\label{lem:actual-jet-complete-forcing}
Fix $k\ge4$ and an actual-jet $H^k$ ball with the stated chart and foliation
margins. Let $\mathscr F_{\rm err}$ denote the sum of the residual and
gauge-forcing terms on the right of \eqref{eq:actual-jet-writer}. Then
\begin{align}
 \norm{\mathscr F_{\rm err}}_{H_x^k}
 \le C_M\bigl(&\norm{\mathcal R_{\rm B}}_{H_x^k}
        +\norm{\mathcal R_{\rm D}}_{H_x^{k+1}}\notag\\
        &+\norm C_{H_x^{k+1}}+\norm{\partial_tC}_{H_x^k}\bigr).
 \label{eq:actual-jet-complete-forcing}
\end{align}
\end{lemma}
\begin{proof}
The state determines $g,\partial g$, so every zeroth-order coefficient in
\eqref{eq:actual-jet-writer} has a bounded $H^k$ norm on this ball. Smooth
composition and the $H^k$ algebra estimate on the three-dimensional closed
slices control products with the undifferentiated residuals. The matrices
multiplying $\partial_j r_D$ and $\partial_\alpha C$ depend only on $g$;
these products require respectively $r_D\in H^{k+1}$,
$C\in H^{k+1}$ and $\partial_tC\in H^k$. Lowering the index of $C$ adds only
products of $C$ with first metric jets. The normalized Euler densities and
physical rows are related by smooth invertible algebraic coefficient maps
on this chart. Summing these estimates proves the bound.
\end{proof}

\begin{theorem}[Coupled actual-jet stability and first-exit closure]
\label{prop:coupled-bootstrap}
Let $z_*$ be a regular exact Einstein--Standard-Model solution on $Q=[0,T]\times\mathbb T^3$ in the preceding gauges with $\mathcal U_*\in C_tH_x^{k+1}$, $k\ge4$.  Let $\widehat z$ be a smooth actual field tuple on the same slab and
foliated chart, in exact temporal internal gauge. Both spinor and dual-spinor
blocks are included. Harmonic coordinates are required for $z_*$ but not
for $\widehat z$. Define
\[
 d=\norm{\widehat{\mathcal U}(0)-\mathcal U_*(0)}_{H^k}
 +\norm{\mathcal R_{\rm B}(\widehat z)}_{L^2_tH_x^k}
 +\norm{\mathcal R_{\rm D}(\widehat z)}_{L^2_tH_x^{k+1}}
 +\norm{C(\widehat g)}_{L^2_tH_x^{k+1}}
 +\norm{\partial_tC(\widehat g)}_{L^2_tH_x^k}.
\]
There are constants $d_*,C_*>0$, depending only on the reference and fixed chart margins, such that $d\le d_*$ implies
\begin{align}
 \norm{\widehat{\mathcal U}-\mathcal U_*}_{C_tH_x^k}&\le C_*d,\label{eq:bootstrap-state}\\
 \norm{\Riem(\widehat g)-\Riem(g_*)}_{L^2_tH_x^{k-1}}
 +\norm{T^{\rm SM}(\widehat z)-T^{\rm SM}(z_*)}_{L^2_tH_x^k}
 &\le C_*d.\label{eq:bootstrap-curvature}
\end{align}
No all-time $H^k$ bound for $\widehat{\mathcal U}$ is assumed.  Existence of $\widehat z$ on the slab is assumed.
\end{theorem}
\begin{proof}
Choose a radius $\rho>0$ such that the $H^k$ tube of radius $\rho$ around
the reference state stays inside the common nondegenerate foliated chart.
This is possible by $H^k(\mathbb T^3)\hookrightarrow C^1$ and the fixed
reference margins. On an interval on which
$\norm{\widehat{\mathcal U}-\mathcal U_*}_{H^k}<\rho$, put
$W=\widehat{\mathcal U}-\mathcal U_*$ and subtract the two identities of
\Cref{prop:actual-jet-writer}. The reference has zero physical and harmonic
residuals. With the approximate-state principal matrices on the left,
\[
 \partial_tW+\mathcal A^j(\widehat g)\partial_jW
 =\mathcal F(\widehat{\mathcal U})-\mathcal F(\mathcal U_*)
  +\bigl(\mathcal A^j(g_*)-\mathcal A^j(\widehat g)\bigr)
          \partial_j\mathcal U_*+\mathscr F_{\rm err}.
\]
The first term is locally Lipschitz in $H^k$. The second is bounded by
$C\norm W_{H^k}$ because the reference has one additional spatial derivative.
For $|\alpha|\le k$ the commutator estimate
\[
 \norm{[\partial^\alpha,a]\partial_jW}_2
 \le C_k\bigl(\norm{\nabla a}_\infty\norm W_{H^k}
                +\norm a_{H^k}\norm{\nabla W}_\infty\bigr)
\]
and $H^k\hookrightarrow W^{1,\infty}$ control the differentiated principal
terms. Symmetry and integration by parts on the closed slice give a positive
energy $E_k^{1/2}\simeq\norm W_{H^k}$ with
\begin{equation}
 \begin{aligned}
 \frac{\dd}{\dd t}E_k^{1/2}&\le C E_k^{1/2}+C b_k(t),\\
 b_k&=\norm{\mathcal R_{\rm B}}_{H^k}
      +\norm{\mathcal R_{\rm D}}_{H^{k+1}}
      +\norm C_{H^{k+1}}+\norm{\partial_tC}_{H^k}.
 \end{aligned}
 \label{eq:actual-jet-gauge-energy}
\end{equation}
At a zero of the energy use $(E_k+\epsilon^2)^{1/2}$ and pass to the limit.
The forcing estimate here is precisely
\Cref{lem:actual-jet-complete-forcing}; no exactly harmonic equation has
been substituted for the approximate metric.

Integration and Gronwall give
\begin{equation}
 \sup_{0\le t\le T}\norm W_{H^k}
 \le C_T\left(\norm{W(0)}_{H^k}+\int_0^T b_k(t)\,\dd t\right)
 \le C_*d
 \label{eq:actual-jet-integrated-stability}
\end{equation}
as long as the tube condition holds. Choose $d_*$ so that the right side is
less than $\rho/2$ and the initial state is inside this smaller tube.
Continuity then excludes a first exit before $T$, proving
\eqref{eq:bootstrap-state}. Existence of the smooth approximate record on
the slab is used here; no uniform all-time Sobolev bound for it is assumed.

The equation bounds $\partial_tW$ in $L_t^2H_x^{k-1}$ by the same state and
forcing quantities. Recover ordinary first metric derivatives algebraically
from $(g,p,q)$, and differentiate these identities once, using the metric
first-jet rows for the time derivative. All second metric derivatives then
differ by $O(d)$ in $L_t^2H_x^{k-1}$. The coordinate curvature formula,
linear in these second derivatives and quadratic in first derivatives,
proves the curvature estimate. For the complete stress, normal-spinor
elimination gives smooth algebraic maps with
\[
 T^{\rm SM}(z)=\mathscr T_0(\mathcal U)
                  +\mathscr T_1(\mathcal U)(r_D,\bar r_D).
\]
Consequently
\[
 \norm{T^{\rm SM}(\widehat z)-T^{\rm SM}(z_*)}_{L_t^2H_x^k}
 \le C\left(\sqrt T\norm W_{C_tH_x^k}
                   +\norm{(r_D,\bar r_D)}_{L_t^2H_x^k}\right)
 \le C_*d.
\]
This includes the off-shell normal-spinor contribution and proves
\eqref{eq:bootstrap-curvature}.
\end{proof}

\subsection{Common-slab metric stability}
\label{sec:propagation}

We now describe a way of obtaining geometric convergence from a reduced evolution. This result is intentionally local in physical time. It assumes a common hyperbolic slab and a strongly controlled \emph{complete} source; it does not manufacture that source from the algebraic duality. The energy argument belongs to the standard symmetric-hyperbolic framework \cite{Kato1975,FriedrichRendall2000}.

Let $s\ge5$ be an integer and let smooth metrics $g_h$ on $[0,T]\times\Sigma$ solve, in one background-harmonic coordinate/frame realization,
\begin{equation}
 g_h^{\alpha\beta}\partial_\alpha\partial_\beta g_h
 =\mathcal N(g_h,\partial g_h)+\mathcal S_h.
 \label{eq:reduced-wave}
\end{equation}
Fixed background-connection terms are included in $\mathcal N$. Its dependence on first derivatives is the usual smooth quadratic one on a compact nondegenerate metric chart. Here $\mathcal S_h$ is a tensor source, not the action $\Sact_h$.

\begin{theorem}[Common-slab metric stability]
\label{thm:hyperbolic}
Assume uniform hyperbolicity with a common time function, uniform inverse-metric bounds, and
\begin{equation}
 \sup_h\left(\norm{g_h}_{L^\infty_tH^s_x}
       +\norm{\partial_tg_h}_{L^\infty_tH^{s-1}_x}\right)<\infty.
 \label{eq:hyperbolic-bound}
\end{equation}
Suppose the initial data are Cauchy in $H^{s-1}\times H^{s-2}$ and the sources are Cauchy in $L^1_tH^{s-2}_x$. Then
\begin{align}
 g_h&\longrightarrow g\quad\text{in }C_tH^{s-1}_x,\notag\\
 \partial_tg_h&\longrightarrow\partial_tg\quad\text{in }C_tH^{s-2}_x,\notag\\
 \Riem(g_h)&\longrightarrow\Riem(g)
       \quad\text{in }L^1_tH^{s-3}_x.
 \label{eq:hyperbolic-limit}
\end{align}
In particular, compatible coframes chosen on one compact frame chart converge in the geometric topology of \eqref{eq:strong-geometry}. A vanishing reduced-equation residual can be included in $\mathcal S_h$ with the same source norm.
\end{theorem}

\begin{proof}
Subtract the equations for $h$ and $j$, and put $w=g_h-g_j$. Use the principal wave operator with coefficient $g_h$ on the left. The forcing consists of $\mathcal S_h-\mathcal S_j$, the difference of the quadratic first-derivative terms, and
\[
 (g_j^{\alpha\beta}-g_h^{\alpha\beta})\partial_\alpha\partial_\beta g_j.
\]
At one derivative below the a priori bound, $H^{s-2}(\Sigma)$ is an algebra. Spatial second derivatives of $g_j$ and mixed time--space derivatives are uniformly bounded in $H^{s-2}$. Solving \eqref{eq:reduced-wave} for $\partial_t^2g_j$ gives an $L^1_tH^{s-2}$ bound with size controlled by $1+\norm{\mathcal S_j}_{L^1_tH^{s-2}}$. Consequently the forcing, apart from $\mathcal S_h-\mathcal S_j$, has $H^{s-2}$ norm bounded by
\[
 a_{h,j}(t)\bigl(\norm w_{H^{s-1}}+\norm{\partial_tw}_{H^{s-2}}\bigr),
 \qquad \sup_{h,j}\int_0^T a_{h,j}(t)\dd t<\infty.
\]
The differentiated wave energy $E_{h,j}$ is uniformly equivalent to the square of the displayed difference norm. Its square root obeys the integrated estimate
\begin{equation}
 E_{h,j}(t)^{1/2}
 \le C\left(E_{h,j}(0)^{1/2}
   +\int_0^t\norm{\mathcal S_h-\mathcal S_j}_{H^{s-2}}\dd u\right)
   \exp\!\left(C\int_0^t(1+a_{h,j}(u))\dd u\right).
 \label{eq:hyperbolic-energy}
\end{equation}
For completeness, commute spatial derivatives through the principal operator, use the differentiated normal derivative as energy multiplier, and integrate on the closed $\Sigma$. The uniformly spacelike slices and bounded lapse and shift give a positive wave energy; commutators are bounded by the coefficient norms in \eqref{eq:hyperbolic-bound} and the Sobolev product estimate above. Apply the scalar differential inequality to $(E_{h,j}+\delta)^{1/2}$ and let $\delta\downarrow0$ to obtain \eqref{eq:hyperbolic-energy}. This form uses the $L^1$ source norm directly; it does not replace it by an unjustified squared-source bound.

Cauchy initial and source data prove the first two limits. The reduced equation then gives convergence of $\partial_t^2g_h$ in $L^1_tH^{s-3}$. All other second derivatives converge in $C_tH^{s-3}$, and the inverse metrics and quadratic first-derivative terms converge by Sobolev product estimates. The coordinate formula for curvature gives the last limit. Smooth dependence of a locally fixed coframe on the metric gives the coframe assertion.
\end{proof}

The loss of one derivative in the difference estimate avoids treating a quasilinear coefficient difference times a second derivative as if it had top-order Lipschitz regularity. Likewise, $L^1$ control of the source yields the stated $L^1$-in-time curvature conclusion, not continuity in time of every curvature component. Stronger conclusions require stronger source and multiplier estimates.

\subsection{The Ward residual forces the harmonic constraint}

For a harmonic gauge covector $C_\mu$ define
\begin{equation}
 \widehat\Eins_{\mu\nu}
 =\Eins_{\mu\nu}-\nabla_{(\mu}C_{\nu)}
        +\tfrac12g_{\mu\nu}\nabla^\alpha C_\alpha.
 \label{eq:reduced-Einstein}
\end{equation}
Suppose a smooth reconstructed field satisfies
\[
 \widehat\Eins_h+\Lambda_hg_h=\kappa_h T_h+r_h,
 \qquad
 \mathcal W_{h,\nu}=\kappa_h\nabla_h^\mu T_{h,\mu\nu}
                         +\nabla_h^\mu r_{h,\mu\nu},
\]
with constant parameters on spacetime. The source $\mathcal W_h$ includes every retained sector and the reduced-equation residual.

\begin{proposition}[Subsidiary equation and constraint stability]
\label{prop:subsidiary}
With the curvature convention fixed by \eqref{eq:reduced-Einstein},
\begin{equation}
 \Box_{g_h}C_{h,\nu}+\Ric(g_h)_\nu{}^\mu C_{h,\mu}
 =-2\mathcal W_{h,\nu}.
 \label{eq:subsidiary}
\end{equation}
On a common closed spatial slab with uniform wave-energy coefficients and a uniform $L^1_tL^\infty_x$ bound for the curvature potential,
\begin{equation}
 \norm{(C_h,\nabla C_h)}_{L^\infty_tL^2_x}
 \le C_T\left(
 \norm{(C_h,\nabla C_h)(0)}_{L^2_x}
 +\norm{\mathcal W_h}_{L^1_tL^2_x}\right).
 \label{eq:constraint-bound}
\end{equation}
Vanishing initial constraint data and Ward residual therefore imply $C_h\to0$ in the subsidiary energy norm.
\end{proposition}
\begin{proof}
Take the covariant divergence of the reduced Einstein equation. Contracted Bianchi removes the Einstein tensor and the constant cosmological term. The identity
\[
 2\nabla^\mu\nabla_{(\mu}C_{\nu)}
   -\nabla_\nu\nabla^\alpha C_\alpha
 =\Box C_\nu+\Ric_\nu{}^\mu C_\mu
\]
gives \eqref{eq:subsidiary}. For the estimate, use the positive wave energy of a one-form relative to the common time foliation, including its $L^2$ term. Integration by parts has no spatial boundary flux. Coefficient derivatives and the curvature potential multiply the energy by an integrable function, and the forcing pairs with the time derivative. The square-root energy inequality and Gronwall give \eqref{eq:constraint-bound}.
\end{proof}

This proposition is not a consequence of the finite commutant identity. Exact classical gauge and relabeling invariance provide Noether identities, but their continuum use requires the whole source and the physical test realization. In an evolution-based construction, the subsidiary residual must be monitored along with the reduced equation. On the variational branch of \Cref{thm:main-limit}, by contrast, independent physical variations already impose the full Euler system. Conservation is asserted at the regularity at which the corresponding covariant distributional identity is defined, not by multiplying rough distributions formally.

A physical timelike boundary adds a boundary energy flux to \eqref{eq:constraint-bound}. Maximal dissipation for the main system does not automatically preserve the subsidiary constraint, and dissipation alone does not identify a boundary domain of the same closed action. These are separate extensions; the present closed-spatial-cylinder result avoids, rather than solves, that boundary problem.

\section{Controlled comparison and generated finite families}
\label{app:comparison-families}

These are conditional nonemptiness and calibration constructions.
Smooth sampling verifies consistency on prescribed smooth fields;
nondegenerate finite shadowing additionally requires a uniform linearized
inverse. The generated evolution below instead starts from constrained
initial data and retains its separate initial-solver hypotheses.

\subsection{Smooth calibration and finite critical-point shadowing}

\begin{proposition}[Smooth sampling is a consistency realization]
\label{prop:smooth-sampling}
Let $z_*$ be a smooth classical configuration on a compact slab with a fixed nondegenerate coframe chart. On the prescribed shape-regular reconstruction and local action of \Cref{app:reconstruction}, smooth sampling gives $z_h\to z_*$ in the strong packet and
\begin{equation}
 c_h(K)\le C_Kh.
 \label{eq:sampling-consistency}
\end{equation}
If $z_*$ solves the classical Euler equations, then on the same smooth family
\begin{equation}
 \varepsilon_h(K)\le C'_Kh.
 \label{eq:sampling-stationarity}
\end{equation}
Under the nondegeneracy hypotheses of \Cref{prop:critical-shadowing}, these approximate critical points are shadowed by exact finite critical points.
\end{proposition}
\begin{proof}
Apply \Cref{prop:mesh-consistency} for strong-packet convergence and
$c_h(K)=O(h)$. If $z_*$ is critical, \Cref{prop:variation-continuity}
gives $\varepsilon_h(K)\le c_h(K)+C_Kd_K(z_h,z_*)=O(h)$.
The final assertion is \Cref{prop:critical-shadowing}.
\end{proof}

\begin{proposition}[Nondegenerate critical-point shadowing]
\label{prop:critical-shadowing}
Let $z_*$ be a smooth classical critical configuration and $x_h^0$ its sampled finite configuration in a gauge-fixed finite-dimensional physical slice $X_h$.  Represent the finite Euler covector on that slice by a map $F_h:X_h\to X_h$ using a fixed positive slice inner product.  Suppose, on a ball $B_r(x_h^0)$ and uniformly for small $h$,
\begin{enumerate}[label=\textup{(S\arabic*)},leftmargin=3.2em]
\item $\norm{F_h(x_h^0)}\le C_0h$;
\item $D F_h(x_h^0)$ is invertible and $\norm{D F_h(x_h^0)^{-1}}\le M$;
\item $\norm{D F_h(x)-D F_h(y)}\le L\norm{x-y}$;
\item the reconstruction is locally Lipschitz, $d_K(\mathcal R_hx,\mathcal R_hy)\le C_R\norm{x-y}$.
\end{enumerate}
Then, for sufficiently small $h$, there is a unique exact finite critical point $x_h^*\in B_r(x_h^0)$ in the Newton ball
\begin{equation}
 \norm{x_h^*-x_h^0}\le 2MC_0h,
 \label{eq:critical-shadowing}
\end{equation}
and its reconstruction satisfies $d_K(\mathcal R_hx_h^*,z_*)=O(h)$.  Hence a nondegenerate smooth Einstein--Standard-Model solution generates a nearby exact finite stationary branch whenever the finite linearized physical Euler operator has a cutoff-uniform inverse and the second variation is uniformly controlled.
\end{proposition}
\begin{proof}
Put $B_h=D F_h(x_h^0)^{-1}$ and define the frozen-Newton map
\[
 \mathcal T_h(x)=x-B_hF_h(x).
\]
Let $\rho_h=2MC_0h$.  Taylor's formula and (S3) give, for $x,y\in B_{\rho_h}(x_h^0)$,
\[
 \norm{\mathcal T_h(x)-\mathcal T_h(y)}
 \le ML\rho_h\norm{x-y}.
\]
For small $h$, $ML\rho_h<1/2$.  The same Taylor estimate gives
\[
 \norm{\mathcal T_h(x)-x_h^0}
 \le MC_0h+\tfrac12ML\rho_h^2\le\rho_h,
\]
again for small $h$.  Thus $\mathcal T_h$ is a contraction of the Newton ball into itself and has a unique fixed point there.  Since $B_h$ is invertible, the fixed-point equation is $F_h(x_h^*)=0$.  The reconstruction estimate follows from (S4), \eqref{eq:critical-shadowing}, and the $O(h)$ sampling convergence of \Cref{prop:smooth-sampling}.
\end{proof}

If the positive slice norm used to represent $F_h$ is uniformly equivalent to the physical dual norm induced by the test lifts, condition (S1) follows from \eqref{eq:sampling-stationarity}.  The substantive additional hypothesis is then the cutoff-uniform invertibility and local Lipschitz control of the linearized physical Euler map.

This establishes a nonempty classical comparison class once a smooth solution is specified, and it upgrades that class to exact finite stationarity on every locally nondegenerate branch satisfying \Cref{prop:critical-shadowing}.  It is not a proof that every renewal history selects this discretization or that every classical solution has a uniformly invertible physical linearization.  The following explicit reduction supplies a backreacting comparison field without appealing to a general coupled existence theorem.

\begin{proposition}[A nonvacuum Einstein--Higgs comparison family]
\label{prop:homogeneous}
Let $\Sigma=\mathbb T^3$, $\kappa,\lambda_H>0$, and take
\begin{equation}
 g=-\dd t^2+a(t)^2\delta_{ij}\dd x^i\dd x^j,
 \quad H=\frac{\phi(t)}{\sqrt2}h_0,
 \quad A=0,\quad\Psi=\bPsi=0,
 \label{eq:homogeneous-fields}
\end{equation}
where $h_0$ is a fixed unit Higgs vector and $\phi$ is real. Put
\[
 U(\phi)=\lambda_H(\phi^2/2-v_H^2)^2.
\]
For initial data $a_0>0$, $\phi_0$, $\pi_0$, and $\mathcal H_0$ satisfying
\begin{equation}
 3\mathcal H_0^2=\Lambda+\kappa\left(\tfrac12\pi_0^2+U(\phi_0)\right),
 \label{eq:Friedmann-initial}
\end{equation}
the system
\begin{equation}
 \dot a=a\mathcal H,\qquad
 \dot\phi=\pi,\qquad
 \dot\pi=-3\mathcal H\pi-U'(\phi),\qquad
 \dot{\mathcal H}=-\frac\kappa2\pi^2
 \label{eq:homogeneous-ODE}
\end{equation}
has a smooth local solution with $a>0$, and \eqref{eq:homogeneous-fields} solves the torsion-free Einstein--Standard-Model equations. Choosing $\pi_0\ne0$ gives a genuinely time-dependent matter stress and backreaction. Smooth sampling yields the finite residual conclusion of \Cref{prop:smooth-sampling}.
\end{proposition}
\begin{proof}
The right side of \eqref{eq:homogeneous-ODE} is smooth, so local ordinary differential equation existence and uniqueness apply. Since $a(t)=a_0\exp(\int_0^t\mathcal H)$, the scale factor stays positive on its existence interval. The Higgs kinetic term is $-\tfrac12g^{\mu\nu}\partial_\mu\phi\partial_\nu\phi$, giving energy density $\rho=\tfrac12\pi^2+U$ and pressure $p=\tfrac12\pi^2-U$. Its Euler equation is the third equation of \eqref{eq:homogeneous-ODE}.

The gauge current vanishes: $H$ and $\dot H$ have the same fixed internal direction and a real relative coefficient, whereas all infinitesimal unitary generators are skew-adjoint. All fermionic equations are satisfied by zero spinors. Differentiating the Friedmann residual gives
\begin{align*}
 \frac{\dd}{\dd t}\left(3\mathcal H^2-\Lambda
             -\kappa(\tfrac12\pi^2+U)\right)
 &=6\mathcal H\left(-\tfrac\kappa2\pi^2\right)
   -\kappa\left(\pi(-3\mathcal H\pi-U')+U'\pi\right)=0.
\end{align*}
Thus the constraint persists. Together with $\dot{\mathcal H}=-(\kappa/2)(\rho+p)$ it gives both independent Einstein equations for the flat homogeneous metric. If $\pi_0\ne0$, then $\dot{\mathcal H}(0)<0$, so the example is not merely a vacuum metric relabeled by an internal carrier. On a smaller compact existence slab all required smoothness and noncollapse bounds hold, and \Cref{prop:smooth-sampling} applies.
\end{proof}

\begin{corollary}[A nonempty calibration class for the local action]
\label{cor:local-calibration-nonempty}
Let $z_*$ be a smooth solution on a compact slab, with compatible fixed
gauges and a smooth periodic extension to the buffered comparison box that
remains in the chosen field-value chart.  Nodal sampling followed by
trigonometric reconstruction yields records for $S_h^{\rm loc}$ with
$\sigma_h=O(h)$ and with the budgets of \Cref{thm:native-closure} tending
to zero for a suitable $K_h\to\infty$.  Thus the quantitative hypotheses
are nonempty.  A backreacting Einstein--Higgs example is supplied by
\Cref{prop:homogeneous}, after a regular gauge change on a smaller slab.
\end{corollary}
\begin{proof}
Smooth Fourier decay and the aliasing formula give uniform $H^q$ bounds
and convergence of each fixed number of reconstructed derivatives.
Fixed-scale consistency yields $\sigma_h=O(h)$ because the reference Euler
row vanishes; derivative convergence controls the initial and harmonic
mismatches. Choose $K_h\asymp h^{-\beta}$ with
$0<\beta<1/(k+4)$. The bound
$\tau_{h,j}\le C_qK_h^{2-q}$ from \eqref{eq:app-sharp-tail}, for
sufficiently large $q>k+5$, makes both $hK_h^{k+4}$ and
$K_h^{k+5-q}$ vanish, including the logarithmic factors in the forcing
budget. This constructs a calibration family, not a microscopic selector.
\end{proof}

\subsection{Generated finite dynamics from constrained data}
\label{app:generated-dynamics}

This subsection gives an independent nonemptiness construction for the controlled finite class.  It is local in time and deliberately high-regularity; sharp thresholds are not claimed.  Its role is not to replace the native source theorem of \Cref{thm:native-selected-closure}, but to show that the quantitative hypotheses admit a fully finite realization generated from constrained finite initial data.  The construction uses the positive symmetric actual-jet system of \Cref{app:coupled-stability} and the standard conformal strategy for the Einstein constraints \cite{ChoquetBruhatIsenbergPollack2007,IsenbergMaxwell2025}.

\begin{assumption}[Compatible elliptic and conformal initial solver]
\label{ass:constrained-initial-solver}
The seed family lies on fixed gauge and conformal slices for which the
Gauss and vector-momentum sources are orthogonal to the corresponding
elliptic kernels.  Bounded inverses on these compatible subspaces are
supplied, with cutoff-uniform estimates at every Sobolev order used below.
The chosen conformal weights and matter data give the scalar equation
\eqref{eq:generated-Hamiltonian}, with the source bound
\eqref{eq:generated-source-small}.  The finite solvers preserve compatibility
and have the approximation estimates in \eqref{eq:generated-initial-rate}.
A constraint-compatible analytic core with the same gauge completion is
available for the physical-identification argument.  On this core the
constrained physical Cauchy problem admits locally unique analytic
solutions with continuation controlled by the declared Sobolev bounds
and field-chart margins.
\end{assumption}

The following construction is conditional on this initial-solver and
physical-Cauchy interface.  The principal variational closure and the
deterministic a-posteriori stability theorems do not use this supplementary
existence input.
It does not assert invertibility of an unrestricted Gauss or momentum
operator on a torus.  In particular, a real standing-spinor profile alone
is not a certificate of vanishing gauge charge or of compatibility with
all elliptic zero modes.  The explicit homogeneous zero-spinor family in
\Cref{prop:homogeneous} supplies a nonempty comparison class independently
of any additional charged-spinor seed construction.

\begin{theorem}[Fully finite approximation branch]
\label{thm:generated-dynamics}
Fix integers $s\ge5$, $p\ge1$, and put $q=s+p+8$.  Under \Cref{ass:constrained-initial-solver}, for the constraint-compatible initial class below, each spatial cutoff $N$ has finite initial data $U_{0,N}$ satisfying
\begin{equation}
 \sup_N\norm{U_{0,N}}_{H^q}\le M_q,
 \qquad
 \norm{U_{0,N}-U_0}_{H^{s+4}}\le C N^{-p}.
 \label{eq:generated-initial-rate}
\end{equation}
Let
\begin{equation}
 G_N(U)=P_N\{F(U)-A^i(U)\partial_iU\},
 \qquad
 U^{j+1}=U^j+\tau G_N\!\left(\frac{U^{j+1}+U^j}{2}\right),
 \quad U^0=U_{0,N}.
 \label{eq:generated-midpoint}
\end{equation}
If
\begin{equation}
 \tau(N+1)\le c_*,\qquad \tau\le\tau_*,
 \label{eq:generated-CFL}
\end{equation}
then the fully finite recursion exists on a common interval $[0,T]$ and
\begin{equation}
 \max_j\norm{U^j}_{H^q}\le R_q.
 \label{eq:generated-uniform}
\end{equation}
Its cubic-Hermite physical reconstruction $z_{N,\tau}$ satisfies
\begin{align}
 \norm{z_{N,\tau}-z}_{C_tH^{s+2}}
 +\norm{\partial_tz_{N,\tau}-\partial_tz}_{C_tH^{s+1}}
 &\le C(N^{-p}+\tau^2),\label{eq:generated-first}\\
 \norm{\Riem(g_{N,\tau})-\Riem(g)}_{L^\infty_tH^s}
 +\norm{\mathcal R_{\rm B}(z_{N,\tau})}_{L^\infty_tH^s}
 &\le C(N^{-p}+\tau),\label{eq:generated-bosonic}\\
 \norm{\mathcal R_{\rm D}(z_{N,\tau})}_{C_tH^{s+1}}
 &\le C(N^{-p}+\tau^2).\label{eq:generated-Dirac}
\end{align}
Moreover the finite comparison action on the same generated record obeys
\begin{equation}
 |D S^{\rm cmp}_{N,\tau}[V]|
 \le C(N^{-p}+\tau)\norm V_{0,\tau}.
 \label{eq:generated-stationarity}
\end{equation}
The continuum solution $z$ is used only for the a-posteriori estimate; the finite initial solver and recursion do not sample its future values.
\end{theorem}

\subsubsection{A constrained initial class and its finite approximation}

Work on $\mathbb T^3$ near a smooth expanding constant background.  Fix a constant Higgs position $H_*$, a nonzero radial Higgs normal momentum $p$, and a spatially homogeneous gauge background $A^*$, so that the scalar energy coefficients below are constant.  Put
\begin{equation}
 A_*^0=\kappa p^2,
 \qquad G_*^0=\kappa|D_{A^*}H_*|^2,
 \qquad B_*^0=\kappa|B_{A^*}|^2,
 \label{eq:generated-background-coeffs}
\end{equation}
and choose the constant mean curvature so that the unperturbed Hamiltonian constraint holds.  In the conformal scalar equation the derivative at the unperturbed solution $w=1$ is
\begin{equation}
 L_*=-8\Delta+\mu_*,
 \qquad
 \mu_*=12A_*^0+4G_*^0+8B_*^0>0.
 \label{eq:generated-Lstar}
\end{equation}
Thus $L_*:H^{r+2}\to H^r$ is an isomorphism including the zero Fourier mode.

The Gauss and vector-momentum equations are solved first by the compatible elliptic inverses of \Cref{ass:constrained-initial-solver}.  For a small transverse--traceless tensor seed $\sigma$ and a small standing spinor seed $\chi=\varepsilon f\chi_*$ with real $f$, assume the complete momentum and gauge-charge compatibility conditions of \Cref{ass:constrained-initial-solver}; the resulting source coefficients in the scalar conformal equation obey
\begin{equation}
 \norm{\delta\mathsf A}_{H^r}
 +\norm{\delta\mathsf B}_{H^r}
 +\norm{\mathsf Y}_{H^r}
 \le C\bigl(\norm\sigma_{H^{r+1}}^2+\norm\chi_{H^{r+1}}^2\bigr).
 \label{eq:generated-source-small}
\end{equation}
The Hamiltonian equation can be written
\begin{equation}
 -8\Delta w-\mathsf A w^{-7}-G_*^0w
 -\mathsf B w^{-3}-\mathsf Y w^{-1}+C_*^0w^5=0,
 \label{eq:generated-Hamiltonian}
\end{equation}
where $C_*^0$ is fixed by the background constraint.

\begin{proposition}[Constructive constrained initial data]
\label{prop:generated-initial}
Under \Cref{ass:constrained-initial-solver}, for every integer $r\ge2$, sufficiently small compatible $\sigma,\chi$ as above determine a unique solution $w$ of \eqref{eq:generated-Hamiltonian} in a fixed $H^{r+2}$ neighborhood of one, with
\begin{equation}
 \norm{w-1}_{H^{r+2}}
 \le C\bigl(\norm\sigma_{H^{r+1}}^2+\norm\chi_{H^{r+1}}^2\bigr),
 \qquad \tfrac12<w<\tfrac32.
 \label{eq:generated-w-bound}
\end{equation}
Together with the Gauss and vector solutions this gives exact Einstein, gauge and matter initial constraints.  Nonzero-spinor data are included only when their complete sources satisfy the stated compatibility and conformal hypotheses; their existence is not inferred from the standing-wave ansatz alone.  Nonconstant admissible source coefficients may produce inhomogeneous metric response.

If the seeds have $p$ additional derivatives, Fourier projection of the seeds followed by the same finite elliptic construction gives finite initial tuples $U_{0,N}$ satisfying \eqref{eq:generated-initial-rate}; the Gauss and vector constraints have the accuracy supplied by that compatible finite solver, and the scalar Hamiltonian residual is $O(N^{-p})$ in the indicated lower Sobolev norm.  No spacetime solution is used in this finite initial solver.
\end{proposition}
\begin{proof}
Write $u=w-1$ and subtract the linear term $L_*u$ from \eqref{eq:generated-Hamiltonian}.  The remaining pure-background nonlinearities are quadratic in $u$, while the perturbed source coefficients satisfy \eqref{eq:generated-source-small}.  Hence
\[
 u=L_*^{-1}\bigl[\delta\mathsf A(1+u)^{-7}
 +\delta\mathsf B(1+u)^{-3}+\mathsf Y(1+u)^{-1}-R_0(u)\bigr].
\]
On a small $H^{r+2}$ ball, Sobolev multiplication and the convergent power series for the negative integer powers give
\[
 \norm{\mathcal T(0)}_{H^{r+2}}\le C\eta,
 \qquad
 \norm{\mathcal T(u)-\mathcal T(v)}_{H^{r+2}}
 \le C(\eta+R)\norm{u-v}_{H^{r+2}},
\]
where $\eta=\norm\sigma^2+\norm\chi^2$.  Banach contraction proves \eqref{eq:generated-w-bound}.  The Gauss and vector equations use bounded elliptic inverses of their compatible sources, so all initial constraints hold after substitution into the conformal identities.  For the finite construction, project the finite seeds, solve those elliptic equations exactly, and apply the projected contraction to the scalar equation.  The contraction constant remains uniform because $P_N$ is a contraction in Sobolev space and commutes with $L_*$.  Subtracting the exact and projected fixed-point equations and using the ordinary Fourier tail estimate gives $O(N^{-p})$.  The same argument applies to the derived harmonic initial completion and actual first jets.
\end{proof}

\begin{lemma}[Physical identification of the symmetric extension]
\label{lem:generated-physical-identification}
Let $U_0$ belong to the constrained Sobolev class of \Cref{prop:generated-initial}, with sufficiently many finite derivatives for the estimates below.  The independent symmetric system
\begin{equation}
 U_t+A^i(U)\partial_iU=F(U)
 \label{eq:generated-extension}
\end{equation}
has a common local Sobolev solution, and on a possibly smaller interval its head fields satisfy the complete torsion-free Einstein--Standard-Model equations and the harmonic/temporal-gauge constraints.  Thus the physical solution $z$ used in \Cref{thm:generated-dynamics} is identified from the constrained initial data; it is not an external trajectory supplied to the finite recursion.
\end{lemma}
\begin{proof}
Approximate the finite-Sobolev seeds by Fourier-polynomial seeds and apply the exact, unprojected contraction of \Cref{prop:generated-initial} to each approximation.  This gives a sequence of analytic initial tuples satisfying all physical constraints exactly and converging to $U_0$ in the Sobolev topology.  For each tuple in the declared constraint-compatible analytic core, use its physical analytic germ in harmonic coordinates and temporal internal gauge.  Its actual first jets satisfy \eqref{eq:generated-extension} by \Cref{prop:actual-jet-writer}.  Uniqueness for the symmetric system therefore identifies the independent analytic solution with those actual jets on their common local interval.

The symmetric energy estimate for \eqref{eq:generated-extension} depends only on a bounded Sobolev ball and the fixed nondegenerate chart, not on the analytic radius.  Together with the physical-Cauchy continuation input in \Cref{ass:constrained-initial-solver}, this gives a common smaller interval for the physical analytic approximants.  Their actual-jet states converge in one lower Sobolev order to the symmetric-system solution with initial datum $U_0$.  The defining-jet identities, Einstein and matter residuals, Gauss constraint and harmonic constraint vanish on every analytic approximant.  The convergence of fields, first derivatives and the equation then passes these identities distributionally to the Sobolev limit; the algebraic reconstruction identifies the actual connection, curvature and complete stress.  This proves physical identification.  It is an analytic-core argument for the constrained class, not a claim that arbitrary auxiliary initial jets are physical.
\end{proof}

\subsubsection{Finite symmetric evolution and midpoint construction}

Let the coupled actual-jet system of \Cref{prop:actual-jet-writer} be written on its fixed chart as
\begin{equation}
 U_t+A^i(U)\partial_iU=F(U),
 \qquad G(U)=F(U)-A^i(U)\partial_iU.
 \label{eq:generated-symmetric-system}
\end{equation}
For the finite initial data above define the semidiscrete system
\begin{equation}
 \dot U_N=P_NG(U_N),\qquad U_N(0)=U_{0,N}.
 \label{eq:generated-Galerkin}
\end{equation}
The symmetry of the principal matrices implies that the $H^q$ energy estimate for \eqref{eq:generated-Galerkin} has constants independent of $N$; the orthogonal projection disappears from the energy pairing.  A first-exit argument therefore gives a common existence interval and a uniform $H^q$ bound.  On that interval,
\begin{equation}
 \max_{0\le a\le2}
 \norm{\partial_t^a(U_N-U)}_{C_tH^{s+4-a}}
 \le C N^{-p},
 \label{eq:generated-Galerkin-rate}
\end{equation}
where $U$ is the physical Sobolev solution with initial datum $U_0$.  This comparison identifies the limit only after the finite ODE has been constructed.

For the fully finite evolution use \eqref{eq:generated-midpoint}.  On $\operatorname{Ran}P_N$ the inverse inequality gives
\[
 \norm{G_N(V)-G_N(W)}_{H^q}
 \le C_R(N+1)\norm{V-W}_{H^q}.
\]
Thus \eqref{eq:generated-CFL} makes the midpoint fixed-point map a contraction.  The exact discrete symmetric energy identity then removes the apparent $N$-dependence from the propagated bound and discrete Gronwall gives \eqref{eq:generated-uniform}.  If two midpoint trajectories have a normalized forcing difference $r_j$, the same energy calculation at any $3\le m\le q-1$ gives
\begin{equation}
 \max_{j\tau\le T}\norm{U^j-V^j}_{H^m}
 \le C_T\left(\norm{U^0-V^0}_{H^m}
 +\tau\sum_j\norm{r_j}_{H^m}\right),
 \label{eq:generated-causal}
\end{equation}
with $C_T$ independent of $N,\tau$.  Applying this estimate to the local truncation residual of the exact semidiscrete solution yields
\begin{equation}
 \max_j\norm{U^j-U_N(t_j)}_{H^{s+4}}\le C\tau^2,
 \qquad
 \max_j\norm{U^j-U(t_j)}_{H^{s+4}}
 \le C(\tau^2+N^{-p}).
 \label{eq:generated-time-rate}
\end{equation}

\subsubsection{Physical Hermite readout and equation residuals}

On each interval $[t_j,t_{j+1}]$, reconstruct the physical head by the cubic Hermite polynomial having endpoint values equal to the generated physical heads and endpoint slopes equal to the corresponding components of $G_N(U^j)$ and $G_N(U^{j+1})$.  The resulting $z_{N,\tau}$ is $C^1$ in time.  The Hermite basis, \eqref{eq:generated-time-rate}, and the uniform higher time derivatives of the semidiscrete system give
\begin{align}
 \norm{z_{N,\tau}-z}_{C_tH^{s+2}}
 +\norm{\partial_tz_{N,\tau}-\partial_tz}_{C_tH^{s+1}}
 &\le C(N^{-p}+\tau^2),\notag\\
 \norm{\partial_t^2z_{N,\tau}-\partial_t^2z}_{L^\infty_tH^s}
 &\le C(N^{-p}+\tau).
 \label{eq:generated-Hermite}
\end{align}
The Riemann tensor is linear in second metric derivatives plus smooth quadratic first-derivative terms, giving \eqref{eq:generated-bosonic}.  The wave-type Yang--Mills and Higgs rows have the same derivative count.  The Dirac row is first order and therefore inherits the sharper estimate \eqref{eq:generated-Dirac}.  The harmonic, Gauss and normal Einstein constraints use only the controlled first jets and spatial derivatives and are $O(N^{-p}+\tau^2)$ in the corresponding lower norms.

Finally let $S^{(1)}$ be the complete first-derivative torsion-free Einstein--Standard-Model action representative with fixed-collar boundary completion, and let $\mathcal R$ be the Hermite reconstruction from independent nodal values and slopes.  Define
\begin{equation}
 S^{\rm cmp}_{N,\tau}=S^{(1)}\circ\mathcal R,
 \qquad
 \norm V_{0,\tau}^2=\tau\sum_j
   \left(\norm{v_j}_2^2+\tau^2\norm{\dot v_j}_2^2\right).
 \label{eq:generated-comparison-action}
\end{equation}
At the generated record, integration by parts expresses $D S^{\rm cmp}_{N,\tau}$ through the physical rows just estimated.  Since $\norm{D\mathcal R[V]}_{L^2}\le C\norm V_{0,\tau}$ and the temporal gauge and all metric entries remain independent variations before evaluation, this gives \eqref{eq:generated-stationarity}.  Keeping the first-derivative variational representative and using \eqref{eq:generated-first} gives the sharper $O(N^{-p}+\tau^2)$ continuum first-variation comparison.  In particular the complete symmetric spinor stress is included and converges strongly.

\section{Finite source budgets and same-record closure}
\label{sec:native-selection}

The finite-action theorem above is deterministic and applies to any complete record satisfying the displayed physical budgets.  We now give a finite source-level entrance that produces those budgets on one \emph{actually observed} record.  The construction does not replace a record by its expectation, by a projected state, or by a trajectory of an inserted evolution equation.

At spacing $h$, let $X_0,\ldots,X_J$ be the records of an actual finite accepted process.  The state may retain the entire finite past, so no Markov assumption is required.  Let $K^{\rm leg}_{h,j}(x,y)$ denote the transition probabilities restricted by a predeclared legal-edge mask for the raw-action chart and action identification; an attempted illegal transition is counted as failure and the diagnostic stopped process retains its last legal action value.  Write $\mu_j$ for the resulting unconditioned surviving mass and $p_h^{\rm exit}=1-\sum_x\mu_J(x)$.  Every legal state has a complete physical record $u_h(x)$ and define
\begin{equation}
 g_h(x)=\norm{E_h^{\rm raw}(u_h(x))}_{0,h},\qquad
 W_h(x)=\bigl(i_{h,k}(x)+\gamma_{h,k}(x)\bigr)^2,
 \label{eq:source-gW}
\end{equation}
together with a nonnegative native energy readout $V_h(x)$.  Put
\begin{equation}
 M_{h,J}=\frac1J\sum_{j<J,x}\mu_j(x)V_h(x),\qquad
 A_{h,J}=\frac1J\sum_{j<J,x}\mu_j(x)W_h(x).
 \label{eq:source-occupations}
\end{equation}
Let $a_h(x)=S_h^{\rm loc}(u_h(x))$ and suppose $A_h^-\le a_h(x)\le A_h^+$ on the legal chart, with action span $D_h^A=A_h^+-A_h^-$.  The stopped conditional action increment is
\begin{equation}
 \mathfrak d_{h,j}(x)=\sum_yK^{\rm leg}_{h,j}(x,y)\,[a_h(y)-a_h(x)].
 \label{eq:source-action-drift}
\end{equation}

\begin{theorem}[Same-record selection from the actual signed action drift]
\label{thm:source-selection}
Assume there are $\delta_h>0$, $\kappa_A>0$ independent of $h$, and nonnegative finite remainders $r_{h,j}$ such that every active legal row satisfies
\begin{equation}
 \mathfrak d_{h,j}(x)
 \le-\kappa_A\delta_h g_h(x)^2+\delta_h r_{h,j}(x).
 \label{eq:source-drift-ineq}
\end{equation}
Set
\begin{equation}
 \mathcal R_{h,J}=\frac1J\sum_{j<J,x}\mu_j(x)r_{h,j}(x).
 \label{eq:source-work-occupation}
\end{equation}
Then
\begin{equation}
 \frac1J\sum_{j<J,x}\mu_j(x)g_h(x)^2
 \le \frac{D_h^A}{\kappa_A\delta_hJ}
       +\frac{\mathcal R_{h,J}}{\kappa_A}.
 \label{eq:source-slope-occupation}
\end{equation}
For thresholds $s_h,R_h,\alpha_h>0$, select on every surviving prefix an observed index minimizing
\begin{equation}
 \frac{g_h(X_j)^2}{s_h^2}
 +\frac{V_h(X_j)}{R_h^2}
 +\frac{W_h(X_j)}{\alpha_h^2}.
 \label{eq:source-score}
\end{equation}
The probability of failing to output one observed record with
\begin{equation}
 g_h\le s_h,\qquad V_h\le R_h^2,\qquad
 i_{h,k}+\gamma_{h,k}\le\alpha_h
 \label{eq:source-selected}
\end{equation}
is at most
\begin{equation}
 \boxed{
 p_h^{\rm exit}
 +\frac{D_h^A}{\kappa_A\delta_hJ s_h^2}
 +\frac{\mathcal R_{h,J}}{\kappa_A s_h^2}
 +\frac{M_{h,J}}{R_h^2}
 +\frac{A_{h,J}}{\alpha_h^2}.}
 \label{eq:source-selection-failure}
\end{equation}
No invariant distribution, mixing estimate, independence between records, or uniform high-order bound on every member of the prefix is required.
\end{theorem}
\begin{proof}
The stopped-action telescope and the deterministic score argument are proved in \Cref{app:native-source-selection}.  The key point is that killing is recorded as failure rather than being credited as a decrease of the signed action.
\end{proof}

The native energy must control the \emph{physical reader actually entering} \eqref{eq:native-tail}.  This relation has an exact finite-dimensional criterion.

\begin{proposition}[Exact native-energy to physical-tail certificate]
\label{prop:native-reader}
Let $\mathsf H_h^{\rm src}$ be a finite real source-coordinate Hilbert space, let $\mathsf Q_h\succeq0$ be a positive native energy form, and let $\xi_h(x)\in\mathsf H_h^{\rm src}$ be the source coordinate of a legal record.  Suppose the physical reader has the affine form
\begin{equation}
 u_h(x)=u_h^{\rm soft}+\mathsf R_h\xi_h(x).
 \label{eq:source-reader}
\end{equation}
Let $\mathsf T_{h,K,m}$ map a source vector to the high-frequency
Fourier coefficients of $\mathsf R_h\xi$, with the weights of
\eqref{eq:native-tail}.  A finite constant $c_{h,K,m}$ such that
\begin{equation}
 \norm{\mathsf T_{h,K,m}\xi}_{\ell^1(w)}
 \le c_{h,K,m}\,\ip{\xi}{\mathsf Q_h\xi}^{1/2}
 \qquad(\xi\in\mathsf H_h^{\rm src})
 \label{eq:source-reader-linear}
\end{equation}
exists if and only if
\begin{equation}
 \mathsf T_{h,K,m}\ker\mathsf Q_h=0.
 \label{eq:source-reader-kernel}
\end{equation}
When this holds, the optimal constant in
\eqref{eq:source-reader-linear} is
\begin{equation}
 c_{h,K,m}
 =\norm{\mathsf T_{h,K,m}\mathsf Q_h^{\dagger/2}}_{2\to\ell^1(w)}.
 \label{eq:source-reader-constant}
\end{equation}
In particular every legal record obeys the affine estimate
\begin{equation}
 \tau_{h,m}(K;u_h(x))
 \le \tau_{h,m}(K;u_h^{\rm soft})
      +c_{h,K,m}\,\ip{\xi_h(x)}{\mathsf Q_h\xi_h(x)}^{1/2}.
 \label{eq:source-reader-tail}
\end{equation}
The equivalence concerns the linear estimate on the entire source space;
it is not inferred from testing only a finite list of legal states.
Thus the leakage assumption is a finite, directly auditable incidence condition between the native energy and the declared physical reader; finite-dimensional norm equivalence alone is not used.
\end{proposition}
\begin{proof}
See \Cref{app:native-reader-proof}.
\end{proof}

\begin{corollary}[Sobolev positive-form reader criterion]
\label{cor:native-reader-sobolev}
Let $\mathsf L_h$ be the diagonal multiplier $(1+|\ell|_2^2)^{1/2}$ on the all-mode Fourier reader.  If, on the retained source coordinates,
\begin{equation}
 \mathsf R_h^*\mathsf L_h^{2q}\mathsf R_h\preceq C_R\mathsf Q_h,
 \qquad \norm{u_h^{\rm soft}}_{H^q}\le C_R,
 \label{eq:source-reader-LMI}
\end{equation}
then for every $0\le j\le m<q-2$,
\begin{equation}
 \norm{u_h(x)}_{H^q}\le C\bigl(1+\sqrt{V_h(x)}\bigr),\qquad
 \tau_{h,j}(K;u_h(x))
 \le C K^{2-q}\bigl(1+\sqrt{V_h(x)}\bigr),
 \label{eq:source-reader-Sobolev}
\end{equation}
whenever $\ip{\xi_h}{\mathsf Q_h\xi_h}\le V_h$.  The power $K^{2-q}$ is sharp for this four-dimensional Sobolev-to-weighted-tail route; in particular the composition with the $k$th-order forcing estimate requires $q>k+5$ if no additional cancellation is used.
\end{corollary}
\begin{proof}
The positive-form inequality gives the $H^q$ estimate.  The sharp four-dimensional weighted Fourier bound is proved in \Cref{app:native-reader-proof}; inserting $m=k+2$ into the forcing budget yields the stated threshold.
\end{proof}


\begin{corollary}[Same-record action, reader, and initial-constraint selection]
\label{cor:source-selected-shadow-budgets}
Assume the signed action-drift hypotheses of \Cref{thm:source-selection}.  In
addition, let every legal record carry a nonnegative composite initial
constraint residual
\begin{equation}
 c_h(x)=\abs{\Phi_h(x,0)},
 \label{eq:source-composite-charge}
\end{equation}
where \Cref{def:exact-initial-reduction} holds for the normalized actual
initial record. Thus $\Phi_h$ is the retained exact continuum constraint
map, and its complementary constraints vanish throughout the correction
family. A certified upper bound such as
$|\widetilde\Phi_h(x,0)|+\zeta_{0,h}(x)$ from
\Cref{cor:validated-initial-retraction} may replace $c_h(x)$ in both the
score and its occupation bound. Put
\begin{equation}
 C_{h,J}=\frac1J\sum_{j<J,x}\mu_j(x)c_h(x)^2.
 \label{eq:source-charge-occupation}
\end{equation}
For thresholds $s_h,R,\beta_h>0$, select on each surviving prefix an observed
index minimizing
\begin{equation}
 \frac{g_h(X_j)^2}{s_h^2}
 +\frac{V_h(X_j)}{R^2}
 +\frac{c_h(X_j)^2}{\beta_h^2}.
 \label{eq:source-shadow-score}
\end{equation}
Then the probability of failing to output one actual unfiltered record with
\begin{equation}
 g_h\le s_h,
 \qquad
 V_h\le R^2,
 \qquad
 c_h\le\beta_h
 \label{eq:source-shadow-selected}
\end{equation}
is at most
\begin{equation}
 \boxed{
 p_h^{\rm exit}
 +\frac{D_h^A}{\kappa_A\delta_hJ s_h^2}
 +\frac{\mathcal R_{h,J}}{\kappa_A s_h^2}
 +\frac{M_{h,J}}{R^2}
 +\frac{C_{h,J}}{\beta_h^2}.}
 \label{eq:source-shadow-selection-failure}
\end{equation}
No independence, stationarity, or conditioning on successful selection is
used.
\end{corollary}
\begin{proof}
Equation \eqref{eq:source-slope-occupation} bounds the occupation of $g_h^2$.
On survival, the minimum of the nonnegative score
\eqref{eq:source-shadow-score} is at most its arithmetic mean along the
observed prefix.  If that mean is at most one, each of the three selected
terms is at most one.  Markov's inequality bounds failure on survival by the
expectation of the mean score, and adding the original exit event gives
\eqref{eq:source-shadow-selection-failure}.
\end{proof}

\begin{theorem}[Same-source selected exact physical shadow]
\label{thm:native-selected-shadow}
Assume \Cref{ass:physical-Cauchy-tube}, the hypotheses of
\Cref{cor:source-selected-shadow-budgets}, and the reader criterion
\Cref{cor:native-reader-sobolev} for some indices satisfying
\eqref{eq:shadow-regularity-indices}.  Require every legal record used by the
selector to lie in the common buffered field-value and inverse-metric chart
and to admit the regular harmonic/temporal normalization of
\Cref{thm:aposteriori-physical-shadow}; violations are included in
$p_h^{\rm exit}$.  Assume also the exact family and basepoint identification of
\Cref{def:exact-initial-reduction}, the uniform initial-constraint margins
\eqref{eq:constraint-linear-margin}--\eqref{eq:constraint-derivative-margin}
and the reconstruction Lipschitz constant $B_*$.  Require the corrected
data to remain in $\mathcal C_k^{\rm reg}\cap\mathcal O_k$ with the certified
initial regularity; any failure of this condition is included in the legal
domain exit event.

On the success event of \Cref{cor:source-selected-shadow-budgets}, define
\begin{equation}
 \varepsilon_h=C_R(s_h+h),
 \qquad
 \eta_{h,k}=C_R\left(
 \varepsilon_h^{1-k/s}
 +\varepsilon_h^{1-(k+1)/s}\right),
 \label{eq:selected-shadow-forcing} 
\end{equation}
and
\begin{equation}
 \bar d_h
 =2B_*\gamma_*^{-1}\beta_h
  +C_R\sqrt T\,\eta_{h,k}.
 \label{eq:selected-shadow-budget}
\end{equation}
If $\bar d_h$ is below the uniform Cauchy-tube stability radius, then the
selected actually observed record admits the unique exact physical
Einstein--Standard-Model solution $z_{*,h}$ with the corrected initial datum
from \Cref{prop:initial-constraint-retraction}, on the common slab, and
\begin{align}
 &\norm{\mathcal U(\widehat z_{h,j_*})-\mathcal U(z_{*,h})}_{C_tH_x^k}
 +\norm{\Riem(\widehat g_{h,j_*})-\Riem(g_{*,h})}_{L_t^2H_x^{k-1}}
 \notag\\
 &\qquad
 +\norm{T^{\rm SM}(\widehat z_{h,j_*})-T^{\rm SM}(z_{*,h})}_{L_t^2H_x^k}
 \le C_R\bar d_h.
 \label{eq:selected-exact-shadow}
\end{align}
The probability that the selection stage fails before this deterministic
conclusion is bounded by \eqref{eq:source-shadow-selection-failure}.  Thus the
same native-law record supplies the finite Euler row, the regularity reserve,
and the near-compatible initial datum from which its exact physical comparison
future is constructed.
\end{theorem}
\begin{proof}
The selected record satisfies \eqref{eq:source-shadow-selected}.
\Cref{cor:native-reader-sobolev,prop:shadow-residual-upgrade} give the
forcing budget, and \Cref{prop:initial-constraint-retraction} gives the
exact compatible initial datum. Apply
\Cref{thm:aposteriori-physical-shadow}; the probability bound is
\eqref{eq:source-shadow-selection-failure}.
\end{proof}

\begin{theorem}[Same-source selected finite-action Einstein--Standard-Model closure]
\label{thm:native-selected-closure}
Let $k\ge4$ and fix the same regular reference solution as in \Cref{thm:native-closure}.  Assume the hypotheses of \Cref{thm:source-selection}.  Every legal record used here is required to exist on the common slab, to be in temporal internal gauge, and to satisfy the buffered field-value, inverse-metric and reconstruction-chart conditions of \Cref{thm:native-source}.  The legal-domain failure probability includes violations of these conditions.  Suppose the same legal records and the same native energy satisfy, for $j=2,m=k+2$,
\begin{equation}
 \tau_{h,j}(K_h;u_h(x))\le t_{h,j}+c_{h,j}\sqrt{V_h(x)},
 \label{eq:selected-reader-general}
\end{equation}
with certified finite reader constants.  Put $T_{h,j}=t_{h,j}+c_{h,j}R_h$ and define
\begin{align}
 \bar\varepsilon_{c,h}&=s_h+hK_h^3+K_h^2T_{h,2},\notag\\
 \bar L_{h,k}&=1+\log\left(2+\frac{C_kK_h^{k+5}}{\bar\varepsilon_{c,h}}\right),\notag\\
 \bar F_{h,k}&=\bar\varepsilon_{c,h}K_h^{k+1}\bar L_{h,k}^{k+1}
                   +K_h^{k+2}T_{h,m},\qquad
 \bar D_{h,k}&=\alpha_h+\bar F_{h,k}.
 \label{eq:selected-composed-budget}
\end{align}
Assume $K_h\to\infty$, $hK_h\to0$, $T_{h,m}\to0$, $\bar\varepsilon_{c,h}\to0$ and $\bar D_{h,k}\to0$.  Outside the event in \eqref{eq:source-selection-failure}, the selected actually observed record satisfies
\begin{align}
 &\norm{\mathcal U_{h,j_*}-\mathcal U_*}_{C_tH_x^k}
 +\norm{\Riem(g_{h,j_*})-\Riem(g_*)}_{L^2_tH_x^{k-1}}\notag\\
 &\hspace{4em}
 +\norm{T^{\rm SM}(z_{h,j_*})-T^{\rm SM}(z_*)}_{L^2_tH_x^k}
 \le C\bar D_{h,k}.
 \label{eq:selected-closure}
\end{align}
Its zero-order physical residual is at most $C(s_h+hK_h^3)$ and its literal Cartan discrepancy is at most $ChK_h^3$.  If the failure bounds and $\bar D_{h_n,k}$ are summable along a common cofinal chain, then under any common coupling the selected successful records eventually have summable adjacent state, curvature and stress differences almost surely.  No independence between cutoffs is required.
\end{theorem}
\begin{proof}
Selection gives $\sigma_h\le s_h$, initial/gauge mismatch at most
$\alpha_h$, and $V_h\le R_h^2$ on the same record. The reader certificate
then gives $\tau_{h,j}\le T_{h,j}$. Apply
\Cref{thm:native-source,thm:native-closure} with these upper budgets.
The first Borel--Cantelli lemma and comparison through the common reference
prove eventual success and summable adjacent differences.
\end{proof}

\begin{corollary}[An explicit polynomial selected-source regime]
\label{cor:selected-polynomial}
Under the signed-drift, common-slab and legal-chart hypotheses of \Cref{thm:native-selected-closure}, fix $k\ge4$, put $m=k+2$, and suppose \Cref{cor:native-reader-sobolev} holds for some $q>k+5$.  Assume, for positive $a,b,c,\epsilon$, $d,d_A\ge0$,
\begin{align}
 \delta_h&\ge c_\delta h^d,& D_h^A&\le Ch^{-d_A},\notag\\
 J_h&\ge c_Jh^{-(d+d_A+2b+2\epsilon)},&
 \mathcal R_{h,J_h}&\le Ch^{2b+2\epsilon},\notag\\
 M_{h,J_h}&\le C,& A_{h,J_h}&\le Ch^{2c+2\epsilon},\notag\\
 p_h^{\rm exit}&\le Ch^{2\epsilon}.&&
 \label{eq:selected-polynomial-source}
\end{align}
Choose
\begin{equation}
 s_h=h^b,\qquad R_h=h^{-a},\qquad \alpha_h=h^c,
 \qquad K_h\asymp h^{-\beta},
 \label{eq:selected-polynomial-thresholds}
\end{equation}
where
\begin{equation}
 0<\beta<\min\left\{\frac{b}{k+1},\frac1{k+4}\right\},
 \qquad 0<a<\beta(q-k-5).
 \label{eq:selected-polynomial-window}
\end{equation}
Then the selection failure probability is $O(h^{2a}+h^{2\epsilon})$ and, on success,
\begin{equation}
 \text{left side of \eqref{eq:selected-closure}}
 \le C h^\rho(1+|\log h|)^{k+1},
 \quad
 \rho=\min\{c,b-\beta(k+1),1-\beta(k+4),\beta(q-k-5)-a\}>0.
 \label{eq:selected-polynomial-rate}
\end{equation}
Along $h_n\asymp2^{-n}$ the successful selected branch is eventually almost sure and has summable adjacent differences whenever the displayed probability bound is summable.
\end{corollary}
\begin{proof}
See \Cref{app:native-polynomial-proof}.
\end{proof}

\begin{corollary}[Weak-sector example of a reader certificate]
\label{cor:gauge-robust-reader}
On a nonvanishing Higgs chart of the weak $SU(2)$ sector, a cutoff-uniform bound on the finite covariant-jet energy of \Cref{app:gauge-reader} yields, after algebraic Higgs normalization and a causal finite temporal gauge, a bounded all-mode $H^q$ reader and hence the tail estimate
\begin{equation}
 \tau_{h,m}(K)\le C K^{2-q},\qquad q>m+2,
 \label{eq:gauge-robust-tail}
\end{equation}
without requiring high ordinary derivatives of the original site frame or small original links.  Exact finite gauge invariance preserves the norm of the \emph{full} physical action covector under this normalization, and logarithmic link coordinates are uniformly equivalent to the physical link-tangent norm on the normalized chart.  This result supplies one source-auditable realization of the reader hypothesis for the weak sector; color, determinant-line, coframe and other independent reader components require their own corresponding bounds.
\end{corollary}
\begin{proof}
The finite covariant-jet, temporal-normalization and full-covector identities are proved in \Cref{app:gauge-reader}.
\end{proof}

\paragraph{Alternative same-action prefix balance.}
The direct signed action-drift entrance above is not the only way to obtain a small native Euler row.  A strongly convex operational potential gives the following independent sufficient criterion.  It is useful when its force-balance discrepancy is directly available, but is not imposed on source laws for which action drift can be verified more economically.

Give the free physical nodal variables the mass norm
\begin{equation}
 \norm v_{0,h}^2=h^4\sum_{x\in Q'_h}|v(x)|^2.
 \label{eq:prefix-mass}
\end{equation}
Let $\mathcal A_h=S_h^{\rm loc}$ on a compact convex field chart $\Theta_h$.  Suppose a positive operational potential $\Psi_h$ has
\begin{equation}
 g_*I\preceq \nabla_h^2\Psi_h(y)\preceq g^*I,
 \qquad 0<g_*\le g^*<\infty,
 \label{eq:prefix-Fisher}
\end{equation}
and the local action obeys
\begin{equation}
 \norm{\nabla_h^2\mathcal A_h(y)}_{0,h\to0,h}\le K_0h^{-2},
 \qquad \operatorname{osc}_{\Theta_h}\mathcal A_h\le D_0.
 \label{eq:prefix-Hessian}
\end{equation}

\begin{theorem}[Alternative finite-prefix stationarity criterion]
\label{thm:prefix-stationarity}
Let $\theta_0,\ldots,\theta_J\in\Theta_h$ be an accepted finite sequence whose connecting segments stay in the chart.  Put $E_h=\nabla_h\mathcal A_h$, choose $\upsilon_h=\upsilon_0h^2$ with $\upsilon_0K_0<g_*/2$, and define
\begin{equation}
 b_j=E_h(\theta_{j+1})
 +\upsilon_h^{-1}\{\nabla_h\Psi_h(\theta_{j+1})-\nabla_h\Psi_h(\theta_j)\}.
 \label{eq:prefix-balance}
\end{equation}
Then
\begin{equation}
 \frac1J\sum_{j=0}^{J-1}\norm{E_h(\theta_j)}_{0,h}^2
 \le C_0h^{-2}\frac{D_0}{J}
 +D_1\frac1J\sum_{j=0}^{J-1}\norm{b_j}_{0,h}^2.
 \label{eq:prefix-average}
\end{equation}
Consequently the smallest actually observed full native slope satisfies
\begin{equation}
 \norm{E_h(\theta_{j_*})}_{0,h}
 \le C\left[h^{-1}\sqrt{D_0/J}
 +\left(J^{-1}\sum_{j<J}\norm{b_j}_{0,h}^2\right)^{1/2}\right].
 \label{eq:prefix-selected}
\end{equation}
In particular $D_0/J=O(h^4)$ and $J^{-1}\sum\norm{b_j}_{0,h}^2=O(h^2)$ give $\sigma_h=O(h)$ without an invariant probability distribution or an assumed small selected action covector.
\end{theorem}
\begin{proof}
The same-action descent calculation is given in \Cref{app:prefix-stationarity}.
\end{proof}

\begin{corollary}[Finite-table certification for the prefix-balance entrance]
\label{cor:prefix-table}
Suppose the accepted transitions are represented by finite tables, a safe set carries the chart and balance hypotheses above, and $p_h^{\rm exit}$ is the probability of leaving that safe set before $J$ accepted comparisons.  Put
\begin{equation}
 \mathcal Q_{h,J}^{\rm bal}
 =\frac1J\sum_{j<J}\mathbb E\!\left[
   \mathbf 1_{\rm safe}\norm{b_j}_{0,h}^2\right].
 \label{eq:prefix-table-balance}
\end{equation}
For any $0<\alpha\le1$, the probability that the best safe observed record fails a fixed-constant $O(h^\alpha)$ full native-row bound is at most
\begin{equation}
 p_h^{\rm exit}
 +C\frac{D_0}{J h^{2+2\alpha}}
 +C h^{-2\alpha}\mathcal Q_{h,J}^{\rm bal}.
 \label{eq:prefix-table-probability}
\end{equation}
If this bound is summable along a dyadic cutoff family, eventual certification holds almost surely without an independence assumption between cutoffs.
\end{corollary}
\begin{proof}
See \Cref{app:prefix-stationarity}.
\end{proof}

\begin{corollary}[Nonemptiness of the controlled finite class]
\label{cor:generated-nonempty}
A separate comparison-action approximation is available under the explicit constrained-initial-solver hypotheses of \Cref{app:generated-dynamics}.  For that class, the fully finite spectral--midpoint evolution there exists on a common local slab and produces records $z_{N,\tau}$ satisfying
\begin{equation}
 \norm{z_{N,\tau}-z}_{C_tH^{s+2}}
 +\norm{\Riem(g_{N,\tau})-\Riem(g)}_{L^\infty_tH^s}
 +\norm{\mathcal R_{\rm B}(z_{N,\tau})}_{L^\infty_tH^s}
 \le C(N^{-p}+\tau),
 \label{eq:generated-nonempty}
\end{equation}
with the Dirac residual $O(N^{-p}+\tau^2)$ and summable dyadic transport.  The construction starts from finite constrained initial data and does not sample a prescribed future solution.  Its purpose is a generated comparison-action realization; it does not by itself establish the raw local-action budgets or the source-drift conditions.  Nonemptiness for the raw local action is treated in \Cref{cor:local-calibration-nonempty}, and the conditional source-native theorem is \Cref{thm:native-selected-closure}.
\end{corollary}
\begin{proof}
This is the content of \Cref{app:generated-dynamics}.
\end{proof}

\begin{remark}[What remains outside the quantitative branch]
\label{rem:native-boundary}
For one actual accepted finite law, \Cref{thm:native-selected-closure} derives the selected Euler-row, leakage and initial/gauge budgets jointly from explicit source-level conditions.  What is not proved is that an arbitrary independently specified microscopic source satisfies the required signed action-drift, reader-energy, legal-domain and full coupled gravity--gauge--chiral-matter conditions.  In the Renewal Geometry realization assembled from the Lorentzian construction \cite{PelissierRenewalSpacetime2026}, the differential/spin construction \cite{PelissierSpectral2026}, and the structural gravity--matter construction \cite{PelissierDuality2026}, this becomes a separate provenance and source-selection problem; the complete microscopic-source Einstein--Standard-Model implication is not assumed here.
\end{remark}

\begin{remark}[All-sector reader hypothesis]
\label{rem:all-sector-reader}
The tails in \eqref{eq:native-tail} concern the full tuple
$(e,A,H,\Psi,\bPsi)$ on a fixed calibrated chart, including color,
weak, hypercharge and spin-geometric components.  The weak-sector
normalization of \Cref{app:gauge-reader} is only an example for the
$\SU(2)$ block.  Neither it nor finite gauge invariance supplies an estimate
for the other blocks.  One must verify the full-reader certificate
\Cref{prop:native-reader}, or independent component certificates combined
as in \Cref{prop:component-reader}.  No all-sector normalization theorem
is inferred from the Higgs frame.
\end{remark}

\section{Source-selection proofs and physical-reader certificates}
\label{app:source-proof-package}

\subsection{Native source selection and physical reader certificates}
\label{app:native-source-selection}

This subsection proves the finite source-level results used in \Cref{sec:native-selection}.  All probabilities refer to the actual accepted table; failures are retained, and no transition law is conditioned on success.

\subsubsection{Stopped signed-action drift and same-record selection}

Let $Z_j$ denote the diagnostic stopped action value: on a legal transition it equals $a_h(X_j)$, while after the first attempted illegal transition it retains the last legal value.  Shifting the action by the constant $A_h^-$ gives $0\le Z_j\le D_h^A$ without changing any increment or Euler covector.  Conditional on an active legal row,
\begin{equation}
 \mathbb E[Z_{j+1}-Z_j\mid X_j=x]=\mathfrak d_{h,j}(x).
 \label{eq:app-stopped-drift}
\end{equation}
Summing \eqref{eq:source-drift-ineq} against the surviving subprobability masses and telescoping gives
\begin{equation}
 \kappa_A\delta_h\sum_{j<J,x}\mu_j(x)g_h(x)^2
 \le D_h^A+\delta_h\sum_{j<J,x}\mu_j(x)r_{h,j}(x),
 \label{eq:app-slope-telescope}
\end{equation}
which is \eqref{eq:source-slope-occupation}.

For the selection step define
\[
 \mathcal S_j=\frac{g_h(X_j)^2}{s_h^2}
 +\frac{V_h(X_j)}{R_h^2}
 +\frac{W_h(X_j)}{\alpha_h^2}.
\]
On a surviving prefix, if $j_*$ minimizes $\mathcal S_j$ and $J^{-1}\sum_{j<J}\mathcal S_j\le1$, then every term in \eqref{eq:source-selected} is bounded by its threshold.  Hence selection failure on survival implies that the mean score exceeds one.  Markov's inequality, \eqref{eq:source-slope-occupation}, and the definitions \eqref{eq:source-occupations} give exactly \eqref{eq:source-selection-failure}.  Notice that no union bound over all candidates is used; the high-order energy and alignment conditions are required only of the selected record.

\subsubsection{Exact reader constant and sharp Fourier transfer}
\label{app:native-reader-proof}

For the affine reader \eqref{eq:source-reader}, collect the high-frequency coefficient maps in
\[
 \mathsf T_{h,K,m}\xi=(T_{h,\ell}\xi)_{|\ell|_\infty>K},
 \qquad w_\ell=(1+|\ell|_1/K)^m.
\]
If $\xi\in\ker\mathsf Q_h$, the energy side of \eqref{eq:source-reader-linear} vanishes; hence \eqref{eq:source-reader-kernel} is necessary.  A bound only on a prescribed finite set of record coordinates would not imply this universal kernel condition.  Conversely restrict to $(\ker\mathsf Q_h)^\perp$ and write $\xi=\mathsf Q_h^{\dagger/2}y$.  Then
\[
 \norm y^2=\ip{\xi}{\mathsf Q_h\xi},
\]
and weighted block-$\ell^1$ duality gives
\begin{align}
 \sup_{\ip{\xi}{\mathsf Q_h\xi}\le1}
 \sum_{|\ell|_\infty>K}w_\ell\norm{T_{h,\ell}\xi}
 &=\norm{\mathsf T_{h,K,m}\mathsf Q_h^{\dagger/2}}_{2\to\ell^1(w)}.
 \label{eq:app-reader-dual}
\end{align}
This proves \Cref{prop:native-reader}; the affine offset is handled by the triangle inequality.

For completeness, on the normalized four-torus Parseval and Cauchy--Schwarz give
\begin{equation}
 \tau_{h,m}(K;u)
 \le A_{N,K,m,q}\norm u_{H^q},\qquad
 A_{N,K,m,q}^2=
 \sum_{K<|\ell|_\infty\le N}
 \frac{(1+|\ell|_1/K)^{2m}}{(1+|\ell|_2^2)^q}.
 \label{eq:app-exact-tail-sum}
\end{equation}
Dyadic shells in four dimensions contain $O(R^4)$ lattice points.  On a shell of radius $R\asymp2^aK$, the contribution to \eqref{eq:app-exact-tail-sum} is bounded by $CK^{-2m}R^{2m-2q+4}$.  Thus
\begin{equation}
 \tau_{h,m}(K;u)\le C_{q,m}K^{2-q}\norm u_{H^q}
 \qquad(q>m+2),
 \label{eq:app-sharp-tail}
\end{equation}
and the exponent is sharp on shell-supported examples.  At $q=m+2$ the dyadic contributions are comparable and the constant diverges logarithmically with the ultraviolet cutoff.  Combining \eqref{eq:source-reader-LMI} with \eqref{eq:app-sharp-tail} proves \Cref{cor:native-reader-sobolev}.

\subsubsection{A directly verifiable action-drift mechanism}

The drift hypothesis \eqref{eq:source-drift-ineq} is itself testable against the actual transition law and the same action.  Let $E=\nabla\mathcal A$ in the physical mass norm and $\Delta=Y-x$ for one legal transition.  Suppose
\begin{equation}
 m(x):=\mathbb E[\Delta\mid x]
 =-\delta M(x)E(x)+\delta e(x),
 \qquad \frac{M+M^*}{2}\succeq\lambda I,
 \label{eq:app-mean-incidence}
\end{equation}
with $\lambda>0$.  Taylor's integral identity gives
\begin{equation}
 \mathbb E[\mathcal A(Y)-\mathcal A(x)\mid x]
 =\ip{E(x)}{m(x)}+\mathcal W(x),
 \quad
 \mathcal W(x)=\mathbb E\!\left[\int_0^1(1-t)D^2\mathcal A(x+t\Delta)[\Delta,\Delta]\,dt\,\middle|\,x\right].
 \label{eq:app-signed-work}
\end{equation}
Therefore
\begin{equation}
 \mathfrak d(x)\le-\frac\lambda2\delta\norm E^2
 +\delta\left(\frac{\norm e^2}{2\lambda}+\frac{[\mathcal W(x)]_+}{\delta}\right).
 \label{eq:app-work-drift}
\end{equation}
Only the signed action work of the innovation is charged.  In particular, if the actual transition factors as $K(x,dy)=\int D(x,dz)N_x(z,dy)$ and the second kernel is action-neutral,
\begin{equation}
 \int\mathcal A(y)N_x(z,dy)=\mathcal A(z),
 \label{eq:app-action-neutral}
\end{equation}
then its stochastic variance contributes nothing to the action-drift remainder.  The selected record itself, including that retained noise, is still tested by $V_h$ and $W_h$.  These identities do not authorize inserting a gradient step or an action-neutral component into a source that does not actually possess them.

\subsubsection{Polynomial regime}
\label{app:native-polynomial-proof}

Under \eqref{eq:selected-polynomial-source}--\eqref{eq:selected-polynomial-window}, the five terms in \eqref{eq:source-selection-failure} are respectively $O(h^{2\epsilon})$, $O(h^{2\epsilon})$, $O(h^{2a})$, $O(h^{2\epsilon})$, and the declared exit term.  Hence the failure probability is $O(h^{2a}+h^{2\epsilon})$.  The Sobolev reader gives
\begin{equation}
 T_{h,j}\le C h^{-a}K_h^{2-q},\qquad j=2,k+2.
 \label{eq:app-polynomial-tail}
\end{equation}
Consequently
\[
 \bar\varepsilon_{c,h}
 \le C\bigl(h^b+hK_h^3+h^{-a}K_h^{4-q}\bigr),
 \qquad \bar L_{h,k}\le C(1+|\log h|),
\]
and the terms in $\bar F_{h,k}$ have orders
\[
 h^bK_h^{k+1},\qquad hK_h^{k+4},\qquad
 h^{-a}K_h^{k+5-q},\qquad h^{-a}K_h^{k+4-q},
\]
up to the fixed logarithmic factor.  The fourth exponent is larger than the third by $\beta$; adding $\alpha_h=h^c$ gives precisely \eqref{eq:selected-polynomial-rate}.  Positive powers of a dyadic spacing times a fixed logarithmic power are summable.

\begin{proposition}[Assembly of full-reader bounds from component certificates]
\label{prop:component-reader}
In a fixed compatible atlas write the physical reader as the finite tuple
$u=(u^{(1)},\ldots,u^{(B)})$, with blocks for the independent geometric,
gauge and matter components.  Suppose on the same source record
\[
 \tau_{h,m}(K;u^{(b)})\le t_b+c_b\sqrt{V_b},\qquad
 V=\sum_{b=1}^BV_b.
\]
Then, up to the fixed equivalence constants of the component norms,
\begin{equation}
 \tau_{h,m}(K;u)
 \le\sum_{b=1}^Bt_b+
       \left(\sum_{b=1}^Bc_b^2\right)^{1/2}\sqrt V.
 \label{eq:component-reader}
\end{equation}
The conclusion applies to local Lie-algebra components of
$S(\U(3)\times\U(2))$ provided their bundle transitions and gauges have
already been identified consistently.  It neither chooses those gauges
nor establishes any missing component estimate.
\end{proposition}
\begin{proof}
The Euclidean norm of a Fourier coefficient of the tuple is at most the
sum of its component norms.  Sum with the nonnegative tail weights, apply
the component estimates, and use Cauchy--Schwarz in the finite block index.
The finite central quotient changes no local Lie-algebra estimate; global
compatibility remains the declared bundle hypothesis.
\end{proof}

\subsection{Gauge-covariant realization of the weak-sector reader}
\label{app:gauge-reader}

The Fourier reader used in the quantitative theorem is a coordinate object, but its control need not originate from ordinary derivatives of an arbitrary gauge frame.  We record a finite gauge-covariant entrance on a nonvanishing Higgs chart.  This appendix concerns the weak $SU(2)$ sector only; it is not a substitute for independent color, determinant-line or gravitational reader estimates.

On a periodic grid let $U_i(x)\in SU(2)$ and define covariant shifts and differences
\begin{equation}
 T_i^UF(x)=\rho(U_i(x))F(x+e_i),\qquad
 D_i^UF=h^{-1}(T_i^UF-F),\qquad D_I^U=D_{i_1}^U\cdots D_{i_\ell}^U.
 \label{eq:gauge-covdiff}
\end{equation}
For a fixed integer $L$ put
\begin{equation}
 \mathcal J_{h,L}(F;U)^2
 =\sum_{\ell=0}^L\ \sum_{|I|=\ell}\norm{D_I^UF}_{2,h}^2.
 \label{eq:gauge-jet-energy}
\end{equation}
Expanding the ordered product $h^{-\ell}(T_{i_1}^U-I)\cdots(T_{i_\ell}^U-I)$ gives, for each word $I$,
\begin{equation}
 D_I^UF(x)=h^{-\ell}\sum_{S\subset\{1,\ldots,\ell\}}(-1)^{\ell-|S|}F_S(x),
 \label{eq:gauge-path-expansion}
\end{equation}
where all transported endpoints $F_S(x)$ transform by the same unitary at $x$.  Hence $|D_I^UF(x)|^2$ is a fixed contraction of the positive mixed path Gram $G_I(x)_{ST}=\langle F_S(x),F_T(x)\rangle$.  The energy \eqref{eq:gauge-jet-energy} is therefore gauge invariant and can be audited from same-record mixed path data.

The discrete Kato inequality
\begin{equation}
 |\delta_i|F|(x)|\le |D_i^UF(x)|
 \label{eq:gauge-Kato}
\end{equation}
and ordinary Sobolev inequalities on the fixed grid imply, for $d\le4$,
\begin{equation}
 \norm F_{\infty,h}\le C_d\mathcal J_{h,3}(F;U),\qquad
 \max_{|I|\le L-3}\norm{D_I^UF}_{\infty,h}
 \le C_{d,L}\mathcal J_{h,L}(F;U),
 \label{eq:gauge-cov-embed}
\end{equation}
with constants independent of the links.

Assume now $|H_x|\ge\rho_*>0$.  Let $Q_x\in SU(2)$ be the algebraic frame sending $e_1$ to $H_x/|H_x|$ and put
\begin{equation}
 W_i(x)=Q_x^*U_i(x)Q_{x+e_i},\qquad
 a_i=h^{-1}(W_i-I),\qquad J_I=Q_x^*D_I^UH.
 \label{eq:gauge-normalized}
\end{equation}
The exact identities
\begin{equation}
 \delta_iJ_I=J_{iI}-a_iT_iJ_I,
 \qquad |H_{x+e_i}|=\bigl||H_x|e_1+hJ_i(x)\bigr|
 \label{eq:gauge-closed}
\end{equation}
close the ordinary difference hierarchy in terms of covariant jets.  On a fixed lower-radius chart, repeated finite-difference chain rules therefore give
\begin{equation}
 \mathcal J_{h,s+4}(H;U)+\sum_b\mathcal J_{h,s+3}(F_b;U)\le R
 \Longrightarrow
 \norm{|H|}_{C_h^{s+1}}
 +\sum_i\norm{h^{-1}\log W_i}_{C_h^s}
 +\sum_b\norm{Q^*F_b}_{C_h^s}
 \le C_{s,\rho_*,R}.
 \label{eq:gauge-normalized-reader}
\end{equation}
No derivative of the original $Q_x$ and no smallness of the original links is assumed.

The algebraic Higgs frame need not be temporal gauge.  Starting with $g_{0,x}=I$ and defining causally
\begin{equation}
 g_{j+1,x}=g_{j,x}W_0(j,x)
 \label{eq:gauge-temporal}
\end{equation}
sets the transformed temporal links to the identity.  Iterated product rules and discrete Gronwall show that one additional covariant derivative controls the transformed spatial logarithmic links through order $s$ on a fixed physical-time interval.  All transformed matter fields retain the same order of control.  Parseval and the discrete-difference symbol estimate then yield an all-mode $H^q$ bound and, by \eqref{eq:app-sharp-tail}, the tail estimate \eqref{eq:gauge-robust-tail}.

Finally, exact finite gauge invariance prevents this normalization from hiding the action row.  If $E_h$ is the Riesz gradient of the \emph{full} physical action in the mass norm, then for every fixed site gauge transformation $g$,
\begin{equation}
 E_h(g\cdot z)=T_gE_h(z),\qquad \norm{E_h(g\cdot z)}_{0,h}=\norm{E_h(z)}_{0,h}.
 \label{eq:gauge-full-covector}
\end{equation}
When $g=g(z)$ is the algebraic Higgs normalization, differentiating $g(z)\cdot z$ adds a vertical gauge tangent; the exact finite Ward identity annihilates that tangent, so the complete first variation is unchanged.  On the normalized logarithm chart $W=e^{h\mathsf A}$ with $h\norm{\mathsf A}\le c$, the physical left link variation and $\dot{\mathsf A}$ are related by
\begin{equation}
 J_h(\mathsf A)\dot{\mathsf A}
 =\int_0^1\operatorname{Ad}_{e^{th\mathsf A}}\dot{\mathsf A}\,dt,
 \qquad \norm{J_h-I}\le Cc,
 \label{eq:gauge-log-source}
\end{equation}
so the logarithmic and physical link-source norms are uniformly equivalent for fixed sufficiently small $c$.  A derivative restricted to the radial Higgs slice alone is not used as a substitute for this full covector.

\subsection{Finite-prefix stationarity from the same local action}
\label{app:prefix-stationarity}

This subsection proves \Cref{thm:prefix-stationarity,cor:prefix-table}.  The argument does not define a new field evolution and does not replace the indefinite Lorentzian action by a least-squares cost.  It studies any accepted finite prefix using the same local action \eqref{eq:native-local-action} and a positive operational potential supplied by the source geometry.

\begin{lemma}[Shifted-first-jet normal form of the unchanged local action]
\label{lem:native-firstjet-normal-form}
On a bounded nondegenerate coframe/logarithm chart, the action \eqref{eq:native-local-action} has the exact interior representation
\begin{equation}
 S_h^{\rm loc}[y]
 =h^4\sum_{x\in Q'_h}F_h^{(1)}\!\left(\Xi_h^{(1)}y(x)\right)
 +\mathscr B_h[y],
 \label{eq:native-firstjet-normal-form}
\end{equation}
where $\Xi_h^{(1)}y$ consists of a fixed finite set of shifted field values and normalized first differences, $F_h^{(1)}$ is smooth with fixed-order derivatives bounded uniformly for small $h$, and $\mathscr B_h$ depends only on the fixed collars for the interior variation problem.  Therefore the interior Euler covector and all of its mixed physical derivatives are exactly those of the shifted-first-jet bulk representative; no continuum approximation is used in this identity.
\end{lemma}
\begin{proof}
Only the gravitational plaquette needs a reduction.  Write its curvature logarithm as the exact discrete curl of the coframe connection plus a remainder that is a regular function of shifted connections.  If $C^{\mu\nu}(e)$ is the Palatini coframe coefficient, the discrete product identity
\begin{equation}
 \ip{C}{\delta_\mu^+W}
 =-\ip{\delta_\mu^-C}{W}
   +\delta_\mu^-\ip{C}{T_\mu W}
 \label{eq:native-SBP}
\end{equation}
converts both curvature-curl terms into first-jet bulk expressions plus an exact lattice divergence.  Moreover, with $e_-=T_{-\mu}e$ and $p_-=T_{-\mu}\delta_\mu^+e$,
\begin{equation}
 \delta_\mu^-C(e)
 =\int_0^1 DC(e_-+thp_-)[p_-]\,\dd t,
 \label{eq:native-C-firstjet}
\end{equation}
which is a smooth shifted-first-jet function with no inverse power of $h$.  Summing the divergence leaves only the fixed collar functional $\mathscr B_h$.

The Yang--Mills plaquette is already a uniformly regular function of shifted gauge values and first differences by \Cref{lem:native-plaquette}; the Higgs and Dirac link terms in \eqref{eq:native-matter-links}--\eqref{eq:native-spin-differences} are shifted first-jet functions directly.  Thus all sectors share the representation \eqref{eq:native-firstjet-normal-form}.  Since the transformation is an exact discrete summation-by-parts identity, interior first and higher physical variations are unchanged.
\end{proof}

\begin{lemma}[Uniform upper Hessian scale of the local action]
\label{lem:native-action-Hessian}
On a compact convex subchart of the finite field coordinates, with all shifted first-jet arguments of \eqref{eq:native-densities} remaining in a common compact coefficient set, there are constants $C_1,K_0,D_0$ independent of $h$ and of the number of nodes such that
\begin{align}
 |D^2\mathcal A_h(y)[u,v]|&\le C_1\norm u_{1,h}\norm v_{1,h},\label{eq:action-Hessian-bilinear}\\
 \norm{\nabla_h^2\mathcal A_h(y)}_{0,h\to0,h}&\le K_0h^{-2},\label{eq:action-Hessian-scale}\\
 \operatorname{osc}_{\Theta_h}\mathcal A_h&\le D_0,\label{eq:action-span}
\end{align}
where
\[
 \norm u_{1,h}^2=\norm u_{0,h}^2
 +\sum_{\mu=0}^3\norm{\delta_\mu^+u}_{0,h}^2.
\]
These are upper bounds only; no positivity or inverse for the Lorentzian action Hessian is asserted.
\end{lemma}
\begin{proof}
By \Cref{lem:native-firstjet-normal-form}, the interior local density is exactly a smooth function of finitely many shifted first jets.  Differentiating the finite sum twice gives $h^4$ times a sum of uniformly bounded coefficient bilinear forms applied to shifted values and first differences of $u,v$.  Each shift occurs only a fixed number of times.  Cauchy--Schwarz gives \eqref{eq:action-Hessian-bilinear}.  The elementary inverse estimate $\norm{\delta_\mu^+u}_{0,h}\le2h^{-1}\norm u_{0,h}$ gives \eqref{eq:action-Hessian-scale}.  Finally the density is uniformly bounded on the compact chart and $h^4$ times the number of lattice cells in the fixed cylinder is uniformly bounded, giving \eqref{eq:action-span}.  The coframe/spin-connection chain is contained in the smooth density and is differentiated throughout.
\end{proof}

\begin{proof}[Proof of \Cref{thm:prefix-stationarity}]
Let $x=\theta_j$, $y=\theta_{j+1}$, $d=y-x$, and write $p_h=\nabla_h\Psi_h$.  By \eqref{eq:action-Hessian-scale}, Taylor's integral formula gives
\begin{equation}
 \mathcal A_h(x)-\mathcal A_h(y)
 \ge -(E_h(y),d)_{0,h}-\tfrac12K_0h^{-2}\norm d_{0,h}^2.
 \label{eq:prefix-descent1}
\end{equation}
The lower bound in \eqref{eq:prefix-Fisher} yields
\[
 (p_h(y)-p_h(x),d)_{0,h}\ge g_*\norm d_{0,h}^2.
\]
Substitute \eqref{eq:prefix-balance} and put
\[
 a_h=g_*\upsilon_h^{-1}-\tfrac12K_0h^{-2}=a_0h^{-2}>0,
 \qquad
 \ell_h=g^*\upsilon_h^{-1}+K_0h^{-2}=\ell_0h^{-2}.
\]
Young's inequality gives
\begin{equation}
 \mathcal A_h(x)-\mathcal A_h(y)
 \ge \tfrac12a_h\norm d_{0,h}^2
 -\frac1{2a_h}\norm{b_j}_{0,h}^2.
 \label{eq:prefix-descent2}
\end{equation}
Moreover
\[
 E_h(x)=b_j-\upsilon_h^{-1}(p_h(y)-p_h(x))
                 +(E_h(x)-E_h(y)),
\]
so
\[
 \norm{E_h(x)}_{0,h}\le\norm{b_j}_{0,h}+\ell_h\norm d_{0,h}.
\]
Combining with \eqref{eq:prefix-descent2}, squaring, and summing over $j$ gives
\[
 \sum_{j<J}\norm{E_h(\theta_j)}_{0,h}^2
 \le \frac{4\ell_h^2}{a_h}
       \{\mathcal A_h(\theta_0)-\mathcal A_h(\theta_J)\}
 +\left(2+\frac{2\ell_h^2}{a_h^2}\right)
       \sum_{j<J}\norm{b_j}_{0,h}^2.
\]
Since $4\ell_h^2/a_h=C_0h^{-2}$, the second coefficient is a cutoff-independent constant, and the telescoped action decrement is at most the uniform oscillation $D_0$, this is \eqref{eq:prefix-average}.  Taking the minimum of the finite list gives \eqref{eq:prefix-selected}.
\end{proof}

\begin{proof}[Proof of \Cref{cor:prefix-table}]
Let $G$ be the event that the accepted prefix stays inside the declared safe set.  On $G$, \eqref{eq:prefix-average} holds pathwise.  Multiplying by $\mathbf1_G$, taking expectations, and enlarging the nonnegative balance sum to all safe occupied edges gives
\[
 \mathbb E\!\left[\mathbf1_G\frac1J\sum_{j<J}
     \norm{E_h(\theta_j)}_{0,h}^2\right]
 \le C_0h^{-2}\frac{D_0}{J}+D_1\mathcal Q_{h,J}^{\rm bal}.
\]
If the best safe slope exceeds $h^\alpha$, the mean square exceeds $h^{2\alpha}$.  Markov's inequality therefore gives the last two terms of \eqref{eq:prefix-table-probability}; adding $\mathbb P(G^c)=p_h^{\rm exit}$ proves the bound.  Summability along dyadic cutoffs and the first Borel--Cantelli lemma yield eventual certification without independence.
\end{proof}

\subsection{Complete physical variations on constrained source charts}
\label{app:complete-source}

The source estimates must concern the same complete action covector as the
continuum variation.  A restriction to a constraint surface, a selected
principal block, or a partial reader does not by itself provide this
identification.  The following finite statements separate these requirements.

\begin{proposition}[Tangential stationarity and the missing normal row]
\label{prop:normal-source}
Let $S_h$ be differentiable on a finite positive Hilbert space $X_h$, and let
$C_h:X_h\to Z_h$ be differentiable with surjective derivative at $x_h$.
Assume $B_h:Z_h\to X_h$ is a right inverse of $DC_h(x_h)$ and put
$P_h=I-B_hDC_h(x_h)$.  Suppose every lifted unit physical test satisfies
\[
 \norm{P_hI_hv}\le L_h^{\rm tan},\qquad
 \norm{DC_h(x_h)I_hv}\le L_h^{\rm nor}.
\]
If
\[
 |DS_h(x_h)[w]|\le a_h\norm w\quad(w\in\ker DC_h(x_h)),
 \qquad
 |DS_h(x_h)[B_hq]|\le b_h\norm q,
\]
then
\begin{equation}
 \sup_{\norm v\le1}|DS_h(x_h)[I_hv]|
 \le a_hL_h^{\rm tan}+b_hL_h^{\rm nor}.
 \label{eq:normal-source-bound}
\end{equation}
In particular exact stationarity of $S_h$ restricted to $C_h=0$ does not
imply complete physical stationarity unless the normal contribution also
vanishes or all physical lifts are tangent.
\end{proposition}
\begin{proof}
Since $DC_hB_h=I$, $P_hI_hv$ lies in the tangent kernel and
$I_hv=P_hI_hv+B_hDC_hI_hv$.  Apply the two covector bounds.
For the failure of the unrestricted implication, take
$S(x,y)=x^2+y$ and $C(x,y)=y$.  The restriction is stationary at $(0,0)$,
whereas $DS(0,0)[(0,1)]=1$.
\end{proof}

\begin{proposition}[Residual-preserving extension of a constrained identity]
\label{prop:constraint-extension}
Let a finite-dimensional regular constraint chart have coordinates $(s,c)$,
with the constraint surface $c=0$.  Let $J(s,c)$ be a $C^1$ matrix-valued map.
Then
\begin{equation}
 J(s,c)=J(s,0)+\int_0^1D_cJ(s,tc)[c]\,\dd t.
 \label{eq:constraint-extension}
\end{equation}
If $c_a(t)$ satisfies the exact tracking equation
$\dot c_a=A_a(t)c_a+r_a(t)$, with
$\int_0^T\norm{A_a(t)}\dd t\le M$ and
$\norm{c_a(0)}+\norm{r_a}_{L^1(0,T)}\le Ca^p$, then
\begin{equation}
 \sup_{0\le t\le T}\norm{c_a(t)}\le Ce^M a^p.
 \label{eq:constraint-tracking}
\end{equation}
A uniformly bounded $D_cJ$ consequently gives an $O(a^p)$ difference
between the full and constrained principal matrices on these paths.
This conclusion gives no estimate for an independently occurring
state-only source $f(s,c)$ without a bound or an identity for that source.
\end{proposition}
\begin{proof}
The first identity is the fundamental theorem of calculus in the normal
coordinate.  The variation-of-constants inequality followed by Gronwall
gives \eqref{eq:constraint-tracking}.  Substitution into
\eqref{eq:constraint-extension} gives the matrix estimate.
For the last assertion take $J(s,c)$ constant while $f(s,c)$ is any fixed
nonzero covector, or any cutoff-dependent multiple of it.  The principal
identity then holds exactly and imposes no restriction on the source.
\end{proof}

These observations apply before invoking either variational closure or the
coupled stability estimate.  A constraint-compatible principal operator
must be accompanied by control of the full source, including its normal
and state-only rows.  Likewise, an exact algebraic decision procedure for
stationarity does not establish the existence of a stationary record until
its actual coefficients, physical-domain conditions and complete covector
have been supplied and evaluated.  No such source-existence conclusion is
used in the continuum theorems.


\clearpage
\begin{thebibliography}{99}
\small
\setlength{\itemsep}{2pt}

\bibitem{ConnesLott1991}
A.~Connes and J.~Lott,
Particle models and noncommutative geometry,
\emph{Nucl. Phys. B Proc. Suppl.} \textbf{18B} (1991), 29--47.
\href{https://doi.org/10.1016/0920-5632(91)90120-4}{doi:10.1016/0920-5632(91)90120-4}.

\bibitem{Krajewski1998}
T.~Krajewski,
Classification of finite spectral triples,
\emph{J. Geom. Phys.} \textbf{28} (1998), 1--30.
\href{https://doi.org/10.1016/S0393-0440(97)00068-5}{doi:10.1016/S0393-0440(97)00068-5}.

\bibitem{ChamseddineConnes1997}
A.~H.~Chamseddine and A.~Connes,
The spectral action principle,
\emph{Commun. Math. Phys.} \textbf{186} (1997), 731--750.
\href{https://doi.org/10.1007/s002200050126}{doi:10.1007/s002200050126}.

\bibitem{ChamseddineConnesMarcolli2007}
A.~H.~Chamseddine, A.~Connes, and M.~Marcolli,
Gravity and the standard model with neutrino mixing,
\emph{Adv. Theor. Math. Phys.} \textbf{11} (2007), 991--1089.
\href{https://doi.org/10.4310/ATMP.2007.v11.n6.a3}{doi:10.4310/ATMP.2007.v11.n6.a3}.

\bibitem{Barrett2007}
J.~W.~Barrett,
A Lorentzian version of the non-commutative geometry of the standard model of particle physics,
\emph{J. Math. Phys.} \textbf{48} (2007), 012303.
\href{https://doi.org/10.1063/1.2408400}{doi:10.1063/1.2408400}.

\bibitem{Finster2016}
F.~Finster,
\emph{The Continuum Limit of Causal Fermion Systems:
From Planck Scale Structures to Macroscopic Physics},
Springer, Cham, 2016.
\href{https://doi.org/10.1007/978-3-319-42067-7}{doi:10.1007/978-3-319-42067-7}.

\bibitem{Regge1961}
T.~Regge,
General relativity without coordinates,
\emph{Nuovo Cimento} \textbf{19} (1961), 558--571.
\href{https://doi.org/10.1007/BF02733251}{doi:10.1007/BF02733251}.

\bibitem{CheegerMullerSchrader1984}
J.~Cheeger, W.~M\"uller, and R.~Schrader,
On the curvature of piecewise flat spaces,
\emph{Commun. Math. Phys.} \textbf{92} (1984), 405--454.
\href{https://doi.org/10.1007/BF01210729}{doi:10.1007/BF01210729}.

\bibitem{ChristiansenRegge2011}
S.~H.~Christiansen,
On the linearization of Regge calculus,
\emph{Numer. Math.} \textbf{119} (2011), 613--640.
\href{https://doi.org/10.1007/s00211-011-0394-z}{doi:10.1007/s00211-011-0394-z}.

\bibitem{ChristiansenRegge2024}
S.~H.~Christiansen,
On the definition of curvature in Regge calculus,
\emph{IMA J. Numer. Anal.} \textbf{44} (2024), 2698--2715.
\href{https://doi.org/10.1093/imanum/drad095}{doi:10.1093/imanum/drad095}.

\bibitem{GawlikNeunteufelScalar2025}
E.~S.~Gawlik and M.~Neunteufel,
Finite element approximation of scalar curvature in arbitrary dimension,
\emph{Math. Comp.} \textbf{94} (2025), 2685--2722.
\href{https://doi.org/10.1090/mcom/4038}{doi:10.1090/mcom/4038}.

\bibitem{GawlikNeunteufelEinstein2026}
E.~S.~Gawlik and M.~Neunteufel,
Finite element approximation of the Einstein tensor,
\emph{IMA J. Numer. Anal.} \textbf{46} (2026), 500--539.
\href{https://doi.org/10.1093/imanum/draf004}{doi:10.1093/imanum/draf004}.

\bibitem{Wilson1974}
K.~G.~Wilson,
Confinement of quarks,
\emph{Phys. Rev. D} \textbf{10} (1974), 2445--2459.
\href{https://doi.org/10.1103/PhysRevD.10.2445}{doi:10.1103/PhysRevD.10.2445}.

\bibitem{MarsdenPatrickShkoller1998}
J.~E.~Marsden, G.~W.~Patrick, and S.~Shkoller,
Multisymplectic geometry, variational integrators, and nonlinear PDEs,
\emph{Commun. Math. Phys.} \textbf{199} (1998), 351--395.
\href{https://doi.org/10.1007/s002200050505}{doi:10.1007/s002200050505}.

\bibitem{ArnoldFalkWinther2006}
D.~N.~Arnold, R.~S.~Falk, and R.~Winther,
Finite element exterior calculus, homological techniques, and applications,
\emph{Acta Numer.} \textbf{15} (2006), 1--155.
\href{https://doi.org/10.1017/S0962492906210018}{doi:10.1017/S0962492906210018}.

\bibitem{ArnoldFalkWinther2010}
D.~N.~Arnold, R.~S.~Falk, and R.~Winther,
Finite element exterior calculus: from Hodge theory to numerical stability,
\emph{Bull. Amer. Math. Soc.} \textbf{47} (2010), 281--354.
\href{https://doi.org/10.1090/S0273-0979-10-01278-4}{doi:10.1090/S0273-0979-10-01278-4}.

\bibitem{ChristiansenHalvorsen2009}
S.~H.~Christiansen and T.~G.~Halvorsen,
Convergence of lattice gauge theory for Maxwell's equations,
\emph{BIT Numer. Math.} \textbf{49} (2009), 645--667.
\href{https://doi.org/10.1007/s10543-009-0242-z}{doi:10.1007/s10543-009-0242-z}.

\bibitem{ChristiansenHalvorsen2011}
S.~H.~Christiansen and T.~G.~Halvorsen,
Discretizing the Maxwell--Klein--Gordon equation by the lattice gauge theory formalism,
\emph{IMA J. Numer. Anal.} \textbf{31} (2011), 1--24.
\href{https://doi.org/10.1093/imanum/drp019}{doi:10.1093/imanum/drp019}.

\bibitem{ChristiansenHalvorsenScheid2023}
S.~H.~Christiansen, T.~G.~Halvorsen, and C.~Scheid,
Convergence of a discretization of the Maxwell--Klein--Gordon equation based on finite element methods and lattice gauge theory,
\emph{Numer. Methods Partial Differential Equations} \textbf{39} (2023), 3271--3308.
\href{https://doi.org/10.1002/num.23008}{doi:10.1002/num.23008}.

\bibitem{ChristiansenHalvorsen2012}
S.~H.~Christiansen and T.~G.~Halvorsen,
A simplicial gauge theory,
\emph{J. Math. Phys.} \textbf{53} (2012), 033501.
\href{https://doi.org/10.1063/1.3692167}{doi:10.1063/1.3692167}.

\bibitem{ShaliziCrutchfield2001}
C.~R.~Shalizi and J.~P.~Crutchfield,
Computational mechanics: Pattern and prediction, structure and simplicity,
\emph{J. Stat. Phys.} \textbf{104} (2001), 817--879.
\href{https://doi.org/10.1023/A:1010388907793}{doi:10.1023/A:1010388907793}.

\bibitem{ChiribellaEtAl2009}
G.~Chiribella, G.~M.~D'Ariano, and P.~Perinotti,
Theoretical framework for quantum networks,
\emph{Phys. Rev. A} \textbf{80} (2009), 022339.
\href{https://doi.org/10.1103/PhysRevA.80.022339}{doi:10.1103/PhysRevA.80.022339}.

\bibitem{PollockEtAl2018}
F.~A.~Pollock, C.~Rodr\'iguez-Rosario, T.~Frauenheim,
M.~Paternostro, and K.~Modi,
Non-Markovian quantum processes: Complete framework and efficient characterization,
\emph{Phys. Rev. A} \textbf{97} (2018), 012127.
\href{https://doi.org/10.1103/PhysRevA.97.012127}{doi:10.1103/PhysRevA.97.012127}.

\bibitem{PelissierRenewalSpacetime2026}
A.~Pelissier,
\newblock Renewal Geometry and the emergence of Lorentzian spacetime,
\newblock HAL preprint hal-05744772, 2026.
\newblock \url{https://hal.science/hal-05744772}

\bibitem{PelissierSpectral2026}
A.~Pelissier,
\newblock From Predictive Dynamics to Spectral Geometry,
\newblock HAL preprint hal-05752672, 2026.
\newblock \url{https://hal.science/hal-05752672}

\bibitem{PelissierDuality2026}
A.~Pelissier,
\newblock A Duality Between Spacetime--Gauge--Incidence Actions and Matter Multiplicity,
\newblock HAL preprint hal-05755429, 2026.
\newblock \url{https://hal.science/hal-05755429}

\bibitem{Uhlenbeck1982}
K.~K.~Uhlenbeck,
Connections with $L^p$ bounds on curvature,
\emph{Commun. Math. Phys.} \textbf{83} (1982), 31--42.
\href{https://doi.org/10.1007/BF01947069}{doi:10.1007/BF01947069}.

\bibitem{Sedlacek1982}
S.~Sedlacek,
A direct method for minimizing the Yang--Mills functional over $4$-manifolds,
\emph{Commun. Math. Phys.} \textbf{86} (1982), 515--527.
\href{https://doi.org/10.1007/BF01214887}{doi:10.1007/BF01214887}.

\bibitem{Lions1985}
P.-L.~Lions,
The concentration-compactness principle in the calculus of variations.
The limit case, Part~1,
\emph{Rev. Mat. Iberoam.} \textbf{1} (1985), no.~1, 145--201.
\href{https://doi.org/10.4171/RMI/6}{doi:10.4171/RMI/6}.

\bibitem{DiPernaMajda1987}
R.~J.~DiPerna and A.~J.~Majda,
Oscillations and concentrations in weak solutions of the incompressible fluid equations,
\emph{Commun. Math. Phys.} \textbf{108} (1987), 667--689.
\href{https://doi.org/10.1007/BF01214424}{doi:10.1007/BF01214424}.

\bibitem{AlibertBouchitte1997}
J.~J.~Alibert and G.~Bouchitt\'e,
Non-uniform integrability and generalized Young measure,
\emph{J. Convex Anal.} \textbf{4} (1997), 129--147.

\bibitem{Tartar1990}
L.~Tartar,
H-measures, a new approach for studying homogenisation, oscillations and concentration effects in partial differential equations,
\emph{Proc. Roy. Soc. Edinburgh Sect. A} \textbf{115} (1990), 193--230.
\href{https://doi.org/10.1017/S0308210500020606}{doi:10.1017/S0308210500020606}.

\bibitem{Gerard1991}
P.~G\'erard,
Microlocal defect measures,
\emph{Comm. Partial Differential Equations} \textbf{16} (1991), 1761--1794.
\href{https://doi.org/10.1080/03605309108820822}{doi:10.1080/03605309108820822}.

\bibitem{Isaacson1968}
R.~A.~Isaacson,
Gravitational radiation in the limit of high frequency. II. Nonlinear terms and the effective stress tensor,
\emph{Phys. Rev.} \textbf{166} (1968), 1272--1280.
\href{https://doi.org/10.1103/PhysRev.166.1272}{doi:10.1103/PhysRev.166.1272}.

\bibitem{Burnett1989}
G.~A.~Burnett,
The high-frequency limit in general relativity,
\emph{J. Math. Phys.} \textbf{30} (1989), 90--96.
\href{https://doi.org/10.1063/1.528594}{doi:10.1063/1.528594}.

\bibitem{HuneauLuk2024}
C.~Huneau and J.~Luk,
Burnett's conjecture in generalized wave coordinates,
\href{https://arxiv.org/abs/2403.03470}{arXiv:2403.03470 [math.AP]}, 2024.

\bibitem{BourguignonGauduchon1992}
J.-P.~Bourguignon and P.~Gauduchon,
Spineurs, op\'erateurs de Dirac et variations de m\'etriques,
\emph{Commun. Math. Phys.} \textbf{144} (1992), 581--599.
\href{https://doi.org/10.1007/BF02099184}{doi:10.1007/BF02099184}.

\bibitem{MullerNowaczyk2017}
O.~M\"uller and N.~Nowaczyk,
A universal spinor bundle and the Einstein--Dirac--Maxwell equation as a variational theory,
\emph{Lett. Math. Phys.} \textbf{107} (2017), 933--961.
\href{https://doi.org/10.1007/s11005-016-0929-4}{doi:10.1007/s11005-016-0929-4}.

\bibitem{Kibble1961}
T.~W.~B.~Kibble,
Lorentz invariance and the gravitational field,
\emph{J. Math. Phys.} \textbf{2} (1961), 212--221.
\href{https://doi.org/10.1063/1.1703702}{doi:10.1063/1.1703702}.

\bibitem{HehlEtAl1976}
F.~W.~Hehl, P.~von der Heyde, G.~D.~Kerlick, and J.~M.~Nester,
General relativity with spin and torsion: Foundations and prospects,
\emph{Rev. Mod. Phys.} \textbf{48} (1976), 393--416.
\href{https://doi.org/10.1103/RevModPhys.48.393}{doi:10.1103/RevModPhys.48.393}.

\bibitem{NielsenNinomiya1981}
H.~B.~Nielsen and M.~Ninomiya,
A no-go theorem for regularizing chiral fermions,
\emph{Phys. Lett. B} \textbf{105} (1981), 219--223.
\href{https://doi.org/10.1016/0370-2693(81)91026-1}{doi:10.1016/0370-2693(81)91026-1}.

\bibitem{Friedan1982}
D.~Friedan,
A proof of the Nielsen--Ninomiya theorem,
\emph{Commun. Math. Phys.} \textbf{85} (1982), 481--490.

\bibitem{LigterinkWaletBishop2000}
N.~E.~Ligterink, N.~R.~Walet, and R.~F.~Bishop,
Toward a many-body treatment of Hamiltonian lattice $\mathrm{SU}(N)$ gauge theory,
\emph{Ann. Phys.} \textbf{284} (2000), 215--262.
\href{https://doi.org/10.1006/aphy.2000.6070}{doi:10.1006/aphy.2000.6070}.

\bibitem{DalMaso1993}
G.~Dal Maso,
\emph{An Introduction to $\Gamma$-Convergence},
Progress in Nonlinear Differential Equations and Their Applications,
vol.~8, Birkh\"auser, Boston, 1993.
\href{https://doi.org/10.1007/978-1-4612-0327-8}{doi:10.1007/978-1-4612-0327-8}.

\bibitem{Braides2002}
A.~Braides,
\emph{Gamma-Convergence for Beginners},
Oxford Lecture Series in Mathematics and Its Applications, vol.~22,
Oxford University Press, Oxford, 2002.
\href{https://doi.org/10.1093/acprof:oso/9780198507840.001.0001}{doi:10.1093/acprof:oso/9780198507840.001.0001}.

\bibitem{JerrardSternberg2009}
R.~L.~Jerrard and P.~Sternberg,
Critical points via $\Gamma$-convergence: general theory and applications,
\emph{J. Eur. Math. Soc.} \textbf{11} (2009), 705--753.
\href{https://doi.org/10.4171/JEMS/164}{doi:10.4171/JEMS/164}.

\bibitem{Attouch1984}
H.~Attouch,
\emph{Variational Convergence for Functions and Operators},
Applicable Mathematics Series, Pitman, Boston, 1984.

\bibitem{SandierSerfaty2004}
E.~Sandier and S.~Serfaty,
Gamma-convergence of gradient flows with applications to Ginzburg--Landau,
\emph{Comm. Pure Appl. Math.} \textbf{57} (2004), 1627--1672.
\href{https://doi.org/10.1002/cpa.20046}{doi:10.1002/cpa.20046}.

\bibitem{Simon1986}
J.~Simon,
Compact sets in the space $L^p(0,T;B)$,
\emph{Ann. Mat. Pura Appl.} \textbf{146} (1986), 65--96.
\href{https://doi.org/10.1007/BF01762360}{doi:10.1007/BF01762360}.
% The publisher's issue metadata gives December 1986; the work is often cited as 1987.

\bibitem{GerochTraschen1987}
R.~Geroch and J.~Traschen,
Strings and other distributional sources in general relativity,
\emph{Phys. Rev. D} \textbf{36} (1987), 1017--1031.
\href{https://doi.org/10.1103/PhysRevD.36.1017}{doi:10.1103/PhysRevD.36.1017}.

\bibitem{LeFlochMardare2007}
P.~G.~LeFloch and C.~Mardare,
Definition and stability of Lorentzian manifolds with distributional curvature,
\emph{Portugaliae Math.} \textbf{64} (2007), 535--573.
\href{https://doi.org/10.4171/PM/1794}{doi:10.4171/PM/1794}.

\bibitem{ErlerFreitas2024}
J.~Erler and A.~Freitas,
Electroweak model and constraints on new physics,
in S.~Navas et al. (Particle Data Group), Review of Particle Physics,
\emph{Phys. Rev. D} \textbf{110} (2024), 030001, Sec.~10.
\url{https://pdg.lbl.gov/2024/reviews/rpp2024-rev-standard-model.pdf}.

\bibitem{Luscher1999}
M.~L\"uscher,
Abelian chiral gauge theories on the lattice with exact gauge invariance,
\emph{Nucl. Phys. B} \textbf{549} (1999), 295--334.
\href{https://doi.org/10.1016/S0550-3213(99)00115-7}{doi:10.1016/S0550-3213(99)00115-7}.

\bibitem{FriedrichRendall2000}
H.~Friedrich and A.~D.~Rendall,
The Cauchy problem for the Einstein equations,
in B.~G.~Schmidt (ed.), \emph{Einstein's Field Equations and Their Physical Implications},
Lecture Notes in Physics \textbf{540}, Springer, Berlin, 2000, pp.~127--223.
\href{https://arxiv.org/abs/gr-qc/0002074}{arXiv:gr-qc/0002074}.

\bibitem{Kato1975}
T.~Kato,
The Cauchy problem for quasi-linear symmetric hyperbolic systems,
\emph{Arch. Ration. Mech. Anal.} \textbf{58} (1975), 181--205.
\href{https://doi.org/10.1007/BF00280740}{doi:10.1007/BF00280740}.

\bibitem{ChoquetBruhatIsenbergPollack2007}
Y.~Choquet-Bruhat, J.~Isenberg, and D.~Pollack,
The constraint equations for the Einstein--scalar field system on compact manifolds,
\emph{Class. Quantum Grav.} \textbf{24} (2007), 809--828.
\href{https://arxiv.org/abs/gr-qc/0610045}{arXiv:gr-qc/0610045}.

\bibitem{IsenbergMaxwell2025}
J.~Isenberg and D.~Maxwell,
A phase space approach to the conformal construction of non-vacuum initial data sets in general relativity,
\emph{Ann. Henri Poincar\'e} \textbf{26} (2025), 2505--2555.
\href{https://doi.org/10.1007/s00023-024-01492-5}{doi:10.1007/s00023-024-01492-5}.
\end{thebibliography}


\end{document}
